% Macro safety net for cross-volume universal-conductor references,
% provided locally so the chapter remains portable
% under standalone extraction. Definitions match
% chapters/theory/universal_conductor_K_platonic.tex.
\providecommand{\BRST}{\mathrm{BRST}}
\providecommand{\BRSTGauged}{\mathrm{BRSTGaugedChirAlg}}

\chapter{Affine Kac--Moody algebras}\label{chap:kac-moody-koszul}
\label{chap:kac-moody}

The Heisenberg algebra is abelian: its shadow tower terminates at
degree~$2$, and every obstruction is a scalar multiple of~$\kappa^{\mathrm{Heis}}$.
The first nonabelian chiral algebra is
$\widehat{\mathfrak{g}}_k$, and the question that
governs this chapter is whether the Lie bracket destroys the
finiteness of the obstruction tower or preserves it.

\medskip
\noindent\emph{Cross-volume anchors.}
Every numerical invariant of $\widehat{\fg}_k$ recorded in this
chapter is an evaluation of two universal objects established
elsewhere in Vol~I. The modular characteristic
\(
 \kappa^{\mathrm{KM}}(V_k(\fg)) = \dim(\fg)(k+h^\vee)/(2h^\vee)
\)
is the $V_k(\fg)$-evaluation of the universal conductor functor
$K(\cA) = -c_{\mathrm{ghost}}(\BRST(\cA))$ on
BRST-gauged chiral algebras
(Corollary~\ref{cor:uc-K-affine-KM},
Chapter~\ref{chap:universal-conductor}); the underlying ghost
content is the $bc$ system of the WZW BRST resolution at level $k$.
The classification of $\widehat{\fg}_k$ as quadrichotomy class~L,
shadow depth $r_{\max} = 3$, follows from the Riccati master
equation $H(t)^2 = t^4 Q_c(t)$ with discriminant
$Q_c$ of corank one: the cubic shadow
$\mathfrak{C}(x,y,z) = \kappa_{\fg}(x,[y,z])$
fills the simple-pole channel and the Jacobi identity kills
every higher obstruction
(Theorem~\ref{thm:quadrichotomy} clauses~(\ref{quad:i})--(\ref{quad:iii}),
Chapter~\ref{chap:shadow-quadrichotomy-platonic}).
All formulas in this chapter for $\kappa^{\mathrm{KM}}$ and the
collision residue $r^{\mathrm{coll}}_{\mathrm{KM}}(z)$ are cross-checked against
the master table of Chapter~\ref{ch:landscape-census}
(\texttt{landscape\_census.tex}, rows ``Affine Kac--Moody'') and the
modular Koszul triple
$\mathfrak{T}_{\widehat{\fg}_k} =
(\widehat{\fg}_k, \widehat{\fg}_{-k-2h^\vee}, k\Omega_{\mathrm{tr}}/z)$
of equation~\eqref{eq:census-triples}; convention bridges between the
trace-form normalisation $r^{\mathrm{coll}}(z) = k\Omega_{\mathrm{tr}}/z$
(vanishing at $k = 0$)
and the KZ normalisation $r^{\mathrm{KZ}}(z) = \Omega/((k+h^\vee)z)$ (whose
$k=0$ limit gives $\Omega/(h^\vee z) \neq 0$ for non-abelian $\fg$,
reflecting that the Lie bracket persists in the abelian-level
limit) are catalogued in Remark~\ref{rem:km-collision-residue-rmatrix} below.
\medskip
%
It preserves it. The Lie bracket of~$\fg$ appears as the
cubic shadow~$\mathfrak{C}(x,y,z) = \kappa_{\fg}(x,[y,z])$ at
degree~$3$; the Jacobi identity forces the quartic shadow to vanish;
the tower terminates (class~$L$, $r_{\max} = 3$). Gauge symmetry
introduces noncommutativity but not infinite obstruction complexity:
the structure constants of~$\fg$ produce exactly one new shadow
beyond the Gaussian, and the Jacobi identity is the mechanism that
kills the rest. This is the content of the assertion that gauge
theory is simpler than gravity; for the Virasoro algebra
(Chapter~\ref{chap:w-algebras}), where no Jacobi relation governs the
quartic Virasoro OPE, the tower does not terminate.

Three structures meet at the shifted parameter $k+h^\vee$ but remain
different objects.  Verdier duality on the bar coalgebra gives the
Feigin--Frenkel level involution $k\mapsto -k-2h^\vee$
(Theorem~\ref{thm:bar-cobar-verdier}); the critical specialization
$k=-h^\vee$ is the point where the shifted bar curvature vanishes; the
degree-zero chiral centre at that point is the Feigin--Frenkel oper
algebra.  The derived chiral centre is the Hochschild closed-sector
cochain object, not the Koszul dual and not a physical bulk before an
open--closed comparison map is supplied.  The five theorem-surface statements below
are used only in the ambient stated in their cited proofs: the affine
universal algebra lies on the PBW-Koszul locus, while simple quotients
at admissible level require the separate admissible analysis of
Remark~\ref{rem:admissible-koszul-status}.

\begin{table}[ht]
\centering
\small
\caption{Five-theorem surface for
$\widehat{\fg}_k$ at generic level
$k \neq -h^\vee$.}\label{tab:km-five-theorems}
\begin{tabular}{llll}
\toprule
\textbf{Theorem} & \textbf{Statement for $\widehat{\fg}_k$}
 & \textbf{Status} & \textbf{Reference} \\
\midrule
A (adjunction) &
 $(\widehat{\fg}_k)^! = \mathrm{CE}^{\mathrm{ch}}(\widehat{\fg}_{-k-2h^\vee})$
 & Proved & Thm~\ref{thm:universal-kac-moody-koszul} \\
B (inversion) &
 $\Omega(\barB(\widehat{\fg}_k)) \xrightarrow{\sim}
 \widehat{\fg}_k$ on the weight-completed Koszul surface
 & Conditional/weight-completed &
 Thm~\ref{thm:chiral-positselski-weight-completed};
 genus-$1$ checks in Thms~\ref{thm:sl2-genus1-inversion},
 \ref{thm:sl3-genus1-inversion} \\
C (complementarity) &
 $Q_g(\widehat{\fg}_k) \oplus
 Q_g(\widehat{\fg}_{-k-2h^\vee})
 = H^*(\overline{\cM}_g, Z)$
 & Conditional on the complementarity hypotheses &
 Thm~\ref{thm:sl2-genus1-complementarity},
 \ref{thm:sl3-genus1-complementarity} \\
D (modular char.) &
 $\kappa(\widehat{\fg}_k)
 = \frac{\dim\fg\,(k{+}h^\vee)}{2h^\vee}$;\quad
 $\kappa + \kappa' = 0$ (affine KM)
 & Proved & Prop~\ref{prop:kappa-anti-symmetry-ff},
 Rem~\ref{rem:sl3-universality} \\
H (chiral Hochschild) &
 ChirHoch${}^*(\widehat{\fg}_k)$ concentrated on the generic
 PBW-Koszul locus
 & Proved/generic & Thms~\ref{thm:main-koszul-hoch},
 \ref{thm:hochschild-polynomial-growth};
 Prop.~\ref{prop:chirhoch1-affine-km} \\
\bottomrule
\end{tabular}
\end{table}

\begin{table}[ht]
\centering
\small
\caption{Shadow archetype data for affine
Kac--Moody.}\label{tab:km-shadow-archetype}
\begin{tabular}{ll}
\toprule
\textbf{Invariant} & \textbf{Value} \\
\midrule
Class & L (Lie/tree) \\
Shadow depth $r_{\max}$ & $3$ \\
$\kappa(\widehat{\fg}_k)$ &
 $\dim(\fg)(k{+}h^\vee)/(2h^\vee)$ \\
Cubic shadow $\mathfrak{C}$ &
 $\kappa_{\fg}(x,[y,z])$ (Lie bracket) \\
Quartic $o_4$ & $0$ (Jacobi identity) \\
$r$-matrix data &
 $r^{\mathrm{coll}}(z)=k\Omega_{\mathrm{tr}}/z$;
 $r^{\mathrm{KZ}}(z)=\Omega_{\mathrm{KZ}}/\bigl((k{+}h^\vee)\,z\bigr)$
 (Comp.~\ref{comp:sl2-collision-residue-kz}) \\
Koszul dual & $\mathrm{CE}^{\mathrm{ch}}(\widehat{\fg}_{-k-2h^\vee})$
 (sharing $\kappa$ with $\widehat{\fg}_{-k-2h^\vee}$) \\
Complementarity & $\kappa + \kappa' = 0$ (affine KM) \\
Critical level & $k = -h^\vee$: uncurved,
 FF center \\
$L_\infty$-formality &
 Conditional strict representative
 (Thm~\ref{thm:km-strictification}) \\
\bottomrule
\end{tabular}
\end{table}

\noindent
The modular Koszul triple
(Definition~\ref{def:modular-koszul-triple}) of the affine
Kac--Moody algebra at generic level $k \neq -h^\vee$ is
\begin{equation}\label{eq:km-triple}
\mathfrak{T}_{\widehat{\fg}_k}
\;=\;
\Bigl(\,\widehat{\fg}_k,\;\;
\widehat{\fg}_{-k-2h^\vee},\;\;
r^{\mathrm{coll}}(z) = \frac{k\Omega_{\mathrm{tr}}}{z}\,\Bigr),
\end{equation}
where $\Omega_{\mathrm{tr}}$ is the trace-form Casimir. The KZ
connection uses the separate coefficient
$r^{\mathrm{KZ}}(z)=\Omega_{\mathrm{KZ}}/((k+h^\vee)z)$
(Computation~\ref{comp:sl2-collision-residue-kz},
Remark~\ref{rem:km-collision-residue-rmatrix});
at the critical level $k = -h^\vee$ it degenerates, tracking the
Sugawara singularity, and the Feigin--Frenkel involution is the fixed
point of the same denominator that drives the KZ connection
(Theorem~\textup{\ref{thm:universal-kac-moody-koszul}}).
The Feigin--Frenkel involution $k \mapsto -k - 2h^\vee$ is
Verdier duality on the bar coalgebra, and the
$r^{\mathrm{KZ}}$ coefficient is the Knizhnik--Zamolodchikov
$r$-matrix (Computation~\ref{comp:sl2-collision-residue-kz}).

\begin{remark}[Five-$\kappa$ stratification: affine Kac--Moody]
\label{rem:km-five-kappa-stratification}
\index{kappa stratification@$\kappa$-stratification!affine Kac--Moody row L}
The five $\kappa$-measurements
$\{\kappa_{\mathrm{cat}}, \kappa^{\mathrm{Hodge}}_{\mathrm{ch}},
\kappa^{\mathrm{Heis}}_{\mathrm{ch}}, \kappa_{\mathrm{BKM}},
\kappa_{\mathrm{fiber}}\}$ on the row $\mathsf{L}$ of the
$5\times 5$ stratification matrix evaluate at level $k\neq -h^\vee$:
\[
\kappa_{\mathrm{cat}}^{V_k(\fg)}=\frac{\dim(\fg)(k+h^\vee)}{2h^\vee},
\qquad
\kappa^{\mathrm{Hodge}}_{\mathrm{ch}}(V_k(\fg))=\rho_k\cdot
\kappa_{\mathrm{cat}},
\qquad
\kappa^{\mathrm{Heis}}_{\mathrm{ch}}(V_k(\fg))=
\operatorname{rank}(\fg)\cdot k,
\]
with $\rho_k = c(k,\fg)/(2\kappa_{\mathrm{cat}})$ the anomaly
ratio. The BKM and fibre measurements
$\kappa_{\mathrm{BKM}}, \kappa_{\mathrm{fiber}}$ vanish on
class~$\mathsf{L}$ outside the lattice locus where
$\widehat{\fg}_k$ embeds in a Borcherds--Kac--Moody algebra.
This is an Open-quadrant chart-level scalar tuple at Beilinson
level~$5$. It uses five distinct measurements:
$\kappa_{\mathrm{cat}}$ the trace-form scalar,
$\kappa^{\mathrm{Hodge}}_{\mathrm{ch}}$ the chiral Hodge supertrace,
$\kappa^{\mathrm{Heis}}_{\mathrm{ch}}$ the Heisenberg projection
through the Cartan oscillator subalgebra,
$\kappa_{\mathrm{BKM}}, \kappa_{\mathrm{fiber}}$ the BKM and fibre
data; the collapse pattern is $r_{\max}=3$.
The five entries are five distinct invariants per family
and their collapse pattern under averaging is the
classification axis, not any single value.
\end{remark}

\begin{remark}[Endpoint hypotheses for the affine
Kac--Moody surface]
\label{rem:km-endpoint-hypotheses}
\index{endpoint hypothesis!affine Kac--Moody}
Several quantitative results recorded for $\widehat{\fg}_k$ in this
chapter promote a finite-spin / quadratic / classical input to a
chiral statement; each carries an endpoint hypothesis package
(quadratic-candidate discipline):
\begin{enumerate}[label=\textup{(\arabic*)}]
\item \emph{Sugawara central charge} $c(k,\fg) =
k\dim(\fg)/(k+h^\vee)$: holds at every $k\neq -h^\vee$ on the
Sugawara conformal vector, conditional on the conformal-weight
grading existing (no completion needed at generic level).
\item \emph{KZ $r$-matrix} $r^{\mathrm{KZ}}(z)=
\Omega_{\mathrm{KZ}}/((k+h^\vee)z)$: classical-to-quantum
promotion conditional on the KZ analytic Strong Deformation
Retract (Etingof--Kazhdan; CKL analytic SDR).
\item \emph{Quantum group equivalence}
(Theorem~\ref{thm:km-quantum-groups}): finite-spin lift to
$U_q(\fg)$ at $q=e^{\pi i/(k+h^\vee)}$ conditional on the
Kazhdan--Lusztig analytic equivalence at non-integer levels.
\item \emph{Wakimoto BRST resolution}
(Theorem~\ref{thm:wakimoto-brst-full-nondegenerate}): screening-charge
realisation conditional on the generic nonresonant locus, with
explicit Morita obstruction off the locus.
\item \emph{Drinfeld--Sokolov reduction}
(Proposition~\ref{prop:ds-admissible}): $V_k(\fg)\to
\mathcal{W}_k(\fg)$ conditional on Arakawa exactness at admissible
level (Yamada / Penkava--Vanhaecke quadratic-to-PBW endpoint).
\end{enumerate}
None of (1)--(5) is asserted as an unconditional finite-spin /
classical theorem.
\end{remark}

Affine Kac--Moody algebras occupy the Lie/tree archetype
($r_{\max} = 3$), between the Gaussian families
(Heisenberg, lattice, free fermion; $r_{\max} = 2$) and
the contact/quartic archetype of $\beta\gamma$
($r_{\max} = 4$). At generic level $k \neq -h^\vee$, the
scalar curvature
$\kappa = \dim(\mathfrak{g}) \cdot (k + h^\vee)/(2h^\vee)$
controls the genus tower, and one-channel scalar saturation
(Theorem~\ref{thm:theta-direct-derivation}) identifies the
minimal scalar package
$\Theta^{\min} = \kappa \cdot \eta \otimes \Lambda$. The strict model $\Convstr$ carries a nonzero cubic shadow at degree~$3$
from the Lie bracket, but the Jacobi identity forces $o_4 = 0$
and the tower terminates at $r = 3$. Under the conditional
strictification theorem, $\Definfmod(\widehat{\fg}_k)$ has a
strict representative
(Theorem~\ref{thm:km-strictification}): the canonical class is
represented without higher homotopy corrections, making Kac--Moody
the archetype of strictification
(Corollary~\ref{cor:strictification-comparison}(iii)). At critical level
$k = -h^\vee$, the shifted curvature vanishes and the bar complex
becomes uncurved.  The Feigin--Frenkel center is seen on the
completed critical bar-Ext/Hochschild comparison surface, not in the
raw bar-cobar counit.

\begin{remark}[Chart-class enumeration: affine Kac--Moody]
\label{rem:km-chart-class-enumeration}
\index{chart-class!affine Kac--Moody}
This is an Open-quadrant chart-presentation statement at Beilinson
level~$1$. The hypotheses are a factorisation dg-category
$\mathcal{C}^{\mathrm{op}}$ on $(X, D, \tau)$, vacuum $b$,
$A_b = \mathrm{End}_\mathcal{C}(b)$; archetype label invariant
under Morita equivalence on the level-$1\to 0$ fibre.
At generic level $k\neq -h^\vee$ the chart algebra
\[
A_b^{\mathrm{KM}, k} = \mathrm{End}_{\mathcal{C}^{\mathrm{op}}_X}(b^{\mathrm{KM},k})
\;\simeq\; V_k(\fg)\otimes_{\mathcal{D}_X}(\,\cdot\,)
\]
on the standard vacuum $b^{\mathrm{KM},k}$ realises the affine
universal Kac--Moody chiral algebra at level~$k$.
The five-archetype label is~$\mathsf{L}$
(shadow depth $r_{\max}=3$, Lie/tree),
controlled by the scalar curvature
$\kappa^{\mathrm{cat}}_{V_k(\fg)} = \dim(\fg)(k+h^\vee)/(2h^\vee)$
on the trace-form normalisation. Two Morita-distinct chart
presentations exist:
\begin{enumerate}[label=\textup{(L\arabic*)}]
\item \emph{Universal current chart} ($V_k(\fg)$): the standard
chart from current generators $J^a$ with double-pole OPE
coefficient $k$. The KZ chart $\Phi_{\mathrm{KZ}}$ of the
landscape census embeds here as the universal generator-presentation.
\item \emph{Wakimoto free-field chart} ($W_k$): the screening-charge
realisation with $\beta\gamma$ + free-boson generators
(Theorem~\ref{thm:w3-wakimoto-sl3}). Wakimoto and current charts
are Morita-equivalent on the generic nonresonant locus
(Theorem~\ref{thm:wakimoto-brst-full-nondegenerate} below), so the
class-$\mathsf{L}$ label is chart-independent on the generic locus.
The resonant locus carries an explicit Morita obstruction.
\end{enumerate}
At critical level $k=-h^\vee$ the chart algebra $V_{-h^\vee}(\fg)$
falls into class~$\mathsf{G}$ on the scalar lane (because
$\kappa(V_{-h^\vee}(\fg))=0$) but the simple-pole Lie bracket
persists in the OPE: the chart presentation is class-$\mathsf{L}$
at the operator level and class-$\mathsf{G}$ on the scalar projection.
This stratification of the critical level is the FF
(Feigin--Frenkel) chart-class anomaly. Within Vol~III the matching
$\Phi_d^{\mathrm{FA}}$-evaluation at the same vacuum yields the
holographic CY companion (Vol~III, level~$1\leftrightarrow 1$).
\end{remark}

\begin{remark}[Theorem A in four lanes: affine Kac--Moody]
\label{rem:km-theorem-A-four-lanes}
\index{Theorem A!four lanes!affine Kac--Moody}
On the chart algebra $A_b^{\mathrm{KM},k}=V_k(\fg)$, Theorem~A
(Theorem~\ref{thm:A-infinity-2} in
Chapter~\ref{chap:theorem-A-infinity-2}) appears in the four lanes
$\mathrm{L}1,\mathrm{L}2,\mathrm{L}3,\mathrm{L}4$
of the Theorem~A four-lane discipline:
\begin{itemize}
\item \emph{$\mathrm{L}1$ (chain level, $\mathrm{Ch}(\mathrm{Vect})$).}
The bar--cobar adjunction
$\Omegach_X\dashv\Bbarch_X$ is realised by the curved chiral
Chevalley--Eilenberg complex
$\bar{B}^{\mathrm{ch}}(V_k(\fg))\simeq
\mathrm{CE}^{\mathrm{ch},c}(\fg_{-k-2h^\vee})$
of equation~\eqref{eq:kac-moody-koszul-basic}; the witness chain
homotopy is the Sugawara-corrected one
$h_L = h_{\mathrm{LV}}/(2(k+h^\vee))$
for $k\neq -h^\vee$; the generic denominator has a pole at
$k=-h^\vee$, where the critical Sugawara construction itself is
undefined.
\item \emph{$\mathrm{L}2$ (Quillen-equivalence,
$(\Fact(X),\star)$).}
The Vallette bar--cobar Quillen equivalence
\cite[Theorem~2.1]{Val16} applies in the Francis--Gaitsgory
factorisation ambient on the generic Koszul-locus stratum
(Theorem~\ref{thm:universal-kac-moody-koszul}); the
Hinich~\cite{Hinich2003} symmetric-monoidal promotion is
unobstructed on this locus.
\item \emph{$\mathrm{L}3$ ($(\infty,1)$-categorical).}
The $(\infty,1)$-localisation of the L2 Quillen pair gives the
adjoint equivalence between
$\Fact^{\aug}(X)$-objects of class~$\mathsf{L}$ type and
their $\CoFact^{\conil,\comp}$-counterparts; the Lurie HA \S5.5
parallel topological case applies after the chiral-to-topological
specialisation of Vol II.
\item \emph{$\mathrm{L}4$ ($(\infty,2)$-properad).}
Theorem~$A^{\infty,2}$ (Theorem~\ref{thm:A-infinity-2})
inscribes the $(\infty,2)$-categorical adjoint equivalence on the
properadic envelope generated by $V_k(\fg)$; the Hackney--Robertson
\cite[Theorem~5.12]{HackneyRobertson2019} lift identifies the
class-$\mathsf{L}$ properadic shadow with the chiral
associative bar--cobar pair on $V_k(\fg)$.
\end{itemize}
The four lanes coincide on the generic Koszul locus (where
hypothesis (H3) holds for the conformal-weight grading);
they diverge at the critical level and at the admissible-level
resonant locus, where (H3) fails on the raw direct-sum ambient and
the Positselski coderived ambient is mandatory
(completed-ambient ambient discipline).
\end{remark}

The Feigin--Frenkel shift
$k \mapsto -k - 2h^\vee$ is Verdier duality on
configuration spaces
(Theorem~\ref{thm:bar-cobar-verdier}).
In the WZW model, this shift exchanges holomorphic
and anti-holomorphic sectors; geometrically, it reverses
worldsheet orientation. The same involution governs modular
$S$-transformations of $\widehat{\mathfrak{g}}_k$-characters, the
$q \mapsto q^{-1}$ symmetry of $U_q(\mathfrak{g})$ at
$q = e^{\pi i/(k+h^\vee)}$
(Theorem~\ref{thm:km-quantum-groups}), and 4d S-duality on circles.
Each is a projection of Verdier duality on the Fulton--MacPherson
compactification.

\begin{remark}[Shadow obstruction tower for Kac--Moody]\label{rem:km-master-mc}
For $\widehat{\mathfrak{g}}_k$ at generic level, the shadow
obstruction tower in the modular convolution algebra
$\mathfrak{g}^{\mathrm{mod}}_\cA$
(Definition~\ref{def:modular-convolution-dg-lie}) has scalar
minimal package $\Theta^{\min} = \kappa \cdot \eta \otimes \Lambda$
with
$\kappa = \dim(\mathfrak{g})\cdot(k+h^\vee)/(2h^\vee)$, while
$\Theta^{\leq 3}$ already carries the cubic Lie/tree shadow.
The bar complex, Koszul duals, genus tower, and complementarity
splitting computed in this chapter are projections of both the
scalar minimal package and the nontrivial cubic shadow
\textup{(}Remark~\textup{\ref{rem:unifying-principle}}\textup{)}.
At critical level $k = -h^\vee$, $\kappa = 0$ and the shifted bar
curvature vanishes. The Feigin--Frenkel center and opers live in the
completed critical center/Hochschild lane.
\end{remark}

\begin{remark}[Three-pillar interpretation: Kac--Moody]
\label{rem:km-three-pillar}
\index{Kac--Moody!three-pillar interpretation}
In the three-pillar architecture
(\S\ref{sec:concordance-three-pillars}):
(i)~the \v{C}ech complex of affine Kac--Moody is a \emph{non-strict}
homotopy chiral algebra with $F_3 \neq 0$
(Example~\ref{ex:cech-hca-sl2}); the nontrivial $F_3$ is the Jacobiator
homotopy witnessing the failure of the Jacobi identity at the chain level.
This is the Lie/tree secondary Borcherds archetype:
$j'_{(p,q,r)} \neq 0$ at degree~$3$ but all higher degrees vanish;
(ii)~the convolution $sL_\infty$-algebra
$\operatorname{hom}_\alpha(\cC^{\textup{!\textasciigrave}}, \widehat{\mathfrak{g}}_k)$
has a conditional strict representative
(Theorem~\ref{thm:km-strictification}): on the stated strictification
surface the transferred higher brackets vanish
($\ell_k^{\mathrm{tr}} = 0$ for $k \geq 3$), and the strict dg~Lie
model represents the $L_\infty$-structure at every genus
\textup{(}; all currents $J^a$ have conformal weight~$1$, so
the scalar formula $\mathrm{obs}_g = \kappa\cdot\lambda_g$ is the
Theorem~D input on this surface\textup{)};
(iii)~the Sugawara construction at non-critical level gives a Virasoro
element in the PVA, so Khan--Zeng~\cite{KZ25} topological resonance
reduction applies and the shadow obstruction tower simplifies.
\end{remark}

\begin{remark}[Genus-two shell profile: Kac--Moody]
\label{rem:km-genus-two-shells}
\index{Kac--Moody!genus-two shells}
\index{genus-two shells!Lie/tree archetype}
In the genus-two shell decomposition
(Construction~\ref{const:vol1-genus-two-shells}),
the Lie/tree archetype sees two of the three shells:
$\Theta^{(2)}_{\mathrm{loop} \circ \mathrm{loop}}$ (iterated
non-separating loop) and
$\Theta^{(2)}_{\mathrm{sep} \circ \mathrm{loop}}$ (mixed
separating/non-separating clutching). The planted-forest
shell $\Theta^{(2)}_{\mathrm{pf}}$ is absent. The conditional
strict representative of Theorem~\ref{thm:km-strictification} kills
the transferred higher brackets at tree level: $\ell_4^{(0)} = 0$
implies no quartic contact contribution, and the planted-forest
correction terms from codimension-${\geq}2$ log-FM boundary strata
vanish identically. The surviving genus-two terms arise from
iterated application of the BV operator~$\Delta$ and the
binary sewing differential, both of which are nonzero because
the structure constants $f_{abc}$ produce nonvanishing Casimir
traces.
\end{remark}

\begin{computation}[Genus-$1$ Hessian for $\widehat{\mathfrak{sl}}_2$;
\ClaimStatusProvedHere]
\label{comp:km-genus1-hessian}
\index{Kac--Moody!genus-1 Hessian}
\index{spectral discriminant!Kac--Moody}
At level~$k$ with $\kappa=3(k+2)/4$, the genus-$1$ deformation
space is one-dimensional (generated by the Sugawara variation
$\partial_k T_{\mathrm{Sug}}$) and the genus-$1$ Hessian is
the scalar
\[
H^{(1)}_{\widehat{\mathfrak{sl}}_2{}_k}
\;=\;
\kappa
\;=\;
\frac{3(k+2)}{4}\,.
\]
The spectral discriminant is
$\Delta_{\widehat{\mathfrak{sl}}_2}(x)
= (1-kx)(1-(k+4)x)/(1-2x)$
(Definition~\ref{def:spectral-branch-object},
Computation~\ref{comp:spectral-discriminants-standard}).
The two zeros $x=1/k$ and $x=1/(k+4)$ are the branch points
at levels $k$ and $-k-4$ (Feigin--Frenkel dual); the pole
at $x=1/2$ is the critical level $k=-2$.
\end{computation}

\begin{remark}[Four-level structure: Kac--Moody]
\label{rem:km-four-level}
The affine algebra exhibits all four levels of the modular engine:
(i)~the WZW model / affine current algebra (geometric/physical);
(ii)~the CE complex and bar resolution (strict model $\Convstr$);
(iii)~a conditional strict representative by
Theorem~\ref{thm:km-strictification}: the archetype of
strictification on the stated surface, with no transferred higher
homotopy corrections
(Corollary~\ref{cor:strictification-comparison}(iii));
(iv)~genus tower with $\kappa = \dim(\mathfrak{g})\cdot(k + h^\vee)/(2h^\vee)$, Lie/tree shadow
archetype, terminates at $r = 3$.
\end{remark}

\section{The Kac--Moody bar complex}%
\label{sec:km-bar-complex}

Let $\mathfrak{g}$ be a finite-dimensional simple Lie algebra with
normalized invariant form $(\cdot,\cdot)$, long roots of square~$2$,
and dual Coxeter number~$h^\vee$. The Killing trace form is
$K_{\fg}=2h^\vee(\cdot,\cdot)$ in this convention.
The affine Kac--Moody chiral algebra $\widehat{\mathfrak{g}}_k$
at level $k \neq -h^\vee$ is generated by currents
$J^a(z)$, $a = 1, \ldots, \dim \mathfrak{g}$, with OPE
\[
 J^a(z)\,J^b(w) \;=\; \frac{k\,\kappa_{\fg,ab}}{(z-w)^2}
 \;+\; \frac{f^{ab}{}_c \, J^c(w)}{z-w}
 \;+\; \text{regular},
\]
where $\kappa_{\fg,ab}=(T^a,T^b)$ and $f^{ab}{}_c$ are the structure
constants. The double pole encodes the level; the simple pole encodes
the Lie bracket. Both contribute to the bar differential.

The bar complex $\barB^{\mathrm{ch}}(\widehat{\mathfrak{g}}_k)$ is
the curved chiral Chevalley--Eilenberg coalgebra at the shifted level:
\index{Kac--Moody algebra!as modular Koszul family}
\begin{align}\label{eq:kac-moody-koszul-basic}
\bar{B}^{\mathrm{ch}}(\widehat{\mathfrak{g}}_k) &\simeq \mathrm{CE}^{\mathrm{ch},c}(\mathfrak{g}_{-k-2h^\vee}), \nonumber\\
(\widehat{\mathfrak{g}}_k)^! &\simeq
\mathrm{CE}^{\mathrm{ch}}(\widehat{\mathfrak{g}}_{-k-2h^\vee})
\end{align}
Explicitly, the bar construction gives the curved chiral
Chevalley--Eilenberg coalgebra
\[
\mathrm{CE}^{\mathrm{ch},c}(\mathfrak{g}_{-k-2h^\vee})
= \bigl(\mathrm{Sym}^c(\mathfrak{g}^*[-1]),\, d_{\mathrm{CE}},\, m_0\bigr)
\]
with curvature~$m_0$ encoding the shifted level; applying cobar
to the bar complex gives the inversion
$\Omega(\bar{B}^{\mathrm{ch}}(\widehat{\mathfrak{g}}_k))
\simeq \widehat{\mathfrak{g}}_k$, while the Koszul dual is the
chiral CE algebra of the level-reflected current presentation
(Theorem~\ref{thm:universal-kac-moody-koszul}).
This is \emph{curved} Koszul duality when $k \neq -h^\vee$.
At the critical level $k = -h^\vee$, the curvature vanishes.

\begin{remark}[Physical context]
The physical and mathematical manifestations of the level shift
$k \mapsto -k - 2h^\vee$ as Verdier duality (WZW orientation
reversal, modular $S$-transformations, quantum group inversion,
and 4d S-duality) are discussed in the chapter introduction above.
\end{remark}

\subsection{The critical level as pivot point}

The critical level $k = -h^\vee$ is the fixed point of $k \mapsto -k-2h^\vee$:

\begin{theorem}[Feigin--Frenkel center at critical level \cite{Feigin-Frenkel}; \ClaimStatusProvedElsewhere]\label{thm:critical-level-structure}
\index{critical level|textbf}
\index{Langlands dual}
At $k = -h^\vee$, the center of the affine Kac--Moody algebra
$\widehat{\mathfrak{g}}_{-h^\vee}$ is, on the local formal-disc surface
relevant here,
\begin{equation}
Z(\widehat{\mathfrak{g}}_{-h^\vee}) \cong
\mathrm{Fun}(\mathrm{Op}_{\mathfrak{g}^\vee}(D)),
\end{equation}
the algebra of functions on $\mathfrak{g}^\vee$-opers on the formal
disc.
\end{theorem}

\begin{remark}[Related external critical-level structures]
Other external critical-level structures, such as affine-flag
cohomological models and Wakimoto/free-field realizations, are
important background but are not part of the specific
Feigin--Frenkel-center theorem surface recorded in
Theorem~\ref{thm:critical-level-structure}.
\end{remark}

\subsection{Local computational surface}

\begin{remark}[Generators, residues, relations, genus correction]
The explicit Kac--Moody calculations use four data:
\begin{enumerate}[label=(\roman*)]
\item current generators and conformal weights (\S\ref{sec:km-bar-complex});
\item OPE multi-residues on configuration spaces
 (Theorem~\ref{thm:universal-kac-moody-koszul}).
\item quadratic relations from OPE associativity
 (\S\ref{sec:km-koszul-abstract}).
\item higher-genus correction classes
 (Theorem~\ref{thm:km-higher-genus-corrections}).
\end{enumerate}
They are evaluated for $\widehat{\mathfrak{sl}}_2$,
$\widehat{\mathfrak{sl}}_3$, and general $\widehat{\mathfrak{g}}$.
\end{remark}

\section{Affine Kac--Moody algebras: precise setup}

\subsection{Loop algebras and central extensions}

\begin{definition}[Loop algebra]\label{def:loop-algebra}
Let $\mathfrak{g}$ be a simple finite-dimensional Lie algebra with
normalized invariant form $\kappa_{\mathfrak{g}}$, normalized so that
$(\theta|\theta) = 2$ where $\theta$ is the highest root. The actual
Killing trace form is \(2h^\vee\kappa_{\mathfrak g}\). The
\emph{loop algebra} is:
\begin{equation}
L\mathfrak{g} := \mathfrak{g} \otimes \mathbb{C}[t,t^{-1}]
\end{equation}
with bracket:
\begin{equation}
[x \otimes t^m, y \otimes t^n] = [x,y] \otimes t^{m+n}, \quad x,y \in \mathfrak{g}, \; m,n \in \mathbb{Z}
\end{equation}
\end{definition}

\begin{definition}[Affine Kac--Moody Lie algebra]\label{def:affine-kac-moody}
\index{Kac--Moody algebra!affine|textbf}
\index{dual Coxeter number}
The (untwisted) affine Kac--Moody algebra $\widehat{\mathfrak{g}}$ is the central extension:
\begin{equation}
\widehat{\mathfrak{g}} = L\mathfrak{g} \oplus \mathbb{C} K
\end{equation}
with bracket:
\begin{equation}
[x \otimes t^m, y \otimes t^n] = [x,y] \otimes t^{m+n} + m \delta_{m+n,0} \cdot \kappa_{\mathfrak{g}}(x,y) \cdot K
\end{equation}
and $[K, \widehat{\mathfrak{g}}] = 0$ (central element).
\end{definition}

\begin{remark}[Cocycle interpretation]
The central extension is classified by $H^2(L\mathfrak{g}, \mathbb{C})$. The cocycle is:
\begin{equation}
\nu(x \otimes f, y \otimes g) = \kappa_{\mathfrak{g}}(x,y) \cdot \mathrm{Res}_{t=0}(g \, df)
\end{equation}
This residue pairing is the algebraic shadow of the geometric residue pairing on configuration spaces developed below.
\end{remark}

\subsection{Vertex algebra and chiral algebra presentations}

\begin{definition}[Vertex algebra perspective]
The \emph{universal affine vertex algebra} $V_k(\mathfrak{g})$ at level $k$ is generated by fields:
\begin{equation}
J^a(z) = \sum_{n \in \mathbb{Z}} J^a_n z^{-n-1}, \quad a = 1,\ldots,\dim(\mathfrak{g})
\end{equation}
satisfying the OPE:
\begin{equation}\label{eq:affine-kac-moody-ope}
J^a(z) J^b(w) \sim
\frac{k \kappa_{\fg,ab}}{(z-w)^2}
+ \frac{{f^{ab}}_c J^c(w)}{z-w}
\end{equation}
where $\kappa_{\fg,ab} = \kappa_{\mathfrak{g}}(T^a, T^b)$ is the
normalized invariant bilinear form on $\mathfrak{g}$ (normalized so
that long roots have squared length $2$), ${f^{ab}}_c$ are the
structure constants, and $\sim$ means ``has singular part.'' Its
inverse matrix is denoted $\kappa_{\fg}^{ab}$. In an orthonormal
basis with $\kappa_{\fg,ab} = \delta_{ab}$, this reduces to the
familiar form $k\delta_{ab}/(z-w)^2$.
\end{definition}

\begin{definition}[Chiral algebra perspective]\label{def:chiral-kac-moody}
Following Beilinson--Drinfeld, the affine Kac--Moody chiral algebra $\widehat{\mathfrak{g}}_k$ at level $k$ on a smooth curve $X$ is the $\mathcal{D}_X$-module:
\begin{equation}
\widehat{\mathfrak{g}}_k = \mathcal{U}_k(\mathfrak{g}) := \left(\mathfrak{g} \otimes_{\mathbb{C}} \mathcal{D}_X\right) / \langle [x \otimes P, y \otimes Q] - [x,y] \otimes PQ - k \cdot \kappa_{\mathfrak{g}}(x,y) \cdot P(Q) \rangle
\end{equation}
where $P, Q \in \mathcal{D}_X$ are differential operators and $P(Q)$ denotes the action of $P$ on $Q$ as a function.
\end{definition}

\begin{theorem}[Equivalence of perspectives \cite{FBZ04, BD04}; \ClaimStatusProvedElsewhere]\label{thm:vertex-chiral-equivalence}
For $X = \mathbb{A}^1$ with coordinate $z$, the vertex algebra $V_k(\mathfrak{g})$ and the chiral algebra $\widehat{\mathfrak{g}}_k$ encode the same mathematical structure. The dictionary is:
\begin{align}
J^a_n &\longleftrightarrow x^a \otimes t^n \in L\mathfrak{g} \quad \text{(loop algebra element)} \\
T(z) = \sum_n L_n z^{-n-2} &\longleftrightarrow \text{Sugawara stress tensor}
\end{align}
where the Sugawara construction gives:
\begin{equation}
T = \frac{1}{2(k+h^\vee)} \sum_a {:}J^a J^a{:}
\end{equation}
\end{theorem}

\subsection{The level and its meaning}

\begin{definition}[Level and Sugawara central charge]
At non-critical level $k\neq -h^\vee$, the \emph{level} $k$
determines the central charge of the Virasoro algebra via the
Sugawara construction:\index{Sugawara construction}
\begin{equation}
c(k, \mathfrak{g}) = \frac{k \cdot \dim(\mathfrak{g})}{k + h^\vee}
\end{equation}
where $h^\vee$ is the dual Coxeter number. At $k=-h^\vee$ the
Sugawara tensor is undefined, while the affine current OPE still has
the trace-form double pole. For $\mathfrak{sl}_n$ one has
$h^\vee = n$ (so $h^\vee(\mathfrak{sl}_2) = 2$,
$h^\vee(\mathfrak{sl}_3) = 3$); in general
$h^\vee = 1 + \sum_i a_i^\vee$ where
$\theta = \sum_i a_i \alpha_i$ is the highest root and $a_i^\vee$
are the dual Kac labels, or equivalently
$h^\vee = (\rho|\theta^\vee) + 1$ with
$\theta^\vee = 2\theta/(\theta|\theta)$.
\end{definition}

\begin{computation}[Sugawara central charge from $TT$-OPE]
\label{comp:sugawara-central-charge-from-tt}
\index{Sugawara construction!derivation from $TT$-OPE}
The Virasoro algebra inside $V_k(\fg)$ is generated by the Sugawara
stress tensor $T = \frac{1}{2(k+h^\vee)}\sum_a :J^aJ^a:$ with
currents normalised to $J^aJ^b \sim k\delta^{ab}/(z-w)^2
+ f^{ab}{}_c J^c(w)/(z-w)$. The $TT$-OPE is computed by two
successive Wick contractions:
\begin{equation}\label{eq:sugawara-tt-ope}
T(z)T(w) \;\sim\;
\underbrace{\frac{c/2}{(z-w)^4}}_{\text{double contraction}}
\;+\;\frac{2T(w)}{(z-w)^2}
\;+\;\frac{\partial T(w)}{z-w}.
\end{equation}

\emph{Double-contraction channel.} Both $J^a$ factors of $T(z)$
contract with both $J^b$ factors of $T(w)$. Each contraction
produces $k\delta^{ab}/(z-w)^2$. The two contractions pair as
$(1{\leftrightarrow}1',\,2{\leftrightarrow}2')$ and
$(1{\leftrightarrow}2',\,2{\leftrightarrow}1')$, giving a symmetry
factor of~$2$:
\[
\frac{1}{4(k+h^\vee)^2}\cdot 2\cdot (k\delta^{ab})^2\cdot
\frac{1}{(z-w)^4}
\cdot\dim\fg
\;=\;
\frac{k^2 \dim\fg}{2(k+h^\vee)^2}\cdot\frac{1}{(z-w)^4}.
\]

\emph{Single-contraction channel.} One $J^a(z)$ contracts with
one $J^b(w)$, producing the simple pole $f^{ab}{}_c J^c$ term.
This contributes to the $2T/(z-w)^2$ and $\partial T/(z-w)$ terms
via the adjoint-Casimir identity
$\sum_{a,b} f^{ab}{}_c f^{ab}{}_d = 2h^\vee\delta_{cd}$. Summing
gives an additional $h^\vee\dim\fg/(k+h^\vee)^2$ contribution to
the quartic pole via the normal-ordering relation
$:J^aJ^a:(z){:}J^bJ^b{:}(w)$ double-contraction of one pair with the
simple pole.

\emph{Combining channels.} The total coefficient of
$1/(z-w)^4$ is $c/2$ with
\begin{equation}\label{eq:sugawara-c-derivation}
\frac{c}{2}
\;=\;
\frac{k^2\dim\fg + 2kh^\vee\dim\fg\cdot 0 + h^\vee\cdot(\text{adj-Casimir correction})}{2(k+h^\vee)^2}.
\end{equation}
Executing the full double-contraction trace (Kac--Raina
\cite[Proposition~10.1]{KacRaina} or Frenkel--Ben-Zvi
\cite[Theorem~3.3.4]{FBZ04}) yields the closed form
\begin{equation}\label{eq:sugawara-c-closed}
\boxed{\;c(k, \fg) \;=\; \frac{k\,\dim\fg}{k + h^\vee}.\;}
\end{equation}

\emph{Consistency.} At $\mathfrak{sl}_2$, $k=1$: $c = 3/3 = 1$,
matching the WZW level-$1$ $\mathfrak{sl}_2$ identification with
the free-boson rank-$1$ lattice $\Lambda = A_1$. At $k = -h^\vee$:
$c\to\infty$, confirming the Sugawara degeneration at critical
level (the bar curvature element
$m_0=(k+h^\vee)\mathsf C_{\fg}/(2h^\vee)$ vanishes, while its
scalar trace is
$\kappa=(k+h^\vee)\dim\fg/(2h^\vee)$; the Sugawara central charge
diverges, and the limits commute only after separate
regularisation). The modular characteristic
$\kappa = (k+h^\vee)\dim\fg/(2h^\vee)$
\textup{(}Computation~\textup{\ref{comp:km-kappa-first-principles}}\textup{)}
and the Sugawara central charge $c = k\dim\fg/(k+h^\vee)$ are
\emph{distinct} invariants related by the anomaly-ratio bridge
$\varrho = \kappa/c \neq \mathrm{const}$ for Kac--Moody, in
contrast to the level-independence $\varrho = \mathrm{const}$ for
principal $\mathcal W_N$ (Corollary~\textup{\ref{cor:anomaly-ratio-ds}}).
\end{computation}

\begin{remark}[Critical level]
At $k = -h^\vee$ (the \emph{critical level}) three independent
phenomena meet at the same shifted parameter.  The Sugawara tensor
and the KZ coefficient are undefined because their denominator
$k+h^\vee$ vanishes.  The shifted bar curvature
$m_0=(k+h^\vee)\mathsf C_{\fg}/(2h^\vee)$ vanishes, so the bar
coalgebra is uncurved.  The Feigin--Frenkel centre enlarges and
identifies with functions on opers.  None of these statements says
that the critical affine vertex algebra is commutative: the OPE still
has the Lie-bracket simple pole and the trace-form double pole
$-h^\vee\Omega_{\mathrm{tr}}/(z-w)^2$.  The fixed point of
$k\mapsto -k-2h^\vee$ is a fixed \emph{level}, not an identification
of $A$, $\barB(A)$, $A^i$, $A^!$, and
$Z^{\mathrm{der}}_{\mathrm{ch}}(A)$.
\end{remark}

\begin{proposition}[Critical-level separation of affine invariants;
\ClaimStatusProvedHere]\label{prop:km-critical-separation}
\index{critical level!separation of invariants}
\index{Sugawara construction!undefined versus zero curvature}
\index{Koszul dual!object firewall!Kac--Moody}
Let $\fg$ be simple and let the invariant form be normalized by
$(\theta|\theta)=2$. At the critical level $k=-h^\vee$:
\begin{enumerate}[label=\textup{(\roman*)}]
\item The trace-form double pole in the current OPE is not zero:
\[
J^a(z)J^b(w)\sim
\frac{-h^\vee(J^a,J^b)}{(z-w)^2}
+\frac{f^{ab}{}_{c}J^c(w)}{z-w}.
\]
\item The shifted bar curvature is zero:
\[
m_0(k)=\frac{k+h^\vee}{2h^\vee}\,\mathsf C_{\fg},
\qquad
\kappa^{\mathrm{KM}}(V_k(\fg))
=\operatorname{tr}_{\fg}(m_0(k))
=\frac{(k+h^\vee)\dim\fg}{2h^\vee},
\]
so $m_0(-h^\vee)=0$ and
$\kappa^{\mathrm{KM}}(V_{-h^\vee}(\fg))=0$.
\item The Sugawara tensor and the KZ coefficient are undefined,
not zero:
\[
T_{\mathrm{Sug}}=
\frac{1}{2(k+h^\vee)}\sum_a {:}J^aJ_a{:},
\qquad
r^{\mathrm{KZ}}(z)=
\frac{\Omega_{\mathrm{KZ}}}{(k+h^\vee)z}.
\]
The trace-form collision residue remains
$r^{\mathrm{coll}}(z)=-h^\vee\Omega_{\mathrm{tr}}/z$.
\item The Feigin--Frenkel involution fixes the level:
$-(-h^\vee)-2h^\vee=-h^\vee$.  The fixed level does not identify
the four Koszul objects or the bulk:
\[
\begin{array}{rcl}
A&=&V_{-h^\vee}(\fg),\\[2pt]
\barB(A)&=&\barB\bigl(V_{-h^\vee}(\fg)\bigr),\\[2pt]
A^i&=&H^\bullet\!\bigl(\barB(A)\bigr),\\[2pt]
A^!&=&\mathrm{CE}^{\mathrm{ch}}\bigl(\widehat{\fg}_{-h^\vee}\bigr),\\[2pt]
Z^{\mathrm{der}}_{\mathrm{ch}}(A)
&=&C^\bullet_{\mathrm{ch}}(A,A).
\end{array}
\]
Bar-cobar inversion recovers $A$ from $\barB(A)$; the Koszul dual
is $A^!$, obtained from the bar cohomology coalgebra by Verdier
duality and completed linear duality.  The derived centre is the
Hochschild closed-sector cochain object; at critical level its degree-zero
cohomology is the Feigin--Frenkel centre, not the Koszul dual
algebra.  Equality of the scalar characteristic
$\kappa^{\mathrm{KM}}=0$ is therefore not self-duality of chiral
algebras.
\end{enumerate}
\end{proposition}

\begin{proof}
Clause~(i) is the affine OPE with $k=-h^\vee$.  Clause~(ii) is
Proposition~\ref{prop:km-bar-curvature} specialized to this level;
the trace gives the modular characteristic formula recorded in
Chapter~\ref{ch:landscape-census}.  Clause~(iii) follows from the
same denominator $k+h^\vee$ in the Sugawara construction and in the
Knizhnik--Zamolodchikov connection, while the collision residue uses
the unshifted trace-form coefficient~$k$.  Clause~(iv) is the
involutivity calculation for $k\mapsto -k-2h^\vee$ together with
Convention~\ref{conv:bar-coalgebra-identity}: $\barB(A)$, $A^i$,
and $A^!$ are distinct objects, and
Definition~\ref{def:thqg-chiral-derived-center} defines
$Z^{\mathrm{der}}_{\mathrm{ch}}(A)$ as the Hochschild cochain
object.  Theorem~\ref{thm:bar-cobar-verdier} identifies the dual by
Verdier duality rather than by replacing $A^!$ with~$A$.
\end{proof}

\section{Configuration space realization}

\subsection{Currents as differential forms}

\begin{construction}[Current fields on configuration space]
A current $J^a \in \mathfrak{g}$ at level $k$ is realized as a section:
\begin{equation}
J^a \in \Gamma(\overline{C}_2(X), \mathfrak{g} \boxtimes \omega_X \otimes \Omega^1_{\log})
\end{equation}
Explicitly, in coordinates $(z_1, z_2)$ on $\overline{C}_2(X)$:
\begin{equation}
J^a = x^a(z_1) \otimes d\log(z_1-z_2)
\end{equation}
where $x^a$ is a basis element of $\mathfrak{g}$.
\end{construction}

\begin{remark}[Choice of bundle]
Currents $J^a$ are spin-1 fields, hence sections of $\omega_X$. The level $k$ enters through the double-pole OPE coefficient, not the bundle twist; the logarithmic form $d\log(z_1-z_2)$ captures the $1/(z-w)$ singularity.
\end{remark}

\subsection{OPEs via multi-residue calculus}

\begin{theorem}[Geometric OPE formula; \ClaimStatusProvedHere]\label{thm:geometric-ope-kac-moody}
The binary collision of two affine currents on
\(\overline{C}_2(X)\) has two distributional carriers:
\begin{equation}\label{eq:geometric-ope-distributional}
J^a(z)J^b(w)
\;\sim\;
\frac{k\,\kappa_{\fg,ab}}{(z-w)^2}
+\frac{f^{ab}{}_{c}J^c(w)}{z-w},
\qquad
\{J^a{}_\lambda J^b\}
= f^{ab}{}_{c}J^c+k\,\kappa_{\fg,ab}\lambda .
\end{equation}
Equivalently, in the delta-function kernel convention used for
Poisson vertex algebras,
\begin{equation}\label{eq:geometric-ope-delta-carriers}
\{J^a(z),J^b(w)\}
= f^{ab}{}_{c}J^c(w)\,\delta(z-w)
+k\,\kappa_{\fg,ab}\,\partial_w\delta(z-w).
\end{equation}
Thus the simple pole gives the zeroth product
\(J^a_{(0)}J^b=f^{ab}{}_{c}J^c\), while the double pole gives the
first product \(J^a_{(1)}J^b=k\,\kappa_{\fg,ab}\); the two carriers
are not interchangeable.
\end{theorem}

\begin{proof}
Write \(u=z-w\). The OPE-to-\(\lambda\)-bracket dictionary is
\[
a(z)b(w)\sim \sum_{n\ge0}\frac{(a_{(n)}b)(w)}{u^{n+1}},
\qquad
\{a_\lambda b\}=\sum_{n\ge0}\frac{\lambda^n}{n!}\,a_{(n)}b .
\]
For affine currents,
\[
J^a_{(0)}J^b=f^{ab}{}_{c}J^c,\qquad
J^a_{(1)}J^b=k\,\kappa_{\fg,ab},\qquad
J^a_{(n)}J^b=0\quad(n\ge2),
\]
which gives~\eqref{eq:geometric-ope-distributional}. Passing from the
\(\lambda\)-bracket to the distributional kernel sends
\(\lambda\) to \(\partial_w\delta(z-w)\), giving
\eqref{eq:geometric-ope-delta-carriers}. The logarithmic residue
\(d\log(z-w)=du/u\) extracts the \(u^{-1}\) term for the
current-valued collision; the \(u^{-2}\) term is the central
first-product carrier and feeds the curvature/level package.
\end{proof}

\section{Koszul duality: abstract theory}
\label{sec:km-koszul-abstract}

The level shift $k \mapsto -k - 2h^\vee$ is Verdier duality on configuration spaces: the Serre dual of the level-$k$ invariant form on $\overline{C}_{n+1}(X)$ is the level-$(-k-2h^\vee)$ form, the shift $2h^\vee$ arising from the adjoint Casimir eigenvalue.

\subsection{The general pattern}

\begin{theorem}[Level-shifting duality, abstract form; \ClaimStatusProvedHere]\label{thm:level-shifting-abstract}
\index{Feigin--Frenkel duality|textbf}
\index{level!shifted}
\index{Kac--Moody algebra!Koszul dual}
For the Heisenberg algebra, Koszul duality sends the Lie-type
$\mathcal{H}_k$ to the Com-type
$\mathrm{Sym}^{\mathrm{ch}}(V^*)$
(Theorem~\ref{thm:frame-heisenberg-koszul-dual}): the dual algebra
changes type.
For non-abelian Kac--Moody algebras on the non-critical completed
universal surface, the dual bar construction shifts the level.  The
Koszul dual is
\[
(\widehat{\mathfrak{g}}_k)^!
= \mathrm{CE}^{\mathrm{ch}}(\widehat{\mathfrak{g}}_{-k-2h^\vee}),
\]
the chiral Chevalley--Eilenberg algebra built from the level-reflected
current presentation.  This CE object has the same scalar
$\kappa$-value as the affine current algebra
$\widehat{\mathfrak{g}}_{-k-2h^\vee}$ and is a different chiral
algebra.  Bar-cobar inversion recovers $\widehat{\mathfrak g}_k$
from its own bar coalgebra.  The Feigin--Frenkel center belongs to the
critical degree-zero center surface.

For any simple Lie algebra $\mathfrak{g}$ and level $k \neq -h^\vee$, there exists a quasi-isomorphism of complexes:
\begin{equation}
\Omega\bigl(\bar{B}^{\mathrm{geom}}(\widehat{\mathfrak{g}}_k)^{\vee_{\mathrm{V}}}\bigr)
\simeq
\mathrm{CE}^{\mathrm{ch}}(\widehat{\mathfrak{g}}_{-k-2h^\vee})
\end{equation}
where $\vee_{\mathrm{V}}$ denotes the completed Verdier-linear dual
on the configuration-space bar coalgebra.  This is the chiral analogue
of the symmetric/exterior Koszul exchange, with curvature governed by
the shifted level $k+h^\vee$.
\end{theorem}

\begin{proof}
This is proved as Theorem~\ref{thm:universal-kac-moody-koszul} below:
the bar complex of $\widehat{\mathfrak{g}}_k$ is computed via residues
on FM compactifications (Step~1), the curvature
$m_0 = (k+h^\vee)\mathsf C_{\fg}/(2h^\vee)$ is identified from the
double-pole OPE (Step~2), the cobar OPE recovers the level
$-k - 2h^\vee$ current presentation underlying the chiral CE algebra
(Step~3), and the result is characterized by functorial uniqueness
(Step~4).
\end{proof}

\begin{remark}[Central charge sum; \ClaimStatusProvedHere]\label{rem:km-central-charge-sum}
The central charges satisfy $c(\widehat{\mathfrak{g}}_k) + c(\widehat{\mathfrak{g}}_{-k-2h^\vee}) = 2\dim\mathfrak{g}$, independent of $k$ (Theorem~\ref{thm:central-charge-complementarity}).
\end{remark}

\begin{definition}[Curved Koszul complex; \ClaimStatusDefinitional]\label{def:curved-koszul-km}
For $k \neq -h^\vee$, the $A_\infty$ structure has curvature:
\begin{equation}
m_0 =
\frac{k+h^\vee}{2h^\vee}\cdot \mathsf C_{\fg},
\qquad
m_1^2(a) = [m_0, a] = m_2(m_0, a) - m_2(a, m_0),
\end{equation}
where $\mathsf C_{\fg} = \sum_a J^a \otimes J_a$ is the Casimir
element defined by the normalized invariant form.  The scalar
modular characteristic is its trace:
\[
\kappa^{\mathrm{KM}}(V_k(\fg))
=\operatorname{tr}_{\fg}(m_0)
=\frac{(k+h^\vee)\dim\fg}{2h^\vee}.
\]
The curvature vanishes precisely at $k = -h^\vee$ (critical level),
where $m_1$ becomes a genuine differential.
\end{definition}

\subsection{The Wakimoto perspective}

\begin{definition}[Wakimoto module]\label{def:wakimoto}
\index{Wakimoto representation|textbf}
The Wakimoto module $\mathcal{M}_{\mathrm{Wak}}$ at critical level is the free field algebra
\begin{equation}
\mathcal{M}_{\mathrm{Wak}} = \mathrm{Free}[\beta_\alpha, \gamma_\alpha, \phi_i]
\end{equation}
generated by a $\beta$-$\gamma$ system $(\beta_\alpha, \gamma_\alpha)$ of conformal weights $(1, 0)$ for each positive root $\alpha \in \Delta_+$, together with $r = \mathrm{rank}(\mathfrak{g})$ free bosons $\phi_i$ for the Cartan generators. The currents $J^a = f^a(\beta, \gamma, \phi, \partial\phi)$ are explicit differential polynomials determined by the Wakimoto construction.
\end{definition}

\begin{theorem}[Critical Wakimoto realization and bar comparison;
\ClaimStatusConditional]\label{thm:wakimoto-koszul}
At critical level $k = -h^\vee$, the Wakimoto free field realization
gives a completed free-field model for the critical affine algebra and
its Feigin--Frenkel center. The bar comparison used in this chapter is
the following conditional identity:
\begin{equation}
\bar{B}(\widehat{\mathfrak{g}}_{-h^\vee}) \;\simeq\; \bar{B}(\mathcal{M}_{\mathrm{Wak}})^{Q_{\mathrm{DS}}}
\end{equation}
provided the bar functor commutes with the completed
Drinfeld--Sokolov BRST totalization in the Wakimoto category.
Concretely, generators $J^a$ of $\widehat{\mathfrak{g}}_{-h^\vee}$
correspond to composite operators in Wakimoto, while the center maps to
the commutative algebra of local $\fg^\vee$-opers
(Theorem~\ref{thm:critical-level-structure}).
\end{theorem}

\begin{proof}
The external Wakimoto construction
\cite{Wakimoto86,FeiginFrenkel90bosonization} realizes the critical
currents as differential polynomials in the
$\beta\gamma$-fields and Cartan bosons, and
Theorem~\ref{thm:critical-level-structure} identifies the critical
center with functions on local opers. The displayed bar identity is
therefore not a formal consequence of the center theorem alone. Its
extra hypothesis is the exactness of the completed BRST totalization
after applying $\bar B$. Under that hypothesis,
$Q_{\mathrm{DS}}$ is a derivation of the vertex algebra structure and
is compatible with the bar differential. The Wakimoto module is a
tensor product of free fields:
$\mathcal{M}_{\mathrm{Wak}} =
\bigotimes_{\alpha \in \Delta_+}\mathcal{F}_{\beta_\alpha\gamma_\alpha}
\otimes \mathcal{H}^{\otimes r}$. The bar complexes of these free
fields are the Heisenberg and $\beta\gamma/bc$ bar complexes computed
in Chapters~\ref{ch:heisenberg-frame} and~\ref{chap:free-fields}.
K\"unneth then identifies $\bar B(\mathcal M_{\mathrm{Wak}})$ with the
tensor product of those free-field bar complexes, and the completed
BRST totalization gives the displayed comparison.
\end{proof}

\begin{remark}[Wakimoto triangle: three realizations of Koszul duality]\label{rem:wakimoto-triangle}
\index{Wakimoto representation!Koszul triangle}
Under the bar--BRST comparison hypothesis in
Theorem~\ref{thm:wakimoto-koszul}, one obtains a commutative triangle
of Koszul-duality models:
\begin{center}
\begin{tikzcd}[row sep=large, column sep=large]
\widehat{\mathfrak{g}}_{-h^\vee} \arrow[rr, "\text{chiral CE}"] \arrow[dr, "\text{Wakimoto}"'] & &
\mathrm{CE}^{\mathrm{ch}}(\widehat{\mathfrak{g}}_{-h^\vee}) \\
& \beta\gamma^{\otimes |\Delta_+|} \otimes \mathcal{H}^{\otimes r} \arrow[ur, "Q_{\mathrm{DS}}"'] &
\end{tikzcd}
\end{center}
At critical level, the fixed-point / uncurved package for
$\widehat{\mathfrak{g}}_{-h^\vee}$ factors through the Wakimoto free
field realization under the completed bar--BRST comparison hypothesis.
Each edge of the triangle corresponds to a different manifestation of
the same critical-level structure: algebraic (bar-cobar inversion),
analytic (free field realization), and homological (BRST reduction).
The triangle commutes under the conditional comparison of
Theorem~\ref{thm:wakimoto-koszul}.
\end{remark}

\subsection{Wakimoto BRST exactness on the generic nonresonant locus}
\label{subsec:wakimoto-brst-nondegenerate}

The Koszul-dual triangle at $k = -h^\vee$ is a shadow of a stronger
generic statement away from the Kac--Kazhdan resonant hyperplanes. Let
\[
\beta\coloneqq k + h^\vee, \qquad \nu(k)\coloneqq\sqrt{\beta}
\]
parametrise the deformation away from the critical point. Feigin and
Frenkel constructed a one-parameter family of Wakimoto modules
$W_k(\mathfrak{g})$ by twisting the Cartan bosons with the level:
\[
W_k(\mathfrak{g})
\;=\;
\bigotimes_{\alpha\in\Delta_+}\mathcal{F}_{\beta_\alpha\gamma_\alpha}
\;\otimes\;
\Pi^{\nu(k)}_{\mathfrak{h}},
\]
where $\Pi^{\nu}_{\mathfrak{h}}$ is the rank-$r$ Heisenberg vertex
algebra of Cartan bosons $\phi_i$ with OPE
$\phi_i(z)\phi_j(w)\sim \delta_{ij}\nu^{-2}\log(z-w)$.
The currents $J^a(z)$ are the level-$k$ Wakimoto polynomials
\cite{Wakimoto86, FeiginFrenkel90bosonization}, rescaled so that
the residue OPE $J^a(z)J^b(w)$ reproduces the affine Kac--Moody
algebra $\widehat{\mathfrak{g}}_k$ at level $k$.

The screening operators
$S_{\alpha_i}\in\mathrm{End}(W_k(\mathfrak{g}))$ ($i=1,\ldots,r$) are
the contour integrals
\begin{equation}
\label{eq:wakimoto-screenings}
S_{\alpha_i}
\;=\;
\oint \bigl({:}\,\beta_{\alpha_i}\,e^{-\alpha_i(\phi)/\nu(k)}\,{:}\bigr)(z)\,dz,
\qquad i=1,\ldots,r,
\end{equation}
one per simple root. The BRST differential is their sum:
\begin{equation}
\label{eq:wakimoto-brst}
Q_{\mathrm{BRST}}^{(k)}
\;\coloneqq\;
\sum_{i=1}^{r} S_{\alpha_i}\,\otimes\,c_i,
\end{equation}
where $c_i$ is the ghost of charge $+1$ associated to the simple root
$\alpha_i$. The identities $[S_{\alpha_i},S_{\alpha_j}]=0$ (quantum
Serre relations for screenings, Feigin--Frenkel
\cite{FeiginFrenkel90bosonization}) imply
$(Q_{\mathrm{BRST}}^{(k)})^2 = 0$ on the enlarged complex
$W_k(\mathfrak{g})\otimes\Lambda^\bullet[c_1,\ldots,c_r]$ for every
$k\ne -h^\vee$.

The question the next theorem resolves: does the BRST cohomology of
$(W_k(\mathfrak{g}),Q_{\mathrm{BRST}}^{(k)})$ compute
$V_k(\mathfrak{g})$ (the universal Kac--Moody chiral algebra at
level $k$) on the nonresonant locus where the vacuum
Kac--Shapovalov determinants are nonzero? At resonant and admissible
levels the answer requires the simple-quotient analysis of
Remark~\ref{rem:admissible-koszul-status} and
Theorem~\ref{thm:admissible-sl3-non-koszul-qge3}.

\begin{definition}[Non-critical and Wakimoto-nonresonant loci]
\label{def:wakimoto-nondegenerate-locus}
\index{non-degenerate locus!Wakimoto}
The \emph{non-critical locus} of Kac--Moody levels is
\(
\mathcal{L}^{\mathrm{nd}}_{\mathfrak{g}}
\;\coloneqq\;\{k\in\mathbb{C}\;:\;k\ne -h^\vee\},
\)
the complement of the critical level in the complex $k$-line. The
\emph{Wakimoto-nonresonant locus}
\(\mathcal L^{\mathrm{Wak}}_{\fg}\subset\mathcal L^{\mathrm{nd}}_{\fg}\)
is the complement of the Kac--Kazhdan determinant hyperplanes for the
vacuum module. Admissible levels $k+h^\vee=p/q$ are resonant
rational points; they are not covered by the generic theorem below.
\end{definition}

\begin{theorem}[Wakimoto BRST exactness on the generic nonresonant
locus; \ClaimStatusProvedHere]
\label{thm:wakimoto-brst-full-nondegenerate}
\index{Wakimoto BRST resolution!generic nonresonant locus}
\index{Feigin--Frenkel!BRST resolution}
For every $k\in\mathcal{L}^{\mathrm{Wak}}_{\mathfrak{g}}$, the
Wakimoto BRST complex
$(W_k(\mathfrak{g})\otimes\Lambda^\bullet[c_1,\ldots,c_r],
Q_{\mathrm{BRST}}^{(k)})$ is a resolution of $V_k(\mathfrak{g})$
concentrated in ghost degree zero:
\begin{equation}
\label{eq:wakimoto-brst-resolution-nondegenerate}
H^{n}\bigl(W_k(\mathfrak{g})\otimes\Lambda^\bullet[c],
Q_{\mathrm{BRST}}^{(k)}\bigr)
\;=\;
\begin{cases}
V_k(\mathfrak{g}), & n = 0,\\
0, & n\ne 0.
\end{cases}
\end{equation}
Equivalently, the complex
\begin{equation}
\label{eq:wakimoto-brst-complex-nondegenerate}
0 \to W_k(\mathfrak{g}) \xrightarrow{\;Q_{\mathrm{BRST}}^{(k)}\;}
W_k(\mathfrak{g})\otimes\mathbb{C}^r
\xrightarrow{\;Q_{\mathrm{BRST}}^{(k)}\;} \cdots
\xrightarrow{\;Q_{\mathrm{BRST}}^{(k)}\;}
W_k(\mathfrak{g})\otimes\Lambda^r(\mathbb{C}^r)\to 0
\end{equation}
is a right resolution of $V_k(\mathfrak{g})$ as an
$\widehat{\mathfrak{g}}_k$-module, and the higher cohomology groups
$H^n$ ($n>0$) vanish. The kernel of the first screening,
\(
V_k(\mathfrak{g})=\bigcap_{i=1}^{r}\ker S_{\alpha_i}
\;\subset\; W_k(\mathfrak{g}),
\)
realises $V_k(\mathfrak{g})$ as the Heisenberg--$\beta\gamma$
bosonisation of $\widehat{\mathfrak{g}}_k$ chain-level on the
generic nonresonant locus.
\end{theorem}

\begin{proof}
The proof combines three ingredients: (i) Feigin--Frenkel's
asymptotic-level BRST exactness at $|k+h^\vee|\to\infty$, (ii) the
Kac--Kazhdan module filtration on Wakimoto modules, and (iii) a
spectral-sequence argument on the open locus where the
Kac--Shapovalov determinants remain nonzero.

\emph{Step 1: Verma filtration on the Wakimoto module.}
The Wakimoto module $W_k(\mathfrak{g})$ admits a canonical filtration
by affine Verma submodules. As a bigraded vector space it decomposes
\(
W_k(\mathfrak{g})
\;=\;
\bigoplus_{\lambda\in\mathbb{Z}_{\geq 0}^{|\Delta_+|}\times\mathbb{Z}^r}
W_k(\mathfrak{g})_\lambda
\)
with the first grading the monomial degree in the $\beta_\alpha$
(ghost number in the $\beta\gamma$ sector) and the second the
$\mathfrak{h}$-weight. On the associated graded
$\mathrm{gr}\, W_k(\mathfrak{g})$ the screening charges
$S_{\alpha_i}$ act by the classical boundary differential of the
nilpotent Lie algebra cohomology complex
$C^\bullet(\mathfrak{n}_+;\mathrm{Fun}(N_+))$
(Feigin--Frenkel \cite{FeiginFrenkel90bosonization},
Theorem~3; see also \cite[Chapter~15]{FBZ04}). The classical limit
is acyclic in positive degrees by the affine analogue of the
Koszul resolution
$\mathrm{Sym}(\mathfrak{n}_+)\otimes\Lambda(\mathfrak{n}_+^\vee)
\xrightarrow{\sim}\mathbb{C}$.

\emph{Step 2: Kac--Kazhdan module filtration.}
Away from the Kac--Kazhdan resonant hyperplanes, every
$\widehat{\mathfrak{g}}_k$-submodule of a Verma module is generated
by Kac--Kazhdan singular vectors
\cite[Theorem~2]{KacKazhdan79}. The
Kac--Kazhdan determinant
\[
\det_n(k)
\;=\;
\prod_{\lambda}\;\prod_{m=1}^{\infty}
\bigl(k + h^\vee - \frac{p(\lambda)}{q(\lambda,m)}\bigr)^{P(n-m)}
\]
vanishes on the Kac--Kazhdan resonant hyperplanes; admissible levels
$k+h^\vee=p/q$ are a distinguished rational subfamily of this
resonant locus. On the complement of this countable union of
hyperplanes the vacuum Verma module is irreducible, and the Wakimoto
module $W_k(\mathfrak{g})$ is isomorphic to a Verma module of
highest weight~$0$. For every $k$ with
$k\in\mathcal L^{\mathrm{Wak}}_{\fg}$, the map
$V_k(\mathfrak{g})\hookrightarrow W_k(\mathfrak{g})$ realises
$V_k(\mathfrak{g})$ as the irreducible quotient / Verma module,
and the quotient $W_k(\mathfrak{g})/V_k(\mathfrak{g})$ carries an
acyclic filtration by Kac--Kazhdan singular vectors.

\emph{Step 3: Asymptotic-level BRST exactness.}
Feigin and Frenkel proved
\cite[Theorem~4.3.2]{FeiginFrenkel90bosonization} that for every
$k\in\mathbb{C}$ with $|k+h^\vee|$ sufficiently large,
\[
H^n\bigl(W_k(\mathfrak{g})\otimes\Lambda^\bullet[c],
Q_{\mathrm{BRST}}^{(k)}\bigr)
\;=\;\delta_{n,0}\,V_k(\mathfrak{g}).
\]
The argument works because the leading-order screening action
reproduces the classical limit of Step~1, and the subleading
corrections are bounded by $O(\beta^{-1})$.

\emph{Step 4: Spectral-sequence propagation on
$\mathcal{L}^{\mathrm{Wak}}_{\mathfrak{g}}$.}
Filter the BRST complex by conformal weight:
\(
F^{\leq N} W_k(\mathfrak{g})
\;\coloneqq\;
\bigoplus_{h\leq N} W_k(\mathfrak{g})_h.
\)
Each $F^{\leq N}$ is finite-dimensional, the BRST differential
preserves the filtration strictly, and the spectral sequence
\(
E_1^{p,q} = H^{p+q}(\mathrm{gr}^p W_k(\mathfrak{g})\otimes
\Lambda^\bullet, Q_{\mathrm{BRST}}^{(k)})
\;\Longrightarrow\;
H^{p+q}(W_k(\mathfrak{g})\otimes\Lambda^\bullet,
Q_{\mathrm{BRST}}^{(k)})
\)
converges at finite stage for each $N$. The $E_1$ page is computed
by the classical Koszul resolution of Step~1, tensored with a
finite-dimensional $\mathfrak{h}$-character. The differentials
$d_r$ ($r\geq 1$) are polynomials in $\beta = k + h^\vee$ with
coefficients independent of $k$ (the shape of the Wakimoto
recursion is universal in $\beta$). At each fixed conformal weight
$N$, the cohomology $H^n(F^{\leq N})$ is a polynomial in $k$
divided by a power of $\beta$; the vanishing at
$|k+h^\vee|\to\infty$ (Step~3) forces the numerator to vanish
identically on each nonresonant component. Therefore
$H^n(F^{\leq N}) = 0$ for $n>0$ and every
$k\in\mathcal{L}^{\mathrm{Wak}}_{\mathfrak{g}}$.

\emph{Step 5: Passage to the direct limit.}
The cohomology of $W_k(\mathfrak{g})\otimes\Lambda^\bullet$ is the
direct limit of the cohomologies of $F^{\leq N}$. Direct limits
preserve exactness, so
$H^n(W_k(\mathfrak{g})\otimes\Lambda^\bullet, Q_{\mathrm{BRST}}^{(k)}) = 0$
for $n>0$ and every $k\in\mathcal L^{\mathrm{Wak}}_{\fg}$.
The $n = 0$ cohomology is
$\bigcap_i \ker S_{\alpha_i}$, and by Feigin--Frenkel Step~1 this
intersection realises precisely the image of $V_k(\mathfrak{g})$
inside $W_k(\mathfrak{g})$. The identification
$H^0 = V_k(\mathfrak{g})$ follows.
\end{proof}

\begin{remark}[Admissible levels are a separate simple-quotient problem]
\label{rem:wakimoto-admissible-robust}
The generic theorem deliberately stops before the admissible
hyperplanes $k+h^\vee=p/q$. On those hyperplanes the
Kac--Kazhdan determinant vanishes, the universal vacuum module
contains null vectors, and the simple quotient
$L_k(\fg)=V_k(\fg)/I_k$ need not inherit the PBW-Koszul property.
The local source of truth is
Remark~\ref{rem:admissible-koszul-status}: $L_k(\mathfrak{sl}_2)$ is
chirally Koszul at admissible level, while
Theorem~\ref{thm:admissible-sl3-non-koszul-qge3} gives
non-Koszulness for $L_k(\mathfrak{sl}_3)$ on the denominator
$q\geq 3$ non-degenerate admissible surface.
\end{remark}

\begin{corollary}[Rank-one admissible BRST resolution]
\label{cor:wakimoto-brst-admissible}
\ClaimStatusConditional
\index{admissible level!Wakimoto resolution}
For $\widehat{\mathfrak{sl}}_2$ at every admissible level
$k=-2+p/q$ with $p\geq2$, $q\geq1$, and $\gcd(p,q)=1$, the
admissible Wakimoto screening complex resolves the simple quotient
$L_k(\mathfrak{sl}_2)$:
\[
H^n\bigl(W_k(\mathfrak{sl}_2)\otimes\Lambda^\bullet[c],
Q_{\mathrm{BRST}}^{(k)}\bigr)
\;=\;
\delta_{n,0}\;L_k(\mathfrak{sl}_2).
\]
In particular, for $\widehat{\mathfrak{sl}}_2$ at $k=-1/2$:
\[
L_{-1/2}(\mathfrak{sl}_2)
\;=\;
\bigl(\ker S_{\alpha_1}\bigr) \;\subset\;
\mathcal{F}_{\beta\gamma}\otimes\Pi^{\sqrt{3/2}}_{\mathfrak{h}},
\]
and the full complex is exact in positive ghost degrees.
No analogous assertion is made here for higher-rank simple quotients.
For $\mathfrak{sl}_3$, the value $k=-1$ has $k+h^\vee=2=2/1$ and is
not admissible under the local convention $p\geq h^\vee=3$; on the
non-degenerate admissible denominator-$q\geq3$ surface,
Theorem~\ref{thm:admissible-sl3-non-koszul-qge3} proves that
$L_k(\mathfrak{sl}_3)$ is not chirally Koszul.
\end{corollary}

\begin{proof}
The rank-one statement is the admissible $\mathfrak{sl}_2$ case of
Remark~\ref{rem:admissible-koszul-status}, with the explicit
$k=-1/2$ character calculation supplied by
Theorem~\ref{thm:kw-bar-spectral}. The screening relation
$(Q_{\mathrm{BRST}}^{(k)})^2=0$ follows from the rank-one quantum
Serre relation, and the Kac--Wakimoto character formula gives the
positive-degree exactness of the quotient bar spectral sequence.
\end{proof}

\begin{corollary}[Wakimoto presentation preserves the
class $\mathsf{L}$ carrier]
\label{cor:class-L-reduces-to-G-chain-level}
\ClaimStatusConditional
\index{Wakimoto presentation!class L carrier}
On the Wakimoto-nonresonant locus
$\mathcal{L}^{\mathrm{Wak}}_{\mathfrak{g}}$, every affine Kac--Moody
vertex algebra \(V_k(\mathfrak g)\) has the chain-level
Wakimoto free-field presentation
\[
V_k(\mathfrak{g})
\;=\;
H^0_{Q_{\mathrm{BRST}}^{(k)}}
\bigl(\mathcal{F}_{\beta\gamma}^{\otimes|\Delta_+|}
\otimes \Pi^{\nu(k)}_{\mathfrak{h}}\bigr).
\]
This is a resolution of the affine current algebra, not a reduction
of its carrier to class~\(\mathsf G\). The cohomology algebra remains
class~\(\mathsf L\), with
\[
\kappa^{\mathrm{KM}}(V_k(\fg))
=\frac{\dim(\fg)(k+h^\vee)}{2h^\vee},
\qquad
\mathfrak C(x,y,z)=\kappa_{\fg}(x,[y,z]),
\qquad
r_{\max}=3 .
\]
The free fields are the presentation; the BRST differential is part of
the structure that reconstructs the non-abelian simple-pole carrier.
\end{corollary}

\begin{proof}
The statement $V_k(\mathfrak{g}) = H^0$ is
Theorem~\ref{thm:wakimoto-brst-full-nondegenerate}. The source
complex is built from Heisenberg and \(\beta\gamma\) free fields, whose
bare OPEs have their own shadow data.  The affine OPE is the
cohomology OPE after imposing the screening differential
\(Q_{\mathrm{BRST}}^{(k)}\). In that cohomology, the quadratic
coefficient is the level-normalized invariant form
\(k\,\kappa_{\fg}(J^a,J^b)\): in the Wakimoto realization it is the
sum of the \((k+h^\vee)\) Cartan-Heisenberg contraction and the
\((-h^\vee)\) adjoint \(\beta\gamma\) contraction. The simple-pole
carrier is reconstructed by the screening charges and equals
\(J^a_{(0)}J^b=f^{ab}{}_{c}J^c\). Therefore the cyclic cubic tensor is
\(\kappa_{\fg}(x,[y,z])\), and the Jacobi identity kills the quartic
obstruction. By Theorem~\ref{thm:quadrichotomy}, the resulting
cohomology algebra is class~\(\mathsf L\) with \(r_{\max}=3\).
\end{proof}

\begin{remark}[Critical-level degeneration]
\label{rem:wakimoto-brst-critical-limit}
Theorem~\ref{thm:wakimoto-brst-full-nondegenerate} does not extend
to $k = -h^\vee$: at the critical level the rescaling
$\nu(k) = \sqrt{\beta}$ becomes singular, the Heisenberg module
$\Pi^{\nu}_{\mathfrak{h}}$ degenerates to a classical Poisson
module, and the screening charges lose their vertex-operator
structure. Under the completed bar--BRST comparison hypothesis of
Theorem~\ref{thm:wakimoto-koszul}, the critical-level BRST
cohomology contains the Feigin--Frenkel center
$\mathfrak{z}(\widehat{\mathfrak{g}})$, an enlargement of the
completed critical center by classical opers. The
generic nonresonant theorem recovers the universal chiral algebra
$V_k(\mathfrak{g})$; the critical-level comparison records the
enlarged Feigin--Frenkel center in the center/Hochschild ambient. They
are two related BRST
constructions, distinguished by whether
the Wakimoto pairing is quantum (non-critical) or classical
(critical).
\end{remark}

\begin{remark}[Chain-level and $(\infty,1)$ comparison]
\label{rem:wakimoto-brst-infty-lane}
Theorem~\ref{thm:wakimoto-brst-full-nondegenerate} is stated and
proved chain-level: every cohomology group is computed from an
explicit complex of Heisenberg--$\beta\gamma$ modules, with named
BRST generators
\eqref{eq:wakimoto-brst} and named screening contours
\eqref{eq:wakimoto-screenings}. The $(\infty,1)$-categorical
form is the assertion that the BRST resolution is a
quasi-isomorphism of $E_1$-algebra objects in the
$\infty$-category of factorisation $D$-modules on the curve $X$:
\[
W_k(\mathfrak{g})\otimes\Lambda^\bullet[c]
\xrightarrow{\;\simeq\;}
V_k(\mathfrak{g})
\qquad \text{in } \mathrm{Alg}^{E_1}(\mathrm{Fact}_X).
\]
The chain-level statement and the $(\infty,1)$-statement are
independent theorems with independent content: the chain-level
version is what lets one compute $\kappa^{\mathrm{KM}}$, the
Kac--Kazhdan determinant, and the shadow tower; the
$(\infty,1)$-version is what lets one assert universality in
the derived category (Vol II Chapter~\ref{chap:factorization-swiss-cheese}).
\end{remark}

\begin{remark}[Master Gravity comparison and class-$\mathsf{L}$ domain]
\label{rem:wakimoto-brst-gravity-upgrade}
Corollary~\ref{cor:class-L-reduces-to-G-chain-level} supplies the
chain-level Heisenberg--$\beta\gamma$ resolution needed by the Master
Gravity comparison on the Wakimoto-nonresonant locus. It does not
lower the affine carrier from class~\(\mathsf L\) to
class~\(\mathsf G\), and it does not make admissible or critical
cases to unconditional statements:
admissible simple quotients are governed by
Remark~\ref{rem:admissible-koszul-status} and
Theorem~\ref{thm:admissible-sl3-non-koszul-qge3}, while the critical
case is controlled by the Feigin--Frenkel center and the conditional
bar--BRST comparison of Theorem~\ref{thm:wakimoto-koszul}.
\end{remark}

\section{\texorpdfstring{Explicit computation: $\widehat{\mathfrak{sl}}_2$}{Explicit computation: sl-2}}

\subsection{Setup and generators}

For $\mathfrak{g} = \mathfrak{sl}_2$, the data are: $h^\vee = 2$, $\dim(\mathfrak{sl}_2) = 3$, with the standard basis $\{e, f, h\}$ satisfying $[e,f] = h$, $[h,e] = 2e$, $[h,f] = -2f$. The normalized invariant form is $(h,h) = 2$, $(e,f) = 1$ (the Killing trace pairing is $4 \cdot \mathrm{tr}_{\mathrm{fund}}$, but we use the standard normalization $(\theta|\theta) = 2$).

\begin{definition}[\texorpdfstring{$\widehat{\mathfrak{sl}}_2$}{sl-hat_2} at level \texorpdfstring{$k$}{k}]
The affine $\mathfrak{sl}_2$ vertex algebra has generators:
\begin{align}
e(z) &= \sum_n e_n z^{-n-1}, \quad \text{conformal weight } h_e = 1 \\
f(z) &= \sum_n f_n z^{-n-1}, \quad \text{conformal weight } h_f = 1 \\
h(z) &= \sum_n h_n z^{-n-1}, \quad \text{conformal weight } h_h = 1
\end{align}
with OPEs:
\begin{align}
e(z) f(w) &\sim \frac{k}{(z-w)^2} + \frac{h(w)}{z-w} \\
h(z) e(w) &\sim \frac{2e(w)}{z-w} \\
h(z) f(w) &\sim \frac{-2f(w)}{z-w} \\
h(z) h(w) &\sim \frac{2k}{(z-w)^2}
\end{align}
\end{definition}

\subsection{\texorpdfstring{Critical level: $k = -2$}{Critical level: k = -2}}

\begin{theorem}[Critical level simplification for \texorpdfstring{$\mathfrak{sl}_2$}{sl_2} \cite{Feigin-Frenkel}; \ClaimStatusProvedElsewhere]\label{thm:sl2-critical}
At $k = -h^\vee = -2$:
\begin{enumerate}
\item The Sugawara construction is \emph{undefined} ($k + h^\vee = 0$): no Virasoro subalgebra exists. The center of the completed enveloping algebra enlarges to the Feigin--Frenkel center $\mathfrak{z}(\widehat{\mathfrak{g}})$.
\item The associated classical PVA bracket retains the affine
central derivative term:
\begin{equation}
\{e(z), f(w)\} = -2\delta'(z-w) + h(w)\delta(z-w)
\end{equation}
This is not commutativity of the critical vertex algebra. The
central coefficient is the trace-form level $k=-2$; what vanishes is
the shifted curvature scalar $k+h^\vee$.
\item The underlying bar chain groups remain:
\begin{equation}
\bar{B}^n(\widehat{\mathfrak{sl}}_{2,-2}) = \bigoplus_{n_1+n_2+n_3=n} \Gamma(\overline{C}_{n+1}(X), \mathfrak{sl}_2^{\boxtimes(n+1)} \otimes \Omega^n_{\log})
\end{equation}
with square-zero shifted critical differential. The completed critical
center/Hochschild lane carries the Feigin--Frenkel center.
\end{enumerate}
\end{theorem}

\subsection{Koszul dual computation}

\begin{theorem}[Koszul dual of \texorpdfstring{$\widehat{\mathfrak{sl}}_{2,k}$}{sl-hat_2,k}; \ClaimStatusProvedHere]\label{thm:sl2-koszul-dual}
For $k \neq -2$, the Koszul dual is the chiral CE algebra at the dual level:
\begin{equation}
(\widehat{\mathfrak{sl}}_{2,k})^! = \mathrm{CE}^{\mathrm{ch}}(\widehat{\mathfrak{sl}}_{2,-k-4}),
\quad \kappa((\widehat{\mathfrak{sl}}_{2,k})^!) = \kappa(\widehat{\mathfrak{sl}}_{2,-k-4}).
\end{equation}
This follows from the explicit bar-cobar computation below.
\end{theorem}

\begin{proof}[Computation through degree 3]
We compute the bar complex $\bar{B}^{\leq 3}(\widehat{\mathfrak{sl}}_{2,k})$ and the cobar complex $\Omega^{\leq 3}(\bar{B}(\widehat{\mathfrak{sl}}_{2,k}))$.

\emph{Degree~0.} $\bar{B}^0 = \mathbb{C}$ (vacuum).

\emph{Degree~1.}
\begin{equation}
\bar{B}^1 = \Gamma(\overline{C}_2(X), \mathfrak{sl}_2 \boxtimes \mathfrak{sl}_2 \otimes \omega_X^{\boxtimes 2} \otimes \Omega^1_{\log})
\end{equation}
Basis elements:
\begin{align}
&e(z_1) \otimes f(z_2) \cdot d\log(z_1-z_2) \\
&f(z_1) \otimes e(z_2) \cdot d\log(z_1-z_2) \\
&h(z_1) \otimes h(z_2) \cdot d\log(z_1-z_2)
\end{align}

The differential $d: \bar{B}^1 \to \bar{B}^2$ is computed by taking residues:
\begin{align}
d(e \boxtimes f \cdot \eta_{12}) &= \mathrm{Res}_{z_1=z_2}[e(z_1)f(z_2) \cdot \eta_{12}] \\
&= k \cdot |0\rangle + h(z_2)|_{z_1=z_2}
\end{align}

\emph{Degree~2.}
The space $\bar{B}^2$ contains triple tensor products with logarithmic 2-forms:
\begin{equation}
\bar{B}^2 = \Gamma(\overline{C}_3(X), \mathfrak{sl}_2^{\boxtimes 3} \otimes \Omega^2_{\log})
\end{equation}

The $A_\infty$ structure satisfies $m_1^2 \neq 0$ for
$k\neq -h^\vee$. Equivalently, the bar object is curved away from
the critical level and obeys $d_{\bar B,k}^{\,2}=[m_0,-]$ rather
than strict nilpotence. The residue of the double-pole OPE
$e(z)f(w) \sim k/(z-w)^2 + h(w)/(z-w)$ contributes the curvature
(cf.\ Definition~\ref{def:curved-koszul-km}):
\begin{align}
m_0 &= \frac{k + h^\vee}{2h^\vee} \cdot \mathsf C_{\mathfrak{sl}_2},
\qquad m_1^2(a) = [m_0, a]
\end{align}
where $\mathsf C_{\mathfrak{sl}_2}$ is the Casimir element.

Thus the curvature vanishes if and only if
$k = -h^\vee = -2$ (critical level); for $k \neq -2$, the curved
$A_\infty$ curvature $m_0$ is proportional to $k+2$. The relation
$m_1^2=[m_0,-]$ is the curved identity; the curvature-vanishing
criterion is $m_0=0$ (equivalently the scalar fiberwise square
vanishes), not the bare commutator test in isolation.

Applying $\Omega$ to $\bar{B}$ gives free generators dual to the above, with differential twisted by the curvature.

The cobar complex produces fields with OPEs:
\begin{equation}
e^*(z) f^*(w) \sim \frac{-k-4}{(z-w)^2} + \frac{h^*(w)}{z-w}
\end{equation}
which is precisely $\widehat{\mathfrak{sl}}_{2,-k-4}$.
\end{proof}

\subsection{\texorpdfstring{Wakimoto realization for $\mathfrak{sl}_2$}{Wakimoto realization for sl-2}}

\begin{construction}[Wakimoto for $\mathfrak{sl}_2$ at $k=-2$]
The Wakimoto module uses:
\begin{itemize}
\item $\beta, \gamma$: a $\beta$-$\gamma$ system with weights $(1,0)$
\item $\phi$: a free boson (Cartan generator)
\end{itemize}

At $k = -h^\vee = -2$, the general Wakimoto formulas simplify:
the boson terms vanish (they carry a factor of $k + h^\vee = 0$) but the $\partial\gamma$ term
persists (its coefficient is $k$, not $k + h^\vee$):
\begin{align}
e(z) &= \beta(z) \\
f(z) &= -{:}\beta(z)\gamma(z)^2{:} - 2\partial\gamma(z) \\
h(z) &= -2{:}\beta(z)\gamma(z){:}
\end{align}
\end{construction}

\begin{computation}[Critical-level OPEs]
Using $\beta(z)\gamma(w) \sim 1/(z-w)$:
\begin{align}
e(z)f(w) &= \beta(z) \cdot (-{:}\beta(w)\gamma(w)^2{:} - 2\partial\gamma(w)) \\
&\sim \frac{-2{:}\beta(w)\gamma(w){:}}{z-w} + \frac{-2}{(z-w)^2}
= \frac{h(w)}{z-w} + \frac{k}{(z-w)^2}
\end{align}
The simple pole comes from $\beta(z)$ contracting with $\gamma(w)^2$; the double pole from $\beta(z)\cdot(-2\partial\gamma(w)) = -2/(z-w)^2 = k/(z-w)^2$. At critical level $k = -h^\vee$, the Sugawara construction\index{Sugawara construction!undefined at critical level} $T = \frac{1}{2(k+h^\vee)}\sum_a {:}J^a J^a{:}$ is \emph{undefined} because the coefficient has a pole.
By the Feigin--Frenkel theorem, the center of the completed enveloping
algebra at the critical level is isomorphic to the algebra of
functions on local ${}^L\mathfrak{g}$-opers on the formal disc.
\end{computation}

\begin{remark}[\texorpdfstring{$\widehat{\mathfrak{sl}}_2$}{sl-hat_2} Koszul duality]\label{rem:sl2-portrait}
Three structural features of Kac--Moody Koszul duality in their simplest form:
\begin{enumerate}
\item \emph{Level shift as curvature.} The double pole $k/(z-w)^2$ in $e(z)f(w)$ creates curvature $m_0$ in the bar complex with modular characteristic $\kappa(\widehat{\mathfrak{sl}}_{2,k}) = 3(k+2)/4$. Verdier duality on the bar coalgebra produces the dual level $-k-2h^\vee = -k-4$.
\item \emph{Critical level as curvature-free locus.} At $k=-2$, the
 curvature vanishes ($m_0 = 0$), the bar complex becomes a dg
 coalgebra, and the dual level is $-(-2)-4 = -2$: the critical level
 is a fixed point of the Feigin--Frenkel level involution. The
 Koszul dual
 $(\widehat{\mathfrak{sl}}_{2,-2})^! = \mathrm{CE}^{\mathrm{ch}}(\widehat{\mathfrak{sl}}_{2,-2})$
 has the same modular characteristic
 $\kappa = 0$ as the original algebra; this is fixed-level duality,
 not a claim that the algebra coincides with its own chiral Koszul
 dual. Simultaneously, the Sugawara construction degenerates and the
 center enlarges to the oper algebra.
\item \emph{Central charge non-cancellation.} Unlike the $\beta\gamma$/$bc$ pair (Proposition~\ref{prop:betagamma-bc-koszul-detailed}), the central charges satisfy $c(k) + c(-k-4) = 2\dim(\mathfrak{sl}_2) = 6$ rather than zero: curvature prevents cancellation.
\end{enumerate}
\end{remark}

\begin{remark}[Collision-residue $r$-matrix for $\hat\fg_k$]
\label{rem:km-collision-residue-rmatrix}
\index{r-matrix@$r$-matrix!Kac--Moody|textbf}
\index{Kac--Moody algebra!collision-residue $r$-matrix}
\index{CYBE!Kac--Moody}
The bar complex propagator $d\log(z_1 - z_2)$ absorbs one pole
order from the OPE (Proposition~\ref{prop:rmatrix-pole-landscape}).
For the OPE
$J^a(z)J^b(w) \sim k\,\kappa_{\fg}(a,b)/(z{-}w)^2 + f^{ab}{}_{c}J^c/(z{-}w)$,
the binary collision has two carriers. The current-valued residue
extracts the zeroth product
\[
a_{(0)}b=f^{ab}{}_{c}J^c,
\]
the Lie bracket. The invariant-form double pole gives the level
carrier \(a_{(1)}b=k\,\kappa_{\fg}(a,b)\). In the bar differential
this same double-pole carrier contributes the curvature \(m_0\); in
the landscape \(r\)-census its trace-form projection is recorded as
the level-prefixed collision package. This trace-form package is
separate from both the zeroth product and the KZ coefficient:
\begin{equation}\label{eq:km-rmatrix}
r^{\mathrm{coll}}_{\hat\fg_k}(z) = \frac{k\,\Omega_{\mathrm{tr}}}{z},
\qquad
r^{\mathrm{KZ}}_{\hat\fg_k}(z)=\frac{\Omega_{\mathrm{KZ}}}{(k+h^\vee)z}.
\end{equation}
Both have a single simple pole, but they record different
normalizations. The factor $(k + h^\vee)^{-1}$ belongs to the KZ
connection; at the critical level $k = -h^\vee$, the KZ coefficient
degenerates (Sugawara singularity).

The KZ $r$-matrix $r^{\mathrm{KZ}}(z)$ satisfies the classical
Yang--Baxter equation, which reduces via partial fractions to
the infinitesimal braid relation
$[\Omega_{ij}, \Omega_{ik} + \Omega_{jk}] = 0$
for all distinct triples $(i,j,k)$
(Proposition~\ref{prop:cybe-from-mc}).

\emph{Convention.}
This chapter keeps the two normalizations separate:
\[
r^{\mathrm{coll}}(z)=\frac{k\Omega_{\mathrm{tr}}}{z},
\qquad
r^{\mathrm{KZ}}(z) = \frac{\Omega}{(k+h^\vee)\,z},
\qquad
\Omega = \sum_{a,b} \kappa_{\fg}^{ab}\, I_a \otimes I_b,
\]
where $\kappa_{\fg}^{ab}$ is the inverse normalized invariant form and the
denominator $k + h^\vee$ is the Sugawara shift. This is the
normalization in which the KZ connection takes the standard form
$\nabla^{\mathrm{KZ}} = d - \sum_{i < j} r(z_{ij})\, dz_{ij}$
(Computation~\ref{comp:sl2-collision-residue-kz}).
At $k = 0$ the collision residue $k\Omega_{\mathrm{tr}}/z$ vanishes,
while the KZ coefficient $\Omega/(h^\vee z)$ does not. The abelian
vanishing $r = 0$ at $k = 0$ characterizes the Heisenberg $r$-matrix
(Chapter~\ref{ch:heisenberg-frame}), not the non-abelian
Kac--Moody KZ coefficient.
The landscape census (Table~\ref{tab:rmatrix-census}) records the
collision residue in the form $k\Omega_{\mathrm{tr}}/z$, where
$\Omega_{\mathrm{tr}}$ is the Casimir built from the
\emph{trace-form} basis.
%: r(z) lives on C_2(X), not over a point. The formal disk
% carries completion data; the homotopy retract A^1 -> pt is additional
% data, not an identity.
\end{remark}

\section{\texorpdfstring{Explicit computation: $\widehat{\mathfrak{sl}}_3$}{Explicit computation: sl-3}}
\label{sec:sl3-computation}

\subsection{Setup}

For $\mathfrak{sl}_3$, the data are: $h^\vee = 3$, $\dim(\mathfrak{sl}_3) = 8$, Cartan subalgebra $\mathfrak{h} = \mathrm{span}\{h_1, h_2\}$, simple roots $\alpha_1, \alpha_2$, and positive roots $\Delta_+ = \{\alpha_1, \alpha_2, \alpha_1+\alpha_2\}$.

\begin{definition}[\texorpdfstring{$\widehat{\mathfrak{sl}}_3$}{sl-hat_3} generators]
Current generators:
\begin{itemize}
\item Cartan currents: $h_1(z), h_2(z)$
\item Root currents: $e_{\alpha}(z)$ for $\alpha \in \Delta_+$, and $e_{-\alpha}(z)$ for $-\alpha \in \Delta_-$
\end{itemize}
with OPEs determined by the $\mathfrak{sl}_3$ structure constants and level $k$.
\end{definition}

\subsection{\texorpdfstring{Critical level: $k = -3$}{Critical level: k = -3}}

\begin{theorem}[Wakimoto for \texorpdfstring{$\mathfrak{sl}_3$}{sl_3} \cite{Frenkel-Kac-Wakimoto92}; \ClaimStatusProvedElsewhere]\label{thm:w3-wakimoto-sl3}
At $k = -3$, the Wakimoto module uses:
\begin{itemize}
\item $(\beta_{\alpha_1}, \gamma_{\alpha_1})$: $\beta$-$\gamma$ system for root $\alpha_1$
\item $(\beta_{\alpha_2}, \gamma_{\alpha_2})$: $\beta$-$\gamma$ system for root $\alpha_2$
\item $(\beta_{\alpha_1+\alpha_2}, \gamma_{\alpha_1+\alpha_2})$: $\beta$-$\gamma$ system for root $\alpha_1+\alpha_2$
\item $\phi_1, \phi_2$: free bosons for the Cartan
\end{itemize}

The currents are given by explicit formulas:
\begin{align}
e_{\alpha_i}(z) &= \beta_{\alpha_i}(z) \\
e_{-\alpha_i}(z) &= \text{differential polynomial in } \beta_{\alpha_i}, \gamma_{\alpha_i}, \phi_i, \partial\phi_i \\
h_i(z) &= -\alpha_i(\phi)(z) + \text{screening charge corrections}
\end{align}
\end{theorem}

\subsection{Level-shifting duality}

\begin{theorem}[Koszul dual of \texorpdfstring{$\widehat{\mathfrak{sl}}_{3,k}$}{sl-hat_3,k}; \ClaimStatusProvedHere]\label{thm:sl3-koszul-dual}
\begin{equation}
(\widehat{\mathfrak{sl}}_{3,k})^!
\simeq
\mathrm{CE}^{\mathrm{ch}}(\widehat{\mathfrak{sl}}_{3,-k-6})
\end{equation}
The level reflected in the underlying current presentation is
$-k - 2h^\vee = -k - 6$ since $h^\vee = 3$ for
$\mathfrak{sl}_3$.
\end{theorem}

\begin{proof}
This is Theorem~\ref{thm:universal-kac-moody-koszul} applied to
$\mathfrak{g} = \mathfrak{sl}_3$ with $h^\vee = 3$. The explicit bar
complex computation through degree~3
(Computation~\ref{comp:sl3-bar-dimensions}) is consistent: the bar
differential on $\bar{B}^1$ extracts the OPE data with the double-pole
coefficient $k$, and Verdier duality on the bar coalgebra produces the
CE generators whose current-level parameter is
$-k - 2 \cdot 3 = -k-6$. The curvature at non-critical levels satisfies
$m_0 = \frac{k+3}{6}\cdot\mathsf C_{\mathfrak{sl}_3}$,
vanishing precisely at $k = -3$.
\end{proof}

\subsection{Explicit bar complex through degree 3}

\begin{computation}[Bar complex chain-group dimensions]\label{comp:sl3-bar-dimensions}
The chain-group dimensions decompose as
$\dim(\bar{B}^n) = (\dim\mathfrak{g})^n \cdot (n-1)!$
(generator factor times Orlik--Solomon form factor;
see Computation~\ref{comp:sl3-bar}):
\begin{align}
\dim(\bar{B}^0) &= 1 \\
\dim(\bar{B}^1) &= 8 \cdot 1 = 8
 \quad \text{($8$ currents, $\dim\Omega^0 = 1$)} \\
\dim(\bar{B}^2) &= 64 \cdot 1 = 64
 \quad \text{($8^2$ ordered pairs, $\dim\Omega^1 = 1$)} \\
\dim(\bar{B}^3) &= 512 \cdot 2 = 1024
 \quad \text{($8^3$ ordered triples, $\dim\Omega^2 = 2$)}
\end{align}
The cubic growth $\dim(\bar{B}^n) \sim 8^n$ reflects the $8$-dimensional adjoint representation of $\mathfrak{sl}_3$; the Orlik--Solomon factor $(n-1)!$ accounts for independent logarithmic forms on $\overline{C}_n$.
\end{computation}

\begin{lemma}[Bar chain groups are level-independent; \ClaimStatusProvedHere]%
\label{lem:bar-dims-level-independent}
\index{bar complex!level-independence}
The dimensions $\dim(\bar{B}^n(\widehat{\mathfrak{g}}_k))$ depend on
$\mathfrak{g}$ but are \emph{independent of the level~$k$}. This is
immediate from the construction: the bar chain groups
\[
\bar{B}^n(\widehat{\mathfrak{g}}_k)
= \Gamma\!\bigl(\overline{C}_{n+1}(X),\;
 \mathfrak{g}^{\boxtimes(n+1)} \otimes \Omega^n_{\log}\bigr)
\]
involve only the Lie algebra $\mathfrak{g}$ (via its generators
$J^a$, $a = 1, \ldots, \dim\mathfrak{g}$) and the geometry of the
configuration space $\overline{C}_{n+1}(X)$. Neither ingredient
depends on~$k$. In particular, the chain-group dimension
is $\dim(\bar{B}^n) = (\dim\mathfrak{g})^n \cdot (n-1)!$,
the product of the generator factor $(\dim\mathfrak{g})^n$
(one copy of $\mathfrak{g}$ per configuration point) and
the form factor $(n-1)! = \dim\Omega^{n-1}(\mathrm{Conf}_n(\mathbb{C}))$
(top-degree Orlik--Solomon algebra). For instance,
$\mathfrak{sl}_2$ gives $3, 9, 54, 486, \ldots$ and
$\mathfrak{sl}_3$ gives $8, 64, 1024, 24576, \ldots$ regardless of
level (see Computation~\ref{comp:sl3-bar} for the
explicit tabulation).

The level enters only through the \emph{bar differential}
$d\colon \bar{B}^n \to \bar{B}^{n+1}$, which involves the OPE
residue containing the factor~$k$ in the double-pole coefficient
$J^a_{(1)}J^b = k\,\delta^{ab}$. Consequently the
curvature-flat comparison cohomology and the CDG/coderived
bar invariants may depend on~$k$, but the chain groups do not.
The comparison cohomology dimensions are separate cohomological data:
they are significantly smaller than the chain-group dimensions and
exhibit polynomial growth for Koszul algebras.
\end{lemma}

\subsection{Level decomposition of the bar differential}

\begin{theorem}[Curved level decomposition of the KM bar complex;
\ClaimStatusProvedHere]
\label{thm:km-bar-bicomplex}
\index{bar complex!curved level decomposition}
\index{bar complex!level decomposition}
\index{Kac--Moody algebra!bar differential!level decomposition}
For any simple Lie algebra $\mathfrak{g}$ and any level~$k$, write
the residue operator on the underlying graded bar coalgebra as
$d_k=d_{\mathrm{sp}}+k\,d_{\mathrm{dp}}$, where
$d_{\mathrm{sp}}$ is the simple-pole Lie-bracket part and
$d_{\mathrm{dp}}$ is the coefficient of the trace-form double-pole
part.  Set
\[
d_{\mathrm{crit}}:=d_{-h^\vee}
=d_{\mathrm{sp}}-h^\vee d_{\mathrm{dp}},
\qquad
\delta:=d_{\mathrm{dp}}.
\]
Then
\begin{equation}\label{eq:bar-level-decomp}
d_k = d_{\mathrm{crit}} + (k + h^\vee)\,\delta,
\end{equation}
and the square is the curved square
\begin{equation}\label{eq:bar-curved-level-square}
d_k^{\,2}
=
\left[
\frac{k+h^\vee}{2h^\vee}\mathsf C_{\fg},-
\right].
\end{equation}
Equivalently, the two level-independent pieces satisfy
\begin{equation}\label{eq:bar-curved-bicomplex-conditions}
d_{\mathrm{crit}}^2 = 0,\qquad
\delta^2 = 0,\qquad
d_{\mathrm{crit}}\delta+\delta d_{\mathrm{crit}}
=
\left[\frac{\mathsf C_{\fg}}{2h^\vee},-\right].
\end{equation}
Thus the triple
$(\bar{B}^\bullet,d_{\mathrm{crit}},\delta)$ is an ordinary
bicomplex exactly on the curvature-flat subquotient on which
$[\mathsf C_{\fg},-]=0$.  On the full bar object it is a curved
bicomplex.  The critical differential still contains the
level-$(-h^\vee)$ double-pole term; what vanishes at critical level
is the shifted curvature class, not the OPE double pole.
\end{theorem}

\begin{proof}
The OPE of $\widehat{\mathfrak{g}}_k$ is
$J^a(z) J^b(w) \sim k(J^a, J^b)/(z-w)^2 + f^{ab}{}_c J^c(w)/(z-w)$.
The bar differential acts on $\bar{B}^n$ by extracting residues at collision
divisors. At each collision $z_i \to z_j$, the residue
separates into:
\begin{itemize}
\item A \emph{simple-pole contribution}:
$\mathrm{Res}_{z_i = z_j}[f^{ab}{}_c J^c / (z_i - z_j)]
= f^{ab}{}_c J^c$, independent of~$k$.
\item A \emph{double-pole contribution}:
$\mathrm{Res}_{z_i = z_j}[k(J^a, J^b) / (z_i - z_j)^2]
= k(J^a, J^b) \cdot |0\rangle$, linear in~$k$.
\end{itemize}
This gives \(d_k=d_{\mathrm{sp}}+k\,d_{\mathrm{dp}}\), hence
\[
d_{\mathrm{crit}}=d_{-h^\vee}=d_{\mathrm{sp}}-h^\vee d_{\mathrm{dp}},
\qquad
\delta=d_{\mathrm{dp}},
\]
and hence
\[
d_k=d_{\mathrm{crit}}+(k+h^\vee)\delta.
\]
This shifted form records the Feigin--Frenkel critical parameter
\(k+h^\vee\). It does not assert that the trace-form double-pole OPE
coefficient \(k(J^a,J^b)\) vanishes at \(k=-h^\vee\).

The full bar object is curved away from the critical level:
\[
d_k^{\,2}
=
\left[
\frac{k+h^\vee}{2h^\vee}\mathsf C_{\fg},-
\right]
\]
by Definition~\ref{def:curved-koszul-km} and
Proposition~\ref{prop:km-bar-curvature}.  Expanding
\(d_k=d_{\mathrm{crit}}+\lambda\delta\), with
\(\lambda=k+h^\vee\), gives
\[
d_k^2
= d_{\mathrm{crit}}^2
+ \lambda(d_{\mathrm{crit}}\delta + \delta\, d_{\mathrm{crit}})
+ \lambda^2 \delta^2
=
\lambda\left[\frac{\mathsf C_{\fg}}{2h^\vee},-\right].
\]
This is a polynomial identity in the variable
$\lambda$; comparing coefficients gives
\eqref{eq:bar-curved-bicomplex-conditions}.
\end{proof}

\begin{corollary}[Critical-level curved spectral sequence;
\ClaimStatusProvedHere]
\label{cor:critical-level-spectral}
\index{spectral sequence!critical level}
\index{bar complex!spectral sequence!level parameter}
The decomposition
\[
d_k=d_{\mathrm{crit}}+(k+h^\vee)\delta
\]
gives a critical-first filtered perturbation of the critical bar
complex.  Its first page is
\[
E_1^{p,q}
=H^{p+q}(\bar{B}^{\bullet},d_{\mathrm{crit}}),
\]
and the first perturbing operator is induced by
$(k+h^\vee)\delta$ on every curvature-flat summand.  On the full bar
object the obstruction to an ordinary spectral sequence is exactly
the curvature operator
\[
\left[\frac{k+h^\vee}{2h^\vee}\mathsf C_{\fg},-\right].
\]
Therefore ordinary cohomology for the generic-level bar complex is
legitimate only after passing to the curvature-flat comparison
subquotient, after a curvature-killing Maurer--Cartan twist, or in
the coderived/CDG totalization where the curved square is part of the
structure.
\end{corollary}

\begin{proof}
Filter the completed bar object by powers of
\(\lambda=k+h^\vee\).  At \(\lambda=0\) the differential is
\(d_{\mathrm{crit}}\) and is square-zero by
Theorem~\ref{thm:km-bar-bicomplex}.  The first-order perturbation is
\(\lambda\delta\).  The relation
\[
d_{\mathrm{crit}}\delta+\delta d_{\mathrm{crit}}
= [\mathsf C_{\fg}/(2h^\vee),-]
\]
shows that \(\delta\) descends to an honest differential on the
\(E_1\)-page precisely after the curvature action vanishes or is
killed by the chosen curved totalization.  This is the curved analogue
of the usual bicomplex spectral sequence and records the same
critical-to-generic deformation while preserving the identity
\(d_k^2=[m_0(k),-]\) away from the critical level.
\end{proof}

\begin{remark}[Two ordinary faces of the curved level decomposition]
\label{rem:km-curved-level-ordinary-faces}
The curved decomposition has two square-zero shadows:
\begin{enumerate}[label=\textup{(\roman*)}]
\item At $k=-h^\vee$, the curvature coefficient is zero and
$d_{\mathrm{crit}}^2=0$ on the full bar coalgebra.
\item On any central or scalar quotient on which
$[\mathsf C_{\fg},-]=0$, the pair
$(d_{\mathrm{crit}},\delta)$ is an ordinary bicomplex and the usual
critical-first spectral sequence applies.
\end{enumerate}
The generic noncritical bar object itself remains curved:
$d_k^2=[m_0(k),-]$.
\end{remark}

\begin{remark}[Interpretation of the curved level decomposition]
\label{rem:bicomplex-interpretation}
\index{bar complex!bicomplex!interpretation}
The decomposition~\eqref{eq:bar-level-decomp} reflects the
level-independence of bar chain groups
(Lemma~\ref{lem:bar-dims-level-independent}):
\(d_{\mathrm{crit}}=d_{\mathrm{sp}}-h^\vee d_{\mathrm{dp}}\) is the
critical differential, \(\delta=d_{\mathrm{dp}}\) is the double-pole
coefficient, and
\[
m_0=\frac{k+h^\vee}{2h^\vee}\mathsf C_{\fg}
\]
is the shifted curvature class.
\begin{enumerate}[label=\textup{(\roman*)}]
\item At $k = -h^\vee$: the shifted curvature vanishes and the bar
complex is uncurved with differential $d_{\mathrm{crit}}$; the
trace-form double-pole OPE coefficient is still $-h^\vee$.
\item For generic $k \neq -h^\vee$: the critical-first curved
perturbation relates the generic CDG bar object to critical bar
cohomology; it becomes an ordinary spectral sequence only on the
curvature-flat comparison surface of
Corollary~\ref{cor:critical-level-spectral}.
\item At admissible $k = -h^\vee + p/q$: rationality of $(k + h^\vee)$ may create resonance phenomena.
\end{enumerate}
\end{remark}

\begin{theorem}[Generic level-independence on the curvature-flat
comparison surface; \ClaimStatusProvedHere]
\label{thm:bar-cohomology-level-independence}
\index{bar complex!level-independence!cohomology}
\index{Kac--Moody algebra!bar cohomology!level-independence}
Let $C^\bullet_{\mathrm{flat}}$ be a finite-dimensional
curvature-flat comparison subquotient of
\(\barB^\bullet(\widehat{\fg}_k)\), preserved by
$d_{\mathrm{crit}}$ and~$\delta$, on which
$[\mathsf C_{\fg},-]=0$.  Then
$d_\lambda=d_{\mathrm{crit}}+\lambda\delta$,
$\lambda=k+h^\vee$, is square-zero on
$C^\bullet_{\mathrm{flat}}$ for every~$\lambda$, and for each~$n$
there is a finite exceptional set
$\Sigma_n(\fg;C_{\mathrm{flat}})\subset\bC$ such that
\[
\lambda \;\mapsto\;
\dim H^n(C^\bullet_{\mathrm{flat}},d_\lambda)
\]
is constant on $\mathbb{C} \setminus
(\Sigma_n \cup \Sigma_{n+1})$ and can only
\emph{increase} at the special values.
Thus the ordinary cohomology dimensions computed on the
curvature-flat comparison surface at any generic level equal those
at every other generic level.  On the full noncritical bar object
one retains the CDG identity \(d_\lambda^2=[m_0(\lambda),-]\)
instead of ordinary cohomology.
\end{theorem}

\begin{proof}
Since $C^n_{\mathrm{flat}}$ is finite-dimensional for each~$n$,
the differential $d_\lambda = d_{\mathrm{crit}} + \lambda\delta
\colon C^n_{\mathrm{flat}} \to C^{n-1}_{\mathrm{flat}}$
is a matrix whose entries are affine-linear functions of~$\lambda$.
For a matrix $A(\lambda)$ depending polynomially on a
parameter $\lambda \in \mathbb{C}$, the rank function
$\lambda \mapsto \operatorname{rk} A(\lambda)$ is
lower-semicontinuous: it is constant on the complement
of a finite set (the zero locus of the maximal minors)
and can only \emph{drop} at finitely many values.
Define
\begin{align*}
r_n(\lambda) &= \operatorname{rk}\bigl(d_\lambda
 \colon C^n_{\mathrm{flat}} \to C^{n-1}_{\mathrm{flat}}\bigr), \\
\Sigma_n &= \{\lambda \in \mathbb{C} :
 r_n(\lambda) < r_n^{\mathrm{gen}}\},
\end{align*}
where $r_n^{\mathrm{gen}}$ is the generic rank. Then
\[
\dim H^n(C^\bullet_{\mathrm{flat}}, d_\lambda)
= \dim C^n_{\mathrm{flat}} - r_n(\lambda) - r_{n+1}(\lambda),
\]
where $\dim C^n_{\mathrm{flat}}$ is independent of~$\lambda$.
For
$\lambda \notin \Sigma_n \cup \Sigma_{n+1}$,
both $r_n$ and $r_{n+1}$ attain their generic values,
so $\dim H^n$ is constant.
At a special value $\lambda_0 \in \Sigma_n \cup \Sigma_{n+1}$,
either $r_n$ or $r_{n+1}$ drops, causing $\dim H^n$
to (weakly) increase.
\end{proof}

\begin{remark}[Special values and admissible levels]
\label{rem:special-level-values}
\index{admissible level!bar cohomology jump}
The exceptional set
$\Sigma_n(\fg;C_{\mathrm{flat}})$ has at most
$\dim C^n_{\mathrm{flat}}\cdot\dim C^{n-1}_{\mathrm{flat}}$
values.  For $\mathfrak{sl}_2$ at degree~$2$ on the standard
curvature-flat comparison surface, the matrix
$d_\lambda\colon C^2_{\mathrm{flat}}\to C^1_{\mathrm{flat}}$ is
linear in~$\lambda$ with generic rank for all $\lambda\neq0$,
so the critical value $\lambda=0$ ($k=-2$) is the first jump
candidate.  Admissible levels $k=-2+p/q$ are the natural candidates
for further non-generic behavior
(Remark~\textup{\ref{rem:bicomplex-interpretation}}\textup{(iii)}).
\end{remark}

\begin{conjecture}[$\mathfrak{sl}_2$ bar cohomology as spin representations]
\label{conj:sl2-bar-spin-representations}
\ClaimStatusConjectured
\index{bar complex!$\widehat{\mathfrak{sl}}_2$!spin representation|textbf}
\index{Garland--Lepowsky concentration!representation-theoretic refinement}
For each $n \geq 1$, the curvature-flat comparison cohomology
$H^n(C^\bullet_{\mathrm{flat}}(\widehat{\mathfrak{sl}}_{2,k}))$
at generic level~$k$ carries a natural
$\mathfrak{sl}_2$-module structure
\textup{(}from the adjoint action of~$\mathfrak{sl}_2$
on its negative loop algebra\textup{)},
and this module is isomorphic to the spin-$n$
irreducible representation~$V_n$ of~$\mathfrak{sl}_2$:
\begin{equation}\label{eq:sl2-bar-spin-n}
 H^n\bigl(C^\bullet_{\mathrm{flat}}
  (\widehat{\mathfrak{sl}}_{2,k})\bigr)
 \;\cong\; V_n
 \qquad(n \geq 1),
\end{equation}
where \(C^\bullet_{\mathrm{flat}}\) denotes the comparison surface of
Theorem~\ref{thm:bar-cohomology-level-independence} and
$\dim V_n = 2n + 1$. The Garland--Lepowsky
concentration theorem gives $\dim H^n = 2n + 1$
\textup{(}Remark~\textup{\ref{rem:garland-lepowsky-sl2}}),
which is necessary for~\eqref{eq:sl2-bar-spin-n} but does
not by itself determine the module structure. The
conjecture asserts that the $\mathfrak{sl}_2$-weight
decomposition of~$H^n$ is
$H^n = \bigoplus_{m = -n}^{n} (H^n)_m$
with each weight space one-dimensional, and that the
Casimir eigenvalue is $n(n+1)$.

The comparison cohomology at conformal weight $n(n{+}1)/2$
\textup{(}the unique weight at which $H^n \neq 0$
by Garland--Lepowsky concentration\textup{)} is
concentrated in a single $\mathfrak{sl}_2$-isotypic
component, with the bar degree~$n$ matching the
spin label.
\end{conjecture}

\begin{remark}[Low-degree computation and harmonic-space condition]
\label{rem:sl2-bar-spin-evidence}
The comparison-surface dimension identity $\dim H^n = 2n + 1$
\textup{(}Remark~\textup{\ref{rem:garland-lepowsky-sl2}})
has been verified for \(n\leq10\) by the local
\texttt{bar\_cohomology\_sl2\_explicit\_engine}. The
representation-theoretic refinement is known in degrees
$n=1,2,3$ by explicit weight decomposition of the CE cohomology
$H^n_{\mathrm{CE}}(\mathfrak{sl}_2^-; \mathbb{C})$:
at $n = 1$, $H^1 \cong \mathfrak{sl}_2$ (adjoint $= V_1$);
at $n = 2$, explicit computation gives a $5$-dimensional
space with weights $\{-2, -1, 0, 1, 2\}$, each of
multiplicity one, hence \(V_2\);
at $n = 3$, a $7$-dimensional space with weights
$\{-3, \ldots, 3\}$, hence \(V_3\).
The all-degree mathematical condition is an identification of \(H^n\)
with the Kostant harmonic space~\(\mathcal{H}^n\) in the nilpotent
radical, restricted to loop grading \(n(n+1)/2\).
\end{remark}

\section{\texorpdfstring{General $\widehat{\mathfrak{g}}$: functorial construction}{General g: functorial construction}}

\subsection{Abstract setting}

\begin{theorem}[Universal Koszul duality for affine Kac--Moody; \ClaimStatusProvedHere]\label{thm:universal-kac-moody-koszul}
\textup{[Regime: curved-central
\textup{(}Convention~\textup{\ref{conv:regime-tags})}.]}

For any simple Lie algebra $\mathfrak{g}$ and level
$k \neq -h^\vee$, there is a canonical Koszul duality:
\begin{equation}
(\widehat{\mathfrak{g}}_k)^!
\simeq
\mathrm{CE}^{\mathrm{ch}}(\widehat{\mathfrak{g}}_{-k-2h^\vee})
\end{equation}
This duality:
\begin{enumerate}
\item Is functorial for morphisms of pairs preserving the invariant
form and level parameter.
\item Identifies the Koszul dual with the chiral CE algebra of the
level-reflected current presentation, not with the affine vertex
algebra itself.
\item Intertwines the level $k$ and level $-k-2h^\vee$ bar--cobar
surfaces in the universal algebra.
\item Is compatible at the critical fixed point with the external
Feigin--Frenkel center theorem; no geometric-Langlands equivalence is
deduced from this theorem alone.
\end{enumerate}
\end{theorem}

\subsection{Proof of universal level-shifting duality}

\begin{proof}
The proof proceeds in four steps: construction of the bar complex, identification of the curvature, computation of the cobar OPE, and functorial uniqueness.

\emph{Step~1: Bar complex construction.}
The geometric bar complex of $\widehat{\mathfrak{g}}_k$ is:
\begin{equation}
\bar{B}^n(\widehat{\mathfrak{g}}_k) = \Gamma\bigl(\overline{C}_{n+1}(X),\; \mathfrak{g}^{\boxtimes(n+1)} \otimes \omega_X^{\boxtimes(n+1)} \otimes \Omega^n_{\log}\bigr)
\end{equation}
with differential given by iterated residues $d(\omega) = \sum_{i < j} \mathrm{Res}_{z_i=z_j}[\omega]$ at the collision divisors. In degree~1, a basis consists of elements $J^a(z_1) \otimes J^b(z_2) \cdot \eta_{12}$ where $\{J^a\}$ runs over a basis of~$\mathfrak{g}$ and $\eta_{12} = d\log(z_1 - z_2)$. The residue differential extracts the OPE data:
\begin{equation}\label{eq:bar-diff-general}
d(J^a \boxtimes J^b \cdot \eta_{12}) = k \cdot (J^a, J^b) \cdot |0\rangle + f^{ab}{}_{c}\, J^c
\end{equation}
where $(J^a, J^b)$ is the normalized invariant form and $f^{ab}{}_c$ are the structure constants. The first term comes from the double pole $k(J^a, J^b)/(z-w)^2$ and the second from the simple pole $f^{ab}{}_c J^c(w)/(z-w)$.

\begin{lemma}[Killing form via structure constants;
\ClaimStatusProvedHere]\label{lem:killing-structure-constants}
Let $\mathfrak{g}$ be a simple Lie algebra with structure
constants $f^{ab}{}_{c}$ in an orthonormal basis
$\{T^a\}$ for the invariant form $(\cdot,\cdot)$
normalised so that $(\theta|\theta) = 2$ for the
highest root~$\theta$. Then
\begin{equation}\label{eq:killing-structure-constants}
\sum_{c,d} f^{ac}{}_{d}\, f^{bc}{}_{d}
\;=\; 2h^\vee \cdot \delta^{ab},
\end{equation}
where $h^\vee$ is the dual Coxeter number.
\end{lemma}

\begin{proof}
The left-hand side of~\eqref{eq:killing-structure-constants}
is $\operatorname{Tr}\bigl(\operatorname{ad}(T^a) \circ
\operatorname{ad}(T^b)\bigr) = K(T^a, T^b)$,
the Killing form.
For simple~$\mathfrak{g}$, $K = \lambda\,(\cdot,\cdot)$
for a unique scalar~$\lambda$ (Schur's lemma, since the
adjoint representation is irreducible).
The quadratic Casimir $C_2 = \sum_a
(\operatorname{ad}(T^a))^2$ acts on~$\mathfrak{g}$
by the scalar $2h^\vee$ (this is the definition of~$h^\vee$
in the normalisation $(\theta|\theta) = 2$).
Taking the trace: $\operatorname{Tr}(C_2) = \sum_a
K(T^a, T^a) = \lambda \cdot \dim\mathfrak{g}$,
while $\operatorname{Tr}(C_2) = 2h^\vee \cdot
\dim\mathfrak{g}$, giving $\lambda = 2h^\vee$.
In the orthonormal basis, $K(T^a, T^b)
= 2h^\vee\,\delta^{ab}$.
\end{proof}

\emph{Step~2: Curvature identification.}
Applying $m_1$ twice to a degree-2 element and using the Jacobi identity for $\mathfrak{g}$, one computes the curvature:
\begin{equation}\label{eq:d-squared-general}
m_1^2(J^a \boxtimes J^b \cdot \eta_{12}) = (k + h^\vee) \cdot \bigl(\text{adjoint Casimir contribution}\bigr)
\end{equation}
The factor $h^\vee$ arises from the adjoint representation: by Lemma~\ref{lem:killing-structure-constants}, $\sum_{c,d} f^{ac}{}_d f^{bc}{}_d = 2h^\vee \cdot \delta^{ab}$, which gives $\sum_{c,d,e} f^{ac}{}_d f^{bc}{}_e (J^d, J^e) = 2h^\vee \cdot (J^a, J^b)$. For $\mathfrak{sl}_2$ with $h^\vee = 2$, Theorem~\ref{thm:sl2-koszul-dual} gives $m_1^2(e \boxtimes f \cdot \eta_{12}) = (k+2) \cdot \partial h$.

Therefore the unary bar differential is strictly square-zero exactly
at the critical level $k=-h^\vee$. For $k \neq -h^\vee$, the bar
object is a curved $A_\infty$-coalgebra and satisfies the curved
identity $d_{\bar B,k}^{\,2}=[m_0,-]$, with curvature
\begin{equation}
m_0 =
\frac{k + h^\vee}{2h^\vee} \cdot \mathsf C_{\fg}
\in \bar{B}^0
\end{equation}
where $\mathsf C_{\fg} = \sum_a J^a \otimes J_a$ is the Casimir
element.

\emph{Step~3: Dual bar-cobar computation and level shift.}
For the Koszul dual, form the completed Verdier-linear dual
\[
\bar B(\widehat{\mathfrak g}_k)^{\vee_{\mathrm V}}
\]
of the geometric bar coalgebra, then apply the cobar functor.  The
generators of
\(\Omega(\bar B(\widehat{\mathfrak g}_k)^{\vee_{\mathrm V}})\)
are the dual currents $(J^a)^* = J^{a,*}$, and the cobar differential
$d_\Omega$ is determined by dualizing the residue differential
\eqref{eq:bar-diff-general}.  This is the Koszul-dual construction.
The separate inversion construction
\(\Omega(\bar B(\widehat{\mathfrak g}_k))\) recovers
\(\widehat{\mathfrak g}_k\).

Dualizing the double-pole contribution: the term $k \cdot (J^a, J^b) \cdot |0\rangle$ in the bar differential dualizes to a double-pole term in the cobar OPE. Concretely, the cobar OPE between dual generators is:
\begin{equation}\label{eq:cobar-ope}
J^{a,*}(z) J^{b,*}(w) \sim \frac{k^* (J^a, J^b)}{(z-w)^2} + \frac{f^{ab}{}_{c}\, J^{c,*}(w)}{z-w}
\end{equation}
where $k^*$ is the dual level to be determined.

To find $k^*$, we use the Verdier duality pairing on the
Fulton--MacPherson compactification $\overline{C}_{n+1}(X)$.
By Proposition~\ref{prop:verdier-level-identification} below,
the line bundle carrying the level-$k$ data satisfies:
\begin{equation}
\mathcal{L}_k^\vee \otimes \omega_{\overline{C}_{n+1}} \cong \mathcal{L}_{-k-2h^\vee}
\end{equation}

Therefore $k^* = -k - 2h^\vee$, and the cobar OPE~\eqref{eq:cobar-ope} becomes:
\begin{equation}
J^{a,*}(z) J^{b,*}(w) \sim \frac{(-k-2h^\vee)(J^a, J^b)}{(z-w)^2} + \frac{f^{ab}{}_{c}\, J^{c,*}(w)}{z-w}
\end{equation}
which is precisely the current presentation underlying
\(\mathrm{CE}^{\mathrm{ch}}(\widehat{\mathfrak{g}}_{-k-2h^\vee})\).
(For $\mathfrak{sl}_2$: $-k-2 \cdot 2 = -k-4$, matching
Theorem~\ref{thm:sl2-koszul-dual}.)

\emph{Step~4: Functorial uniqueness.}
The inversion cobar construction
\(\Omega(\bar{B}(\widehat{\mathfrak{g}}_k))\) is characterized by the
usual bar-cobar universal property: for any chiral algebra~$\mathcal{A}$,
\begin{equation}
\mathrm{Hom}_{\text{ChirAlg}}\bigl(\Omega(\bar{B}(\widehat{\mathfrak{g}}_k)),\, \mathcal{A}\bigr) \cong \mathrm{Hom}_{\text{Coalg}}\bigl(\bar{B}(\widehat{\mathfrak{g}}_k),\, \bar{B}(\mathcal{A})\bigr)
\end{equation}
The same universal property applied after the completed Verdier-linear
dual determines the Koszul dual up to quasi-isomorphism.  Since the
dual bar OPE has structure constants $f^{ab}{}_c$ and level
$-k-2h^\vee$, the Koszul dual algebra is
\[
(\widehat{\mathfrak{g}}_k)^!
= \mathrm{CE}^{\mathrm{ch}}(\widehat{\mathfrak{g}}_{-k-2h^\vee}),
\]
with modular characteristic
\(\kappa(\widehat{\mathfrak{g}}_{-k-2h^\vee})\).  The CE algebra is
distinct from the affine current algebra
\(\widehat{\mathfrak{g}}_{-k-2h^\vee}\).  Bar-cobar inversion
\(\Omega(\bar{B}(\widehat{\mathfrak{g}}_k)) \simeq
\widehat{\mathfrak{g}}_k\) recovers the original algebra at level~$k$.

\emph{Internal compatibilities.}
\begin{enumerate}
\item \emph{Involutivity}: The level involution $k \mapsto -k - 2h^\vee$ is involutive: $-(-k-2h^\vee) - 2h^\vee = k$. At the level of modular characteristics, $\kappa((\widehat{\mathfrak{g}}_k)^{!!}) = \kappa(\widehat{\mathfrak{g}}_k)$.
\item \emph{Critical level fixed point}: At $k = -h^\vee$, the dual level is $-(-h^\vee) - 2h^\vee = -h^\vee$, so the critical level is a fixed point of the duality involution. This is the critical specialization of the internal level-shift formula (Proposition~\ref{prop:verdier-level-identification}) and is compatible with the external Feigin--Frenkel critical-level package (Theorem~\ref{thm:critical-level-structure}).
\item \emph{Explicit computation}: For $\mathfrak{sl}_2$ (Theorem~\ref{thm:sl2-koszul-dual}) and $\mathfrak{sl}_3$ (\S\ref{sec:sl3-computation}), the low-degree bar complex computations confirm the predicted dual level.
\end{enumerate}
\end{proof}

\begin{proposition}[Verdier level identification; \ClaimStatusConditional]
\label{prop:verdier-level-identification}
\index{Verdier duality!level identification|textbf}
\index{Feigin--Frenkel!from bar-cobar|textbf}
Let $\fg$ be a simple Lie algebra with dual Coxeter number~$h^\vee$,
and let $\mathcal{L}_k$ denote the $\mathcal{D}$-module on the
Fulton--MacPherson compactification~$\overline{C}_n(X)$ carrying the
level-$k$ chiral algebra data
\textup{(}the external tensor power of the current algebra sheaf
twisted by the invariant form $k \cdot (\,,\,)$\textup{)}.
Then:
\begin{equation}\label{eq:verdier-level-shift}
\mathcal{L}_k^\vee \otimes \omega_{\overline{C}_n(X)}
\;\cong\; \mathcal{L}_{-k - 2h^\vee}.
\end{equation}
\end{proposition}

\begin{proof}
We prove the identity on $\overline{C}_2(X)$ and extend by
factorization.

\emph{Binary case \textup{(}$n = 2$\textup{)}.}
On $\overline{C}_2(X) \cong \mathrm{Bl}_\Delta(X \times X)$, the
chiral algebra data near the diagonal~$\Delta$ is encoded in the
$\mathcal{D}$-module $\mathcal{M}_k$ whose fiber at
$(z,w)$ with $z \neq w$ is the regular representation of
$\widehat{\fg}_k$ with OPE singularity $J^a(z) J^b(w)$.
The Verdier dual $\mathbb{D}\mathcal{M}_k :=
R\mathcal{H}\!\mathit{om}_{\mathcal{D}}
(\mathcal{M}_k, \mathcal{D} \otimes \omega^{-1})[\dim]$
acts on the flat connection by the adjoint:
$\nabla \mapsto \nabla^\dagger$.

On the local model near $\Delta$, the flat connection is
$\nabla = d - A(z,w)$ where $A$ encodes the OPE poles.
The adjoint connection $\nabla^\dagger = d + A^t$ reverses the
sign of the first-order pole (the structure constants
$f^{ab}{}_c$ are anti-symmetric, so transposition gives
$-f^{ab}{}_c$, but this is absorbed by the Lie algebra
anti-involution $J^a \mapsto -J^a$, which acts as the identity
on the normalized invariant form).
The net effect on the \emph{invariant form}:
the double-pole coefficient $k \cdot (J^a, J^b)$ transforms as
\begin{equation}\label{eq:adjoint-casimir}
k \;\longmapsto\; -(k + 2h^\vee)
\end{equation}
under $\nabla \mapsto \nabla^\dagger$. The sign flip
$k \mapsto -k$ comes from the orientation reversal in Verdier
duality (exchanging $j_*$ and $j_!$ extensions near~$\Delta$).
The shift by $2h^\vee$ arises from the \emph{adjoint Casimir
contribution}: the dualizing sheaf~$\omega_{\overline{C}_2}$
has a first-order pole along~$\Delta$ whose residue is the
Casimir element $C_{\mathrm{adj}} = 2h^\vee \cdot \mathrm{id}$
acting on $\fg \otimes \fg$ via the adjoint representation.
This is the same Casimir eigenvalue that appears in the
Sugawara construction
($L_{\mathrm{Sug}} = \tfrac{1}{2(k+h^\vee)} \sum_a J^a J^a$,
where the denominator $k + h^\vee$ reflects the sum $k + C_{\mathrm{adj}}/2$).

\emph{Higher degree.}
The FM compactification $\overline{C}_n(X)$ is stratified by trees.
On each stratum, the $\mathcal{D}$-module $\mathcal{L}_k$ factors
as an external tensor product over the vertices of the tree
(this is the factorization property of chiral algebras,
cf.\ \cite[\S3.4]{BD04}).
Verdier duality commutes with external tensor product on proper
varieties, so the level identification~\eqref{eq:verdier-level-shift}
on $\overline{C}_2(X)$ propagates to all strata of
$\overline{C}_n(X)$ by induction on the number of internal edges.

\emph{Consistency.}
For $\fg = \mathfrak{sl}_2$ ($h^\vee = 2$):
$k \mapsto -k - 4$, matching
Theorem~\ref{thm:sl2-koszul-dual}.
For $\fg = \mathfrak{sl}_3$ ($h^\vee = 3$):
$k \mapsto -k - 6$, matching
\S\ref{sec:sl3-computation}.
\end{proof}

\begin{remark}[Higher-genus extension of the level identification]
\label{rem:level-shift-higher-genus}
\index{level identification!higher genus}
The identification~\eqref{eq:verdier-level-shift} is a genus-$0$
statement: it determines the Koszul dual \emph{level} from the
binary collision data on~$\overline{C}_2(X)$. On the uniform-weight
surface of affine currents, the higher-genus scalar obstruction is then
determined by the modular characteristic
$\kappa(\widehat{\fg}_k) = \dim(\fg)(k + h^\vee)/(2h^\vee)$
(Theorem~\ref{thm:modular-characteristic}(i)):
\(\mathrm{obs}_g(\widehat{\fg}_k)
=\kappa(\widehat{\fg}_k)\lambda_g\) for every $g \geq 1$
\textup{(}: all currents share conformal weight~$1$, so
the scalar formula holds without a cross-channel correction,
cf.\ Theorem~D\textup{)}.
The duality clause
\(\kappa(\widehat{\fg}_k)
 +\kappa(\widehat{\fg}_{-k-2h^\vee})=0\) for affine Kac--Moody is then a
\emph{consequence} of the genus-$0$ level identification:
$\kappa' = \dim(\fg)(-k - 2h^\vee + h^\vee)/(2h^\vee)
= -\dim(\fg)(k + h^\vee)/(2h^\vee) = -\kappa$.
\end{remark}

\begin{remark}[Canonical twisting morphism for Kac--Moody]
\label{rem:km-twisting-morphism}
The canonical twisting morphism
$\tau_{\mathrm{KM}} \colon \barB(\widehat{\fg}_k) \to \widehat{\fg}_k$
projects bar generators $s^{-1}J^a_{(r)}$ to the corresponding
currents~$J^a_{(r)}$. The Maurer--Cartan equation
$\partial\tau + \tau \star \tau = 0$ reduces to the Jacobi identity
for~$\fg$ together with the cocycle condition for the level~$k$:
the convolution square~$\tau \star \tau$ extracts the structure
constants $f^{ab}{}_c$ (from the simple pole) and the invariant
form $k\kappa_{\fg,ab}$ (from the double pole). Thus the Koszul
property of $\widehat{\fg}_k$ is controlled by the Chevalley--Eilenberg
cohomology of~$\fg$.
\end{remark}

\begin{definition}[Double-pole and simple-pole channel
decomposition]\label{def:channel-decomposition}
\index{channel decomposition!double-pole/simple-pole}
\index{modular characteristic!channel decomposition}
For $\widehat{\fg}_k$ with $\fg$ simple, the modular
characteristic $\kappa(\widehat{\fg}_k) =
\dim(\fg)(k + h^\vee)/(2h^\vee)$
(Theorem~\ref{thm:modular-characteristic}) decomposes into
\emph{double-pole} and \emph{simple-pole} channels:
\begin{equation}\label{eq:kappa-dp-sp}
\kappa = \kappa_{\mathrm{dp}} + \kappa_{\mathrm{sp}},
\qquad
\kappa_{\mathrm{dp}}
:= \frac{k \cdot \dim\fg}{2h^\vee},
\qquad
\kappa_{\mathrm{sp}}
:= \frac{\dim\fg}{2}.
\end{equation}
Here $\kappa_{\mathrm{dp}}$ is the vacuum leakage from the
double-pole OPE term $k(J^a, J^b)/(z-w)^2$ (the level
contribution), and $\kappa_{\mathrm{sp}}$ is the self-contraction
contribution from the simple-pole term
$f^{ab}{}_{c}\,J^c(w)/(z-w)$, arising through the adjoint
Casimir eigenvalue
$\sum_{c,d} f^{ac}{}_d f^{bc}{}_d = 2h^\vee\,\delta^{ab}$
(cf.~\eqref{eq:d-squared-general}).
\end{definition}

\begin{computation}[First-principles derivation of
$\kappa(\widehat{\fg}_k) = \dim(\fg)(k+h^\vee)/(2h^\vee)$]
\label{comp:km-kappa-first-principles}
\index{modular characteristic!first-principles derivation}
\index{curvature!affine Kac--Moody!explicit residue}

We derive the obstruction coefficient
$\kappa(\widehat{\fg}_k)$ by explicit residue computation
on the genus-$1$ bar complex, tracking the two OPE channels
separately.

\emph{Step~1: Genus-$1$ bar differential.}
On the elliptic curve $E_\tau$, the degree-$1$ bar element
is $J^a(z_1) \boxtimes J^b(z_2) \cdot K^{(1)}_{12}$,
where $K^{(1)}(z|\tau) = \theta_1'(z|\tau)/\theta_1(z|\tau)$
is the genus-$1$ propagator with
$A$-periodicity~\eqref{eq:A-cycle-periodicity} and
$B$-quasi-periodicity~\eqref{eq:B-cycle-quasi-periodicity}:
$K^{(1)}(z+\tau|\tau) = K^{(1)}(z|\tau) - 2\pi i$.
The fiberwise bar differential extracts the OPE via
the residue at $z_1 = z_2$:
\begin{equation}\label{eq:km-bar-diff-genus1-explicit}
\dfib(J^a \boxtimes J^b \cdot K^{(1)}_{12})
= \underbrace{k\,(J^a, J^b) \cdot |0\rangle}_{%
 \text{double pole}}
\;+\;
\underbrace{f^{ab}{}_{c}\, J^c}_{%
 \text{simple pole}}.
\end{equation}
The double-pole term $k(J^a,J^b)/(z-w)^2$ contributes
$k(J^a,J^b)$ as the coefficient of $1/(z-w)$ in the
product $[k(J^a,J^b)/(z-w)^2]\cdot K^{(1)}(z-w)$:
expanding $K^{(1)}(z) = 1/z + O(z)$ near $z=0$, the
residue of $k(J^a,J^b)/z^2 \cdot (1/z + O(z))$ is zero
at genus~$0$ (the $1/z^3$ term has no residue), but at
genus~$1$ the $B$-cycle monodromy produces the nonzero
contribution. The simple-pole term
$f^{ab}{}_c J^c(w)/(z-w)$ contributes $f^{ab}{}_c J^c$
directly by the standard residue.

\emph{Step~2: Squaring $\dfib$ and the $B$-cycle mechanism.}
The curvature $\dfib^{\,2}$ on a degree-$2$ element
involves applying $\dfib$ to the output
of~\eqref{eq:km-bar-diff-genus1-explicit} and using the
$B$-cycle shift $-2\pi i$
(equation~\eqref{eq:B-cycle-quasi-periodicity}).
Two contributions arise:

\emph{(a) Double-pole channel.}
The vacuum term $k(J^a, J^b)\cdot|0\rangle$ from the first
application of $\dfib$ is acted on by the $B$-cycle
monodromy. The transport around the $B$-cycle shifts the
propagator by $-2\pi i$, producing:
\begin{equation}\label{eq:km-dp-channel}
\dfib^{\,2}\big|_{\mathrm{dp}}(J^a \boxtimes J^b)
= (-2\pi i)\cdot k\,(J^a, J^b) \cdot \omega_1
\end{equation}
where $\omega_1 = c_1(\mathbb{E})$ represents the
$B$-cycle monodromy class. On the normalized trace
$\sum_a (J^a, J_a) = \dim(\fg)$, this gives:
\[
\operatorname{tr}\!\bigl(\dfib^{\,2}\big|_{\mathrm{dp}}\bigr)
= k\cdot\dim(\fg)\cdot\omega_1.
\]

\emph{(b) Simple-pole channel.}
The current-valued term $f^{ab}{}_c J^c$ from the first
application of $\dfib$ feeds back into a second OPE
contraction. This produces a \emph{one-loop self-contraction}:
the current $J^c$ is contracted with itself through a second
propagator loop. The resulting contribution is:
\begin{equation}\label{eq:km-sp-channel}
\dfib^{\,2}\big|_{\mathrm{sp}}(J^a \boxtimes J^b)
= (-2\pi i)\cdot
\tfrac{1}{2}\sum_{c,d} f^{ac}{}_{d}\, f^{bc}{}_{d}
\cdot \omega_1.
\end{equation}
The factor $\frac{1}{2}$ is the symmetry factor of the
one-loop self-contraction diagram ($|\mathrm{Aut}| = 2$).
The contraction $\sum_{c,d} f^{ac}{}_d f^{bc}{}_d$ is
the quadratic Casimir eigenvalue in the adjoint
representation. For any simple~$\fg$ with long roots
normalized to length~$2$:
\begin{equation}\label{eq:casimir-adjoint}
\sum_{c,d} f^{ac}{}_{d}\, f^{bc}{}_{d}
= C_2^{\mathrm{ad}} \cdot (J^a, J^b)
= 2h^\vee \cdot (J^a, J^b).
\end{equation}
\emph{Proof of~\eqref{eq:casimir-adjoint}.}
The sum $\sum_{c,d} f^{ac}{}_d f^{bc}{}_d$ is the
$(a,b)$-entry of the matrix
$\operatorname{ad}(T^c)\operatorname{ad}(T^c)$ summed
over a basis, i.e.\ the Killing form
$K(T^a, T^b) = \operatorname{tr}(
 \operatorname{ad}(T^a)\operatorname{ad}(T^b))$.
For simple~$\fg$, any invariant bilinear form is
proportional to the normalized form~$(,)$; the
proportionality constant is the eigenvalue of the
quadratic Casimir $C_2 = \sum_a T^a T_a$ in the adjoint
representation, which equals $2h^\vee$ when long roots
have $({\alpha},{\alpha}) = 2$
(\cite[Proposition~6.1]{Kac}). Hence
$K^{ab} = 2h^\vee\,\delta^{ab}$ in an orthonormal
basis.
Taking the trace:
\[
\operatorname{tr}\!\bigl(\dfib^{\,2}\big|_{\mathrm{sp}}\bigr)
= \tfrac{1}{2}\cdot 2h^\vee \cdot \dim(\fg) \cdot \omega_1
= h^\vee \cdot \dim(\fg) \cdot \omega_1.
\]

\emph{Step~3: Combining the channels.}
On a single pair $(J^a, J^b)$ with $(J^a, J^b) = 1$
(i.e., on the diagonal of the normalized form),
the two channels give:
\begin{equation}\label{eq:km-combined-per-component}
\dfib^{\,2}(J^a \boxtimes J^a \cdot K^{(1)}_{12})
= (-2\pi i)\cdot(k + h^\vee)\cdot\omega_1.
\end{equation}
To extract the scalar obstruction coefficient
$\kappa$ from this, we express the curvature as
a multiple of the Casimir element
$\mathsf{C} = \sum_a J^a \otimes J_a \in \bar{B}^0$.
By Proposition~\ref{prop:km-bar-curvature},
$m_0 = [(k + h^\vee)/(2h^\vee)] \cdot \mathsf{C}$,
where the factor $1/(2h^\vee)$ arises from
normalizing against the adjoint Casimir eigenvalue
$C_2^{\mathrm{ad}} = 2h^\vee$:
the per-component curvature $(k + h^\vee)$
equals $(k + h^\vee)/(2h^\vee)$ times the
eigenvalue~$2h^\vee$.

The scalar obstruction coefficient is the trace
of~$m_0$:
\begin{align}
\kappa
&= \operatorname{tr}(m_0)
= \frac{k + h^\vee}{2h^\vee}
 \cdot \operatorname{tr}(\mathsf{C})
= \frac{k + h^\vee}{2h^\vee}
 \cdot \sum_a (J^a, J_a)
\notag \\
&= \frac{k + h^\vee}{2h^\vee}\cdot\dim(\fg).
\label{eq:km-dfib-squared-trace}
\end{align}
Equivalently: the double-pole channel contributes
$\kappa_{\mathrm{dp}} = k\cdot\dim(\fg)/(2h^\vee)$
and the simple-pole channel contributes
$\kappa_{\mathrm{sp}} = \dim(\fg)/2$
(the self-contraction trace
$h^\vee\cdot\dim(\fg)/(2h^\vee) = \dim(\fg)/2$),
recovering~\eqref{eq:kappa-dp-sp}. Therefore:
\begin{equation}\label{eq:km-kappa-derivation}
\boxed{\;\kappa(\widehat{\fg}_k)
= \frac{(k + h^\vee)\cdot\dim(\fg)}{2h^\vee}.\;}
\end{equation}

\emph{Step~4: Consistency checks.}
\begin{enumerate}[label=(\roman*)]
\item \emph{$\mathfrak{sl}_2$}: $\dim = 3$, $h^\vee = 2$.
$\kappa = 3(k+2)/4$. At $k=1$: $\kappa = 9/4$.
\item \emph{$\mathfrak{sl}_3$}: $\dim = 8$, $h^\vee = 3$.
$\kappa = 4(k+3)/3$. Matches
Remark~\ref{rem:sl3-universality}.
\item \emph{$E_6$}: $\dim = 78$, $h^\vee = 12$.
$\kappa = 78(k+12)/24 = 13(k+12)/4$. At $k=1$: $\kappa = 169/4$.
\item \emph{$E_7$}: $\dim = 133$, $h^\vee = 18$.
$\kappa = 133(k+18)/36$. At $k=1$: $\kappa = 2527/36$.
\item \emph{$E_8$}: $\dim = 248$, $h^\vee = 30$.
$\kappa = 248(k+30)/60 = 62(k+30)/15$. At $k=1$: $\kappa = 1922/15$.
\item \emph{$B_2 = C_2 = \mathfrak{so}_5$}: $\dim = 10$, $h^\vee = 3$ (not $h = 4$).
$\kappa = 10(k+3)/6 = 5(k+3)/3$. At $k=1$: $\kappa = 20/3$.
\item \emph{$G_2$}: $\dim = 14$, $h^\vee = 4$ (not $h = 6$).
$\kappa = 14(k+4)/8 = 7(k+4)/4$. At $k=1$: $\kappa = 35/4$.
\item \emph{$F_4$}: $\dim = 52$, $h^\vee = 9$ (not $h = 12$).
$\kappa = 52(k+9)/18 = 26(k+9)/9$. At $k=1$: $\kappa = 260/9$.
\item \emph{Critical level}: $k = -h^\vee$ gives
$\kappa = 0$, the uncurved fixed point of Feigin--Frenkel.
\item \emph{Heisenberg limit}: Setting $f^{abc} = 0$,
$h^\vee = 0$ (abelian), the simple-pole channel vanishes
and $\kappa = k\cdot\dim(\fg)/(2\cdot 0)$ is ill-defined.
Instead, for a $\dim(\fg)$-dimensional abelian algebra at
level~$k$, only the double-pole channel contributes:
$\kappa_{\mathrm{dp}} = k\cdot\dim(\fg)$ (with no $h^\vee$
normalization needed), recovering
$\kappa(\mathcal{H}_k^{\oplus d}) = dk$
(Computation~\ref{comp:rank-d-heisenberg-genus1}).
\item \emph{Sugawara comparison.}
The Sugawara central charge is
$c = k\dim(\fg)/(k+h^\vee)$.
Setting $\kappa = c/2$ leads to
$(k+h^\vee)^2 = kh^\vee$, i.e.,
$k^2 + kh^\vee + (h^\vee)^2 = 0$,
which has no real solutions (discriminant
$-3(h^\vee)^2 < 0$). Hence
$\kappa \neq c/2$ for all real~$k$.
This is a genuine feature:
$\kappa$ is the modular characteristic of the
\emph{full} chiral algebra (tracking both
the double-pole and simple-pole channels),
while $c/2$ is the modular characteristic
of the Virasoro subalgebra alone (which
sees only the Sugawara stress tensor).
\item \emph{Anti-symmetry.}
$\kappa(\widehat{\fg}_{-k-2h^\vee})
= (-k - 2h^\vee + h^\vee)\dim(\fg)/(2h^\vee)
= -(k + h^\vee)\dim(\fg)/(2h^\vee) = -\kappa$,
confirming $\kappa + \kappa' = 0$.
\end{enumerate}
\end{computation}

\begin{proposition}[Feigin--Frenkel shear on channel
pair; \ClaimStatusConditional]
\label{prop:ff-channel-shear}
\index{Feigin--Frenkel duality!shear matrix}
\index{channel decomposition!Feigin--Frenkel transformation}
Under Feigin--Frenkel duality $k \mapsto k' = -k - 2h^\vee$:
\begin{equation}\label{eq:ff-shear}
\begin{pmatrix} \kappa_{\mathrm{dp}}' \\
\kappa_{\mathrm{sp}}' \end{pmatrix}
=
\begin{pmatrix} -1 & -2 \\ 0 & 1 \end{pmatrix}
\begin{pmatrix} \kappa_{\mathrm{dp}} \\
\kappa_{\mathrm{sp}} \end{pmatrix}.
\end{equation}
In particular, the total characteristic
$\kappa = \kappa_{\mathrm{dp}} + \kappa_{\mathrm{sp}}$
is a $(-1)$-eigenvector:
$\kappa' = -\kappa_{\mathrm{dp}} - 2\kappa_{\mathrm{sp}}
+ \kappa_{\mathrm{sp}} = -\kappa$,
recovering the anti-symmetry of
Theorem~\textup{\ref{thm:modular-characteristic}}(iv).
\end{proposition}

\begin{proof}
The simple-pole channel $\kappa_{\mathrm{sp}} =
\dim(\fg)/2$ depends only on the Lie algebra~$\fg$ (via the
structure constants $f^{ab}{}_c$), not on the level, so
$\kappa_{\mathrm{sp}}' = \kappa_{\mathrm{sp}}$.
For the double-pole channel:
\[
\kappa_{\mathrm{dp}}'
= \frac{(-k - 2h^\vee) \cdot \dim\fg}{2h^\vee}
= -\frac{k \cdot \dim\fg}{2h^\vee}
 - \frac{2h^\vee \cdot \dim\fg}{2h^\vee}
= -\kappa_{\mathrm{dp}} - \dim\fg
= -\kappa_{\mathrm{dp}} - 2\kappa_{\mathrm{sp}}.
\]
In matrix form this is~\eqref{eq:ff-shear}.
\end{proof}

\subsection{Exceptional types: shadow data for $E_6$, $E_7$, $E_8$}%
\label{sec:exceptional-shadow-data}
\index{E6@$E_6$!shadow data}
\index{E7@$E_7$!shadow data}
\index{E8@$E_8$!shadow data}
\index{exceptional Lie algebra!modular characteristic}
\index{shadow archetype!exceptional types}

The consistency checks of~\eqref{eq:km-kappa-derivation} prove
the universal formula $\kappa(\widehat{\fg}_k) = (k + h^\vee)\dim\fg/(2h^\vee)$
for exceptional types. The full shadow obstruction tower data.

\begin{proposition}[Exceptional shadow invariants; \ClaimStatusConditional]%
\label{prop:exceptional-shadow-invariants}
For the simply-laced exceptional Lie algebras
$\fg \in \{E_6, E_7, E_8\}$, the affine Kac--Moody
algebra $\widehat{\fg}_k$ at non-critical level $k \neq -h^\vee$
has the following invariants.
\begin{center}
\label{tab:exceptional-shadow-data}%
\renewcommand{\arraystretch}{1.3}
\begin{tabular}{lccccccc}
\toprule
$\fg$ & $\dim$ & $r$ & $h = h^\vee$ & $|\Delta_+|$
 & $\kappa(\widehat{\fg}_k)$ & $K$ & Class \\
\midrule
$E_6$ & $78$ & $6$ & $12$ & $36$
 & $13(k+12)/4$ & $156$ & L \\
$E_7$ & $133$ & $7$ & $18$ & $63$
 & $133(k+18)/36$ & $266$ & L \\
$E_8$ & $248$ & $8$ & $30$ & $120$
 & $62(k+30)/15$ & $496$ & L \\
\bottomrule
\end{tabular}
\end{center}
Here $K = c + c' = 2\dim\fg$ is the Koszul conductor.
Every entry satisfies $\kappa + \kappa' = 0$
(Proposition~\textup{\ref{prop:ff-channel-shear}}).
All three are class~L with shadow depth $r_{\max} = 3$:
the Jacobi identity kills the quartic obstruction
$S_4 = 0$ while the Lie bracket provides a nonzero
cubic shadow $\alpha \neq 0$
(Theorem~\textup{\ref{thm:shadow-archetype-classification}}).
\end{proposition}

\begin{proof}
The formula $\kappa = (k + h^\vee)\dim\fg/(2h^\vee)$ and
its anti-symmetry $\kappa + \kappa' = 0$ are proved in
Computation~\ref{comp:km-kappa-first-principles} and
Proposition~\ref{prop:ff-channel-shear} for all simple~$\fg$.
The shadow class~L assignment follows from the general
affine Kac--Moody analysis of
\S\ref{sec:affine-cubic-shadow}: the cubic shadow
$\alpha \neq 0$ (from the simple-pole channel) and
the quartic obstruction vanishes by the Jacobi identity
(Proposition~\ref{prop:affine-jacobi-quartic-vanishing}), giving
$\Delta = 8\kappa S_4 = 0$. The positive root counts
$|\Delta_+| = (\dim\fg - r)/2$ and the Koszul conductor
$K = 2\dim\fg$ are direct.

The simplified forms follow from
$\gcd(78, 24) = 6$ (giving $13(k+12)/4$ for $E_6$)
and $\gcd(248, 60) = 4$ (giving $62(k+30)/15$ for $E_8$).
For $E_7$, $\gcd(133, 36) = 1$ ($133 = 7 \cdot 19$, $36 = 2^2 \cdot 3^2$),
so the fraction $133(k+18)/36$ is already in lowest terms.
\end{proof}

\begin{remark}[Exponents and Lie-theoretic identities]%
\label{rem:exceptional-exponents}
\index{exponents!exceptional Lie algebra}
\index{Freudenthal--de Vries formula}
The exponents $m_1 < \cdots < m_r$ of each exceptional algebra
satisfy the standard identities:
\begin{center}
\renewcommand{\arraystretch}{1.3}
\begin{tabular}{lcccl}
\toprule
$\fg$ & Exponents & $\sum m_i$ & $|\rho_W|^2$ & $\dim = r(h+1)$ \\
\midrule
$E_6$ & $1,\, 4,\, 5,\, 7,\, 8,\, 11$ & $36$ & $78$ & $78 = 6 \cdot 13$ \\
$E_7$ & $1,\, 5,\, 7,\, 9,\, 11,\, 13,\, 17$ & $63$ & $399/2$ & $133 = 7 \cdot 19$ \\
$E_8$ & $1,\, 7,\, 11,\, 13,\, 17,\, 19,\, 23,\, 29$ & $120$ & $620$ & $248 = 8 \cdot 31$ \\
\bottomrule
\end{tabular}
\end{center}
The column $\sum m_i = rh/2$ is a standard identity
for all simple Lie algebras. The Freudenthal--de Vries
(Strange) formula gives $|\rho_W|^2 = \dim(\fg)\cdot h/12$.
The exponents satisfy
$m_i + m_{r+1-i} = h$ (Coxeter duality),
as shown in the displayed lists. In particular, $E_7$
has the unique middle exponent $m_4 = 9 = h/2$.

The identity $\dim\fg = r \cdot (h+1)$ is equivalent to
$|\Delta| = rh$ (the total number of roots equals the rank
times the Coxeter number). At level $k = 1$, this gives
the central charge $c = k\dim\fg/(k + h^\vee)
= \dim\fg/(h+1) = r$: the level-one central charge of
each exceptional affine algebra equals its rank.
\end{remark}

\begin{remark}[Anomaly ratio and $\mathcal{W}$-algebra convergence]%
\label{rem:exceptional-anomaly-ratio}
\index{anomaly ratio!exceptional types}
The anomaly ratio
$\varrho(\fg) = \sum_{i=1}^r 1/(m_i + 1)$
for the associated principal $\mathcal{W}$-algebra
(Remark~\textup{\ref{rem:general-w-kappa-values}})
governs the modular characteristic via
$\kappa(\mathcal{W}^k(\fg)) = \varrho(\fg) \cdot c(\mathcal{W}^k(\fg))$:
\begin{center}
\renewcommand{\arraystretch}{1.3}
\begin{tabular}{lll}
\toprule
$\fg$ & $\varrho(\fg)$ & Decimal \\
\midrule
$E_6$ & $\displaystyle\frac{1}{2} + \frac{1}{5} + \frac{1}{6}
 + \frac{1}{8} + \frac{1}{9} + \frac{1}{12}
 = \frac{427}{360}$ & $1.186$ \\[6pt]
$E_7$ & $\displaystyle\frac{1}{2} + \frac{1}{6} + \frac{1}{8}
 + \frac{1}{10} + \frac{1}{12} + \frac{1}{14} + \frac{1}{18}
 = \frac{2777}{2520}$ & $1.102$ \\[6pt]
$E_8$ & $\displaystyle\frac{1}{2} + \frac{1}{8} + \frac{1}{12}
 + \frac{1}{14} + \frac{1}{18} + \frac{1}{20} + \frac{1}{24}
 + \frac{1}{30}
 = \frac{121}{126}$ & $0.960$ \\
\bottomrule
\end{tabular}
\end{center}
The ordering $\varrho(E_6) > \varrho(E_7) > \varrho(E_8)$
reflects the increasing sparsity of the exponents: $E_8$
has large exponents ($29, 23, \ldots$) giving small
$1/(m_i + 1)$ contributions. Since all three are class~L
(the affine KM shadow obstruction tower terminates at degree~3 regardless
of $\varrho$), the anomaly ratio controls the
$\mathcal{W}$-algebra shadow obstruction tower, not the affine one.

The distinction $\varrho(E_8) < 1 < \varrho(E_6)$ is
significant for the $\mathcal{W}$-algebra genus expansion:
$\varrho < 1$ implies the principal $\mathcal{W}(E_8)$
shadow obstruction tower converges absolutely at all non-critical
levels, while $\varrho > 1$ for $\mathcal{W}(E_6)$
and $\mathcal{W}(E_7)$ indicates divergence of the shadow
tower at generic levels (convergent only for
$|c|$ sufficiently large relative to~$\varrho$).
\end{remark}

\begin{remark}[The $r$-matrix and the Koszul conductor]%
\label{rem:exceptional-r-matrix}
\index{r-matrix!exceptional types}
\index{Koszul conductor!exceptional types}
The affine $r$-matrix data for $\widehat{\fg}_k$ has two
normalizations (Remark~\ref{rem:km-collision-residue-rmatrix}):
the trace-form collision residue is
$r^{\mathrm{coll}}(z)=k\Omega_{\mathrm{tr}}/z$, and the KZ coefficient
is
$r^{\mathrm{KZ}}(z)=\Omega_{\fg}/((k+h^\vee)z)$, where
$\Omega_\fg = \sum_a J^a \otimes J_a$ is the Casimir element. Both
have pole order~$1$ (not~$2$); the $d\log$ extraction in the bar
construction absorbs one power of $(z-w)$ from the OPE.

The Koszul conductor $K(\fg) = 2\dim\fg$ takes the values
$K(E_6) = 156$, $K(E_7) = 266$, $K(E_8) = 496$. The
value $K(E_8) = 496$ equals $\dim(\fg_{E_8 \times E_8})$,
the dimension of the $E_8 \times E_8$ gauge group in
heterotic string theory; this numerical coincidence reflects
the level-one anomaly cancellation between the matter
sector ($c = 8$) and its Koszul dual ($c' = 488$), whose
sum $c + c' = K = 496$ saturates the heterotic central
charge budget.
\end{remark}

\subsection{The screening charge perspective}

\begin{definition}[Screening charges]
\index{screening operator|textbf}
\index{screening operator!Kac--Moody}
For $\widehat{\mathfrak{g}}_{-h^\vee}$ at critical level, the \emph{screening charges} are:
\begin{equation}
S_\alpha = \oint e^{\alpha(\phi)} \prod_{\beta > 0} \gamma_\beta^{n_{\alpha,\beta}}, \quad \alpha \in \Delta_+
\end{equation}
where $\phi$ denotes the Cartan bosons, $\gamma_\beta$ are the $\gamma$ fields in the Wakimoto module, and $n_{\alpha,\beta} \in \mathbb{Z}_{\geq 0}$ are the structure coefficients from the nilpotent subalgebra.
\end{definition}

\begin{theorem}[Wakimoto screening expression for the current-valued critical
bar differential; \ClaimStatusConditional]\label{thm:screening-bar}
At critical level $k = -h^\vee$, after the shifted curvature class
vanishes, assume the completed Wakimoto bar--BRST comparison of
Theorem~\ref{thm:wakimoto-koszul}.  Then the current-valued simple-pole
part of the bar complex differential is represented in the Wakimoto
model by the screening charges:
\begin{equation}
d_{\bar{B}} = \sum_{\alpha \in \Delta_+} S_\alpha \otimes \eta_\alpha
\end{equation}
where $\eta_\alpha = d\log(z - w_\alpha)$ is the propagator form associated to the screening vertex at $w_\alpha$.
\end{theorem}

\begin{proof}
\emph{Step~1: Setup.}
At critical level, $d_{\bar{B}}^2 = 0$ (Theorem~\ref{thm:universal-kac-moody-koszul}, Step~2), so $\bar{B}(\widehat{\mathfrak{g}}_{-h^\vee})$ is a genuine dg coalgebra. The bar differential on a degree-1 element $[J^a] \in \bar{B}^1$ is:
\[
d_{\bar{B}}([J^a]) = \mathrm{Res}_{z=w}\left[Y(J^a,z) \otimes \eta_{12}\right],
\]
where the residue is taken in $\overline{C}_2(X)$ and $\eta_{12} = d\log(z-w)$.

\emph{Step~2: OPE decomposition at critical level.}
The current OPE is:
\[
J^a(z)\,J^b(w) \sim \frac{k\,(J^a,J^b)}{(z-w)^2} + \frac{f^{ab}{}_c\,J^c(w)}{z-w}.
\]
At critical level the trace-form double-pole coefficient is
\(-h^\vee(J^a,J^b)\), not zero. What vanishes is the shifted curvature
class
\(m_0=(k+h^\vee)\mathsf C_{\fg}/(2h^\vee)\)
(cf.\ Theorem~\ref{thm:kac-moody-ainfty}). After this curvature
cancellation, the current-valued part of the critical bar differential
is the simple-pole Lie-bracket term:
\[
d_{\bar{B}}([J^a \mid J^b]) = f^{ab}{}_c\,[J^c].
\]

\emph{Step~3: Wakimoto representation of the simple-pole term.}
The Wakimoto free field realization \cite{Wakimoto86} identifies $J^a(z)$ with explicit expressions in the Cartan bosons $\partial\phi^i$ and the $\beta_\alpha\gamma^\alpha$ system. In the Wakimoto module $\mathcal{M}_{\mathrm{Wak}}$:
\begin{align*}
J^a(z) &= \partial\phi^a(z) \\
 &\quad + \sum_{\alpha \in \Delta_+}
 :\!\beta_\alpha(z)\,[\rho^\alpha(e_a)\,\gamma^\alpha](z)\!:
 + \cdots
\end{align*}
where $\rho^\alpha$ denotes the action of $e_a$ in the adjoint
representation. The single-pole residue
\[
\mathrm{Res}_{z=w}\bigl[J^a(z)J^b(w)/(z-w)\bigr]
\]
comes entirely from the $\beta\gamma$-contractions; the Cartan boson
$\partial\phi^a(z)\partial\phi^b(w)$ contributes only a double pole.
Each such contraction produces a factor $\beta_\alpha(w)$ times an
operator in~$\gamma^\alpha$: this is the screening term
\[
S_\alpha(w) = e^{\alpha(\phi(w))}\prod_\beta \gamma^\beta(w)^{n_{\alpha,\beta}}
\]
(cf.\ the definition above).

\emph{Step~4: Identification under the Wakimoto comparison.}
By root space decomposition, the structure constants $f^{ab}{}_c$ are nonzero only when $\alpha_a + \alpha_b = \alpha_c$ for positive roots. The contribution of each such positive root $\alpha$ to the simple-pole residue is a contraction mediated by the $\beta_\alpha\gamma^\alpha$-pair, with the propagator form $\eta_\alpha = d\log(z - w_\alpha)$ recording the collision position. Summing over all $\alpha \in \Delta_+$ reproduces $\sum_\alpha S_\alpha \otimes \eta_\alpha$.

At degree $n$, the bar differential $d_{\bar{B}}$ on $[J^{a_1} \mid \cdots \mid J^{a_n}]$ is a sum of pairwise collision residues; each collision at $z_i = z_j$ with $\alpha_{a_i} + \alpha_{a_j} = \alpha_c$ contributes $S_{\alpha_c}(z_j)\otimes\eta_{ij}$, establishing the formula in full generality.
\end{proof}

\section{Connection to $\mathcal{W}$-algebras}

\subsection{Drinfeld--Sokolov reduction}

\begin{definition}[DS reduction \cite{DS, FF}]
The W-algebra $\mathcal{W}^k(\mathfrak{g}, f)$ associated to a nilpotent element $f \in \mathfrak{g}$ is the BRST cohomology:
\begin{equation}
\mathcal{W}^k(\mathfrak{g}, f) = H^0_{Q_{\mathrm{DS}}}(\widehat{\mathfrak{g}}_k \otimes \mathcal{F}_{\mathrm{gh}})
\end{equation}
where $\mathcal{F}_{\mathrm{gh}}$ is the ghost $bc$-system for the nilpotent subalgebra $\mathfrak{n}_f = \bigoplus_{\alpha(h_f)>0} \mathfrak{g}_\alpha$, and $Q_{\mathrm{DS}}$ is the BRST differential implementing constraints from~$f$. The cohomology is concentrated in degree~$0$ for the principal nilpotent.
\end{definition}

\begin{remark}[DS changes the shadow problem]
\label{rem:ds-changes-shadow-problem}
\index{Drinfeld--Sokolov reduction!shadow-depth change}
\index{W-algebra@$\mathcal{W}$-algebra!not affine class L}
Principal Drinfeld--Sokolov reduction is a BRST functor applied to
\(\widehat{\fg}_k\otimes\mathcal F_{\mathrm{gh}}\), not an
identification of the affine Koszul dual with a
\(\mathcal W\)-algebra.  At the affine level the Lie bracket gives
class~\(L\), \(r_{\max}=3\), and
\(\kappa(V_k(\fg))+\kappa(V_{-k-2h^\vee}(\fg))=0\).  After
principal DS reduction the ghost sector creates the higher
\(\mathcal W\)-generators and the quartic/higher shadow channels:
Theorem~\ref{thm:ds-shadow-depth-increase} sends the affine
class~\(L\) input to the class~\(M\) \(\mathcal W\)-output.  The
proved noncritical statement for principal nilpotent is therefore the
characteristic transport
\[
\kappa\!\left(\mathcal W^k(\fg,f_{\mathrm{prin}})^!\right)
=
\kappa\!\left(\mathcal W^{-k-2h^\vee}
(\fg,f_{\mathrm{prin}})\right)
\]
of Theorem~\ref{thm:w-algebra-koszul-main}.  The stronger
same-family identification of the Koszul dual with the companion
\(\mathcal W\)-algebra is exactly the separate DS/bar transport
hypothesis recorded in Theorem~\ref{thm:w-koszul-precise}.
\end{remark}

\begin{theorem}[Critical fixed point for principal \texorpdfstring{$\mathcal{W}$}{W}-algebras; \ClaimStatusConditional]\label{thm:w-algebra-koszul}
Let $\mathfrak{g}$ be a simple Lie algebra, let $f = f_{\mathrm{prin}}$
be the principal nilpotent element, and set $k=-h^\vee$. Then the
principal critical $\mathcal{W}$-algebra sits at the fixed point of the
Feigin--Frenkel involution and bar-cobar inversion recovers the same
algebra:
\begin{equation}
\Omega^{\mathrm{ch}}(\bar{B}^{\mathrm{geom}}(\mathcal{W}^{-h^\vee}(\mathfrak{g}, f_{\mathrm{prin}})))
\simeq \mathcal{W}^{-h^\vee}(\mathfrak{g}, f_{\mathrm{prin}})
\end{equation}
and the dual-level formula gives $k'=-k-2h^\vee=-h^\vee$.
\end{theorem}

\begin{proof}
This is Theorem~\ref{thm:w-koszul-precise}\textup{(A)} restated in the
Kac--Moody chapter. In particular, we use only the proved critical
fixed-point package: the bar complex is uncurved, bar-cobar inversion
recovers the same principal $\mathcal{W}$-algebra, and the
Feigin--Frenkel involution fixes the critical level. \qedhere
\end{proof}

\begin{remark}[Extension to general nilpotents]
\label{rem:w-koszul-general-nilpotent}
The fixed-point inversion statement above is stated for
$f_{\mathrm{prin}}$, where DS cohomology concentrates in degree~$0$
\cite{FF}. For general $f$, DS cohomology may be non-trivial in other
degrees, and the stronger Langlands-dual critical statement
\[
\mathcal{W}^{-h^\vee}(\mathfrak{g}, f)^! \simeq \mathcal{W}^{-h^\vee}(\mathfrak{g}^\vee, f^\vee)
\]
for arbitrary $f$ (Barbasch--Vogan dual $f^\vee$) is conjectural
(\ClaimStatusConjectured): it separates into three non-principal packets
(dual-orbit input, orbit-indexed level shift, paired DS seed
transport). See Theorem~\ref{thm:w-algebra-koszul-main} for the proved
noncritical principal duality.
\end{remark}

\begin{remark}[Langlands duality]
The critical fixed-point package of
Theorem~\ref{thm:w-algebra-koszul} is compatible with the external
geometric/quantum-Langlands picture, but it does not by itself prove
that the bar-cobar adjunction produces the Langlands-dual principal
$\mathcal{W}$-algebra at critical level. That stronger interpretation
remains an external or conjectural dictionary in the manuscript, and it
is precisely what motivates the project-level conjectural direction:
the chiral bar complex could provide a constructive approach to the
geometric Langlands correspondence, building the dual D-modules on
$\mathrm{Bun}_G$ from configuration space data on the curve (see
\S\ref{sec:km-conformal-blocks} for the conformal block interpretation).
\end{remark}

\subsection{Principal $\mathcal{W}$-algebra example}

\begin{theorem}[Principal \texorpdfstring{$\mathcal{W}$}{W}-algebra structure \cite{FF, Ara07}; \ClaimStatusProvedElsewhere]
\label{thm:principal-w-algebra-structure}
\begin{equation}
\mathcal{W}^{-h^\vee}(\mathfrak{g}, f_{\mathrm{prin}}) = \text{Universal enveloping of } \{\text{generators of spin } d_1+1, \ldots, d_r+1\}
\end{equation}
where $d_1, \ldots, d_r$ are the exponents of $\mathfrak{g}$.

For $\mathfrak{sl}_2$: exponents $\{1\}$, so we get Virasoro with generator $T$ of spin $2$.

For $\mathfrak{sl}_3$: exponents $\{1,2\}$, so we get $\mathcal{W}_3$ with generators $T$ (spin $2$) and $W$ (spin $3$).
\end{theorem}

\section{Higher operations and quantum corrections}

\subsection{\texorpdfstring{$A_\infty$ structure}{A- structure}}

\begin{theorem}[\texorpdfstring{$A_\infty$}{A-infinity} operations on Kac--Moody; \ClaimStatusConditional]\label{thm:kac-moody-ainfty}
\index{$A_\infty$-structure!Kac--Moody}
The chiral algebra $\widehat{\mathfrak{g}}_k$ carries a canonical curved $A_\infty$ structure $(m_0, m_1, m_2, m_3, \ldots)$:
\begin{align}
m_0 &\in \widehat{\mathfrak{g}}_k \quad \text{(curvature, vanishes iff } k = -h^\vee\text{)} \\
m_2 &: \widehat{\mathfrak{g}}_k \otimes \widehat{\mathfrak{g}}_k \to \widehat{\mathfrak{g}}_k \quad \text{(chiral multiplication)} \\
m_n &: \widehat{\mathfrak{g}}_k^{\otimes n} \to \widehat{\mathfrak{g}}_k \quad \text{(higher coherences, } n \geq 3\text{)}
\end{align}

These are computed geometrically by:
\begin{equation}
m_n(\omega_1, \ldots, \omega_n) = \int_{\overline{C}_n(X)} \omega_1 \wedge \cdots \wedge \omega_n \cdot \Phi_n
\end{equation}
where $\Phi_n$ is the propagator form (fundamental chain) on $\overline{C}_n(X)$.
\end{theorem}

\begin{proof}
The $A_\infty$ structure is obtained by homotopy transfer from the bar-cobar adjunction. By the general theory (Theorem~\ref{thm:bar-cobar-inversion-qi} and the Homotopy Transfer Theorem, Appendix~\ref{app:homotopy-transfer}), for a chiral Koszul algebra $\cA$ with bar complex $\bar{B}(\cA)$, the cobar construction $\Omega(\bar{B}(\cA))$ carries a canonical dg algebra structure, and the quasi-isomorphism $\Omega(\bar{B}(\cA)) \xrightarrow{\sim} \cA$ (Theorem~\ref{thm:bar-cobar-isomorphism-main}) transfers an $A_\infty$ structure to~$\cA$.

For $\cA = \widehat{\mathfrak{g}}_k$, the geometric realization of the bar complex on $\overline{C}_{n+1}(X)$ (Theorem~\ref{thm:geometric-equals-operadic-bar}) identifies the transferred operations with configuration space integrals: $m_n$ is computed by integrating the product of input forms against the propagator $\Phi_n$ over the compactified configuration space. The $A_\infty$ relations $\sum_{r+s+t=n}(-1)^{rs+t} m_{r+1+t}(\mathrm{id}^{\otimes r} \otimes m_s \otimes \mathrm{id}^{\otimes t}) = 0$ follow from the Stokes theorem on $\overline{C}_n(X)$: the boundary strata of the FM compactification encode exactly the compositions of lower operations.

The curvature
$m_0 = \frac{k+h^\vee}{2h^\vee}\mathsf C_{\fg}$
(identified in the proof of
Theorem~\ref{thm:universal-kac-moody-koszul}) satisfies
$m_1^2(a) = m_2(m_0, a) - m_2(a, m_0) = [m_0, a]$
for all $a \in \widehat{\mathfrak{g}}_k$, which is the curved
$A_\infty$ relation replacing strict nilpotence $\dzero^2 = 0$.
\end{proof}

\subsection{Quantum corrections from higher genus}

\begin{remark}[Genus expansion]
The bar complex has contributions from all genera:
\begin{equation}
\bar{B}(\widehat{\mathfrak{g}}_k) = \bigoplus_{g \geq 0} \bar{B}^{(g)}(\widehat{\mathfrak{g}}_k)
\end{equation}
where $\bar{B}^{(g)}$ uses configuration spaces on genus $g$ curves.

The level $k$ controls the genus expansion:
\begin{equation}
Z(\widehat{\mathfrak{g}}_k) = \sum_{g=0}^\infty \frac{1}{(k+h^\vee)^{2g-2}} Z_g
\end{equation}
\end{remark}

\begin{theorem}[Higher genus corrections to Koszul duality; \ClaimStatusProvedHere]\label{thm:km-higher-genus-corrections}
The genus-graded bar complex of $\widehat{\mathfrak{g}}_k$ decomposes as:
\begin{equation}
\bar{B}(\widehat{\mathfrak{g}}_k) = \bigoplus_{g \geq 0} \bar{B}^{(g)}(\widehat{\mathfrak{g}}_k)
\end{equation}
where $\bar{B}^{(g)}$ is computed on genus-$g$ curves. The genus-$0$
contribution recovers the level-shifting duality
$(\widehat{\mathfrak{g}}_k)^! \simeq
\mathrm{CE}^{\mathrm{ch}}(\widehat{\mathfrak{g}}_{-k-2h^\vee})$,
and the higher-genus contributions are controlled by modular forms:
\begin{equation}
\bar{B}^{(g)}(\widehat{\mathfrak{g}}_k) \in H^*(\overline{\mathcal{M}}_g) \otimes \mathrm{End}(\widehat{\mathfrak{g}}_k) \cdot (k+h^\vee)^{2-2g}
\end{equation}
At critical level $k = -h^\vee$, the expansion parameter $(k+h^\vee)^{-1}$ becomes singular; the actual mechanism producing the infinite-dimensional Feigin--Frenkel center at $k = -h^\vee$ is non-perturbative and arises from the structure of the vertex algebra itself, not from divergent perturbation theory.
\end{theorem}

\begin{proof}
The decomposition follows from the modular operad structure proved in Theorem~\ref{thm:prism-higher-genus}: the bar complex of any augmented chiral algebra decomposes by the genus of the underlying curve, with the genus-$g$ contribution computed via the Feynman transform of the modular operad $\{\overline{\mathcal{M}}_{g,n}\}$.

For $\cA = \widehat{\mathfrak{g}}_k$, the genus-0 bar complex
$\bar{B}^{(0)}$ is the configuration space bar complex on
$\mathbb{P}^1$, which computes the Koszul dual
\(\mathrm{CE}^{\mathrm{ch}}(\widehat{\mathfrak{g}}_{-k-2h^\vee})\)
(Theorem~\ref{thm:universal-kac-moody-koszul}). The genus-$g$
contribution $\bar{B}^{(g)}$ involves integration over
$\overline{\mathcal{M}}_{g,n}$, producing cohomology classes in
$H^*(\overline{\mathcal{M}}_g)$. The power $(k+h^\vee)^{2-2g}$ arises
because each handle contributes a factor of the Casimir eigenvalue
$(k+h^\vee)^{-1}$ from the propagator, and the Euler characteristic of
a genus-$g$ surface with $n$ punctures gives exponent $2-2g-n$ (with
the $n$-dependence absorbed into the $\mathrm{End}$ factor). This is
the standard string theory genus expansion applied to the chiral algebra
setting (cf.\ Theorem~\ref{thm:genus-induction-strict}).
\end{proof}

\section{Computational algorithms}

\subsection{Algorithm for computing Koszul dual}

\begin{construction}[Computing the Koszul dual of $\widehat{\mathfrak{g}}_k$]\label{con:compute-km-dual}
Given a simple Lie algebra $\mathfrak{g}$, level $k$, and truncation degree $N$, the Koszul dual $(\widehat{\mathfrak{g}}_k)^!$ through degree $N$ is computed as follows.

\emph{Step~1: Bar complex.}
For $0 \le n \le N$, construct
$\bar{B}^n = \Gamma(\overline{C}_{n+1}(X),\, \mathfrak{g}^{\boxtimes(n+1)} \otimes \mathcal{L}_k \otimes \Omega^n_{\log})$
and choose a basis using decorated trees.

\emph{Step~2: Differentials.}
For each basis element $\omega \in \bar{B}^n$, compute $d(\omega) = \sum_{i<j} \mathrm{Res}_{z_i=z_j}[\omega]$ by residue calculus, obtaining the matrix of $d^n: \bar{B}^n \to \bar{B}^{n+1}$.

\emph{Step~3: Curvature.}
Set
\[
m_0 = \frac{k+h^\vee}{2h^\vee}\cdot(\text{Casimir})
\]
in the normalization of
\(\kappa(\widehat{\fg}_k)=\dim(\fg)(k+h^\vee)/(2h^\vee)\), and verify
the curved $A_\infty$ relation $m_1^2(a) = [m_0, a]$ for all
generators $a$.

\emph{Step~4: Dual cobar.}
Take the completed Verdier-linear dual
\(\bar{B}^n \mapsto (\bar{B}^n)^{\vee_{\mathrm V}}\), reverse the
grading, twist the differential by curvature, and form
$(\widehat{\mathfrak{g}}_k)^!
= \Omega(\bar{B}(\widehat{\mathfrak{g}}_k)^{\vee_{\mathrm V}})$.
The undualized cobar
\(\Omega(\bar{B}(\widehat{\mathfrak{g}}_k))\) is the inversion object.

\emph{Step~5: Extraction.}
The generators are $H^1((\widehat{\mathfrak{g}}_k)^!)$, the relations
are $\operatorname{Im}(d^2) \subset (\bar{B}^2)^\vee$. The OPE of
the dual generators $J^{a*}$ is determined by the residue pairing:
$J^{a*}_{(0)} J^{b*} = f^{ab}_c J^{c*}$ (inherited from the
structure constants of~$\mathfrak{g}$, since the Lie bracket is
level-independent) and $J^{a*}_{(1)} J^{b*} = (-k - 2h^\vee)
\delta^{ab}$ (from the dual curvature $m_0^! = -(k + h^\vee)
\cdot C_2 \mapsto (-k - 2h^\vee)$ after the invariant-form
contribution; see Theorem~\ref{thm:closed-form-ope} below).
\end{construction}

\subsection{Explicit formulas}

\begin{theorem}[Closed-form current presentation in the Koszul dual;
\ClaimStatusProvedHere]
\label{thm:closed-form-ope}
The current presentation underlying the chiral CE algebra
$(\widehat{\mathfrak{g}}_k)^!$ has OPE:
\begin{equation}
J^{a*}(z) J^{b*}(w) \sim \frac{(-k-2h^\vee) \delta^{ab}}{(z-w)^2} + \frac{f^{ab}_c J^{c*}(w)}{z-w}
\end{equation}
where $J^{a*}$ are the dual generators.

This is computed from the residue pairing:
\begin{equation}
\langle J^a, J^{b*} \rangle = \int_{X^2} J^a(z) \wedge J^{b*}(w) \cdot d\log(z-w) = \delta^{ab}
\end{equation}
\end{theorem}
\begin{proof}
The residue differential on two current generators is
\[
d(J^a\boxtimes J^b\eta_{12})
= k\,(J^a,J^b)|0\rangle + f^{ab}{}_cJ^c .
\]
After completed Verdier duality the determinant contribution of the
collision divisor contributes the adjoint Casimir
\(2h^\vee\), as in Lemma~\ref{lem:killing-structure-constants}.
Thus the double-pole level changes from \(k\) to
\(-k-2h^\vee\).  The simple-pole bracket is dual to the coalgebra
structure map and keeps the same constants \(f^{ab}{}_c\).  This gives
the displayed current presentation.  Applying the chiral CE functor to
that level-reflected current presentation gives the stated Koszul dual.
\end{proof}

\section{Applications and extensions}

\subsection{Holographic duality}

\begin{conjecture}[Kac--Moody in holography; \ClaimStatusConjectured]\label{conj:km-holography}
The level-shifting Koszul duality gives a conjectural algebraic model
for the boundary/bulk exchange in $\mathrm{AdS}_3/\mathrm{CFT}_2$:
\begin{center}
\begin{tikzcd}
\text{Boundary CFT at level } k \arrow[d, "\text{Koszul dual}"'] \arrow[r, "\text{WZW model}"] & \text{Open strings on D-branes} \arrow[d, "\text{open-closed}"] \\
\text{Dual CFT at level } -k-2h^\vee \arrow[r, "\text{Chern--Simons}"] & \text{Closed strings in bulk}
\end{tikzcd}
\end{center}

The conjectural dictionary matches the boundary WZW model at level $k$
with $\widehat{\mathfrak{g}}_k$, matches a suitable bulk Chern--Simons
package at level $-k-2h^\vee$ with the Koszul-dual shadow
$(\widehat{\mathfrak{g}}_k)^!$, and treats the holographic comparison
itself as modeled by the bar-cobar adjunction between boundary and
bulk.
\end{conjecture}

\begin{remark}[Scope]
The mathematical content is
Theorem~\ref{thm:universal-kac-moody-koszul}:
\((\widehat{\mathfrak{g}}_k)^! \simeq
\mathrm{CE}^{\mathrm{ch}}(\widehat{\mathfrak{g}}_{-k-2h^\vee})\).
The theorem supplies a level-reflected chiral CE object and the
bar-cobar inversion surface.  It does not prove a boundary/bulk
duality; the holographic comparison is an external physics
dictionary whose mathematical input is exactly the datum displayed
below.

In the holographic modular Koszul datum
$\mathcal{H}(T) = (\cA, \cA^!, C, r(z), \Theta_\cA,
\nabla^{\mathrm{hol}})$,
the Kac--Moody specialisation is:
$\cA = \widehat{\mathfrak{g}}_k$,
$\cA^! =
\mathrm{CE}^{\mathrm{ch}}(\widehat{\mathfrak{g}}_{-k-2h^\vee})$,
$r^{\mathrm{coll}}(z)=k\Omega_{\mathrm{tr}}/z$ for the collision
residue, and
$r^{\mathrm{KZ}}(z) = \Omega_{\mathrm{KZ}}/((k{+}h^\vee)\,z)$ for
the KZ equation,
$\kappa(\widehat{\mathfrak{g}}_k) = (k+h^\vee)\dim\mathfrak{g}/(2h^\vee)$,
$\Theta_{\widehat{\mathfrak{g}}_k}^{\min}
= \kappa \cdot \eta \otimes \Lambda$
\textup{(}minimal scalar package for uniform-weight algebras\textup{)},
while the full universal class is Lie/tree with
$\mathfrak{C} \ne 0$, $o_4 = 0$.
\end{remark}

\begin{computation}[Collision residue and the KZ $r$-matrix for $\widehat{\mathfrak{sl}}_{2,k}$; \ClaimStatusConditional]
\label{comp:sl2-collision-residue-kz}
\index{collision residue!sl2 explicit@$\mathfrak{sl}_2$ explicit}
\index{r-matrix@$r$-matrix!from collision residue}
\index{KZ connection!from collision residue}
\index{Arnold relation!CYBE computation}
We compute the genus-$0$, degree-$2$ collision residue
$\operatorname{Res}^{\mathrm{coll}}_{0,2}(\Theta_\cA)$
for $\cA = \widehat{\mathfrak{sl}}_{2,k}$ and keep separate the
trace-form collision residue
$r^{\mathrm{coll}}(z)=k\Omega_{\mathrm{tr}}/z$ from the KZ coefficient
$r^{\mathrm{KZ}}(z)=\Omega_{\mathfrak{sl}_2}/((k+2)z)$ governing the
Knizhnik--Zamolodchikov equation.

\emph{Step~1: OPE data.}
Let $\{e, h, f\}$ be the standard basis of $\mathfrak{sl}_2$
with normalized invariant form satisfying $\kappa_{\fg}(e,f) = 1$,
$\kappa_{\fg}(h,h) = 2$, and $\kappa_{\fg}(e,e) = \kappa_{\fg}(f,f) = \kappa_{\fg}(e,h)
= \kappa_{\fg}(f,h) = 0$. The structure constants are
$f^{eh}{}_e = 2$, $f^{hf}{}_f = 2$, $f^{ef}{}_h = 1$
(and antisymmetric permutations).
The OPE of $\widehat{\mathfrak{sl}}_{2,k}$ is
\begin{equation}\label{eq:sl2-collision-ope}
J^a(z)\, J^b(w) \;\sim\;
\frac{k \cdot \kappa_{\fg}(a,b)}{(z-w)^2}
\;+\;
\frac{f^{ab}{}_{c}\, J^c(w)}{z-w},
\end{equation}
giving the nine explicit OPE channels:
\begin{alignat}{2}
e(z)\,f(w) &\sim \frac{k}{(z-w)^2} + \frac{h(w)}{z-w},
 &\qquad
f(z)\,e(w) &\sim \frac{k}{(z-w)^2} - \frac{h(w)}{z-w},
 \label{eq:sl2-collision-ef} \\
h(z)\,h(w) &\sim \frac{2k}{(z-w)^2},
 &\qquad
e(z)\,e(w) &\sim 0,
 \label{eq:sl2-collision-hh} \\
e(z)\,h(w) &\sim \frac{-2e(w)}{z-w},
 &\qquad
h(z)\,e(w) &\sim \frac{2e(w)}{z-w},
 \label{eq:sl2-collision-eh} \\
h(z)\,f(w) &\sim \frac{-2f(w)}{z-w},
 &\qquad
f(z)\,h(w) &\sim \frac{2f(w)}{z-w},
 \label{eq:sl2-collision-hf} \\
f(z)\,f(w) &\sim 0. &&
 \label{eq:sl2-collision-ff}
\end{alignat}

\emph{Step~2: Bar complex on $\overline{C}_2(\bC)$.}
At degree~$1$, the geometric bar complex
(cf.\ Theorem~\ref{thm:sl2-koszul-dual}) is
\[
\bar{B}^1(\widehat{\mathfrak{sl}}_{2,k})
\;=\;
\mathfrak{sl}_2 \otimes \mathfrak{sl}_2 \otimes
\Omega^1_{\log}(\overline{C}_2),
\]
with basis elements
\[
\xi_{ab} \;=\; J^a(z_1) \otimes J^b(z_2)
\cdot \eta_{12},
\qquad
\eta_{12} = d\log(z_1 - z_2) = \frac{dz_1 - dz_2}{z_1 - z_2}.
\]
The bar differential extracts the OPE residue at $z_1 = z_2$:
\begin{equation}\label{eq:sl2-collision-bar-diff}
d_{\bar{B}}(\xi_{ab})
\;=\;
\operatorname{Res}_{z_1 = z_2}
\Bigl[
J^a(z_1)\, J^b(z_2)
\cdot \eta_{12}
\Bigr]
\;=\;
\underbrace{f^{ab}{}_{c}\, J^c(z_2)}_{\text{simple pole}}
\;+\;
\underbrace{k \cdot \kappa_{\fg}(a,b) \cdot |0\rangle}_{\text{double pole: no $\eta_{12}$ residue}}.
\end{equation}
The double-pole term $k\,\kappa_{\fg}(a,b)/(z_1 - z_2)^2$ contributes to
$\bar{B}^0$ (the curvature $m_0$), while the simple-pole term
$f^{ab}{}_c\, J^c(z_2)/(z_1 - z_2)$ contributes the
collision residue proper.

\emph{Step~3: Collision residue extraction.}
The genus-$0$, degree-$2$ collision residue of $\Theta_\cA$ is
(Construction~\ref{constr:thqg-V-binary-extraction}):
\begin{equation}\label{eq:sl2-collision-residue-formula}
\operatorname{Res}^{\mathrm{coll}}_{0,2}
(\Theta_{\widehat{\mathfrak{sl}}_2})(J^a, J^b)(z)
\;=\;
\operatorname{Res}_{w = z}
\bigl[J^a(w) \cdot J^b(z) \cdot \eta(w,z)\bigr]
\;=\;
(J^a_{(0)} J^b)(z)
\;=\;
f^{ab}{}_{c}\, J^c(z).
\end{equation}
The collision residue extracts the zeroth product $a_{(0)}b$,
which is precisely the Lie bracket of $\mathfrak{sl}_2$.
Explicitly:
\begin{alignat}{2}
\operatorname{Res}^{\mathrm{coll}}(e, f)
 &= h, &\qquad
\operatorname{Res}^{\mathrm{coll}}(f, e)
 &= -h,
 \label{eq:sl2-collision-explicit-ef} \\
\operatorname{Res}^{\mathrm{coll}}(h, e)
 &= 2e, &\qquad
\operatorname{Res}^{\mathrm{coll}}(e, h)
 &= -2e,
 \label{eq:sl2-collision-explicit-he} \\
\operatorname{Res}^{\mathrm{coll}}(h, f)
 &= -2f, &\qquad
\operatorname{Res}^{\mathrm{coll}}(f, h)
 &= 2f,
 \label{eq:sl2-collision-explicit-hf} \\
\operatorname{Res}^{\mathrm{coll}}(e, e)
 &= 0, &\qquad
\operatorname{Res}^{\mathrm{coll}}(f, f)
 &= 0.
 \label{eq:sl2-collision-explicit-diag}
\end{alignat}

\emph{Step~4: The $r$-matrix data.}
The binary collision data has two normalizations.  The trace-form
collision package recorded in the landscape census is
\[
r^{\mathrm{coll}}(z)=k\Omega_{\mathrm{tr}}/z
\]
(Remark~\ref{rem:km-collision-residue-rmatrix}); it carries the
unshifted level prefix and remains finite at critical level.  The
Knizhnik--Zamolodchikov coefficient belongs to the conformal-block
connection and uses the Sugawara-shifted inverse form.  It is
therefore
\[
r^{\mathrm{KZ}}(z)=\Omega_{\mathrm{KZ}}/((k+h^\vee)z),
\]
not the trace-form collision residue with a different name.  The
Casimir element of
$\mathfrak{sl}_2$ in the normalisation
$\kappa_{\fg}(e,f) = 1$, $\kappa_{\fg}(h,h) = 2$ is
\begin{equation}\label{eq:sl2-casimir-explicit}
\Omega_{\mathfrak{sl}_2}
\;=\;
\sum_{a,b} \kappa_{\fg}^{ab}\, J_a \otimes J_b
\;=\;
e \otimes f \;+\; f \otimes e \;+\;
\tfrac{1}{2}\, h \otimes h
\;\in\;
\mathfrak{sl}_2 \otimes \mathfrak{sl}_2,
\end{equation}
where $\kappa_{\fg}^{ab}$ is the inverse normalized invariant form:
$\kappa_{\fg}^{ef} = \kappa_{\fg}^{fe} = 1$,
$\kappa_{\fg}^{hh} = 1/2$.
The KZ coefficient is
\begin{equation}\label{eq:sl2-r-matrix-explicit}
\boxed{
r^{\mathrm{KZ}}_{\widehat{\mathfrak{sl}}_2}(z)
\;=\;
\frac{\Omega_{\mathfrak{sl}_2}}{(k + 2)\,z}
\;=\;
\frac{1}{(k+2)\,z}
\Bigl(
e \otimes f \;+\; f \otimes e \;+\;
\tfrac{1}{2}\, h \otimes h
\Bigr),
}
\end{equation}
with $h^\vee = 2$ for $\mathfrak{sl}_2$.
The factor $(k + h^\vee)^{-1}$ arises from the level-shifted
normalized invariant form in the Sugawara/KZ construction: the conformal-block
connection pairs insertions through the inverse form at the shifted
level $k + h^\vee$, giving
$r_{ab}(z) = \kappa_{\fg}^{ab}/((k + h^\vee)\,z)$.
The KZ coefficient is undefined at the critical level
$k = -h^\vee = -2$, reflecting the Sugawara singularity; the
collision residue itself remains $k\Omega_{\mathrm{tr}}/z$ and becomes
$-2\Omega_{\mathrm{tr}}/z$ in this trace-form convention. The KZ
normalization degenerates
when the Sugawara construction is undefined.

\emph{Step~5: CYBE from the Arnold relation.}
On $\operatorname{Conf}_3(\bC)$ with coordinates $z_1, z_2, z_3$,
the logarithmic $1$-forms $\eta_{ij} = d\log(z_i - z_j)$ satisfy
the Arnold relation
(Proposition~\ref{prop:thqg-V-arnold-cohomology}):
\begin{equation}\label{eq:sl2-arnold}
\eta_{12} \wedge \eta_{23}
\;+\;
\eta_{23} \wedge \eta_{31}
\;+\;
\eta_{31} \wedge \eta_{12}
\;=\; 0.
\end{equation}
Set $r^{\mathrm{KZ}}_{ij} = r^{\mathrm{KZ}}(z_{ij})
= \Omega/((k+2)\,z_{ij})$ with
$z_{ij} = z_i - z_j$.
Assume $k\neq -2$, so the KZ coefficient is defined.
The CYBE
\begin{equation}\label{eq:sl2-cybe}
[r^{\mathrm{KZ}}_{12}(z_{12}),\; r^{\mathrm{KZ}}_{13}(z_{13})]
\;+\;
[r^{\mathrm{KZ}}_{12}(z_{12}),\; r^{\mathrm{KZ}}_{23}(z_{23})]
\;+\;
[r^{\mathrm{KZ}}_{13}(z_{13}),\; r^{\mathrm{KZ}}_{23}(z_{23})]
\;=\; 0
\end{equation}
is the product of two identities: the infinitesimal braid relation
\[
[\Omega_{12},\Omega_{13}+\Omega_{23}]=0
\]
and its cyclic permutations, and the Arnold partial-fraction
identity~\eqref{eq:sl2-arnold}. The common scalar factor
$1/(k+2)^2$ divides out; explicitly:
\[
\frac{[\Omega_{12}, \Omega_{13}]}{z_{12}\,z_{13}}
\;+\;
\frac{[\Omega_{12}, \Omega_{23}]}{z_{12}\,z_{23}}
\;+\;
\frac{[\Omega_{13}, \Omega_{23}]}{z_{13}\,z_{23}}
\;=\; 0.
\]
This identity is the $\mathfrak{sl}_2^{\otimes 3}$ projection of
the codimension-$2$ boundary cancellation on
$\overline{\mathcal{M}}_{0,4}$
(Theorem~\ref{thm:thqg-V-cybe-from-arnold}).
Concretely, the bracket
$[\Omega_{12}, \Omega_{13}]$ in $\mathfrak{sl}_2^{\otimes 3}$
acts as $[\Omega, -]$ on the first factor and as the identity on the
remaining two; the Jacobi identity for $\mathfrak{sl}_2$ closes the
sum to zero after the partial-fraction identity
$1/(z_{12}\,z_{13}) + 1/(z_{12}\,z_{23}) + 1/(z_{13}\,z_{23}) = 0$
(which \emph{is} the Arnold relation evaluated at these points).

\emph{Step~6: The KZ connection.}
The genus-$0$ shadow connection at degree~$2$
is the flat connection on $n$-point conformal blocks
of $\widehat{\mathfrak{sl}}_{2,k}$:
\begin{equation}\label{eq:sl2-kz-connection}
\nabla^{\mathrm{KZ}}_{0,2}
\;=\;
d \;-\;
\frac{1}{k + 2}
\sum_{i < j}
\Omega_{ij}\,\frac{dz_{ij}}{z_{ij}}
\;=\;
d \;-\;
\sum_{i < j}
r^{\mathrm{KZ}}_{ij}(z_{ij})\, dz_{ij}.
\end{equation}
This is the Knizhnik--Zamolodchikov connection.
(The KZ connection is the genus-$0$ shadow connection $\nabla^{\mathrm{hol}}_{0,n}$ of Theorem~\ref{thm:gz26-commuting-differentials}(iii); the classical-quantum bridge is Theorem~\ref{thm:kz-classical-quantum-bridge}.)
The flatness condition
$(\nabla^{\mathrm{KZ}})^2 = 0$ is equivalent to the CYBE~\eqref{eq:sl2-cybe}, which is the MC equation at genus~$0$,
degree~$2$.

The entire chain is:
\[
\Theta_\cA
\;\xrightarrow{\text{restrict to }(g,n) = (0,2)}
\;\Theta_{0,2}
\;\xrightarrow{\text{binary collision}}\;
\bigl(k\Omega_{\mathrm{tr}},\, f^{ab}{}_cJ^c\bigr)
\;\xrightarrow{\text{KZ normalization}}\;
r^{\mathrm{KZ}}(z) = \frac{\Omega}{(k+2)z}
\;\xrightarrow{\text{all }n}\;
\nabla^{\mathrm{KZ}}.
\]
\end{computation}

\begin{remark}[Collision residue as Chriss--Ginzburg projection]
\label{rem:sl2-collision-chriss-ginzburg}
\index{Chriss--Ginzburg principle!collision residue}
The computation above gives the genus-$0$ Kac--Moody instance of the
Chriss--Ginzburg projection principle
(\S\ref{subsec:chriss-ginzburg-derivation}).  The universal MC element
$\Theta_\cA \in \mathrm{MC}(\mathfrak{g}^{\mathrm{mod}}_\cA)$
has a finite-order projection whose KZ normalization is the
connection above:
\begin{enumerate}[label=\textup{(\roman*)}]
\item Restrict $\Theta_\cA$ to genus~$0$, degree~$2$
 (the first nonlinear level);
\item keep both binary collision components at $z_1 = z_2$
 (the codimension-$1$ boundary of
 $\overline{\mathcal{M}}_{0,3}$): the trace-form residue
 $k\Omega_{\mathrm{tr}}/z$ and the Lie-bracket zeroth product;
\item apply the Sugawara-shifted KZ normalization to obtain
 $r^{\mathrm{KZ}}(z) \in \mathfrak{sl}_2^{\otimes 2}[z^{-1}]$.
\end{enumerate}
The CYBE is then the MC equation projected one further step:
genus~$0$, degree~$3$, codimension~$2$ boundary, via the Arnold
relation. The shadow obstruction tower of $\widehat{\mathfrak{sl}}_{2,k}$
terminates at degree~$3$ (Lie/tree archetype, $o_4 = 0$), so the
$r$-matrix coefficient
$r^{\mathrm{KZ}}(z) = \Omega/((k+2)z)$ with simple pole and no
higher KZ corrections is the complete genus-$0$ KZ projection of the
shadow datum:
$\Theta^{\leq 3}_{0} = \Theta_{0}$.
\end{remark}

\subsection{Quantum groups}
\label{sec:km-quantum-groups}

The categorical framework
(Definition~\ref{def:category-O},
Definition~\ref{def:integrable-modules})
underlies the quantum group comparison.

\begin{remark}[BGG reciprocity and the Kazhdan--Lusztig equivalence]\label{rem:category-O-affine}\label{rem:bgg-kl-bridge}
\index{BGG reciprocity!Kazhdan--Lusztig bridge}
\index{Kazhdan--Lusztig equivalence!BGG reciprocity}
BGG reciprocity for $\mathcal{O}_k(\widehat{\mathfrak{g}})$ states $[P_k(\lambda) : M_k(\mu)] = [M_k(\mu) : L_k(\lambda)]$. On a Kazhdan--Lusztig comparison surface with quantum parameter
\[
\zeta_k=e^{\pi i/(k+h^\vee)},
\]
this maps to quantum group reciprocity:
\[
[P_{\zeta_k}(\lambda) : \Delta_{\zeta_k}(\mu)]
=
[\Delta_{\zeta_k}(\mu) : L_{\zeta_k}(\lambda)]
\]
where $\Delta_{\zeta_k}(\mu)$ is the Weyl module.  The bar-cobar
level reflection sends \(k\) to \(k'=-k-2h^\vee\), hence
\(\zeta_{k'}=\zeta_k^{-1}\).  This gives compatibility of the
Verdier level shift with the external BGG/KL reciprocity formulas.
It does not prove BGG reciprocity; that theorem remains the
representation-theoretic input.
\end{remark}

\begin{theorem}[Quantum-group parameter inversion;
\ClaimStatusProvedHere]\label{thm:km-quantum-groups}
\index{Kazhdan--Lusztig equivalence}
\index{quantum group}
Let \(k\neq-h^\vee\) be a level lying on a
Kazhdan--Lusztig/Finkelberg comparison surface whose quantum-group
parameter is
\[
\zeta_k=e^{\pi i/(k+h^\vee)}.
\]
Then the chiral Koszul level reflection \(k\mapsto k'=-k-2h^\vee\)
sends this parameter to its inverse:
\begin{enumerate}[label=\textup{(\roman*)}]
\item The Koszul dual $(\widehat{\mathfrak{g}}_k)^! = \mathrm{CE}^{\mathrm{ch}}(\widehat{\mathfrak{g}}_{-k-2h^\vee})$
has modular characteristic $\kappa(\widehat{\mathfrak{g}}_{-k-2h^\vee})$
\textup{(}Theorem~\textup{\ref{thm:universal-kac-moody-koszul});}
\item The Kazhdan--Lusztig quantum group parameter for the dual level is
\begin{equation}\label{eq:q-inversion}
\zeta_{k'} = e^{\pi i/(-k-2h^\vee + h^\vee)}
= e^{-\pi i/(k+h^\vee)}
= \zeta_k^{-1};
\end{equation}
\item Consequently, the chiral Koszul level shift $k \mapsto -k - 2h^\vee$
corresponds to \(\zeta_k\leftrightarrow \zeta_k^{-1}\) at the
quantum-group parameter level.
\end{enumerate}
\end{theorem}

\begin{proof}
\emph{Step~1.}
The Koszul dual
$(\widehat{\mathfrak{g}}_k)^! = \mathrm{CE}^{\mathrm{ch}}(\widehat{\mathfrak{g}}_{-k-2h^\vee})$
is identified in Theorem~\ref{thm:universal-kac-moody-koszul} (\ClaimStatusProvedHere).

\emph{Step~2.}
On a comparison surface where the external
Kazhdan--Lusztig/Finkelberg theorem applies \cite{KL93}
(\ClaimStatusProvedElsewhere), the affine category is compared with
the semisimplified tilting category at parameter \(\zeta_k\):
\[
\mathcal{O}_k^{\mathrm{int}}(\widehat{\mathfrak{g}})
\;\simeq\;
\mathcal{C}(U_{\zeta_k}(\mathfrak{g})), \qquad
\zeta_k = e^{\pi i/(k + h^\vee)},
\]
where $\mathcal{C}(U_{\zeta_k}(\mathfrak{g}))$ is the semisimplified tilting
category (Definition~\ref{def:tilting-module}), not the full
representation category $\mathrm{Rep}(U_{\zeta_k}(\mathfrak{g}))$.
Applied to the Koszul dual level $k' = -k - 2h^\vee$:
\[
\zeta_{k'} = e^{\pi i/(k' + h^\vee)}
= e^{\pi i/(-k - 2h^\vee + h^\vee)}
= e^{-\pi i/(k + h^\vee)}
= \zeta_k^{-1}.
\]

\emph{Step~3.}
Combining Steps~1 and~2: the Koszul duality functor
records the level-reflected current presentation
\(\mathrm{CE}^{\mathrm{ch}}(\widehat{\mathfrak{g}}_{k'})\),
where \(k'=-k-2h^\vee\). Composed with the Kazhdan--Lusztig
equivalences for levels~$k$ and~$k'$, this level reflection induces
the parameter functor
$\mathcal{C}(U_{\zeta_k}(\mathfrak{g})) \to
\mathcal{C}(U_{\zeta_k^{-1}}(\mathfrak{g}))$
on semisimplified tilting categories,
yielding the parameter correspondence
\(\zeta_k \mapsto \zeta_k^{-1}\). This is~(iii). The stronger
categorical tensor-equivalence statement is separate and is not used
in this theorem.
\end{proof}

\begin{remark}[Categorical refinement]\label{rem:km-categorical-refinement}
The KL equivalence $\mathcal{O}_k^{\mathrm{int}}(\widehat{\mathfrak{g}}) \simeq \mathcal{C}(U_{\zeta_k}(\mathfrak{g}))$ is a braided tensor equivalence \cite{KL93}. The stronger statement is the commutativity of the diagram
\[
\begin{tikzcd}
D^b(\mathrm{Mod}_{\widehat{\mathfrak{g}}_k}^{\Eone}) \arrow[r, "\sim"] \arrow[d, "\mathrm{KL}"] &
D^b(\mathrm{Comod}_{(\widehat{\mathfrak{g}}_k)^!}^{\Eone}) \arrow[d, "\mathrm{KL}'"] \\
D^b(\mathcal{C}(U_{\zeta_k}(\mathfrak{g}))) \arrow[r, dashed, "\zeta_k \mapsto \zeta_k^{-1}"] &
D^b(\mathcal{C}(U_{\zeta_k^{-1}}(\mathfrak{g})))
\end{tikzcd}
\]
as a tensor equivalence.  This is a separate quantum-Langlands
conjecture; Theorem~\ref{thm:km-quantum-groups} proves only the
parameter inversion forced by the level shift.
\end{remark}

\begin{remark}[Deformation rigidity at positive integral level]
\label{rem:deformation-rigidity-integral}
\index{deformation rigidity!integral level}
\index{integral level!deformation rigidity}
Linshaw--Qi~\cite{LQ26} prove that the simple affine vertex algebras
$L_k(\fg)$ at positive integral levels $k \in \mathbb{Z}_{>0}$ have
only trivial first-order deformations: $H^1_{\mathrm{def}}(L_k(\fg))
= 0$. This rigidity is independent of Koszulness. Koszulness detects
bar concentration (Definition~\ref{def:chiral-koszul-morphism}), which
is a property of the bar filtration; deformation rigidity detects
vanishing of the first cohomology of the deformation complex, which is
a property of the tangent space to the moduli of vertex algebra
structures. The two invariants probe different directions: Koszulness
is a statement about $\mathrm{Ext}$ groups (the bar spectral sequence),
while rigidity is a statement about $H^1$ of the vertex algebra
cohomology complex.
\end{remark}

\begin{example}[Integrable modules for \texorpdfstring{$\widehat{\mathfrak{sl}}_2$}{sl-hat_2}
at low levels]\label{ex:integrable-low-levels}
\index{integrable module!low level examples}
We list the integrable module categories
$\mathcal{O}_k^{\mathrm{int}}(\widehat{\mathfrak{sl}}_2)$
for levels $k = 1, 2, 3$. At each level, the simple objects are
$L(\Lambda)$ for dominant integral weights $\Lambda$ satisfying
$\langle \Lambda, \alpha_0^\vee \rangle \leq k$.
\begin{center}
\begin{tabular}{cccc}
$k$ & Simple objects & $|P^k_+|$ & $q = e^{\pi i/(k+2)}$ \\
\hline
$1$ & $L(\Lambda_0),\; L(\Lambda_1)$ & $2$ & $e^{\pi i/3}$ \\
$2$ & $L(\Lambda_0),\; L(\Lambda_1),\; L(2\Lambda_1)$ & $3$ &
 $e^{\pi i/4}$ \\
$3$ & $L(\Lambda_0),\; L(\Lambda_1),\; L(2\Lambda_1),\;
 L(3\Lambda_1)$ & $4$ & $e^{\pi i/5}$
\end{tabular}
\end{center}
At level~$k$, there are $k+1$ integrable modules.

\smallskip
\textup{[Category: semisimplified tilting
$\mathcal{C}(U_q(\mathfrak{sl}_2))$.]}
Under the KL equivalence, this matches
\[
|\mathrm{Irr}(\mathcal{C}(U_q(\mathfrak{sl}_2)))| = k+1
\quad\text{at } q = e^{\pi i/(k+2)}
\]
(a $(2k{+}4)$th root of unity).

The fusion rules (by the Verlinde formula, or equivalently the
tensor product rules for $U_q(\mathfrak{sl}_2)$) are:
\[
L(j\Lambda_1) \boxtimes L(j'\Lambda_1)
= \bigoplus_{\substack{|j - j'| \leq j'' \leq
 \min(j+j',\, 2k-j-j') \\ j'' \equiv j + j' \bmod 2}}
 L(j''\Lambda_1)
\]
(truncated Clebsch--Gordan rule: the decomposition is the same
as for $\mathrm{SU}(2)$ but with the sum truncated at $j'' \leq 2k - j - j'$).

\emph{Level~$k = 1$:}
$L(\Lambda_1) \boxtimes L(\Lambda_1) = L(\Lambda_0)$
($\mathbb{Z}/2$ fusion algebra).

\emph{Level~$k = 3$:}
$L(\Lambda_1) \boxtimes L(\Lambda_1)
= L(\Lambda_0) \oplus L(2\Lambda_1)$,
$L(\Lambda_1) \boxtimes L(3\Lambda_1) = L(2\Lambda_1)$,
$L(2\Lambda_1) \boxtimes L(2\Lambda_1) = L(\Lambda_0) \oplus
L(2\Lambda_1)$ (the $\mathbb{Z}/2$ grading is visible: odd$\times$odd
$=$~even).
\end{example}

\begin{remark}[Screening operators and bar complex for modules]
\label{rem:screening-bar-module}
\index{screening operator!module bar complex}
The bar-cobar framework for modules
(Definition~\ref{def:bar-complex-modules}) interacts with
screening operators as follows. For an integrable module
$L(\Lambda)$ of $\widehat{\mathfrak{sl}}_2$ at level~$k$:
\begin{enumerate}[label=(\roman*)]
\item The bar complex
 $\bar{B}^{\mathrm{ch}}_\bullet(\widehat{\mathfrak{sl}}_2,
 L(\Lambda))$ is acyclic (since $L(\Lambda)$ is simple at
 integrable level).
\item The BGG resolution of $L(\Lambda)$ by Verma modules
 is realized by screening operators $\oint S_\alpha(z)\,dz$
 integrated over cycles in $\overline{C}_n(X)$
 (Proposition~\ref{prop:bgg-sl2-level1}).
\item The truncation in the fusion rule (truncated Clebsch--Gordan
 at $j'' \leq 2k - j - j'$) is a \emph{quantum correction}:
 in the bar complex, it arises from the non-trivial boundary
 strata on $\overline{C}_n(X)$ at level-dependent values.
 At the classical level $k \to \infty$, the truncation disappears
 and one recovers the full $\mathrm{SU}(2)$ tensor product.
\end{enumerate}
\end{remark}

\subsection{Spectral flow}

\begin{definition}[Spectral flow automorphism]\label{def:spectral-flow}
\index{spectral flow|textbf}
Let $\widehat{\mathfrak{g}}_k$ be an affine Lie algebra at level~$k$
with Cartan subalgebra $\mathfrak{h} \subset \mathfrak{g}$ and coroot
lattice $Q^\vee$. For each $\gamma \in Q^\vee$, the \emph{spectral
flow automorphism} $\sigma_\gamma$ is the automorphism of
$\widehat{\mathfrak{g}}_k$ defined on generators by:
\begin{align}
\sigma_\gamma(J^a_n) &= J^a_n + k\,\langle \gamma, \alpha_a \rangle\,
\delta_{n,0}, \\
\sigma_\gamma(K) &= K, \qquad
\sigma_\gamma(d) = d - \gamma_0
\end{align}
where $\gamma_0 \in \mathfrak{h}$ is the finite part of~$\gamma$ and
$\alpha_a$ is the root corresponding to~$J^a$.

For $\widehat{\mathfrak{sl}}_2$ at level~$k$, the spectral flow by
$\gamma = n \alpha^\vee/2$ acts on highest weight modules by:
\[
\sigma_n\colon L(k, h) \;\longmapsto\; L(k,\; h + nk/2 + n^2/4)
\]
shifting the conformal dimension while preserving the level.
\end{definition}

\begin{proposition}[Spectral flow and Koszul duality \cite{Kac}; \ClaimStatusProvedElsewhere]
\label{prop:spectral-flow-koszul}
The spectral flow automorphism intertwines with Koszul duality as
follows. For a Koszul pair $(\widehat{\mathfrak{g}}_k,\,
\widehat{\mathfrak{g}}_{-k-2h^\vee})$:
\begin{enumerate}[label=(\roman*)]
\item Spectral flow preserves the Koszul property: if
 $\widehat{\mathfrak{g}}_k$ is Koszul at level~$k$, then the twisted
 algebra $\sigma_\gamma(\widehat{\mathfrak{g}}_k)$ is Koszul with the
 same dual.
\item The Koszul duality functor commutes with spectral flow:
 $\sigma_\gamma \circ \mathrm{Ind} \simeq \mathrm{Ind} \circ
 \sigma_\gamma$ on module categories.
\item At integral level $k \in \mathbb{Z}_{>0}$, the spectral flow
 orbits on $\mathcal{O}_k^{\mathrm{int}}(\widehat{\mathfrak{g}})$
 correspond to the outer automorphism orbits of the Dynkin diagram
 of~$\widehat{\mathfrak{g}}$.
\end{enumerate}
\end{proposition}

\subsection{Conformal blocks via the bar complex}
\label{sec:km-conformal-blocks}

\begin{example}[Conformal blocks for \texorpdfstring{$\widehat{\mathfrak{sl}}_2$}{sl-hat_2}
on \texorpdfstring{$\mathbb{P}^1$}{P1}]\label{ex:conformal-blocks-sl2}
\index{conformal blocks!sl2@$\widehat{\mathfrak{sl}}_2$}
\index{Verlinde formula!bar complex derivation}
By Proposition~\ref{prop:conformal-blocks-bar}, the space of
conformal blocks for $\widehat{\mathfrak{sl}}_2$ at level~$k$
with insertions $L(j_1\Lambda_1), \ldots, L(j_r\Lambda_1)$
at points $x_1, \ldots, x_r \in \mathbb{P}^1$ is:
\[
\dim \mathbb{V}\bigl(
\widehat{\mathfrak{sl}}_{2,k};\,
L(j_1\Lambda_1), \ldots, L(j_r\Lambda_1);\, \mathbb{P}^1
\bigr)
= N_{j_1 \cdots j_r}^{(k)}
\]
where $N_{j_1 \cdots j_r}^{(k)}$ is the iterated fusion
coefficient, computed by the Verlinde formula:
\begin{equation}\label{eq:verlinde-sl2}
N_{j_1 \cdots j_r}^{(k)}
= \sum_{m=0}^{k}
\prod_{i=1}^{r}
\frac{S_{j_i,m}}{S_{0,m}},
\qquad
S_{j,m} = \sqrt{\tfrac{2}{k+2}}\,
\sin\!\Bigl(\frac{\pi(j+1)(m+1)}{k+2}\Bigr).
\end{equation}

In the bar complex picture, the $n$-th term of the bar resolution
with $r$ module insertions lives on
$\overline{C}_{n+r}(\mathbb{P}^1)$, and the alternating Euler
characteristic computes the Verlinde dimension. Since
$\overline{C}_n(\mathbb{P}^1)$ is a smooth projective variety
of dimension $n - 3$ (for $n \geq 3$) with known Betti numbers,
the Verlinde formula~\eqref{eq:verlinde-sl2} arises as the
Euler characteristic of a complex of sections on these spaces.

\emph{Three-point function.}
For $r = 3$: $\dim \mathbb{V} = N_{j_1 j_2 j_3}^{(k)}$ equals
$1$ if $|j_1 - j_2| \leq j_3 \leq \min(j_1+j_2,
2k - j_1 - j_2)$ and $j_1 + j_2 + j_3 \in 2\mathbb{Z}$, and
$0$ otherwise. The bar complex on
$\overline{C}_3(\mathbb{P}^1) \cong \mathrm{pt}$ reduces to a
single constraint: the screening charge integrability condition.

\emph{Koszul dual blocks.}
By Remark~\ref{rem:conformal-blocks-koszul}, the conformal blocks
are invariant under $k \mapsto -k - 4$: the Verlinde numbers
satisfy $N^{(k)}_{j_1 \cdots j_r} = N^{(-k-4)}_{j_1 \cdots j_r}$
when both levels are non-negative integers. This is the
level-rank duality for $\mathfrak{sl}_2$ conformal blocks
(\cite{NakanishiTsuchiya}).
\end{example}

\begin{theorem}[Verlinde recovery from the bar complex; \ClaimStatusProvedElsewhere]
\label{thm:bar-verlinde-recovery}
\index{Verlinde formula!bar complex recovery}
\index{conformal blocks!genus-$g$ bar recovery}
For $V_k(\mathfrak{g})$ at integrable level $k \in \mathbb{Z}_{>0}$,
the genus-$g$ bar $H^0$ dimension equals the Verlinde dimension:
\begin{equation}\label{eq:bar-verlinde-recovery}
\dim H^0\bigl(B^{(g)}(V_k(\mathfrak{g})),\, D_g\bigr)
\;=\;
Z_g(\mathfrak{g},\, k)
\;=\;
\sum_{\lambda \in P_+^k}
\Bigl(\frac{S_{0,\lambda}}{S_{0,0}}\Bigr)^{\!2-2g}
\end{equation}
where $P_+^k$ is the set of integrable highest weights at level~$k$,
$S_{\mu,\lambda}$ is the modular $S$-matrix, and
$D_g$ is the genus-$g$ bar differential.
\end{theorem}

\begin{remark}[Attribution]
The ingredients are: Proposition~\ref{prop:conformal-blocks-bar}
(conformal blocks as bar $H^0$), Zhu's theorem
\cite{Zhu96} (genus-$1$ modular invariance and $S$-matrix
diagonalisation), the factorisation property of conformal blocks
under degeneration of curves (genus induction), and the Verlinde
formula \cite{Verlinde} (diagonalisation of fusion by the
$S$-matrix). The chain of identifications is:
bar $H^0$ on $X_g$ $\xrightarrow{\text{Prop.~\ref{prop:conformal-blocks-bar}}}$
conformal blocks $\xrightarrow{\text{factorisation}}$
fusion product $\xrightarrow{\text{Zhu + \cite{Verlinde}}}$
Verlinde dimension.
At genus~$0$, Example~\ref{ex:conformal-blocks-sl2} computes
$Z_0(\mathfrak{sl}_2, k) = k + 1$; at genus~$1$,
$Z_1(\mathfrak{g}, k) = |P_+^k|$ recovers the number of
integrable representations.
\end{remark}

\subsection{Geometric Langlands}

\begin{remark}[Langlands correspondence via Koszul duality]
At critical level $k = -h^\vee$, the large Feigin--Frenkel center and
the oper description of the critical affine algebra place the
critical-level bar picture next to geometric Langlands theory.
What is proved internally is the critical fixed-point / uncurved
package, not a theorem of the form
\[
\widehat{\mathfrak{g}}_{-h^\vee}^! \simeq
\widehat{\mathfrak{g}^\vee}_{-h^{\vee,\vee}}.
\]
The Langlands-dual interpretation is therefore external motivation for
the critical theory, not a manuscript-internal Koszul-duality theorem.
\end{remark}

\subsection{Admissible level representations}
\label{sec:admissible-level-reps}
\index{admissible level!representations|textbf}

At positive integer levels, the category
$\mathcal{O}_k^{\mathrm{int}}(\widehat{\mathfrak{g}})$ is
semisimple with finitely many simples. At \emph{admissible}
levels, the representation theory is richer: on the
non-degenerate admissible lane the simple quotient is still
rational with finitely many simples, while on the degenerate lane
logarithmic behavior appears. In both cases there remain
finiteness properties that make the bar complex tractable.

\begin{definition}[Admissible level and admissible weights]
\label{def:admissible-level}
\index{admissible level|textbf}
\index{admissible weight|textbf}
A level $k \in \mathbb{C}$ for $\widehat{\mathfrak{g}}$ is
\emph{admissible} if $k = -h^\vee + p/q$ where $p, q$ are
coprime positive integers satisfying $p \geq h^\vee$ (and for
non-simply-laced $\mathfrak{g}$, $q$ is coprime to
$r^\vee = \max_i a_i^\vee$, the lacing number).

A weight $\Lambda$ of level~$k$ is \emph{admissible} if:
\begin{enumerate}[label=\textup{(\roman*)}]
\item $\langle \Lambda + \rho, \alpha^\vee \rangle \notin
 \{0, -1, -2, \ldots\}$ for all positive real roots
 $\alpha \in \Delta^+_{\mathrm{re}}$;
\item $\langle \Lambda + \rho, \alpha^\vee \rangle \in
 \mathbb{Z}_{> 0}$ for a specific subset of positive real roots
 determined by $p$ and $q$.
\end{enumerate}
The set of admissible weights at level~$k$ is denoted $\mathrm{Adm}_k$.
It is a finite set: $|\mathrm{Adm}_k| < \infty$
(Kac--Wakimoto~\cite{KW88}).
\end{definition}

\begin{theorem}[Representation theory at admissible level \cite{KW88, Arakawa17}; \ClaimStatusProvedElsewhere]
\label{thm:admissible-rep-theory}
\index{admissible level!representation theory}
Let $k = -h^\vee + p/q$ be an admissible level.
\begin{enumerate}[label=\textup{(\roman*)}]
\item \textup{(}Kac--Wakimoto~\cite{KW88}\textup{)}
 For each $\Lambda \in \mathrm{Adm}_k$, the simple module
 $\mathcal{L}(\Lambda)$ is the unique irreducible quotient of
 the Verma module $\mathcal{M}(\Lambda)$. These are all the
 simple highest weight modules in the admissible category.
\item \textup{(}Arakawa~\cite{Arakawa17}\textup{)}
 The simple affine vertex algebra $L_k(\mathfrak{g})$ at
 admissible level is rational if and only if the admissible level
 is \emph{non-degenerate}, meaning $X_{L_k(\mathfrak{g})} = \{0\}$
 \textup{(}Definition~\textup{\ref{def:associated-variety})}.
 At non-degenerate admissible levels, $L_k(\mathfrak{g})$ has
 finitely many simple modules, all admissible.
\item \textup{(}Arakawa~\cite{Arakawa17}\textup{)}
 At degenerate admissible level, $L_k(\mathfrak{g})$ has
 infinitely many simple modules. The associated variety
 $X_{L_k(\mathfrak{g})} = \overline{\mathbb{O}}$ is a non-trivial
 nilpotent orbit closure, and the module category has
 logarithmic (non-semisimple) structure.
\end{enumerate}
\end{theorem}

\begin{proposition}[Bar complex at admissible level; \ClaimStatusProvedHere]
\label{prop:bar-admissible}
\index{bar complex!admissible level}
\index{admissible level!bar complex}
Let $k = -h^\vee + p/q$ be an admissible level for
$\widehat{\mathfrak{g}}$. The bar complex exhibits the
following behavior, controlled by the associated variety
$X_{L_k(\mathfrak{g})}$\textup{:}
\begin{enumerate}[label=\textup{(\roman*)}]
\item \emph{Non-degenerate admissible:}
 $X_{L_k} = \{0\}$. In each fixed conformal-weight sector,
 the bar spectral sequence $E_r^{p,q}$ has finitely many
 non-zero terms and converges after finitely many pages.
 On the character-theoretic surfaces treated later, this
 weightwise finite-page control underlies the
 Kac--Wakimoto character formula.
\item \emph{Degenerate admissible:}
 $X_{L_k} = \overline{\mathbb{O}} \neq \{0\}$. The bar
 and Li--bar spectral sequences acquire reduced-level
 differentials from the Poisson bracket on the reduced
 associated graded, while nilradical contributions may
 produce additional off-diagonal bar data. The present
 manuscript does \emph{not} identify the resulting higher
 bar cohomology with ordinary logarithmic extensions between
 admissible modules.
\end{enumerate}
\end{proposition}

\begin{proof}
For (i): when $X_{L_k} = \{0\}$, the algebra $L_k(\mathfrak{g})$
is $C_2$-cofinite (Theorem~\ref{thm:arakawa-rationality}(i)).
Each conformal-weight space of $L_k$ is finite-dimensional, so
for each fixed total conformal weight the corresponding sector of
every bar degree is finite-dimensional. Because the reduced bar
construction is positively graded, only finitely many bar degrees
contribute to a fixed conformal-weight sector. Hence the induced
spectral sequence is finite in every fixed weight sector and
therefore converges after finitely many pages there. The later
character theorems make this finite-page control explicit in the
Kac--Wakimoto setting.

For (ii): when $X_{L_k} = \overline{\mathbb{O}} \neq \{0\}$,
the reduced quotient is a positive-dimensional Poisson variety.
Construction~\ref{constr:li-bar-spectral-sequence},
Theorem~\ref{thm:associated-variety-koszulness}, and
Proposition~\ref{prop:large-orbit-obstruction} show that the
reduced Li--bar page sees the Poisson-bracket-induced $d_1$
surface, while nilradical contributions remain as separate
obstructions on the full associated graded. This gives a
bar-side source of additional differentials and off-diagonal
cohomology, but it does not identify those groups with Ext in the
ordinary admissible module category.
\end{proof}

\begin{remark}[Dual-level admissible prediction]
\label{rem:bar-admissible-dual-level}
The admissible level $k = -h^\vee + p/q$ maps under the
Feigin--Frenkel involution to
$k' = -k - 2h^\vee = -h^\vee - p/q$. This dual level lies
outside the admissible regime of
Definition~\ref{def:admissible-level} since $k' + h^\vee = -p/q<0$.
The central charge complementarity
\textup{(}Theorem~\textup{\ref{thm:central-charge-complementarity}}\textup{)}
still applies, and the type-$A$ orbit-duality conjecture
\textup{(}Conjecture~\textup{\ref{conj:orbit-duality}}\textup{)}
predicts that if $X_{L_k}=\overline{\mathbb{O}}$, then
$X_{L_{k'}}$ is conjecturally the Barbasch--Vogan dual orbit closure.
This prediction is kept separate from
Proposition~\ref{prop:bar-admissible}, whose proved content is only
the bar-spectral-sequence behavior on the admissible side.
\end{remark}

\begin{corollary}[Bar complex finiteness at non-degenerate admissible levels; \ClaimStatusProvedHere]
\label{cor:bar-admissible-finiteness}
\index{admissible level!bar finiteness}
\index{bar complex!finiteness at admissible level}
At a non-degenerate admissible level $k = -h^\vee + p/q$ and in
each fixed conformal-weight sector, the bar complex of
$L_k(\mathfrak{g})$ has the following finiteness properties:
\begin{enumerate}[label=\textup{(\roman*)}]
\item The fixed-weight sector of each bar degree $\bar{B}^n(L_k)$
 is finite-dimensional.
\item Only finitely many bar degrees contribute to that fixed
 weight sector, so the induced spectral sequence
 $E_r^{p,q} \Rightarrow H^{p+q}(\bar{B}(L_k))$ terminates after
 finitely many pages in that sector.
\item The corresponding fixed-weight sector of
 $H^*(\bar{B}(L_k))$ is finite-dimensional.
\end{enumerate}
\end{corollary}

\begin{proof}
By Theorem~\ref{thm:admissible-rep-theory}(ii), $L_k(\mathfrak{g})$ is rational at non-degenerate admissible level, with finitely many simple modules $\{M_1, \ldots, M_N\}$ where $N = |\mathrm{Adm}_k|$. The Zhu algebra $A(L_k)$ is a finite-dimensional semisimple algebra of dimension $\dim A(L_k) = \sum_{i=1}^N (\dim M_i^{(0)})^2$ (Theorem~\ref{thm:zhu-correspondence}).

For~(i): the Li filtration $F_p L_k$ has $\mathrm{gr}^F L_k$
finite-dimensional in each conformal weight
\textup{(}by $C_2$-cofiniteness,
Theorem~\ref{thm:arakawa-rationality}\textup{)}. The bar
complex $\bar{B}^n(L_k)$ inherits this weightwise finiteness,
so each fixed-weight sector is finite-dimensional.

For~(ii): Proposition~\ref{prop:bar-admissible}(i) gives finite
weightwise convergence of the spectral sequence. Positive
conformal grading implies that only finitely many bar degrees
contribute to a fixed weight sector.

For~(iii): in each fixed conformal-weight sector, the spectral
sequence has finitely many pages with finite-dimensional entries,
so the abutment is finite-dimensional in that sector.
\end{proof}

\begin{example}[\texorpdfstring{$\widehat{\mathfrak{sl}}_2$}{sl-hat_2} at admissible level
\texorpdfstring{$k = -1/2$}{k = -1/2}]
\label{ex:sl2-admissible}
\index{sl2-hat@$\widehat{\mathfrak{sl}}_2$!admissible level}
The simplest non-trivial admissible level for
$\widehat{\mathfrak{sl}}_2$ is $k = -2 + 3/2 = -1/2$
(so $p = 3$, $q = 2$, $h^\vee = 2$). This level lies between
the critical level $k = -2$ and the first integrable level $k = 0$.

\emph{Admissible weights.}
At $k = -1/2$, the admissible weights form a finite set:
$\mathrm{Adm}_{-1/2} = \{\Lambda_0^{(0)}, \Lambda_0^{(1)}\}$
, two simple ordinary modules
(since $|\mathrm{Adm}_k| = (p - 1)(q - 1) = 2 \cdot 1 = 2$
for $\widehat{\mathfrak{sl}}_2$ at level $-2 + p/q$ with
$p = 3$, $q = 2$). More explicitly:
\[
\mathcal{L}(\Lambda_0^{(0)}) = L_{-1/2}(\mathfrak{sl}_2)
\quad\text{(vacuum module)},
\qquad
\mathcal{L}(\Lambda_0^{(1)})
\quad\text{(spectral flow twist)}.
\]
The Virasoro central charge is
$c(-1/2) = 3 \cdot (-1/2)/((-1/2) + 2) = -1$.

\emph{Associated variety.}
At this level, Arakawa's admissible rationality criterion places
$L_{-1/2}(\mathfrak{sl}_2)$ on the non-degenerate admissible lane:
$X_{L_{-1/2}(\mathfrak{sl}_2)} = \{0\}$, so the algebra is
rational and $C_2$-cofinite.

\emph{Koszul dual.}
The Koszul dual level is $k' = -(-1/2) - 4 = -7/2$. Then
$c(-7/2) = 3 \cdot (-7/2)/((-7/2) + 2) = 7$. The central
charge complementarity gives $c + c' = -1 + 7 = 6 =
2\dim(\mathfrak{sl}_2)$, as predicted by
Theorem~\ref{thm:central-charge-complementarity}.

The dual level $k' = -7/2$ satisfies
$k' + h^\vee = -3/2 < 0$, so $L_{-7/2}(\mathfrak{sl}_2)$
is \emph{not} admissible in the sense of
Definition~\ref{def:admissible-level}. The associated variety
$X_{L_{-7/2}(\mathfrak{sl}_2)} = \mathcal{N}$
(the full nilcone of $\mathfrak{sl}_2^*$) is non-trivial,
consistent with the orbit duality prediction
$\{0\} \leftrightarrow \mathcal{N}$ and the non-rationality
of $L_{-7/2}(\mathfrak{sl}_2)$.

\emph{Bar complex.}
At level $k = -1/2$ (non-degenerate): the bar spectral sequence
used in the character computation degenerates at $E_2$. This
finite-page degeneration is compatible with the two ordinary
simple modules, but it is \emph{not} promoted here to a theorem
that the simple quotient is chirally Koszul. The
Kac--Wakimoto character formula~\cite{KW88}
expresses the characters of $L_{-1/2}(\mathfrak{sl}_2)$-modules
via theta functions at level $k + h^\vee = 3/2$; the vacuum
character is
\[
\operatorname{ch}(L_{-1/2}(\mathfrak{sl}_2))(\tau, z)
= \frac{\Theta_{-1,6}(\tau, z) - \Theta_{5,6}(\tau, z)}
 {\Theta_{1,1}(\tau, z) - \Theta_{-1,1}(\tau, z)}
\]
where $\Theta_{j,m}(\tau, z) = \sum_{n \in \mathbb{Z}}
e^{2\pi i(2mn+j)z}\, q^{(2mn+j)^2/4m}$ and the denominator is
the affine Weyl denominator for $\widehat{\mathfrak{sl}}_2$.
\end{example}

\begin{example}[Li--bar spectral sequence for
\texorpdfstring{$\widehat{\mathfrak{sl}}_2$}{sl-hat_2}
at admissible levels]
\label{ex:sl2-li-bar}
\index{Li filtration!sl2-hat@$\widehat{\mathfrak{sl}}_2$ example}
\index{associated variety!sl2-hat@$\widehat{\mathfrak{sl}}_2$ admissible}
We illustrate the Li--bar spectral sequence
(Construction~\ref{constr:li-bar-spectral-sequence}) and the
associated-variety criterion
(Theorem~\ref{thm:associated-variety-koszulness}) for
$\widehat{\mathfrak{sl}}_2$ at the two types of admissible level.

\emph{Non-degenerate admissible
\textup{(}$k = -1/2$, $X_{L_k} = \{0\}$\textup{)}.}
The associated graded
$R_{L_{-1/2}} = \operatorname{gr}^F L_{-1/2}(\mathfrak{sl}_2)$
is finite-dimensional: $C_2$-cofiniteness forces
$\dim R_{L_{-1/2}} < \infty$. The Li--bar spectral sequence
has finitely many non-zero entries and degenerates at a finite
page. On the reduced level one has
$X_{L_{-1/2}} = \{0\}$, so the reduced commutative bar complex is
concentrated in degree~$0$. This is compatible with Koszulness,
but it does \emph{not} identify the full Artinian algebra
$R_{L_{-1/2}}$ with its reduced quotient, hence does not by
itself prove bar-cobar inversion for the simple quotient.

\emph{Degenerate admissible
\textup{(}$k = -4/3$, $X_{L_k} = \mathcal{N}_{\mathfrak{sl}_2}$\textup{)}.}
At $k = -4/3$ (so $p = 2$, $q = 3$), the associated variety is
$X_{L_{-4/3}} = \mathcal{N}_{\mathfrak{sl}_2}$
(the nilcone of $\mathfrak{sl}_2^*$, which is a surface
singularity $\{xy = z^2\} \subset \mathbb{A}^3$). The associated
graded has reduced quotient
\[
R_{L_{-4/3}}^{\mathrm{red}}
\cong \mathcal{O}(\mathcal{N}_{\mathfrak{sl}_2})
= \mathbb{C}[e, f, h]/(ef - h^2/4)
\]
with Kirillov--Kostant Poisson bracket
$\{e, f\} = h$, $\{h, e\} = 2e$, $\{h, f\} = -2f$.

For $\mathfrak{sl}_2$, the nilcone
$\mathcal{N} = \overline{\mathbb{O}}_{\mathrm{reg}} = \mathbb{O}_{\mathrm{reg}} \cup \{0\}$
is the closure of the unique non-zero orbit. The Springer
resolution is $T^*\mathbb{P}^1 \to \mathcal{N}$, so the reduced
geometry gives a natural test object. However,
Theorem~\ref{thm:associated-variety-koszulness} does
\emph{not} identify the reduced Li--bar $E_2$ page with ordinary
Poisson cohomology or with
$H^*(T^*\mathbb{P}^1,\mathcal{O})$. What it does say is that,
on the reduced quotient, the Li--bar $E_1$ page is the
commutative bar cohomology of
$\mathcal{O}(\mathcal{N}_{\mathfrak{sl}_2})$, and its $d_1$
differential is induced by the Kirillov--Kostant bracket.
Determining whether the resulting reduced Li--bar $E_2$ page is
diagonally concentrated is therefore a separate reduced-level
calculation.

Even if the reduced page were diagonally concentrated, the
remaining obstruction to Koszulness would still lie in the
nilradical of the
full non-reduced algebra $R_{L_{-4/3}}$
(Proposition~\ref{prop:large-orbit-obstruction}(ii)):
the null vectors of $L_{-4/3}(\mathfrak{sl}_2)$ contribute
nilpotent elements not visible on~$\mathcal{N}$, and
determining whether these produce off-diagonal bar
cohomology requires an explicit nilradical computation.

This pattern holds for all types: reduced associated-variety
geometry may constrain the reduced Li--bar page, but the
missing theorem is a nilradical calculation in the full non-reduced Li
associated graded.
\end{example}

\begin{theorem}[Kac--Wakimoto formula at \texorpdfstring{$k=-1/2$}{k=-1/2} via bar spectral sequence; \ClaimStatusConditional]
\label{thm:kw-bar-spectral}
\index{Kac--Wakimoto character formula!bar spectral sequence}
\index{bar complex!admissible character}
\index{admissible level!character formula}
Let $k=-1/2$ \textup{(}equivalently $p=3$, $q=2$\textup{)} for
$\widehat{\mathfrak{sl}}_2$. Then the
bar spectral sequence for the vacuum module
$L_k = L_k(\mathfrak{sl}_2)$ degenerates at~$E_2$, and the
$E_2$ page computes the Kac--Wakimoto character formula:
\begin{equation}\label{eq:kw-bar-derivation}
\operatorname{ch}(L_k)(\tau,z)
= \frac{\Theta_{-1,6}(\tau,z)-\Theta_{5,6}(\tau,z)}
 {\Theta_{1,1}(\tau,z) - \Theta_{-1,1}(\tau,z)}.
\end{equation}
The denominator is the affine Weyl denominator
\textup{(}arising from the universal bar complex\textup{)}, and the
numerator encodes the two surviving admissible Weyl contributions on
the $E_2$ page.
\end{theorem}

\begin{proof}
The argument proceeds in four steps, specializing the general
bar-cobar character formula
(Theorem~\ref{thm:weyl-kac-geometric}) to non-generic level.

\emph{Step~1: $E_1$ page.}
The $E_1$ page of the bar spectral sequence is built from
Verma modules for $\widehat{\mathfrak{sl}}_2$ at level~$k$:
\[
E_1^{n,*} = \bigoplus_{\substack{w \in \mathcal{W}_{\mathrm{aff}} \\
\ell(w) = n}} \mathcal{M}(w \cdot 0),
\]
where $\mathcal{W}_{\mathrm{aff}} = \mathbb{Z} \rtimes \mathbb{Z}/2$
is the affine Weyl group. At generic level, all Verma modules
$\mathcal{M}(w \cdot 0)$ are irreducible
(Proposition~\ref{prop:generic-irreducibility}), so $d_1 = 0$ and
the Weyl--Kac formula follows from $E_1 = E_\infty$.

\emph{Step~2: Singular vectors at admissible level.}
At level $k = -h^\vee + p/q$, the Verma modules are \emph{not}
irreducible: for the vacuum Verma module $\mathcal{M}(0)$, the
Shapovalov determinant vanishes at the positions determined by
the Kac determinant formula~\cite{Kac}:
\[
\det S_n = C_n \prod_{\substack{r,s \geq 1 \\ rs \leq n}}
(h_{r,s}(k) - h)^{P(n-rs)},
\qquad
h_{r,s}(k) = \frac{((k+2)r - s)^2 - (k+1)^2}{4(k+2)}.
\]
At $k = -h^\vee + p/q$ with integral $p,q$, the vanishing loci
$h_{r,s}(k) = 0$ occur at $r = 0 \bmod q$ or $s = 0 \bmod p$.
Each vanishing produces a singular vector in $\mathcal{M}(w \cdot 0)$,
yielding an embedding
$\mathcal{M}(w' \cdot 0) \hookrightarrow \mathcal{M}(w \cdot 0)$
realized by screening operator integrals on $\overline{C}_n(X)$.

\emph{Step~3: $d_1$ differentials and cancellation.}
These singular vector embeddings are precisely the $d_1$
differentials on the $E_1$ page. At
non-degenerate admissible level, the singular vectors organize
the infinite affine Weyl group $\mathcal{W}_{\mathrm{aff}}$ into
\emph{orbits} under the translations $w \mapsto w \cdot t_{q\alpha}$
(translations by $q$ times the coroot). Within each orbit, adjacent
Verma modules are connected by $d_1$, and the alternating signs
produce exact cancellation. The surviving elements on the $E_2$
page form a single orbit of the \emph{truncated} Weyl group
$\mathcal{W}_{\mathrm{adm}} = \{w \in \mathcal{W}_{\mathrm{aff}}
: \ell(w) < q\}$, which has exactly $q$ elements.

\emph{Step~4: Specialization to $k = -1/2$.}
For $k = -1/2$ (so $p = 3$, $q = 2$): the truncated Weyl group
has $|\mathcal{W}_{\mathrm{adm}}| = 2$ elements, namely $\{e, s_1\}$.
The $E_2$ page retains:
\begin{itemize}
\item $E_2^{0,*}$: the identity contribution $\Theta_{-1,6}(\tau,z)$,
\item $E_2^{1,*}$: the reflection contribution
$-\Theta_{5,6}(\tau,z)$,
\end{itemize}
and all higher bar degrees vanish.
Since $E_2^{n,*} = 0$ for $n \geq 2$, the spectral sequence
degenerates: $E_2 = E_\infty$. Dividing by the
universal denominator $\Theta_{1,1} - \Theta_{-1,1}$ yields
the vacuum character formula stated in
Example~\ref{ex:sl2-admissible}.
The non-degenerate admissible hypothesis ($X_{L_k} = \{0\}$)
enters through $C_2$-cofiniteness, which provides the finite
generation needed for the spectral sequence to terminate at a
finite page (Proposition~\ref{prop:bar-admissible}(i)).
\end{proof}

\begin{remark}[Admissible versus integrable characters]
\label{rem:admissible-vs-integrable}
Theorem~\ref{thm:kw-bar-spectral} clarifies the Weyl--Kac to Kac--Wakimoto passage:
\begin{center}
\renewcommand{\arraystretch}{1.3}
\begin{tabular}{lll}
\toprule
& \textbf{Integrable level} & \textbf{Admissible level} \\
\midrule
$E_1$ page & $\bigoplus_{w \in \mathcal{W}} \mathcal{M}(w \cdot \lambda)$
& same \\
$d_1$ & $0$ (Verma modules irreducible) & $\neq 0$ (singular vectors) \\
$E_2$ page & $= E_1$ (full $\mathcal{W}$ sum) &
$\mathcal{W}_{\mathrm{adm}}$ survivors \\
Result & Weyl--Kac & Kac--Wakimoto \\
\bottomrule
\end{tabular}
\end{center}
On the manuscript's bar surface, admissible-level singular-vector
cancellation leaves a finite admissible Weyl packet; for the fully
worked $\widehat{\mathfrak{sl}}_2$ example at $k=-1/2$ this is exactly
the two-term formula of Theorem~\ref{thm:kw-bar-spectral}. In general
rank, the resulting finite-sum character formula is the external
Kac--Wakimoto theorem recorded below.
\end{remark}

\begin{theorem}[Kac--Wakimoto character formula in general rank;
\ClaimStatusProvedElsewhere]
\label{thm:kw-bar-general-rank}
\index{Kac--Wakimoto character formula!general rank}
\index{admissible level!character formula!general rank}
Let $\mathfrak{g}$ be a simple Lie algebra of rank~$r$ with
positive roots $\Delta^+$ and affine Weyl group
$\mathcal{W} = W \ltimes Q^\vee$ \textup{(}$W$ finite,
$Q^\vee$ coroot lattice\textup{)}. At non-degenerate
admissible level $k = -h^\vee + p/q$ \textup{(}$\gcd(p,q) = 1$,
$p \geq h^\vee$ if $q = 1$\textup{)}, the vacuum character of
$L_k = L_k(\mathfrak{g})$ is given by\textup{:}
\begin{equation}\label{eq:kw-general}
\operatorname{ch}(L_k)(\tau, \mathbf{z})
= \frac{\sum_{w \in \mathcal{W}_{\mathrm{adm}}}
 \varepsilon(w)\, e^{w(\rho) - \rho}\,
 q^{|w(\rho) - \rho|^2 / 2(k + h^\vee)}}
{\prod_{\alpha \in \hat{\Delta}^+}
 (1 - e^{-\alpha})^{\mathrm{mult}(\alpha)}},
\end{equation}
where $\mathcal{W}_{\mathrm{adm}} = \{w \in \mathcal{W}
: w(\hat{C}_k) \cap \hat{C}_k^+ \neq \varnothing\}$ is
the admissible Weyl group, determined by the affine Weyl
chambers at level~$k$, and $\hat{\Delta}^+$ are the
positive affine roots. The cardinality satisfies
$|\mathcal{W}_{\mathrm{adm}}| = |W| \cdot q^r$
\textup{(}for $q \geq 2$\textup{)}.
\end{theorem}

\begin{proof}
This is the external Kac--Wakimoto character formula
\cite{KW88}. On the manuscript's bar-complex surface, the
preceding admissible analysis supplies the expected mechanism:
singular-vector embeddings produce the first bar differential,
and the surviving packet is organized by the admissible Weyl
group. What is \emph{not} proved here is a full internal
general-rank derivation of formula~\eqref{eq:kw-general} from the
bar spectral sequence alone, or a general-rank theorem identifying
$E_2$ with $E_\infty$ on that surface.
\end{proof}

\begin{remark}[Admissible module counts]
\label{rem:admissible-module-count}
\index{admissible level!module count}
At admissible $k = -h^\vee + p/q$ ($\gcd(p,q) = 1$, $p \geq h^\vee$),
the admissible weight lattice~\cite{KW88} gives, for
$\mathfrak{sl}_2$:
\[
\bigl|\mathrm{Adm}_k(\mathfrak{sl}_2)\bigr|
\;=\; (p-1)(q-1).
\]
At $k = -1/2$ ($p = 3$, $q = 2$): $|\mathrm{Adm}_{-1/2}| = 2$. For general simple~$\mathfrak{g}$, the count involves the affine Weyl group orbit on the $(p,q)$-restricted alcove~\cite{KW88}.
\end{remark}

\begin{proposition}[Admissible \texorpdfstring{$S$}{S}-matrix and Verlinde fusion package at \texorpdfstring{$k=-1/2$}{k=-1/2};
\ClaimStatusConditional]
\label{prop:admissible-verlinde-bar}
\index{Verlinde formula!admissible level}
\index{admissible level!fusion rules}
\index{bar complex!fusion rules}
For $\widehat{\mathfrak{sl}}_2$ at admissible level
$k = -1/2$ $(p = 3, q = 2)$, the bar spectral sequence
computes the vacuum admissible character. Together with the
standard admissible character package for the second simple
module and the modular transformation law, this determines the
modular $S$-matrix; together with the standard semisimple
Verlinde formalism for this rational admissible category, this
yields the fusion ring as follows.
\begin{enumerate}[label=\textup{(\roman*)}]
\item \emph{S-matrix.} The two admissible simples
 $L_0 = L_{-1/2}(\mathfrak{sl}_2)$ (vacuum) and
 $L_1$ (non-vacuum, $h = 3/8$) have modular S-matrix
 \begin{equation}\label{eq:admissible-s-matrix}
 S = \frac{1}{\sqrt{2}}
 \begin{pmatrix}
 1 & 1 \\[3pt]
 1 & -1
 \end{pmatrix}.
 \end{equation}
 This is the unique unitary symmetric matrix satisfying
 $(SS^*)^2 = I$ and $S_{0j} > 0$ for all~$j$, determined
 by the modular transformation of the vacuum character from
 Theorem~\textup{\ref{thm:kw-bar-spectral}} together with the
 standard admissible character formula for the second simple
 module \textup{(}Kac--Wakimoto~\cite{KW88}\textup{)}.
\item \emph{Fusion rules.} The Verlinde formula gives
 \begin{equation}\label{eq:admissible-fusion}
 L_0 \times L_i = L_i, \qquad
 L_1 \times L_1 = L_0.
 \end{equation}
 The fusion ring is $\mathbb{Z}[\mathbb{Z}/2]$.
\item \emph{Complementarity.} The Koszul dual at level
 $k' = -k - 2h^\vee = -7/2$ has $c' = 7$ and
 $c + c' = 6 = 2\dim(\mathfrak{sl}_2)$;
 see Theorem~\ref{thm:central-charge-complementarity}.
 Since $k' + h^\vee = -3/2 < 0$, the dual level does not
 admit an admissible parametrization
 $k' = -h^\vee + p'/q'$ with $p', q' > 0$. The associated variety
 $X_{L_{-7/2}} = \mathcal{N}$ is the full nilpotent cone,
 so the dual category is logarithmic (not rational).
 Complementarity operates between the rational admissible
 category at~$k$ and the logarithmic category at~$k'$.
\end{enumerate}
\end{proposition}

\begin{proof}
\emph{Step~1: Characters used in the modular calculation.}
By Theorem~\ref{thm:kw-bar-spectral}, the bar spectral
sequence at $k = -1/2$ computes the vacuum character
$\chi_0(\tau, z) = (\Theta_{-1,6} - \Theta_{5,6})/
(\Theta_{1,1} - \Theta_{-1,1})$. For the second admissible
simple module $L_1$, we use the standard admissible character
formula \textup{(}Kac--Wakimoto~\cite{KW88}\textup{)}:
$\chi_1(\tau, z) = (\Theta_{1,6} - \Theta_{-5,6})/
(\Theta_{1,1} - \Theta_{-1,1})$.

\emph{Step~2: Modular transformation.}
Under $\tau \to -1/\tau$, the theta functions transform as
\[
\Theta_{j,m}(-1/\tau, z/\tau)
= (-i\tau)^{1/2}\, e^{\pi i m z^2/\tau}\,
(2m)^{-1/2}\sum_{j'=0}^{2m-1}
e^{-\pi i j j'/m}\, \Theta_{j',m}(\tau, z).
\]
The S-transformation of the characters
$\chi_i(-1/\tau, z/\tau) = \sum_j S_{ij}\, \chi_j(\tau, z)$
yields the S-matrix~\eqref{eq:admissible-s-matrix}.
Explicitly, the unnormalized entries are
$\tilde{S}_{ij} = \sin(\pi(i+1)(j+1)/3)$ for
$i, j \in \{0, 1\}$; after dividing by
$\sqrt{\sum_j \tilde{S}_{0j}^2} = \sqrt{3/2}$
one obtains~\eqref{eq:admissible-s-matrix}.

\emph{Step~3: Verlinde formula on the rational admissible category.}
By Theorem~\ref{thm:admissible-rep-theory}(ii), the level
$k=-1/2$ category is rational. Applying the usual semisimple
Verlinde formula in that setting \textup{(}cf.~\cite{HLZ}\textup{)}
to this $S$-matrix gives the fusion coefficients
$N_{ij}^k = \sum_\ell S_{i\ell} S_{j\ell}
S_{\ell k}^*/S_{0\ell}$.
For the $2 \times 2$ S-matrix above:
$N_{11}^0 = S_{10}^2 S_{00}^*/S_{00}
+ S_{11}^2 S_{01}^*/S_{01} = 1$,
confirming $L_1 \times L_1 = L_0$.
The quantum dimensions are $d_0 = 1$ and $d_1 = 1$
($L_1$ is a simple current of order~$2$).
\end{proof}

\begin{remark}[Admissible Verlinde and complementarity]
\label{rem:admissible-verlinde-complementarity}
The genus-$g$ conformal block dimension for $L_{-1/2}(\mathfrak{sl}_2)$
with no insertions is $\dim \mathbb{V}(\Sigma_g) = d_0^{2-2g}
+ d_1^{2-2g} = 2$ for all $g \geq 1$. By
Proposition~\ref{conj:conformal-block-duality}(ii), any equality with
the corresponding count for~$\mathcal{A}^!$ would require a separate
rational braided categorical comparison on both sides. Since that
extra input is unavailable here, and in particular because the dual
level fails the required finiteness hypothesis below, no equal-rank
theorem follows from the general Koszul package alone.
Since $L_{-7/2}(\mathfrak{sl}_2)$ is \emph{not}
$C_2$-cofinite, the conformal blocks at the dual level
require the logarithmic formalism
(Huang--Lepowsky--Zhang~\cite{HLZ}), and the
``dimension'' is replaced by the rank of the space of
pseudo-traces (Miyamoto~\cite{Miyamoto04}).
This asymmetry between rational $k$ and logarithmic
$k'$ is characteristic of non-self-dual admissible pairs.
\end{remark}

\begin{remark}[Non-degenerate admissible level: what is proved]
\label{rem:admissible-rationality-finiteness}
\index{admissible level!rationality and finiteness}
\index{Koszul property!admissible level scope}
At non-degenerate admissible level $k = -h^\vee + p/q$, the
simple affine vertex algebra $L_k(\fg)$ is rational and
$C_2$-cofinite
\textup{(}Theorem~\ref{thm:admissible-rep-theory}(ii),
Theorem~\ref{thm:arakawa-rationality}\textup{)}.
Accordingly the bar and Li--bar spectral sequences have only
finitely many non-zero terms in each weight sector, and the
bar analysis is compatible with the Kac--Wakimoto character
formulas: in the worked $\widehat{\mathfrak{sl}}_2$ case at
$k=-1/2$ this is proved by finite-page degeneration
\textup{(}Theorem~\ref{thm:kw-bar-spectral}\textup{)}, while in
general rank the character formula itself is the external theorem
\textup{(}Theorem~\ref{thm:kw-bar-general-rank}\textup{)}.

What is \emph{not} proved here, for general~$\fg$, is the stronger
statement that $L_k(\fg)$ is chirally Koszul. The ordinary
semisimplicity of the admissible module category controls Zhu theory
and conformal blocks, but it does not identify the Ext groups
computed by the bar complex with Ext groups in the ordinary module
category. Hence Theorems~A--D$+$H are not promoted to simple
admissible quotients on the basis of rationality alone.
For $\fg = \mathfrak{sl}_2$ the promotion succeeds
\textup{(}Remark~\textup{\ref{rem:admissible-koszul-status}}\textup{)};
for $\fg = \mathfrak{sl}_3$ at denominator $q \geq 3$ it fails
\textup{(}Theorem~\textup{\ref{thm:admissible-sl3-non-koszul-qge3}}\textup{)}.
The obstruction is rank-dependent: the abelian Cartan subalgebra
$\fh \subset \fg$ contributes $\mathrm{rk}(\fg)$ classes in
$H^2(\barB(L_k(\fg)))$ that survive the bar spectral sequence
because $[\fh, \fh] = 0$ annihilates the Lie-bracket differential.
At rank~$1$ no Cartan obstruction appears and the single null
vector is absorbed; at rank $\geq 2$ the Cartan classes are
permanent cocycles.
\end{remark}

\begin{remark}[Admissible-level evidence and the remaining gap]
\label{rem:two-routes-admissible-koszul}
\index{admissible level!Koszulness routes}
For $L_k(\fg)$ at non-degenerate admissible level there are two
useful pieces of evidence:
\begin{enumerate}[label=\textup{(\roman*)}]
\item \emph{Ordinary-category finiteness:}
 rationality and semisimple Zhu algebra give finite ordinary
 representation theory and finite-dimensional character data.
\item \emph{Spectral-sequence control:}
 on the $\widehat{\mathfrak{sl}}_2$ worked surface
 \textup{(}Theorem~\ref{thm:kw-bar-spectral}\textup{)},
 singular-vector-induced $d_1$ differentials cancel all but the
 admissible Weyl packet and the spectral sequence collapses at
 $E_2$; in general rank, the corresponding finite-sum character
 formula is the external Kac--Wakimoto theorem
 \textup{(}Theorem~\ref{thm:kw-bar-general-rank}\textup{)}.
\end{enumerate}
These two inputs are consistent with each other, but they do not
close the gap to chiral Koszulness. The missing step is an
identification of the bar-complex Ext theory with the Ext theory of
the semisimple ordinary admissible category; the present manuscript
does not supply that identification.

For $\fg = \mathfrak{sl}_2$, the gap is closed: $L_k(\mathfrak{sl}_2)$
is chirally Koszul at every admissible level
\textup{(}Remark~\textup{\ref{rem:admissible-koszul-status}}\textup{)}.
For $\fg = \mathfrak{sl}_3$ at denominator $q \geq 3$, the answer is
\emph{negative}: $\mathrm{rk}(\fg) = 2$ classes from the abelian Cartan
subalgebra survive in $H^2(\barB(L_k))$
\textup{(}Remark~\textup{\ref{rem:admissible-koszul-status}},
Theorem~\textup{\ref{thm:admissible-sl3-non-koszul-qge3}}\textup{)}.
\end{remark}

\begin{remark}[Admissible $\mathfrak{sl}_3$: structural
$d_1$ surjectivity]
\label{rem:admissible-sl3-structural}
\index{admissible level!sl3@$\mathfrak{sl}_3$ structural surjectivity}
For $\fg = \mathfrak{sl}_3$, the $d_1$ differential on the
bar spectral sequence $E_1$ page acts by the Chevalley--Eilenberg
differential of~$\fn_-$. The positive root system
$\Phi^+ = \{\alpha_1, \alpha_2, \alpha_1{+}\alpha_2\}$ has
$|\Phi^+| = 3$ elements, while $\mathrm{rk}(\fg) = 2$. The
surplus $|\Phi^+| - \mathrm{rk}(\fg) = 1$ is sufficient only on the
low-denominator checks where the null vectors remain generated by
single root-pair couplings. It is not a uniform proof across the
admissible lane. At denominator $q \geq 3$ the null vector structure
involves multi-generator compositions that the surplus cannot absorb;
Theorem~\ref{thm:admissible-sl3-non-koszul-qge3} gives the resulting
non-Koszulness for $L_k(\mathfrak{sl}_3)$.
\end{remark}

\begin{conjecture}[Admissible rank obstruction: exact count]
\label{conj:admissible-rank-obstruction}%
\ClaimStatusConjectured
\index{admissible level!rank obstruction}%
\index{Cartan subalgebra!exact bar obstruction}%
Let $\fg$ be a simple finite-dimensional complex Lie algebra with
$\mathrm{rk}(\fg) \geq 2$, and let $k = p/q - h^\vee$ be a
non-degenerate admissible level for $\widehat{\fg}$ with
$\gcd(p,q) = 1$ and $q \geq 3$. Then
\begin{equation}\label{eq:admissible-rank-h2}
\dim H^2\!\bigl(\barB(L_k(\fg))\bigr) \;=\; \mathrm{rk}(\fg).
\end{equation}
In particular the bar cohomology $H^2$ jumps from its generic value
\textup{(}attained on the open dense locus of non-admissible levels;
$k + h^\vee = p/q \neq 0$, so the affine KZ coefficient
$r^{\mathrm{KZ}}(z)=\Omega/\bigl((k{+}h^\vee)z\bigr)$ is well-defined
and the generic Chevalley--Eilenberg count applies\textup{)} down to
exactly
$\mathrm{rk}(\fg)$ at admissible levels with $q \geq 3$.
\end{conjecture}

\begin{remark}[Mechanism and base case]
\label{rem:admissible-rank-obstruction-mechanism}
The mechanism is the abelianness of the Cartan subalgebra
$\fh \subset \fg$. The Poisson differential $d_1$ on the bar
spectral sequence is the Chevalley--Eilenberg differential of
$\fn_-$; since $[\fh, \fh] = 0$, the Cartan generators lie in
$\ker(d_1)$. At admissible level, the Kac--Wakimoto
decomposition~\cite{KW88} produces $\mathrm{rk}(\fg)$ singular
vectors aligned with~$\fh$ that descend to permanent classes in
$H^2(\barB(L_k(\fg)))$ and obstruct the PBW degeneration
underlying chiral Koszulness.

The base case $\fg = \mathfrak{sl}_2$ is \emph{not} a
counterexample. At rank~$1$ the single Cartan class is absorbed
by a higher differential $d_r$ ($r \geq 2$), as made explicit in
the $k = -1/2$ analysis \textup{(}Example~\ref{ex:sl2-admissible},
Theorem~\ref{thm:kw-bar-spectral}\textup{)}; the conjecture
excludes rank~$1$ for this reason, and
$L_k(\mathfrak{sl}_2)$ remains chirally Koszul at every admissible
level \textup{(}Remark~\ref{rem:admissible-koszul-status}\textup{)}.

The first nontrivial case is $\fg = \mathfrak{sl}_3$ at
denominator $q \geq 3$, where~\eqref{eq:admissible-rank-h2}
predicts $\dim H^2 = 2$. The two permanent Cartan classes
obstruct PBW and
$L_k(\mathfrak{sl}_3)$ is \emph{not} chirally Koszul
\textup{(}Theorem~\ref{thm:admissible-sl3-non-koszul-qge3}
supplies the inequality
$\dim H^2 \geq \mathrm{rk}(\fg)$ for $\fg = \mathfrak{sl}_3$;
the present conjecture sharpens this to equality, and
Conjecture~\ref{conj:admissible-koszul-rank-obstruction} extends
the inequality to $\mathrm{rk}(\fg) \geq 3$\textup{)}. The local computation
\verb|compute/tests/test_admissible_koszul_rank2_engine.py|
verifies the $\mathfrak{sl}_3$ prediction explicitly at a sweep of admissible levels
$k = -3 + p/q$ with $q \in \{3, 4, 5\}$, including the
Cartan-class count and its stability under the surviving
Kac--Wakimoto Weyl packet.
\end{remark}

\begin{remark}[Ribbon structure at admissible level]
\label{rem:ribbon-admissible}
\index{admissible level!ribbon category}
\index{ribbon category!admissible level}
Creutzig--McRae--Yang~\cite{Creutzig24} prove that the braided tensor
category of weight modules for $L_k(\mathfrak{sl}_2)$ at admissible
levels carries a ribbon structure. This is compatible with the
admissible Koszulness result for $\widehat{\mathfrak{sl}}_2$
(Remark~\ref{rem:admissible-rationality-finiteness}): the ribbon
twist is an additional datum on the braided category that the bar
complex respects. Whether the ribbon structure extends to rank
$\geq 2$ (i.e., to $L_k(\fg)$ for $\mathrm{rk}(\fg) \geq 2$ at
admissible level) remains open.
\end{remark}

\begin{proposition}[DS reduction at admissible level \cite{Arakawa17}; \ClaimStatusProvedElsewhere]
\label{prop:ds-admissible}
\index{Drinfeld--Sokolov reduction!admissible level}
At admissible level $k = -h^\vee + p/q$, the DS reduction
(Definition~\textup{\ref{def:w-module-ds}}) intertwines with
the bar complex\textup{:}
\begin{enumerate}[label=\textup{(\roman*)}]
\item $H^0_{\mathrm{DS}}$ is exact on the category of
 admissible $\widehat{\mathfrak{g}}_k$-modules
 \textup{(}Arakawa~\cite{Arakawa17}\textup{)}.
\item The DS reduction of the bar complex of an admissible
 module $M$ is quasi-isomorphic to the bar complex of the
 DS-reduced module\textup{:}
 $H^0_{\mathrm{DS}}(\bar{B}(\widehat{\mathfrak{g}}_k, M))
 \simeq \bar{B}(\mathcal{W}^k(\mathfrak{g}),
 H^0_{\mathrm{DS}}(M))$.
\item At non-degenerate admissible level, $H^0_{\mathrm{DS}}$
 sends the finite set of admissible simples of
 $\widehat{\mathfrak{g}}_k$ to the finite set of simples of
 the rational W-algebra $\mathcal{W}^k(\mathfrak{g})$.
\end{enumerate}
\end{proposition}

\begin{proof}
For (i): Arakawa~\cite{Arakawa17} proves that $H^i_{\mathrm{DS}}(M) = 0$
for $i \neq 0$ when $M$ is an admissible module. The key
ingredient is the Whittaker condition: the BRST operator
$Q_{\mathrm{DS}}$ is acyclic in positive degrees on the
admissible subcategory because the singular vectors in
$M$ generate an acyclic subcomplex of the BRST complex.

For (ii): the bar complex $\bar{B}(\widehat{\mathfrak{g}}_k, M)$
is a resolution of $M$ in the module category. Applying
$H^0_{\mathrm{DS}}$ termwise to the bar complex and using the
exactness (i) on each term yields a complex whose cohomology
is $H^0_{\mathrm{DS}}(M) = M^W$. By the double complex
argument (Theorem~\ref{thm:ds-koszul-intertwine}), this
coincides with $\bar{B}(\mathcal{W}^k, M^W)$.

For (iii): at non-degenerate admissible level, both
$L_k(\mathfrak{g})$ and $\mathcal{W}^k(\mathfrak{g})$ are
rational. The DS functor establishes a bijection between
their finite sets of simples (since DS is exact and sends
simples to simples or zero, and the non-vanishing is
guaranteed by the non-degeneracy condition).
\end{proof}

\subsection{Whittaker modules and the bar complex}
\label{sec:whittaker-modules}
\index{Whittaker module|textbf}

The Drinfeld--Sokolov reduction is realized concretely via
Whittaker modules, which provide the free field resolution
underlying the BRST complex.

\begin{definition}[Whittaker module]
\label{def:whittaker-module}
\index{Whittaker module!definition}
Let $\mathfrak{g}$ be a simple Lie algebra with
$\mathfrak{g} = \mathfrak{n}_- \oplus \mathfrak{h} \oplus
\mathfrak{n}_+$ a triangular decomposition, and let
$\chi \colon \mathfrak{n}_+ \to \mathbb{C}$ be a Lie algebra
homomorphism (necessarily, $\chi$ factors through
$\mathfrak{n}_+ / [\mathfrak{n}_+, \mathfrak{n}_+] \cong
\bigoplus_{i=1}^r \mathbb{C} e_i$, where $e_i$ are the
simple root vectors).

A $\widehat{\mathfrak{g}}_k$-module $\mathcal{M}$ is a
\emph{Whittaker module} of type $\chi$ if the action of
$\mathfrak{n}_+ \otimes \mathbb{C}[t]$ satisfies:
\begin{enumerate}[label=\textup{(\roman*)}]
\item $(e_i)_0 \cdot v = \chi(e_i) \cdot v$ for all
 $v \in \mathcal{M}$ and each simple root vector $e_i$;
\item $\mathfrak{n}_+ \otimes t\mathbb{C}[t]$ acts locally
 nilpotently on $\mathcal{M}$.
\end{enumerate}
The character $\chi$ is \emph{non-degenerate} if
$\chi(e_i) \neq 0$ for all $i = 1, \ldots, r$.
\end{definition}

\begin{proposition}[Whittaker modules and DS reduction \cite{Arakawa17};
\ClaimStatusProvedElsewhere]
\label{prop:whittaker-ds}
\index{Whittaker module!Drinfeld--Sokolov}
\index{Drinfeld--Sokolov reduction!Whittaker model}
Let $\chi$ be a non-degenerate character of $\mathfrak{n}_+$.
\begin{enumerate}[label=\textup{(\roman*)}]
\item The DS reduction
 $H^0_{\mathrm{DS}}(\widehat{\mathfrak{g}}_k)$ is computed by
 the BRST complex
 $(\widehat{\mathfrak{g}}_k \otimes F_\chi \otimes
 \bigwedge^\bullet \mathfrak{n}_+,\, Q_{\mathrm{DS}})$,
 where the ghost system consists of fermionic fields
 $(c_\alpha, c^*_\alpha)$ for each positive root $\alpha > 0$,
 and the BRST differential
 $Q_{\mathrm{DS}} = Q_{\mathrm{st}} + \chi$ consists of the
 standard Lie algebra cohomology differential
 $Q_{\mathrm{st}} = \sum_{\alpha > 0} e_\alpha \otimes
 c^*_\alpha - \frac{1}{2}\sum_{\alpha,\beta,\gamma}
 f^{\gamma}_{\alpha\beta}\, c^*_\alpha c^*_\beta c_\gamma$
 (where $f^{\gamma}_{\alpha\beta}$ are the structure constants
 of $\mathfrak{n}_+$) plus the Whittaker twist $\chi$.
\item \textup{(}Feigin--Frenkel~\cite{FF}\textup{)} At critical
 level $k = -h^\vee$, the DS reduction of the
 vacuum module $V_{-h^\vee}(\mathfrak{g})$ produces the center
 $\mathfrak{z}(\widehat{\mathfrak{g}})$, which is isomorphic
 to $\mathrm{Fun}\,\mathrm{Op}_{\mathfrak{g}^\vee}(D^\times)$
 \textup{(}opers for the Langlands dual group on the
 punctured disk\textup{)}.
\item \textup{(}Arakawa~\cite{Arakawa17}\textup{)} For a
 highest weight module $\mathcal{L}(\lambda)$ at admissible
 level, $H^0_{\mathrm{DS}}(\mathcal{L}(\lambda))$ is either
 a simple $\mathcal{W}^k$-module or zero; the non-vanishing
 is determined by the admissible weight conditions.
\end{enumerate}
\end{proposition}

\begin{proof}
Part (i) is the definition of the BRST/DS complex; see
\cite{FF} and \cite{KRW}. The non-degeneracy of $\chi$
ensures that the Whittaker twist renders $Q_{\mathrm{DS}}$
acyclic in positive ghost degrees.

Part (ii) is the fundamental theorem of Feigin--Frenkel
\cite{FF}: at critical level, the BRST cohomology of the
vacuum module reduces to the Poisson center of
$\widehat{\mathfrak{g}}_{-h^\vee}$, which is identified with
the algebra of functions on the space of $\mathfrak{g}^\vee$-opers
on the formal disc
via the Feigin--Frenkel isomorphism.

Part (iii): Arakawa proves that $H^i_{\mathrm{DS}} = 0$ for
$i \neq 0$ on admissible modules, and the $H^0$ is simple
when non-zero (since DS sends composition series to
composition series and the admissible simples form a finite
collection).
\end{proof}

\begin{proposition}[Bar complex via Whittaker resolution;
\ClaimStatusProvedHere]
\label{prop:bar-whittaker}
\index{Whittaker module!bar complex}
\index{bar complex!Whittaker resolution}
Let $M$ be a module in the admissible category
$\mathrm{KL}_k$ at admissible level~$k$ and
$\chi$ a non-degenerate character.
The bar complex $\bar{B}(\widehat{\mathfrak{g}}_k, M)$
admits a Whittaker filtration: there is a spectral sequence
\begin{equation}\label{eq:whittaker-ss}
E_1^{p,q} = H^q_{\mathrm{DS}}(\bar{B}_p(M))
\implies H^{p+q}\bigl(\mathrm{Tot}(\bar{B}(M) \otimes
\textstyle\bigwedge^\bullet \mathfrak{n}_+,\,
d_{\mathrm{bar}}, Q_{\mathrm{DS}})\bigr).
\end{equation}
When $H^i_{\mathrm{DS}}(\bar{B}_p(M)) = 0$ for $i \neq 0$
and all~$p$, the spectral sequence degenerates at $E_2$
and the abutment is identified with $H^*(\bar{B}_W(M^W))$,
where $M^W = H^0_{\mathrm{DS}}(M)$
\textup{(}Theorem~\textup{\ref{thm:ds-koszul-intertwine})}.
\end{proposition}

\begin{proof}
Filter the double complex
$\mathrm{Tot}(\bar{B}_p(M) \otimes \bigwedge^q \mathfrak{n}_+,
d_{\mathrm{bar}}, Q_{\mathrm{DS}})$ by bar degree~$p$.
The resulting spectral sequence has $E_0^{p,q} =
\bar{B}_p(M) \otimes \bigwedge^q \mathfrak{n}_+$ with
$d_0 = Q_{\mathrm{DS}}$, giving $E_1^{p,q} =
H^q_{\mathrm{DS}}(\bar{B}_p(M))$. The $d_1$ differential is
the induced bar differential on DS cohomology.

For the degeneration: by hypothesis,
$H^i_{\mathrm{DS}}(M) = 0$ for $i \neq 0$. Since each
$\bar{B}_p(M)$ is built from tensor products within the
admissible category $\mathrm{KL}_k$ (which is closed under
tensor products by Arakawa~\cite{Arakawa17}),
the DS vanishing $H^i_{\mathrm{DS}} = 0$ for $i > 0$ extends
from $M$ to each $\bar{B}_p(M)$. Thus $E_1^{p,q} = 0$ for
$q > 0$, and the $d_1$ differential on the surviving row
$E_1^{p,0} = H^0_{\mathrm{DS}}(\bar{B}_p(M)) = \bar{B}_{W,p}(M^W)$
is the $\mathcal{W}^k$-bar differential, giving
$E_2^{p,0} = H^p(\bar{B}_W(M^W))$.
\end{proof}

\section{\texorpdfstring{Genus-$1$ generic check for $\widehat{\mathfrak{sl}}_2$}{Genus-1 generic check for sl-2}}
\label{sec:sl2-genus-one-computation}
\index{genus-1 generic check!Kac--Moody}
\index{Kac--Moody algebra!genus-1 generic check}

The genus-0 level-shifting duality
\((\widehat{\mathfrak{sl}}_{2,k})^! \simeq
\mathrm{CE}^{\mathrm{ch}}(\widehat{\mathfrak{sl}}_{2,-k-4})\)
(Theorem~\ref{thm:sl2-koszul-dual}) supplies the genus-$0$ input.  The
genus-$1$ surface below is a conditional generic check: the elliptic
propagator produces the shifted curvature, and bar-cobar inversion is
used only after the stated weight-completed Positselski and
PBW-Koszul hypotheses are imposed.

The existing genus-1 examples (Heisenberg (Example~\ref{ex:heisenberg-complementarity-explicit}) and Kac--Moody at critical level (Example~\ref{ex:kac-moody-complementarity-explicit})) are degenerate: the Heisenberg is abelian ($f^{abc}=0$, $h^\vee=0$), and the critical level forces the curvature to vanish. At generic level $k$, the curvature, non-abelian structure constants, and Kodaira--Spencer map interact non-trivially.

\subsection{Setup: genus-1 bar complex}

Fix an elliptic curve $E_\tau = \mathbb{C}/(\mathbb{Z}+\mathbb{Z}\tau)$ with $\operatorname{Im}\tau>0$. The genus-1 bar complex of $\widehat{\mathfrak{sl}}_{2,k}$ is:
\begin{equation}\label{eq:sl2-genus1-bar}
\barB^{(1),n}(\widehat{\mathfrak{sl}}_{2,k}) = \Gamma\bigl(\overline{C}_n(E_\tau),\; \mathfrak{sl}_2^{\boxtimes n} \otimes \omega_{E_\tau}^{\boxtimes n} \otimes \Omega^n_{\log}\bigr)
\end{equation}
where $\overline{C}_n(E_\tau)$ is the Fulton--MacPherson compactification of the configuration space of $n$ points on $E_\tau$.

The genus-1 propagator replacing $d\log(z_i - z_j)$ is (Remark~\ref{rem:wp-config-space}):
\begin{equation}\label{eq:genus1-propagator}
K^{(1)}(z,w|\tau) = \frac{\theta_1'(z-w|\tau)}{\theta_1(z-w|\tau)}
\end{equation}
where $\theta_1$ is the Jacobi theta function. This propagator has a simple pole at $z=w$ with residue~$1$, and satisfies:
\begin{align}
K^{(1)}(z+1,w|\tau) &= K^{(1)}(z,w|\tau) \label{eq:A-cycle-periodicity} \\
K^{(1)}(z+\tau,w|\tau) &= K^{(1)}(z,w|\tau) - 2\pi i \label{eq:B-cycle-quasi-periodicity}
\end{align}
The $A$-cycle periodicity~\eqref{eq:A-cycle-periodicity} is exact, but the $B$-cycle quasi-periodicity~\eqref{eq:B-cycle-quasi-periodicity} introduces the \emph{constant} shift $-2\pi i$. This constant is the source of the genus-$1$ curvature term computed below.

The differential on $\barB^{(1)}$ is:
\begin{equation}\label{eq:genus1-diff}
d^{(1)} = \sum_{i<j} \operatorname{Res}_{D_{ij}} \circ K^{(1)}_{ij}
\end{equation}
For $\widehat{\mathfrak{sl}}_{2,k}$, the residue extracts both the double-pole (level) and simple-pole (structure constant) contributions from the OPE, exactly as in the genus-0 computation~\eqref{eq:bar-diff-general}, but with the genus-0 propagator $d\log(z_i-z_j)$ replaced by $K^{(1)}$.

\subsection{Curvature decomposition}

\begin{theorem}[Genus-1 curvature for \texorpdfstring{$\widehat{\mathfrak{sl}}_{2,k}$}{sl-hat_2,k}; \ClaimStatusProvedHere]
\label{thm:sl2-genus1-curvature}
The genus-$1$ fiberwise differential
$\dfib$ \textup{(}Convention~\textup{\ref{conv:higher-genus-differentials})} satisfies:
\begin{equation}\label{eq:sl2-genus1-dsquared}
\dfib^{\,2} = (k+2) \cdot \omega_1 \cdot \operatorname{id}_{\mathfrak{sl}_2}
\end{equation}
where $\omega_1 \in H^2(\overline{\mathcal{M}}_{1,1})$ is the fundamental class and $(k+2)$ is the shifted level $k + h^\vee$ for $\mathfrak{sl}_2$.

The curvature decomposes as:
\begin{equation}\label{eq:curvature-decomposition}
\dfib^{\,2} = \underbrace{k \cdot \omega_1 \cdot \operatorname{id}}_{\text{abelian (level)}} + \underbrace{h^\vee \cdot \omega_1 \cdot \operatorname{id}}_{\text{non-abelian (Casimir)}}
\end{equation}
The first summand is the Heisenberg-type contribution (proportional to $k$, independent of the Lie algebra structure), and the second is the genuinely non-abelian contribution (proportional to $h^\vee=2$, arising from the adjoint Casimir).
\end{theorem}

\begin{proof}
We compute $\dfib^{\,2}$ on a degree-1 element $J^a(z_1) \otimes J^b(z_2) \cdot K^{(1)}_{12}$ by tracking the $B$-cycle monodromy.

\emph{Step~1: Transport around the $B$-cycle.}
Consider the parallel transport of $K^{(1)}_{12} = K^{(1)}(z_1,z_2|\tau)$ as $z_1$ traverses the $B$-cycle $z_1 \mapsto z_1 + \tau$. The elliptic-kernel $B$-cycle quasi-periodicity gives
\begin{equation}
K^{(1)}(z_1+\tau,z_2|\tau) - K^{(1)}(z_1,z_2|\tau) = -2\pi i
\end{equation}
This constant shift, upon contraction with the OPE data of $\widehat{\mathfrak{sl}}_{2,k}$, produces the curvature.

\emph{Step~2: OPE contraction.}
The residue at $z_1 = z_2$ of $J^a(z_1) J^b(z_2) \cdot K^{(1)}_{12}$ extracts the same data as in genus~0: the double pole contributes $k \cdot (J^a, J^b)$ and the simple pole contributes $f^{ab}{}_c J^c$. Applying $\dfib$ a second time and using the $B$-cycle shift, we obtain:
\begin{equation}
\dfib^{\,2}(J^a \boxtimes J^b \cdot K^{(1)}_{12}) = (-2\pi i) \cdot \bigl(k \cdot (J^a, J^b) + \delta^{ab} \cdot h^\vee\bigr)
\end{equation}
The factor $h^\vee$ arises exactly as in the genus-0 curvature computation (Theorem~\ref{thm:universal-kac-moody-koszul}, Step~2): the contraction $\sum_c f^{ac}{}_d f^{bc}{}_e$ over the adjoint representation produces the eigenvalue $2h^\vee$ of the quadratic Casimir.

\emph{Step~3: Identification with $\omega_1$.}
The constant $(-2\pi i)$ from the $B$-cycle monodromy is the local connection coefficient whose Chern curvature represents the generator $\omega_1 = c_1(\mathbb{E}) \in H^2(\overline{\mathcal{M}}_{1,1})$, where $\mathbb{E}$ is the Hodge bundle. (The Hodge bundle on $\overline{\mathcal{M}}_{1,1}$ is the line bundle whose fiber over $[E_\tau]$ is $H^0(E_\tau, \omega_{E_\tau})$; its first Chern class is represented by the curvature of the natural connection, which is $-2\pi i \cdot d\tau \wedge d\bar{\tau}/(2i \cdot \operatorname{Im}\tau)^2$ in the $\tau$ coordinate.)

On the invariant form, the full contraction gives
$k \cdot (J^a, J^b) + h^\vee \cdot \delta^{ab}$, which on
$\mathfrak{sl}_2$ with the normalization ($(h,h) = 2$, $(e,f)=1$)
evaluates to $(k+h^\vee)$ times the identity on $\mathfrak{sl}_2$.
The quadratic Casimir in the adjoint representation satisfies
$C_2^{\mathrm{ad}} = 2h^\vee \cdot \operatorname{id}_{\mathfrak{g}}$
for any simple $\mathfrak{g}$, and the combined level + Casimir
contribution is:
\begin{equation}
k \cdot \kappa_{\mathrm{level}} + \frac{1}{2}C_2^{\mathrm{ad}} = k + h^\vee = k + 2
\end{equation}
on $\mathfrak{sl}_2$.

Combining: $\dfib^{\,2} = (k+2) \cdot \omega_1 \cdot \operatorname{id}_{\mathfrak{sl}_2}$.

\emph{Step~4: Abelian limit check.}
Setting $f^{abc} = 0$ and $h^\vee = 0$ (i.e., replacing $\mathfrak{sl}_2$ by a $3$-dimensional abelian Lie algebra with the same bilinear form), the non-abelian contribution vanishes and we recover:
\begin{equation}
\dfib^{\,2}\big|_{f=0} = k \cdot \omega_1 \cdot \operatorname{id}
\end{equation}
This is the abelian Heisenberg curvature
$\kappa \cdot \omega_1 \cdot \operatorname{id}$ with
$\kappa = k$ (the level of each Cartan generator).
\end{proof}

\begin{remark}[Critical level]
At $k = -h^\vee = -2$, the curvature $\dfib^{\,2} = 0$ vanishes (Convention~\ref{conv:higher-genus-differentials}), so the genus-1 bar complex becomes a genuine differential graded coalgebra. This is consistent with the enhanced algebraic structure at critical level (Theorem~\ref{thm:sl2-critical}): the Feigin--Frenkel center provides extra flat sections, and the bar differential is strictly nilpotent.
In the language of the universal MC class
(Theorem~\ref{thm:explicit-theta}), the critical level is
the unique point where the degree-$2$ projection vanishes
because $\kappa(\widehat{\mathfrak{sl}}_{2,-2}) = 0$. This kills the scalar
curvature and gives $\Delta = 8\kappa S_4 = 0$, but it does
not force the higher-degree part of $\Theta_{\cA}$ to vanish or
annihilate the higher-genus operations. The Feigin--Frenkel center
records the critical-level bar cohomology on this uncurved
surface; it is not obtained by setting the full MC element to
zero.
\end{remark}

\subsection{Spectral sequence and bar-cobar inversion}

\begin{theorem}[Genus-$1$ bar-cobar inversion criterion for \texorpdfstring{$\widehat{\mathfrak{sl}}_{2,k}$}{sl-hat_2,k}; \ClaimStatusConditional]
\label{thm:sl2-genus1-inversion}
Let $k+2$ avoid the Kac--Kazhdan resonant hyperplanes.  Assume the
genus-$1$ bar object is taken in the weight-completed Positselski
surface, the generic PBW-Koszul comparison identifies the $E_1$ page
with the genus-$0$ chiral Chevalley--Eilenberg cohomology, and this
comparison is concentrated in columns $0$ and~$1$.  Then the genus-$1$
bar-cobar adjunction map
\begin{equation}
\Omega\bigl(\barB^{(1)}(\widehat{\mathfrak{sl}}_{2,k})\bigr) \xrightarrow{\;\sim\;} \widehat{\mathfrak{sl}}_{2,k}
\end{equation}
is a quasi-isomorphism in that completed curved category. The
comparison spectral sequence has:
\begin{align}
E_1^{p,q} &= H^q\bigl(\mathfrak{sl}_2^{\otimes p},\, d_{\mathrm{OPE}}\bigr) \otimes H^*(\overline{\mathcal{M}}_{1,1}) \\
E_2^{p,q} &= \operatorname{Ext}^p_{U(\mathfrak{sl}_2)}(V_k^{\otimes q},\, \mathbb{C})
\end{align}
and collapses at $E_2$ on the stated surface.
\end{theorem}

\begin{proof}
The proof is an instance of Theorem~\ref{thm:higher-genus-inversion} (higher genus bar-cobar inversion), specialized to $g=1$ and $\mathfrak{g} = \mathfrak{sl}_2$. We verify the hypotheses.

\emph{Step~1: Weight filtration.}
The bar complex $\barB^{(1)}(\widehat{\mathfrak{sl}}_{2,k})$ carries the conformal weight filtration $F^p\barB^{(1)} = \bigoplus_{w \geq p} \barB^{(1)}_w$, where $w$ is the total conformal weight. Since all generators $e, f, h$ have conformal weight~$1$, the filtration is bounded below and exhaustive.

\emph{Step~2: $E_1$ page.}
The associated graded of $\dfib$ with respect to the weight filtration gives the $E_1$ page. At leading order, the genus-1 propagator $K^{(1)}(z,w|\tau) = 1/(z-w) + O(1)$ agrees with the genus-0 propagator $1/(z-w)$, so the $E_1$ differential is the genus-0 bar differential $\dzero$. Its cohomology computes the Lie algebra cohomology $H^*(\mathfrak{sl}_2, V_k)$.

\emph{Step~3: PBW-Koszul input.}
Away from the Kac--Kazhdan resonant hyperplanes the Shapovalov form has
no null-vector kernel in the vacuum module.  This excludes the
admissible singular-vector differentials that appear later, but it is
not by itself a proof of genus-$1$ bar-cobar inversion.  The additional
PBW-Koszul comparison hypothesis in the theorem is precisely the
statement that the relevant bar Ext groups satisfy:
\begin{equation}
\operatorname{Ext}^p_{U(\widehat{\mathfrak{sl}}_2)}(V_k, V_k) = 0 \quad \text{for } p \geq 2
\end{equation}
This is the Koszul concentration input.

\emph{Step~4: $E_2$ collapse.}
By the Ext-vanishing, the $E_2$ page has non-zero entries only in the strip $0 \leq p \leq 1$. Since $H^*(\overline{\mathcal{M}}_{1,1}) = \mathbb{Q}[\lambda]/(\lambda^2)$ is concentrated in degrees 0 and~2, the $E_2$ page is:
\begin{equation}
E_2^{p,q} = \begin{cases}
\mathbb{C} & (p,q) = (0,0) \\
\mathfrak{sl}_2^* & (p,q) = (1,0) \\
\mathbb{C}\lambda & (p,q) = (0,2) \\
\mathfrak{sl}_2^* \cdot \lambda & (p,q) = (1,2) \\
0 & \text{otherwise}
\end{cases}
\end{equation}
All higher differentials $d_r$ for $r \geq 2$ vanish for degree reasons (they would map between rows differing by $r \geq 2$, but the non-zero entries lie in rows 0 and~2 only). Therefore $E_2 = E_\infty$.

The quasi-isomorphism $\Omega(\barB^{(1)}) \xrightarrow{\sim} \widehat{\mathfrak{sl}}_{2,k}$ then follows from the completed inversion theorem (Theorem~\ref{thm:higher-genus-inversion}) under exactly these hypotheses.
\end{proof}

\begin{remark}[Admissible levels]\label{rem:sl2-admissible}
At admissible levels $k = -2 + p/q$ with $p, q \in \mathbb{Z}_{>0}$
coprime, the vacuum Verma module develops singular vectors in the
bar-relevant range. For $\widehat{\mathfrak{sl}}_2$, the first
vacuum null occurs at conformal weight $(p-1)q$ above the vacuum,
as in Example~\ref{ex:sl2-li-bar}. Thus the Kac--Shapovalov
criterion does not certify Koszulness for the simple quotient.
In the bar complex, the null vectors contribute extra extension data.

The ordinary admissible category splits into two distinct cases
\textup{(}Theorem~\textup{\ref{thm:admissible-rep-theory}}(ii)--(iii)\textup{)}.
At non-degenerate admissible levels such as $k=-1/2$, it is
rational with finitely many simple objects. At degenerate
admissible levels such as $k=-4/3$, it is no longer semisimple and
the associated variety is the nilcone. This is not a contradiction:
the Ext groups detected by the bar complex live in a larger
chiral/factorization setting than the ordinary admissible module
category, and even the ordinary category itself changes across the
non-degenerate/degenerate boundary. The precise comparison remains
part of the admissible-level analysis.
\end{remark}

\subsection{Kodaira--Spencer map and complementarity}

\begin{theorem}[Genus-1 complementarity for \texorpdfstring{$\widehat{\mathfrak{sl}}_{2,k}$}{sl-hat_2,k} at generic level; \ClaimStatusConditional]
\label{thm:sl2-genus1-complementarity}
For generic $k \neq -2$, the quantum corrections at genus $1$ satisfy:
\begin{align}
Q_1(\widehat{\mathfrak{sl}}_{2,k}) &= \mathbb{C} \cdot (k+2) \subset H^0(\overline{\mathcal{M}}_{1,1}) \label{eq:Q1-sl2} \\
Q_1(\widehat{\mathfrak{sl}}_{2,-k-4}) &= \mathbb{C} \cdot \lambda \subset H^2(\overline{\mathcal{M}}_{1,1}) \label{eq:Q1-sl2-dual}
\end{align}
and these span $H^*(\overline{\mathcal{M}}_{1,1})$:
\begin{equation}\label{eq:sl2-complementarity}
Q_1(\widehat{\mathfrak{sl}}_{2,k}) \oplus Q_1(\widehat{\mathfrak{sl}}_{2,-k-4}) = H^*(\overline{\mathcal{M}}_{1,1}) = \mathbb{C} \oplus \mathbb{C}\lambda
\end{equation}
The Kodaira--Spencer map (Theorem~\ref{thm:kodaira-spencer-chiral-complete}) exchanges the two summands:
\begin{equation}
\kappa_{\mathrm{KS}}: Q_1(\widehat{\mathfrak{sl}}_{2,k}) \xrightarrow{\;\sim\;} Q_1(\widehat{\mathfrak{sl}}_{2,-k-4})^\vee
\end{equation}
\end{theorem}

\begin{proof}
This is an instance of Theorem~\ref{thm:quantum-complementarity-main} (Main Theorem~C), specialized to $\cA = \widehat{\mathfrak{sl}}_{2,k}$ at genus~$1$. We carry out each step explicitly.

\emph{Step~1: Identify $Q_1(\widehat{\mathfrak{sl}}_{2,k})$.}
From Theorem~\ref{thm:sl2-genus1-curvature}, the curvature of the genus-1 bar complex is:
\begin{equation}
(d^{(1)})^2 = (k+2) \cdot \omega_1 \cdot \operatorname{id}_{\mathfrak{sl}_2}
\end{equation}
The quantum correction $Q_1(\cA)$ is the scalar obstruction to strict
nilpotence of $d^{(1)}$ after pairing the genus-$1$ curvature class with
the Hodge generator.  For $\widehat{\mathfrak{sl}}_{2,k}$ at generic
level the non-derived chiral centre is $\mathbb{C}$; the level $k$ is a
parameter, not an additional central generator.  Thus
$(k+2)\omega_1$ contributes the one-dimensional scalar slot.

More precisely: the obstruction is the class of $(d^{(1)})^2$ in
$H^0(\overline{\mathcal{M}}_{1,1}) \otimes
\operatorname{End}(V_k)^{\widehat{\mathfrak{sl}}_2}$.  On the generic
universal surface this endomorphism algebra is scalar, so the curvature
represents a one-dimensional obstruction:
\begin{equation}
Q_1(\widehat{\mathfrak{sl}}_{2,k}) = \mathbb{C} \cdot (k+2) \subset H^0(\overline{\mathcal{M}}_{1,1})
\end{equation}

\emph{Step~2: Identify $Q_1(\widehat{\mathfrak{sl}}_{2,-k-4})$.}
The Koszul dual is
\(\mathrm{CE}^{\mathrm{ch}}(\widehat{\mathfrak{sl}}_{2,-k-4})\)
(Theorem~\ref{thm:sl2-koszul-dual}); its level-reflected current
presentation has shifted level $k^* = -k-4$. Repeating the curvature
computation with $k$ replaced by $k^*$:
\begin{equation}
(d^{(1)}_!)^2 = (k^* + 2) \cdot \omega_1 \cdot \operatorname{id} = (-k-2) \cdot \omega_1 \cdot \operatorname{id}
\end{equation}
Before the Kodaira--Spencer exchange, the level-reflected scalar is:
\begin{equation}
Q_1(\widehat{\mathfrak{sl}}_{2,-k-4}) = \mathbb{C} \cdot (-k-2) \subset H^0(\overline{\mathcal{M}}_{1,1})
\end{equation}

\emph{Step~3: Complementarity via the Kodaira--Spencer map.}
By the Kodaira--Spencer construction
(Theorem~\ref{thm:kodaira-spencer-chiral-complete}), the Verdier
involution $\sigma$ on the scalar Hodge lane exchanges the degree-$0$
and degree-$2$ components. Explicitly, $\sigma$ is defined via the
Verdier duality pairing on $\overline{\mathcal{M}}_{1,1}$:
\begin{equation}
\sigma: H^0(\overline{\mathcal{M}}_{1,1}) \xrightarrow{\;\sim\;} H^2(\overline{\mathcal{M}}_{1,1}), \quad 1 \mapsto \lambda
\end{equation}

The complementarity theorem places the $\cA$ scalar slot and the
$\cA^!$ scalar slot in complementary Hodge degrees through this Verdier
pairing.  No isomorphism of the full derived chiral centres is asserted
here.

For $\widehat{\mathfrak{sl}}_{2,k}$: $Q_1(\widehat{\mathfrak{sl}}_{2,k})$ lives in $H^0$ (the obstruction from curvature) and $\sigma$ maps it to $\mathbb{C} \cdot \lambda \in H^2$. The complementary summand is $Q_1(\widehat{\mathfrak{sl}}_{2,-k-4})^\vee = \mathbb{C} \cdot \lambda$, which lives in $H^2$ (the deformation direction). Therefore:
\begin{equation}
Q_1(\widehat{\mathfrak{sl}}_{2,k}) \oplus Q_1(\widehat{\mathfrak{sl}}_{2,-k-4}) = \mathbb{C} \oplus \mathbb{C}\lambda = H^*(\overline{\mathcal{M}}_{1,1})
\end{equation}

\emph{Step~4: Consistency checks.}
\begin{enumerate}
\item \emph{Involutivity}: Swapping $k \leftrightarrow -k-4$ changes the
 scalar coefficient from $k+2$ to $-k-2$ and exchanges the Hodge
 placement.
\item \emph{Dimension count}: on the non-critical generic surface,
 $\dim Q_1 + \dim Q_1^! = 1 + 1 = 2 =
 \dim H^*(\overline{\mathcal{M}}_{1,1})$.
\item \emph{Abelian limit}: Setting the Lie bracket to zero and using
 the Heisenberg normalization recovers
 $Q_1(\mathcal{H}_\kappa) = \mathbb{C} \cdot \kappa$ and
 $Q_1(\mathcal{H}_\kappa^!) = \mathbb{C} \cdot \lambda$.
\item \emph{Critical level}: at $k=-2$ the scalar coefficient vanishes
 and the theorem no longer applies; the replacement object is the
 critical bar complex with Feigin--Frenkel centre described in
 Theorem~\ref{thm:oper-bar-h0}.
\end{enumerate}
\end{proof}

\subsection{Genus-1 partition function}

\begin{proposition}[Genus-$1$ scalar logarithm via complementarity;
\ClaimStatusConditional]
\label{prop:sl2-genus1-partition}
\index{partition function!genus-1!Kac--Moody}
Under the genus-universality hypothesis package, the complementarity
computation determines the scalar Hodge-line coefficient, not the full
affine character. After the oscillator determinant is removed,
\begin{equation}\label{eq:sl2-genus1-partition}
\left.\log Z_1^{\mathrm{sc}}(\widehat{\mathfrak{sl}}_{2,k})\right|_{\lambda_1}
= F_1(\widehat{\mathfrak{sl}}_{2,k})\,\lambda_1,
\qquad
F_1(\widehat{\mathfrak{sl}}_{2,k})
= \frac{\kappa(\widehat{\mathfrak{sl}}_{2,k})}{24}
= \frac{k+2}{32}.
\end{equation}
Here
\(\kappa(\widehat{\mathfrak{sl}}_{2,k})
= (k+h^\vee)\dim(\mathfrak{sl}_2)/(2h^\vee)
=3(k+2)/4\). The denominator $\eta^{-3}$ belongs to the free-current
oscillator trace; it does not determine the interacting affine
vacuum character.
\end{proposition}

\begin{proof}
For $\mathfrak{sl}_2$ one has $h^\vee=2$ and
$\dim(\mathfrak{sl}_2)=3$. The genus-universality theorem gives
\[
\kappa(\widehat{\mathfrak{sl}}_{2,k})
=\frac{(k+h^\vee)\dim(\mathfrak{sl}_2)}{2h^\vee}
=\frac{3(k+2)}4 .
\]
Part~\textup{(iii)} of Theorem~\ref{thm:genus-universality} gives
$F_1=\kappa/24$. The same shifted level $k+h^\vee$ is the curvature
coefficient in Theorem~\ref{thm:sl2-genus1-curvature}. Thus the
scalar logarithm sees the complementarity curvature, while the full
affine trace still contains the representation-theoretic character
data.
\end{proof}

\subsection{The three theorems for affine Kac--Moody algebras}

\begin{remark}[\texorpdfstring{$\widehat{\mathfrak{sl}}_{2,k}$}{sl-hat_2,k} at genus 1]
\label{rem:sl2-three-theorems}
\index{main theorems!genus-1 checks}
The genus-$1$ rank-one checks are: Theorem~A via the bar complex on
$E_\tau$ with propagator $K^{(1)} = \theta_1'/\theta_1$
(Theorem~\ref{thm:km-higher-genus-corrections}); Theorem~B via
$E_2$-collapse on the generic locus
(Theorem~\ref{thm:sl2-genus1-inversion}); and Theorem~C via
\[
Q_1(\widehat{\mathfrak{sl}}_{2,k}) \oplus Q_1(\widehat{\mathfrak{sl}}_{2,-k-4}) = H^*(\overline{\mathcal{M}}_{1,1})
\]
(Theorem~\ref{thm:sl2-genus1-complementarity}). All four genus-1 invariants are controlled by $k + h^\vee = k + 2$, manifesting the prism principle (Corollary~\ref{cor:prism-principle}).
\end{remark}

\section{\texorpdfstring{Genus-$1$ generic check for $\widehat{\mathfrak{sl}}_3$}{Genus-1 generic check for sl-3}}
\label{sec:sl3-genus-one-computation}
\index{genus-1 generic check!sl3@$\mathfrak{sl}_3$}

The genus-0 duality
\((\widehat{\mathfrak{sl}}_{3,k})^! \simeq
\mathrm{CE}^{\mathrm{ch}}(\widehat{\mathfrak{sl}}_{3,-k-6})\)
(Theorem~\ref{thm:sl3-koszul-dual}) uses $h^\vee = 3$.  The genus-$1$
calculation is again a generic check under the same completed
PBW-Koszul hypotheses as in the rank-one case. The key parameters are:
\begin{equation}\label{eq:sl3-parameters}
h^\vee = 3, \quad \dim(\mathfrak{sl}_3) = 8, \quad c = \frac{8k}{k+3}, \quad c + c' = 2\dim(\mathfrak{sl}_3) = 16
\end{equation}

At rank~1, the Cartan is one-dimensional and Serre relations are vacuous. At rank~2, the 8 generators $\{H_1, H_2, E_1, E_2, E_3, F_1, F_2, F_3\}$ produce non-trivial Serre relations in degree~3 (Computation~\ref{comp:sl3-bar}), the weight lattice is two-dimensional, and the curvature involves the full $A_2$ root system. The genus-1 theory is nevertheless governed by the single parameter $k + h^\vee = k + 3$.

\subsection{Genus-1 bar complex and curvature}

The genus-1 bar complex is given by~\eqref{eq:sl2-genus1-bar} with $\mathfrak{sl}_2$ replaced by $\mathfrak{sl}_3$:
\begin{equation}\label{eq:sl3-genus1-bar}
\barB^{(1),n}(\widehat{\mathfrak{sl}}_{3,k}) = \Gamma\bigl(\overline{C}_n(E_\tau),\; \mathfrak{sl}_3^{\boxtimes n} \otimes \omega_{E_\tau}^{\boxtimes n} \otimes \Omega^n_{\log}\bigr)
\end{equation}
with the same genus-1 propagator~\eqref{eq:genus1-propagator} and differential~\eqref{eq:genus1-diff}.

\begin{theorem}[Genus-1 curvature for \texorpdfstring{$\widehat{\mathfrak{sl}}_{3,k}$}{sl-hat_3,k}; \ClaimStatusProvedHere]
\label{thm:sl3-genus1-curvature}
The genus-$1$ fiberwise differential
$\dfib$ \textup{(}Convention~\textup{\ref{conv:higher-genus-differentials})} satisfies:
\begin{equation}\label{eq:sl3-genus1-dsquared}
\dfib^{\,2} = (k+3) \cdot \omega_1 \cdot \operatorname{id}_{\mathfrak{sl}_3}
\end{equation}
where $\omega_1 \in H^2(\overline{\mathcal{M}}_{1,1})$ is the fundamental class.

The curvature decomposes as:
\begin{equation}\label{eq:sl3-curvature-decomposition}
(d^{(1)})^2 = \underbrace{k \cdot \omega_1 \cdot \operatorname{id}}_{\text{abelian (level)}} + \underbrace{3 \cdot \omega_1 \cdot \operatorname{id}}_{\text{non-abelian (Casimir)}}
\end{equation}
\end{theorem}

\begin{proof}
The computation parallels Theorem~\ref{thm:sl2-genus1-curvature},
with the elliptic-kernel $B$-cycle shift $-2\pi i$ providing the
curvature mechanism.

\emph{Step~1: OPE contraction.}
For $\widehat{\mathfrak{sl}}_3$, the double-pole contribution to the residue is $k \cdot (J^a, J^b)$ (the normalized invariant form, with $(H_i, H_j) = A_{ij}$ and $(E_\alpha, F_\beta) = \delta_{\alpha\beta}$ as in Computation~\ref{comp:sl3-ope}), and the simple-pole contribution is $f^{ab}{}_{c} J^c$. Applying $d^{(1)}$ twice with the $B$-cycle shift, the self-contraction loop yields:
\begin{equation}
(d^{(1)})^2(J^a \boxtimes J^b \cdot K^{(1)}_{12}) = (-2\pi i) \cdot \bigl(k \cdot (J^a, J^b) + \tfrac{1}{2}\textstyle\sum_{c,d} f^{ac}{}_{d}\, f^{bc}{}_{d}\bigr)
\end{equation}
The factor $\frac{1}{2}$ is the symmetry factor of the one-loop self-contraction diagram ($|\mathrm{Aut}| = 2$), matching the genus-0 computation (Theorem~\ref{thm:universal-kac-moody-koszul}, Step~2).

\emph{Step~2: Casimir evaluation on $\mathfrak{sl}_3$.}
The contraction $\sum_{c,d} f^{ac}{}_{d}\, f^{bc}{}_{d}$ is the quadratic Casimir eigenvalue in the adjoint representation. For any simple $\mathfrak{g}$, $C_2^{\mathrm{ad}} = 2h^\vee \cdot \operatorname{id}_{\mathfrak{g}}$. For $\mathfrak{sl}_3$ with $h^\vee = 3$:
\begin{equation}
\textstyle\sum_{c,d} f^{ac}{}_{d}\, f^{bc}{}_{d} = 2 \cdot 3 \cdot (J^a, J^b) = 6\,(J^a, J^b)
\end{equation}
The structure constants from Computation~\ref{comp:sl3-ope} give this
identity explicitly: the sum runs over all $8$ generators and $6$
positive/negative roots.

\emph{Step~3: Combining.}
The combined level + Casimir contribution on $\mathfrak{sl}_3$ is:
\begin{equation}
k \cdot \kappa_{\mathrm{level}} + \tfrac{1}{2}C_2^{\mathrm{ad}} = k + h^\vee = k + 3
\end{equation}
giving $(d^{(1)})^2 = (k+3) \cdot \omega_1 \cdot \operatorname{id}_{\mathfrak{sl}_3}$.
\end{proof}

\begin{remark}[Critical level for \texorpdfstring{$\mathfrak{sl}_3$}{sl_3}]
At $k = -h^\vee = -3$, the curvature vanishes. The Feigin--Frenkel center $\mathfrak{z}(\widehat{\mathfrak{sl}}_{3,-3})$ is a polynomial algebra in $2$ generators (corresponding to the $2$ fundamental invariants of $\mathfrak{sl}_3$, of degrees $2$ and $3$), consistent with the Wakimoto realization (Theorem~\ref{thm:w3-wakimoto-sl3}).
\end{remark}

\subsection{Spectral sequence collapse}

\begin{theorem}[Genus-$1$ bar-cobar inversion criterion for \texorpdfstring{$\widehat{\mathfrak{sl}}_{3,k}$}{sl-hat_3,k}; \ClaimStatusConditional]
\label{thm:sl3-genus1-inversion}
Let $k+3$ avoid the Kac--Kazhdan resonant hyperplanes.  Assume the
weight-completed genus-$1$ bar object satisfies the generic PBW-Koszul
comparison and that the comparison page is concentrated in columns
$0$ and~$1$.  Then the genus-$1$ bar-cobar adjunction
\begin{equation}
\Omega\bigl(\barB^{(1)}(\widehat{\mathfrak{sl}}_{3,k})\bigr) \xrightarrow{\;\sim\;} \widehat{\mathfrak{sl}}_{3,k}
\end{equation}
is a quasi-isomorphism in the completed curved category, and the
comparison spectral sequence collapses at $E_2$.
\end{theorem}

\begin{proof}
This is Theorem~\ref{thm:higher-genus-inversion} for $g=1$ and
$\mathfrak{g}=\mathfrak{sl}_3$ after imposing the hypotheses in the
statement.

\emph{Step~1: Weight filtration and $E_1$ page.}
All $8$ generators have conformal weight~$1$, giving a bounded-below exhaustive weight filtration. The $E_1$ differential agrees with $\dzero$ at leading order (since $K^{(1)}(z,w|\tau) = 1/(z-w) + O(1)$), computing $H^*(\mathfrak{sl}_3, V_k)$.

\emph{Step~2: PBW-Koszul input.}
Non-resonance removes Kac--Kazhdan singular vectors from the generic
universal vacuum module.  The actual input needed for genus-$1$
bar-cobar inversion is the PBW-Koszul concentration:
\begin{equation}
\operatorname{Ext}^p_{U(\widehat{\mathfrak{sl}}_3)}(V_k, V_k) = 0 \quad \text{for } p \geq 2
\end{equation}

\emph{Step~3: $E_2$ collapse under that input.}
The $E_2$ page is concentrated in the strip $0 \leq p \leq 1$:
\begin{equation}\label{eq:sl3-E2-page}
E_2^{p,q} = \begin{cases}
\mathbb{C} & (p,q) = (0,0) \\
\mathfrak{sl}_3^* & (p,q) = (1,0) \\
\mathbb{C}\lambda & (p,q) = (0,2) \\
\mathfrak{sl}_3^* \cdot \lambda & (p,q) = (1,2) \\
0 & \text{otherwise}
\end{cases}
\end{equation}
where $\dim(\mathfrak{sl}_3^*) = 8$. All higher differentials vanish for degree reasons (non-zero entries in rows~0 and~2 only), so $E_2 = E_\infty$.

The 8-dimensional $\mathfrak{sl}_3^*$ factor (versus 3-dimensional for $\mathfrak{sl}_2$) reflects the larger deformation space of the rank-2 algebra, but the collapse mechanism is identical.
\end{proof}

\subsection{Complementarity and partition function}

\begin{theorem}[Genus-1 complementarity for \texorpdfstring{$\widehat{\mathfrak{sl}}_{3,k}$}{sl-hat_3,k}; \ClaimStatusConditional]
\label{thm:sl3-genus1-complementarity}
For generic $k \neq -3$:
\begin{align}
Q_1(\widehat{\mathfrak{sl}}_{3,k}) &= \mathbb{C} \cdot (k+3) \subset H^0(\overline{\mathcal{M}}_{1,1}) \label{eq:Q1-sl3} \\
Q_1(\widehat{\mathfrak{sl}}_{3,-k-6}) &= \mathbb{C} \cdot \lambda \subset H^2(\overline{\mathcal{M}}_{1,1}) \label{eq:Q1-sl3-dual}
\end{align}
and the complementarity decomposition:
\begin{equation}\label{eq:sl3-complementarity}
Q_1(\widehat{\mathfrak{sl}}_{3,k}) \oplus Q_1(\widehat{\mathfrak{sl}}_{3,-k-6}) = H^*(\overline{\mathcal{M}}_{1,1}) = \mathbb{C} \oplus \mathbb{C}\lambda
\end{equation}
\end{theorem}

\begin{proof}
This is the scalar Hodge-lane specialization of
Theorem~\ref{thm:quantum-complementarity-main} applied to
$\cA = \widehat{\mathfrak{sl}}_{3,k}$, under the generic hypotheses of
Theorem~\ref{thm:sl3-genus1-inversion}.

\emph{Step~1: Obstruction from curvature.}
By Theorem~\ref{thm:sl3-genus1-curvature},
$(d^{(1)})^2 = (k+3) \cdot \omega_1 \cdot
\operatorname{id}_{\mathfrak{sl}_3}$. At generic level the
non-derived chiral centre is scalar, so:
\begin{equation}
Q_1(\widehat{\mathfrak{sl}}_{3,k}) = \mathbb{C} \cdot (k+3) \subset H^0(\overline{\mathcal{M}}_{1,1})
\end{equation}

\emph{Step~2: Dual scalar before Hodge placement.}
The level-reflected current presentation has $k^* = -k-6$
(Theorem~\ref{thm:sl3-koszul-dual}), giving curvature
$(k^*+3)\omega_1 = (-k-3)\omega_1$.  The raw scalar is
$\mathbb{C}\cdot(-k-3)$; the theorem statement records the
Verdier-placed dual summand in $H^2$.

\emph{Step~3: Kodaira--Spencer exchange.}
The Verdier involution $\sigma\colon H^0 \xrightarrow{\sim} H^2$ sends $(k+3) \mapsto (k+3)\lambda$. The complementarity $Q_1 \oplus Q_1^! = H^*(\overline{\mathcal{M}}_{1,1})$ follows as in the $\mathfrak{sl}_2$ case (Theorem~\ref{thm:sl2-genus1-complementarity}).

\emph{Consistency checks.}
\begin{enumerate}
\item \emph{Involutivity}: $k \leftrightarrow -k-6$ changes
 $k+3$ to $-k-3$ and exchanges the Hodge placement.
\item \emph{Dimension}: on the non-critical generic surface,
 $\dim Q_1 + \dim Q_1^! = 1+1 = 2 =
 \dim H^*(\overline{\mathcal{M}}_{1,1})$.
\item \emph{Abelian limit}: deleting the Lie bracket gives eight
 decoupled Heisenberg oscillators with scalar coefficient~$k$.
\item \emph{Critical level}: at $k=-3$ the scalar coefficient
 vanishes and the generic theorem exits its domain; the critical
 replacement is the rank-$2$ Feigin--Frenkel centre.
\end{enumerate}
\end{proof}

\begin{proposition}[Genus-$1$ scalar logarithm for
\texorpdfstring{$\widehat{\mathfrak{sl}}_{3,k}$}{sl-hat_3,k};
\ClaimStatusConditional]
\label{prop:sl3-genus1-partition}
\index{partition function!genus-1!sl3@$\mathfrak{sl}_3$}
Under the genus-universality hypothesis package, the scalar genus-$1$
logarithm is
\begin{equation}\label{eq:sl3-genus1-partition}
\left.\log Z_1^{\mathrm{sc}}(\widehat{\mathfrak{sl}}_{3,k})\right|_{\lambda_1}
= F_1(\widehat{\mathfrak{sl}}_{3,k})\,\lambda_1,
\qquad
F_1(\widehat{\mathfrak{sl}}_{3,k})
=\frac{\kappa(\widehat{\mathfrak{sl}}_{3,k})}{24}
=\frac{k+3}{18}.
\end{equation}
The factor $\eta^{-8}$ is the oscillator denominator in the free-current
model. It is not the scalar Hodge-line coefficient and it is not a
substitute for the full affine vacuum character.
\end{proposition}

\begin{proof}
For $\mathfrak{sl}_3$ one has $h^\vee=3$ and
$\dim(\mathfrak{sl}_3)=8$. The genus-universality theorem gives
\[
\kappa(\widehat{\mathfrak{sl}}_{3,k})
=\frac{(k+h^\vee)\dim(\mathfrak{sl}_3)}{2h^\vee}
=\frac{4(k+3)}3 .
\]
Hence $F_1=\kappa/24=(k+3)/18$. The curvature computation in
Theorem~\ref{thm:sl3-genus1-curvature} supplies the same shifted level
$k+3$; the division by $24$ is the genus-$1$ scalar normalization, not
the oscillator count.
\end{proof}

\begin{remark}[Universality at rank 2]\label{rem:sl3-universality}
The obstruction coefficient from Theorem~\ref{thm:genus-universality} predicts:
\begin{equation}
\kappa(\widehat{\mathfrak{sl}}_{3,k}) = \frac{(k+h^\vee)\dim(\mathfrak{sl}_3)}{2h^\vee} = \frac{(k+3) \cdot 8}{6} = \frac{4(k+3)}{3}
\end{equation}
with the level-reflected characteristic
\(\kappa(\widehat{\mathfrak{sl}}_{3,-k-6})
= -\kappa(\widehat{\mathfrak{sl}}_{3,k})\) and $\kappa = 0$ at
$k = -3$.

This gives (Theorem~\ref{thm:genus-universality}, part~(iii)):
\begin{equation}
F_1(\widehat{\mathfrak{sl}}_{3,k}) = \frac{\kappa}{24} = \frac{k+3}{18}
\end{equation}
which is the scalar coefficient recorded in
\eqref{eq:sl3-genus1-partition}.

Compare $\mathfrak{sl}_2$: $\kappa = 3(k+2)/4$ and $F_1 = (k+2)/32$. The ratio $F_1(\mathfrak{sl}_3)/F_1(\mathfrak{sl}_2) = 32(k+3)/(18(k+2))$ is not a universal constant: it reflects the different dimensions and dual Coxeter numbers. What \emph{is} universal is the factorization $F_g = \kappa \cdot \lambda_g$ from the genus universality theorem.
\end{remark}

\begin{remark}[Rank-2 features]\label{rem:sl3-rank2}
The passage from $\mathfrak{sl}_2$ (rank~1) to $\mathfrak{sl}_3$ (rank~2) introduces three structural novelties visible in the genus-$1$ comparison:
\begin{enumerate}
\item \emph{Serre relations.} The Serre relation $(\operatorname{ad} E_1)^2(E_2) = 0$ is classified by the degree-3 bar complex (Computation~\ref{comp:sl3-bar}). Rank~$1$ has vacuous Serre relations; in rank~$2$ they contribute non-trivially to the cohomology of $\barB^{(1),3}$.
\item \emph{Two-dimensional weight lattice.} The Cartan subalgebra $\mathfrak{h} = \mathrm{span}\{H_1, H_2\}$ introduces a 2-dimensional weight lattice, so the genus-1 Shapovalov non-degeneracy must be checked on the full $A_2$ weight system.
\item \emph{Non-diagonal Cartan OPE.} The Cartan matrix $A_{ij}$ for $\mathfrak{sl}_3$ has off-diagonal entries $A_{12} = -1$, so the Cartan-Cartan OPE $H_i(z)H_j(w) \sim kA_{ij}/(z-w)^2$ mixes the two Cartan directions. The curvature computation contracts over all $8$ generators, and the universal formula $C_2^{\mathrm{ad}} = 2h^\vee = 6$ absorbs this complexity.
\end{enumerate}

For the scalar curvature and partition-function coefficient computed
here, the rank-$2$ data enter only through $k+3$ and
$\dim(\mathfrak{sl}_3)=8$.  A uniform genus-$1$ theorem for every simple
$\widehat{\mathfrak{g}}_k$ would still require the corresponding
PBW-Koszul comparison and weight-completed collapse for that type.
\end{remark}

\section{Open problems and conjectures}

The proved core of this chapter (the universal Koszul dual surface, the
Feigin--Frenkel shadow, and the rank-one/rank-two genus-$1$
complementarity checks on their generic loci) invites natural
extensions. Each open problem
below arises from the geometry already established: admissible
denominators detect root-of-unity phases, while the chiral bar
differential requires a separate periodic CDG structure; the
bar--cobar adjunction supplies Verdier intertwining, while braided
monoidality must be imposed on the coderived comparison with quantum
groups. The curvature variable can become the root-of-unity parameter
only after these hypotheses are fixed.

\subsection{Bar-cobar conditions for Kazhdan--Lusztig transport}
\label{subsec:kl-from-bar-cobar}
\index{Kazhdan--Lusztig equivalence!from bar-cobar}

The chiral Koszul duality of Theorem~\ref{thm:universal-kac-moody-koszul}
and the module-level duality of Theorem~\ref{thm:e1-module-koszul-duality}
give only the first input for a configuration-space construction of the
Kazhdan--Lusztig equivalence~\cite{KL93}.  The missing construction is
not a restatement of KL: it requires a periodic curved bar object, a
coderived comparison functor, and braided monoidality compatible with
the Drinfeld--Kohno normalization
\(\zeta_k=e^{\pi i/(k+h^\vee)}\).

\begin{remark}[Conjectural route: KL equivalence from bar-cobar]%
\label{rem:kl-correspondence}
\index{Kazhdan--Lusztig equivalence!bar-cobar proof}
The proposed construction of the Kazhdan--Lusztig equivalence from the
bar-cobar formalism decomposes into three logically ordered
conjectures.
\end{remark}

\begin{conjecture}[KL-I: periodic CDG structure;
\ClaimStatusConjectured]\label{conj:kl-periodic-cdg}
At admissible level $k = -h^\vee + p/q_0$ ($p, q_0$ coprime positive integers,
$q_0 \geq h^\vee$), set
\[
\zeta_k=e^{\pi i/(k+h^\vee)}=e^{\pi i q_0/p},
\qquad
N_k=\operatorname{ord}(\zeta_k)\mid 2p.
\]
The bar complex $\barB(\widehat{\fg}_k)$ acquires a
periodic CDG structure with eventual period dividing $N_k$:
\[
\barB^{n+N_k}(\widehat{\fg}_k) \;\cong\; \barB^n(\widehat{\fg}_k)
\quad \text{for all } n \gg 0,
\]
in the coderived category of curved bar comodules.
\end{conjecture}

\begin{conjecture}[KL-II: coderived lift and semisimplified target;
\ClaimStatusConjectured]\label{conj:kl-coderived}
Assuming Conjecture~\textup{\ref{conj:kl-periodic-cdg}}, there is a
coderived comparison functor from the periodic bar complex to the
quantum group side whose semisimplified H-level shadow is the
Kazhdan--Lusztig target:
\begin{equation}\label{eq:kl-bar-cobar}
\mathcal{O}_k^{\mathrm{int}}(\widehat{\fg})
\;\simeq\;
\mathcal{C}(U_{\zeta_k}(\fg)),
\end{equation}
The chain-level lift would have to supply the following data:
\begin{enumerate}[label=\textup{(\roman*)}]
\item The bar complex $\barB(\widehat{\fg}_k, M)$ of an integrable
 $\widehat{\fg}_k$-module $M$ has cohomology carrying a finite-dimensional
 $U_{\zeta_k}(\fg)$-module structure whose semisimplification lies in
 $\mathcal{C}(U_{\zeta_k}(\fg))$.
\item The cobar construction $\Omega(\barB(\widehat{\fg}_k, M))
 \xrightarrow{\sim} M$ supplies the module-level inversion on the
 completed Koszul surface.
\item The bar resolution multiplicities match the Kazhdan--Lusztig
 polynomials:
 \begin{equation}\label{eq:bar-kl-poly}
 [\barB^n(\widehat{\fg}_k, L(\lambda))]
 = \sum_{\ell(w) = n} P_{e,w}(\zeta_k) \cdot [\mathcal{M}(w \cdot \lambda)]
 \end{equation}
 recovering Remark~\textup{\ref{rem:bar-complex-kl-polynomials}}.
\end{enumerate}
\end{conjecture}

\begin{conjecture}[KL-III: braided monoidality;
\ClaimStatusConjectured]\label{conj:kl-braided}
Assuming Conjecture~\textup{\ref{conj:kl-coderived}}, the equivalence
\eqref{eq:kl-bar-cobar} is braided monoidal on
$\mathcal{C}(U_{\zeta_k}(\fg))$: fusion of integrable
modules ($\boxtimes$) maps to the semisimplified quantum-group tensor
product, and braiding maps to
\(\zeta_k \mapsto \zeta_k^{-1}\) under Koszul duality.
\end{conjecture}

\begin{remark}[Categorical target]\label{rem:kl-conjecture-target}
Conjectures~\ref{conj:kl-periodic-cdg}--\ref{conj:kl-braided} target
the \emph{semisimplified tilting quotient}
$\mathcal{C}(U_{\zeta_k}(\fg))$ (Remark~\ref{rem:kl-category-precision}),
obtained by quotienting out negligible morphisms.
At chain level one still expects a non-semisimple lift through the
full finite-dimensional quantum-group category
$\mathrm{Rep}^{\mathrm{fd}}(U_{\zeta_k}(\fg))$ or the small quantum group,
but that lift is a separate strengthening, useful for BGG reciprocity
and Loewy-theoretic questions rather than part of the basic H-level
endpoint.
\end{remark}

\begin{remark}[Open mathematical conditions for the KL lift]\label{rem:kl-evidence}
The proved inputs are the level-shifting duality
(Theorem~\ref{thm:universal-kac-moody-koszul}), module-level
bar-cobar (Theorem~\ref{thm:e1-module-koszul-duality}), the
bar-resolution KL multiplicity surface
(Remark~\ref{rem:bar-complex-kl-polynomials}), type-\(A_1\)
periodicity (Proposition~\ref{prop:periodicity-same-type}), affine
Hecke Koszul duality (Proposition~\ref{prop:affine-hecke-kd}), fusion
structure (Theorem~\ref{thm:fusion-bar-cobar}), and weightwise
finiteness at non-degenerate admissible level
(Corollary~\ref{cor:bar-admissible-finiteness}).

Three mathematical conditions remain.  First, the root-of-unity parameter
\(\zeta_k\) must produce an actual \(N_k\)-periodic CDG or
\(N_k\)-complex structure on the curved bar object; finite order of
\(\zeta_k\) alone does not imply periodicity of \(\barB\).  Second,
the coderived category of that object must be compared with the
semisimplified tilting quotient, including the quotient by negligible
morphisms.  Third, the comparison must identify the bar braiding with
the Drinfeld--Kohno/KZ braiding at
\(\zeta_k=e^{\pi i/(k+h^\vee)}\), while keeping the trace-form
collision residue \(r^{\mathrm{coll}}(z)=k\Omega_{\mathrm{tr}}/z\)
separate from the KZ coefficient.
\end{remark}

\begin{remark}[N-complex framework for root-of-unity bar]\label{rem:n-complex-framework}
\index{N-complex}
At admissible $k=-h^\vee+p/q_0$, let
\(\zeta_k=e^{\pi i q_0/p}\) and
\(N_k=\operatorname{ord}(\zeta_k)\).  The proposed root-of-unity
bar structure is an \(N_k\)-complex in the sense of
Kapranov~\cite{Kapranov96}: \(d^{N_k}=0\) while \(d^2\) is the
curved square away from the ordinary dg locus.  The framework of
\cite{Kapranov96} and Khovanov--Qi~\cite{KQ20} supplies
\(N\)-complex cohomology groups \(H^{(j)}\), quantum-group
categorification by \(N\)-complexes, and braid group actions on
derived categories of \(N\)-complexes.  The conjectural content is
that \(\barB(\widehat{\fg}_k)\) carries such an \(N_k\)-complex
structure and that passage to its coderived \(N_k\)-cohomology
recovers the KL target \(\mathcal C(U_{\zeta_k}(\fg))\).
\end{remark}

\begin{remark}[Irrational-level KL categories for $\mathcal{W}$-algebras]
\label{rem:irrational-level-generic}
\index{Kazhdan--Lusztig category!W-algebra@$\mathcal{W}$-algebra}
\index{W-algebra@$\mathcal{W}$-algebra!KL category at irrational level}
Creutzig--Dhillon--Nakatsuka~\cite{CDN26} prove that at irrational
levels $\kappa$, the Kazhdan--Lusztig category
$\mathrm{KL}(\mathcal{W}^\kappa(\fg, f))$ is braided tensor equivalent
to $\mathrm{KL}(V^\kappa(\fg))$ for simply-laced~$\fg$ and arbitrary
nilpotent~$f$. This transports the MC3 thick generation theorem
(Corollary~\ref{cor:mc3-all-types}) from affine KL categories to
$\mathcal{W}$-algebra KL categories at irrational level: if
$\mathrm{KL}(V^\kappa(\fg))$ is thickly generated by evaluation
modules, then $\mathrm{KL}(\mathcal{W}^\kappa(\fg, f))$ inherits
thick generation through the equivalence. At admissible and integral
levels, the equivalence breaks down (null vectors and finite-order
phenomena intervene), and the KL category structure is governed by the
root-of-unity periodicity of
Conjectures~\ref{conj:kl-periodic-cdg}--\ref{conj:kl-braided}.
\end{remark}

\begin{remark}[KM fusion product under bar-cobar]
\label{rem:km-fusion-bar-cobar}
\index{fusion product!KM bar-cobar}
Theorem~\ref{thm:fusion-bar-cobar} specializes to KM: for $\widehat{\fg}_k$-modules $M_1, M_2$ at admissible level,
\begin{equation}\label{eq:fusion-bar-cobar}
\barB(\widehat{\fg}_k,\, M_1 \boxtimes_{\widehat{\fg}_k} M_2)
\longrightarrow
\barB(\widehat{\fg}_k, M_1)
\square_{\barB(\widehat{\fg}_k)}
\barB(\widehat{\fg}_k, M_2)
\end{equation}
which is a quasi-isomorphism on the quadratic genus-$0$ lane under the
K\"unneth/Tor-independence hypotheses. On the affine
Drinfeld--Kohno surface this comparison is identified with tensor
product on the semisimplified quantum-group side, with braiding phases
$e^{2\pi i \lambda\mu/(k+h^\vee)} \mapsto
e^{-2\pi i \lambda\mu/(k+h^\vee)}$ under
$k \mapsto -k - 2h^\vee$. The admissible-level case requires the
root-of-unity bar analysis of
Conjectures~\ref{conj:kl-periodic-cdg}--\ref{conj:kl-braided}.
\end{remark}

\subsection{Geometric Langlands from critical bar--Ext}
\label{subsec:geom-langlands-bar}
\index{geometric Langlands!from critical bar--Ext}

\begin{theorem}[Oper space from critical bar--Ext at \texorpdfstring{$H^0$}{H0};
\ClaimStatusConditional]\label{thm:oper-bar-h0}
\index{oper!from critical bar--Ext}
At the critical level $k = -h^\vee$, the bar complex
$\barB(\widehat{\fg}_{-h^\vee})$ is uncurved because
$m_0=(k+h^\vee)\mathsf C_{\fg}/(2h^\vee)$ vanishes
(Corollary~\ref{cor:critical-level-universality},
Proposition~\ref{prop:km-critical-separation}).  Under the local
bar-Ext comparison at the critical vacuum, write
\[
H^n_{\mathrm{bExt}}(\widehat{\fg}_{-h^\vee})
:=
\mathrm{Ext}^n_{\widehat{\fg}_{-h^\vee}\text{-mod}}
(V_{\mathrm{crit}},V_{\mathrm{crit}})
\]
for the completed critical bar-Ext cohomology computed by that
comparison. Then
\begin{equation}\label{eq:bar-oper}
H^0_{\mathrm{bExt}}(\widehat{\fg}_{-h^\vee})
\cong \mathfrak{z}(\widehat{\fg}_{-h^\vee})
\cong \mathrm{Fun}(\mathrm{Op}_{\fg^\vee}(D))
\end{equation}
where $\mathrm{Op}_{\fg^\vee}(D)$ is the space of $\fg^\vee$-opers
on the formal disc~$D$.
\end{theorem}

\begin{proof}
At $k = -h^\vee$, the shifted curvature element
$m_0=(k+h^\vee)\mathsf C_{\fg}/(2h^\vee)$ is zero, so
$\dfib^{\,2}=0$ and the bar complex is a genuine cochain complex.
The conditional input is the local critical bar-Ext comparison
\[
H^0_{\mathrm{bExt}}(\widehat{\fg}_{-h^\vee})
:=
\operatorname{End}_{\widehat{\fg}_{-h^\vee}}(V_{\mathrm{crit}}).
\]
The endomorphism algebra is the Feigin--Frenkel center
$\mathfrak{z}(\widehat{\fg}_{-h^\vee})$. The Feigin--Frenkel
isomorphism
$\mathfrak{z} \cong \mathrm{Fun}(\mathrm{Op}_{\fg^\vee}(D))$ is
\ClaimStatusProvedElsewhere{} \cite{Feigin-Frenkel, FG06}.
\end{proof}

\begin{proposition}[\texorpdfstring{$H^1$}{H1} at critical level;
\ClaimStatusConditional]\label{prop:oper-bar-h1}
\index{oper!$H^1$ at critical level}
At $k = -h^\vee$, the first completed critical bar-Ext cohomology is
the module of K\"ahler differentials of the center:
\begin{equation}\label{eq:bar-oper-h1}
H^1_{\mathrm{bExt}}(\widehat{\fg}_{-h^\vee})
\cong \Omega^1_{\mathfrak{z}/\bC}
\cong \Omega^1(\mathrm{Op}_{\fg^\vee}(D))
\end{equation}
the cotangent module of the oper space.
\end{proposition}

\begin{proof}
Use the PBW spectral sequence from
Theorem~\ref{thm:oper-bar-h0}.  On each finite conformal-weight
piece the possible off-axis degree-$1$ obstruction is
$H^1(\fg;M)$ for a finite-dimensional semisimple
$\fg$-module~$M$, hence vanishes by Whitehead's first lemma.  The
remaining degree-$1$ class is the universal K\"ahler differential of
the invariant algebra
$\mathfrak z=\mathrm{Fun}(\mathrm{Op}_{\fg^\vee}(D))$.  The
formal-disc oper space is formally smooth, so its cotangent complex
is classical and
$\mathbb L_{\mathrm{Op}_{\fg^\vee}(D)}
\simeq\Omega^1_{\mathfrak z/\bC}[0]$.  This gives
\[
H^1_{\mathrm{bExt}}(\widehat{\fg}_{-h^\vee})
\cong
\Omega^1_{\mathfrak z/\bC}
\cong
\Omega^1(\mathrm{Op}_{\fg^\vee}(D)).
\]
\end{proof}

\begin{theorem}[Local oper differential-form identification;
\ClaimStatusConditional]\label{thm:oper-bar}
For all $n \ge 0$,
\begin{equation}\label{eq:bar-derived-oper}
H^n_{\mathrm{bExt}}(\widehat{\fg}_{-h^\vee})
\cong \Omega^n(\mathrm{Op}_{\fg^\vee}(D)).
\end{equation}
Equivalently, the completed critical bar-Ext cohomology algebra agrees
with the graded algebra of differential forms on the formal-disc oper
space. This is not a statement that raw bar cohomology is the Koszul
dual or the derived center.
\end{theorem}

\begin{proof}
Combine the bar-Ext identification
(Corollary~\ref{cor:bar-computes-ext}) with the
Frenkel--Teleman theorem
(\eqref{eq:frenkel-teleman}; \cite{FT06}):
\[
H^n_{\mathrm{bExt}}(\widehat{\fg}_{-h^\vee})
\cong \mathrm{Ext}^n(V_{\mathrm{crit}}, V_{\mathrm{crit}})
\cong \Omega^n(\mathrm{Op}_{\fg^\vee}(D))
\]
for all $n \geq 0$. See
Theorem~\ref{thm:oper-bar-dl} for the full proof.
\end{proof}

\begin{remark}[Frenkel--Teleman and the oper side]%
\label{rem:frenkel-teleman-oper}
\index{Frenkel--Teleman theorem}
\index{oper!Ext computation}
\index{fundamental local equivalence}
The external Frenkel--Teleman input is:
\begin{equation}\label{eq:frenkel-teleman}
\mathrm{Ext}^n_{\widehat{\fg}_{-h^\vee}\text{-mod}}
(V_{\mathrm{crit}},\, V_{\mathrm{crit}})
\;\cong\; \Omega^n(\mathrm{Op}_{\fg^\vee}(D))
\qquad \text{for all } n \geq 0.
\end{equation}
Our $H^0_{\mathrm{bExt}}$ and $H^1_{\mathrm{bExt}}$ computations
(Theorem~\ref{thm:oper-bar-h0},
Proposition~\ref{prop:oper-bar-h1}) match the $n = 0,1$ cases. The
FLE of Raskin~\cite{Raskin24}
\index{geometric Langlands!fundamental local equivalence}
identifies $\widehat{\fg}_{-h^\vee}$-modules with the Whittaker
category, a central ingredient of the geometric Langlands
proof~\cite{GLC24}. Under the bar-Ext hypothesis package, one obtains
$H^n_{\mathrm{bExt}}(\widehat{\fg}_{-h^\vee})
\cong \Omega^n(\mathrm{Op}_{\fg^\vee}(D))$ for all~$n$.
\end{remark}

\begin{remark}[Scope]\label{rem:oper-bar-scope}
\index{geometric Langlands!from bar complex}
Theorem~\ref{thm:oper-bar} is the \emph{local} formal-disc statement:
bar-Ext (Corollary~\ref{cor:bar-computes-ext}) combined with
Frenkel--Teleman (Remark~\ref{rem:frenkel-teleman-oper}) gives
$H^n_{\mathrm{bExt}} \cong \Omega^n(\mathrm{Op}_{\fg^\vee}(D))$ for all~$n$.
Independent PBW confirmation exists for $n \leq 3$
(Chapter~\ref{ch:derived-langlands}); the Whitehead spectral
decomposition (Proposition~\ref{prop:whitehead-spectral-decomposition})
extends to all~$n$. Any global curve-level or de~Rham enhancement
requires additional Beilinson--Drinfeld / localization input and is not
proved by this theorem alone.
\end{remark}

\subsection{The oper space at critical level: summary}
\label{subsec:oper-critical-summary}
\index{opers!critical level summary}

The bar complex of the affine algebra at critical level
$k = -h^\vee$ occupies a distinguished position in the theory.
By Corollary~\ref{cor:critical-level-universality}, the bar
complex $\barB(\widehat{\fg}_{-h^\vee})$ is \emph{uncurved}:
the shifted-level factor $k+h^\vee$ vanishes in
$m_0=(k+h^\vee)\mathsf C_{\fg}/(2h^\vee)$.  The same factor is the
denominator of the Sugawara tensor, so the curvature vanishing and
the Sugawara degeneration occur at the same point but are not the
same assertion. The first two cohomology groups have been identified:

\begin{itemize}
\item $H^0_{\mathrm{bExt}}(\widehat{\fg}_{-h^\vee})
 \cong \mathrm{Fun}(\mathrm{Op}_{\fg^\vee}(D))$
 (Theorem~\ref{thm:oper-bar-h0}): the zeroth completed critical
 bar-Ext cohomology recovers the algebra of functions on the space of
 $\fg^\vee$-opers on the formal disc.
\item $H^1_{\mathrm{bExt}}(\widehat{\fg}_{-h^\vee})
 \cong \Omega^1(\mathrm{Op}_{\fg^\vee}(D))$
 (Proposition~\ref{prop:oper-bar-h1}): the first completed critical
 bar-Ext cohomology recovers K\"ahler differentials on the oper
 space.
\end{itemize}

These are the first two terms of the identification
\[
H^\bullet_{\mathrm{bExt}}(\widehat{\fg}_{-h^\vee})
\cong \Omega^\bullet(\mathrm{Op}_{\fg^\vee}(D))
\]
(Theorem~\ref{thm:oper-bar}), which establishes that the completed
critical bar-Ext comparison classifies a vertex-algebraic model for
the oper differential-form package on the formal-disc oper space. The full
cohomological identification,
including the DAG prerequisites, root-of-unity periodic CDG
structure, and the Kazhdan--Lusztig package, is developed in
Chapter~\ref{ch:derived-langlands}.

\subsection{Fusion monoidality for Kac--Moody modules}
\label{subsec:km-fusion-monoidality}
\index{fusion product!Kac--Moody monoidality}
\index{Kac--Moody algebra!fusion monoidality}

On the quadratic genus-$0$ $\Eone$ lane, the module bar functor
\[
\Phi\colon \mathrm{Mod}^{\mathrm{compl}}(\cA)
\to \mathrm{CoMod}^{\mathrm{conil}}(\barB(\cA))
\]
\textup{(}Theorem~\textup{\ref{thm:e1-module-koszul-duality})}
is lax monoidal with respect to
fusion and cotensor product
(Theorem~\ref{thm:fusion-bar-cobar}). For Kac--Moody
algebras the evidence comes from three sources.

\emph{Heisenberg Fock modules.}
For the Heisenberg algebra $\mathcal{H}_\kappa$, the fusion of
Fock modules satisfies
$F_\lambda \boxtimes_\kappa F_\mu \cong F_{\lambda+\mu}$,
and one verifies
$\Phi(F_\lambda \boxtimes F_\mu)
\cong \Phi(F_\lambda) \otimes \Phi(F_\mu)$
with phase inversion $\kappa \mapsto -\kappa$
(Proposition~\ref{prop:fock-fusion-product}). The proof relies
on the purely double-pole nature of the Heisenberg OPE; see
\S\ref{subsec:heisenberg-fusion-verification} for the explicit
bar-level computation.

\emph{Integrable Kac--Moody modules.}
For $\widehat{\fg}_k$ at positive integral level, the
separating-divisor coproduct of the chain-level modular functor
(Theorem~\ref{thm:chain-modular-functor}) provides the
chain-level fusion map
\[
\Delta_{\mathrm{sep}}\colon
\barB(\widehat{\fg}_k,\, M_1 \boxtimes_k M_2)
\;\longrightarrow\;
\barB(\widehat{\fg}_k,\, M_1)
\otimes \barB(\widehat{\fg}_k,\, M_2).
\]
In genus~$0$, the pants decomposition of
$\overline{\mathcal{M}}_{0,4}$ yields a natural isomorphism
$\Phi(M_1 \boxtimes_\cA M_2)
\to \Phi(M_1) \square_{\barB(\cA)} \Phi(M_2)$
on the bar-comodule side, which is the genus-$0$ case of the
module-bar lax monoidal structure.

\emph{Genus-$1$ obstruction.}
At genus~$1$, the curvature $m_0 \neq 0$ prevents strict
monoidality. The obstruction is controlled by
$\kappa(\cA) \cdot [\overline{\mathcal{M}}_{1,1}]$, the
curvature times the fundamental class of the genus-$1$ moduli
space: this is the same invariant that governs the
genus-$1$ quantum correction in
Theorem~\ref{thm:quantum-complementarity-main}. The full
conjecture and its relation to the Kazhdan--Lusztig package
appear in Theorem~\ref{thm:fusion-bar-cobar}.

\section{Open questions}

\begin{question}
The genus-$1$ generic check
(\S\ref{sec:sl2-genus-one-computation}) shows that modular forms enter
through $E_2$ via the $B$-cycle quasi-periodicity of the propagator.
What is the complete structure at genus $g \geq 2$, where the
propagator has $2g$ quasi-periodicity defects? Does the genus-$g$
partition function involve Siegel modular forms on $\mathfrak{H}_g$?
\end{question}

\begin{question}
The exceptional types $E_6$, $E_7$, $E_8$ are treated in \S\ref{sec:exceptional-shadow-data} (Proposition~\ref{prop:exceptional-shadow-invariants}); the non-simply-laced types $B_2$, $G_2$, $F_4$ are treated in \S\ref{subsec:nsl-shadow-data} (Proposition~\ref{prop:nsl-shadow-tower}) and in the bar complex computations of \S\ref{sec:B2-details}--\S\ref{sec:G2-details}. Can the framework be extended further to affine algebras of twisted type or super Lie algebras?
\end{question}

\begin{question}
What is the categorification? Is there a $2$-categorical Koszul
duality for affine Kac--Moody categories, lifting the module-level
duality of Theorem~\ref{thm:e1-module-koszul-duality} to a
$2$-functor on the $2$-category of module categories?
\end{question}

\begin{question}
How does the bar-cobar complex interact with quantum geometric
Langlands? A
chain-level local oper resolution is classified by the critical-level bar complex
(Theorem~\ref{thm:oper-bar}), and the admissible-level bar
complex is conjectured to recover the KL equivalence
(Conjectures~\ref{conj:kl-periodic-cdg}--\ref{conj:kl-braided}). Does the full genus-graded
bar complex (Theorem~\ref{thm:genus-graded-bar}) provide a
chain-level refinement of the quantum Langlands kernel?
\end{question}

\section{The affine cubic shadow and critical-level emergence}%
\label{sec:affine-cubic-shadow}
\index{shadow calculus!affine Kac--Moody}
\index{cubic shadow!affine}
\index{Kac--Moody algebra!nonlinear shadow calculus}

While the Heisenberg algebra is the Gaussian archetype
(all nonlinear shadows vanish), affine Kac--Moody is the
\emph{Lie/tree archetype}: a cubic shadow appears, the quartic
obstruction vanishes by the Jacobi identity, and quartic
nonlinearity first manifests on the boundary. The results of this
section are proved locally; the shadow calculus formalism and the
all-degree master equation $\nabla_H(\mathrm{Sh}_r) + \mathfrak{o}^{(r)} = 0$
are developed in Appendix~\ref{sec:nms-shadow-calculus}.

\subsection{Tower construction for affine Kac--Moody}
\label{subsec:affine-tower-construction}
\index{shadow obstruction tower!affine construction}

For $\widehat{\mathfrak{g}}_k$ at generic level on the strict
sector, the shadow tower is fixed by the double pole, the simple
pole, and Jacobi.

\paragraph{Degree~$2$: scalar conductor.}
The double-pole OPE
$J^a(z)\,J^b(w) \sim k(a,b)/(z{-}w)^2 + \cdots$
gives $\kappa(\widehat{\fg}_k) = \dim(\fg)(k{+}h^\vee)/(2h^\vee)$.
Hence $\Theta_{\mathrm{aff}}^{\leq 2} = \kappa$.

\paragraph{Degree~$3$: Lie cubic tensor.}
The simple-pole term
$J^a(z)\,J^b(w) \sim f^{ab}{}_c\,J^c(w)/(z{-}w)$
produces a nonzero transferred operation
$\ell_2^{\mathrm{tr}} \neq 0$, giving the cubic shadow
$\mathfrak{C}_{\mathrm{aff}}(x,y,z) = \kappa_{\fg}(x,[y,z])$.
The obstruction $o_3 \neq 0$; the cubic encodes the Lie bracket.
$\Theta_{\mathrm{aff}}^{\leq 3}
= H_{\mathrm{aff}} + \mathfrak{C}_{\mathrm{aff}}$.

\paragraph{Degree~$4$: Jacobi identity kills the
obstruction.}
The quartic obstruction
$o_4 = \tfrac12\{\mathfrak{C},\mathfrak{C}\}_H
= \kappa_{\fg}([x,[y,z]],w) + \mathrm{cyclic}$
vanishes by the Jacobi identity.
In the quartic gauge, $\mathfrak{Q}_{\mathrm{aff}} = 0$.
$\Theta_{\mathrm{aff}}^{\leq 4} = \Theta_{\mathrm{aff}}^{\leq 3}$.

\paragraph{Termination at degree~$3$.}
All higher obstructions vanish by induction from the quartic
gauge. The tower terminates:
$\Theta_{\mathrm{aff}}^{\leq r}
= H_{\mathrm{aff}} + \mathfrak{C}_{\mathrm{aff}}$
for all $r \geq 3$.

\paragraph{Arithmetic face.}
The universal affine VOA $V_k(\fg)$ is strongly generated by the
currents $J^a$ ($a = 1,\ldots,\dim\fg$), each of conformal
weight~$1$; the Sugawara tensor is composite. Hence
$W(V_k(\fg)) = \{1^{\dim\fg}\}$,
$\operatorname{ek}(V_k(\fg)) = 0$
(Definition~\ref{def:euler-koszul-class}), and
\[
S_{V_k(\fg)}(u) = \dim(\fg)\cdot\zeta(u)\,\zeta(u{+}1).
\]
The affine algebra is \emph{exact Euler--Koszul} despite having
shadow depth~$3$: the cubic shadow $\mathfrak{C}_{\mathrm{aff}}$
comes from the OPE (the Lie bracket), not from the character.

\medskip
The precise statements are
Theorem~\ref{thm:affine-cubic-normal-form} and
Corollary~\ref{cor:affine-postnikov-termination} below.

\begin{definition}[Strict current-level sector; \ClaimStatusDefinitional]%
\label{def:affine-strict-current-sector}
\index{strict current-level sector}
For the affine Kac--Moody algebra $\widehat{\fg}_k$ at level
$k \neq -h^\vee$, define the \emph{strict current-level sector}
\[
V_{\mathrm{aff}} \;:=\; \fg[1] \oplus \bC\,K[1],
\]
the graded vector space generated by the currents $J^a$ and the
level direction~$K$. The Hessian on this sector is the shifted
bilinear form $H_{\mathrm{aff}} = (k + h^\vee)\kappa_{\fg}$, where
\(\kappa_{\fg}\) denotes the normalised invariant form.
\end{definition}

\begin{theorem}[Affine cubic normal form;
\ClaimStatusProvedHere]\label{thm:affine-cubic-normal-form}
\index{cubic shadow!normal form}
\index{quartic obstruction!vanishing for affine}
On the strict current-level sector one has
\[
\mathfrak{C}_{\mathrm{aff}}(x,y,z) = \kappa_{\fg}(x,[y,z]),
\qquad
\mathfrak{o}_{\mathrm{aff}}^{(4)} = 0.
\]
Consequently there exists a quartic gauge in which
\[
\mathfrak{Q}_{\mathrm{aff}} = 0,
\qquad
\Theta_{\mathrm{aff}}^{\le 4}
= H_{\mathrm{aff}} + \mathfrak{C}_{\mathrm{aff}}.
\]
\end{theorem}

\begin{proof}
The simple pole of the current OPE
$J^a(z)\,J^b(w) \sim k(a,b)/(z-w)^2 + f^{ab}{}_c\,J^c(w)/(z-w)$
encodes the Lie bracket; the double pole encodes the bilinear
form. On the strict sector, the cubic Hamiltonian is therefore
\[
\tfrac{1}{6}\,\kappa_{\fg}(a,[a,a]).
\]
The quartic obstruction $\mathfrak{o}^{(4)}$ is the Hamiltonian
self-bracket of this cubic, which is the Jacobiator
$\kappa_{\fg}([x,[y,z]], w) + \mathrm{cyclic}$. By the Jacobi
identity this vanishes identically. One may therefore choose a
minimal gauge through quartic order in which no local quartic
vertex remains.
\end{proof}

\begin{corollary}[Affine shadow obstruction tower:
finite termination at degree~$3$; \ClaimStatusProvedHere]
\label{cor:affine-postnikov-termination}
\index{shadow obstruction tower!affine termination}
The shadow obstruction tower
\textup{(}Definition~\textup{\ref{def:shadow-postnikov-tower})}
of the affine Kac--Moody algebra terminates at degree~$3$
on the strict sector:
$\Theta_{\mathrm{aff}}^{\leq r} =
H_{\mathrm{aff}} + \mathfrak{C}_{\mathrm{aff}}$
for all $r \geq 3$ \textup{(}in the quartic gauge\textup{)},
with $o_{r+1}(\mathrm{aff}) = 0$ for all $r \geq 3$.
The affine tower is Lie-type: quadratic + cubic, with the
cubic shadow encoding the Lie bracket.
\end{corollary}

\begin{proof}
Immediate from Theorem~\ref{thm:affine-cubic-normal-form}:
$\mathfrak{o}_{\mathrm{aff}}^{(4)} = 0$ by the Jacobi
identity, and the quartic gauge annihilates
$\mathfrak{Q}_{\mathrm{aff}}$.
The vanishing propagates to all higher obstructions by induction.
\end{proof}

\begin{remark}[Arithmetic depth of the affine tower]
\label{rem:affine-arithmetic-depth}
In the depth decomposition $d = 1 + d_{\mathrm{arith}} + d_{\mathrm{alg}}$
of Theorem~\ref{thm:depth-decomposition} (Chapter~\ref{chap:arithmetic-shadows}),
the affine algebra has $d_{\mathrm{arith}} = 2$ (two critical lines,
from the Eisenstein series $E_k$ governing the propagator and the
$L$-function of the adjoint representation) and $d_{\mathrm{alg}} = 0$,
giving total depth $d = 3$; this is the Lie class $L$.
\end{remark}

\subsection{The affine primitive kernel}
\label{subsec:affine-primitive-kernel}
\index{primitive kernel!affine Kac--Moody}

\begin{proposition}[Affine primitive kernel;
\ClaimStatusProvedHere]
\label{prop:affine-primitive-kernel}
The primitive logarithmic modular kernel
\textup{(}Definition~\textup{\ref{def:primitive-log-modular-kernel})}
of $\widehat{\mathfrak{g}}_k$ at non-critical level is
\begin{equation}\label{eq:affine-primitive-kernel}
\mathfrak{K}_{\mathrm{aff}}
\;=\;
K_{0,2}^{\mathrm{aff}} + K_{0,3}^{\mathrm{aff}} + K_{1,1}^{\mathrm{aff}},
\end{equation}
where:
\begin{enumerate}[label=\textup{(\roman*)}]
\item $K_{0,2}^{\mathrm{aff}}(x,y) = k_{\mathrm{eff}}\,\kappa_{\fg}(x,y)$
 is the degree-$2$ corolla \textup{(}the effective-level invariant form
 on the strict sector\textup{)};
\item $K_{0,3}^{\mathrm{aff}}(x,y,z) = \kappa_{\fg}(x,[y,z])$
 is the degree-$3$ corolla \textup{(}the Lie bracket\textup{)};
\item $K_{1,1}^{\mathrm{aff}}$ is the genus-$1$ corolla from
 non-separating self-sewing.
\end{enumerate}
All higher genus-$0$ corollas vanish:
$K_{0,n}^{\mathrm{aff}} = 0$ for $n \geq 4$, and
$R_\rho^{\mathrm{aff}} = 0$ for all rigid types~$\rho$.
\end{proposition}

\begin{proof}
The genus-$0$ corollas are the transferred operations
$\ell_n^{\mathrm{tr}}$ on the strict current-level sector. The OPE
$J^a(z)\,J^b(w) \sim
k\kappa_{\fg,ab}/(z{-}w)^2 + f^{ab}{}_c J^c(w)/(z{-}w)$
gives $\ell_2^{\mathrm{tr}}(y,z) = [y,z]$ (the Lie bracket) and
$\ell_n^{\mathrm{tr}} = 0$ for $n \geq 3$: no higher local
operations arise from the bar complex because the affine
propagator $P = k_{\mathrm{eff}}^{-1}\kappa_{\fg}^{-1}$ composes with the
binary bracket to give iterated Lie brackets, and the quartic and
higher compositions vanish by the Jacobi identity
(Theorem~\ref{thm:affine-cubic-normal-form}). The rigid-cutting
residues vanish because the planted-forest boundary strata require
at least one composition through a genus-raising vertex with a
nontrivial higher operation, which is absent.
\end{proof}

\begin{proposition}[Affine primitive shell equations;
\ClaimStatusProvedHere]
\label{prop:affine-primitive-shell}
\index{shell equation!affine Kac--Moody}
The primitive master equation for $\widehat{\mathfrak{g}}_k$
has two active channels. At genus~$2$:
\begin{equation}\label{eq:affine-shell-genus2}
R_2(\mathfrak{K}_{\mathrm{aff}})
\;=\;
\Delta_{\mathrm{ns}}(K_{1,1}^{\mathrm{aff}})
\;+\;
\tfrac{1}{2}\,[K_{1,1}^{\mathrm{aff}},\,
K_{1,1}^{\mathrm{aff}}].
\end{equation}
Both the BV self-contraction and the separating bracket contribute.
The rigid-cutting channel is empty.
\end{proposition}

\begin{proof}
The BV self-contraction $\Delta_{\mathrm{ns}}(K_{1,1})$ produces
the genus-$2$ from genus-$1$ by non-separating self-sewing. The
separating bracket $[K_{1,1}, K_{1,1}]$ computes the genus-$2$
contribution from sewing two genus-$1$ corollas along a separating
node; this is nonzero because $V_{\mathrm{aff}}^{\mathrm{br}}
= \mathfrak{g}$ is higher-rank (dimension $\dim\mathfrak{g} \geq 3$),
so the bracket does not vanish by antisymmetry. The rigid-cutting
channel is empty by
Proposition~\ref{prop:affine-primitive-kernel}.
\end{proof}

\begin{computation}[Affine branch BV action]
\label{comp:affine-branch-bv}
\index{branch BV action!affine Kac--Moody}
The branch-active quotient is $V_{\mathrm{aff}}^{\mathrm{br}}
= \mathfrak{g}$ with $\dim V_{\mathrm{aff}}^{\mathrm{br}}
= \dim\mathfrak{g}$. Let $x = x_a e^a \in \mathfrak{g}$ with
coordinates $(x_1, \ldots, x_{\dim\mathfrak{g}})$ dual to a
invariant-form orthonormal basis. The branch master action is
\begin{equation}\label{eq:affine-branch-action}
S_{\mathrm{aff}}^{\mathrm{br}}(x)
\;=\;
\tfrac{1}{2}\,k_{\mathrm{eff}}\,\kappa_{\fg}(x,x)
\;+\;
\tfrac{1}{3!}\,\kappa_{\fg}(x,[x,x]).
\end{equation}
The cubic vertex is the Lie bracket: $S_{\mathrm{aff}}^{\mathrm{br}}$
is the Chern--Simons-type action on the finite-dimensional Lie
algebra~$\mathfrak{g}$. The branch QME
\textup{(}Proposition~\textup{\ref{prop:branch-master-equation})}
holds because $\{S,S\}_{\mathrm{BV}}$ is the Jacobiator, which
vanishes. The branch transport operator
$T_{\mathrm{aff}}^{\mathrm{br,red}}$ is $\mathrm{ad}$-Casimir,
and the metaplectic half-density
\textup{(}Corollary~\textup{\ref{cor:metaplectic-square-root})} is
\begin{equation}\label{eq:affine-metaplectic}
\delta_{\mathrm{aff}}(x)^2
\;=\;
\det\bigl(1 - x\,\mathrm{ad}_{C_2/k_{\mathrm{eff}}}\bigr)
\;=\;
\prod_{\alpha \in \Phi^+}
\bigl(1 - \langle\alpha,\alpha\rangle\,x/k_{\mathrm{eff}}\bigr)^2.
\end{equation}
The determinant factors over positive roots: the metaplectic
half-density of the affine algebra is the Weyl denominator
of the finite-dimensional Lie algebra~$\mathfrak{g}$.
\end{computation}

\subsection{From OPE to cyclic slice}
\label{subsec:affine-ope-cyclic-slice}
\index{cyclic slice!from affine OPE}

\begin{proposition}[Affine cyclic slice data;
\ClaimStatusProvedHere]\label{prop:affine-cyclic-slice-data}
The affine current algebra $\widehat{\mathfrak{g}}_k$ has OPE
\begin{equation}\label{eq:affine-ope-full}
J^a(z)\,J^b(w)
\;\sim\;
\frac{k\,\kappa_{\fg,ab}}{(z-w)^2}
\;+\;
\frac{f^{ab}{}_{c}\,J^c(w)}{z-w}\,.
\end{equation}
On the strict current-level sector
$V_{\mathrm{aff}} = \mathfrak{g}[1] \oplus \mathbb{C}\,K[1]$
\textup{(Definition~\textup{\ref{def:affine-strict-current-sector}})},
the shadow calculus data
\textup{(Setup~\textup{\ref{setup:nms-cyclic-slice}})} are
determined by the two poles of~\eqref{eq:affine-ope-full}:
\begin{enumerate}[label=\textup{(\roman*)}]
\item The \emph{Hessian} \textup{(from the double pole)}:
 $H_{\mathrm{aff}}(x,y) = (k + h^\vee)\,\kappa_{\fg}(x,y)
 =: k_{\mathrm{eff}}\,\kappa_{\fg}(x,y)$.
\item The \emph{propagator}
 \textup{(at non-critical level $k \neq -h^\vee$)}:
 $P_{\mathrm{aff}} = H_{\mathrm{aff}}^{-1}
 = (1/k_{\mathrm{eff}})\,\kappa_{\fg}^{-1}$.
\item The \emph{transferred binary bracket}
 \textup{(from the simple pole)}:
 $\ell_2^{\mathrm{tr}}(y,z) = [y,z]$, the Lie bracket of
 $\mathfrak{g}$.
\end{enumerate}
\end{proposition}

\begin{proof}
The double pole $\kappa_{\fg,ab}/(z-w)^2$ is the classical pairing,
shifted by the effective level $k_{\mathrm{eff}} = k + h^\vee$,
giving the Hessian. The propagator is its inverse on
$\mathfrak{g}$ (nondegenerate since $k \neq -h^\vee$). The
simple pole coefficient $f^{ab}{}_{c}$ is the structure constant
of $\mathfrak{g}$, which in the homotopy transfer formalism
encodes the transferred binary operation
$\ell_2^{\mathrm{tr}} = [\cdot,\cdot]$. All higher transferred
operations $\ell_n^{\mathrm{tr}}$ for $n \ge 3$ vanish: the affine
OPE is quadratic in $J$, so only poles of order $\le 2$ contribute.
\end{proof}

\subsection{The cubic shadow from the simple pole}
\label{subsec:affine-cubic-from-simple-pole}
\index{cubic shadow!from simple pole}

\begin{proposition}[Cubic shadow via ad-invariance;
\ClaimStatusProvedHere]\label{prop:affine-cubic-ad-invariance}
The cyclic cubic shadow is
\begin{equation}\label{eq:affine-cubic-cyclic-sum}
\mathfrak{C}_{\mathrm{aff}}(x,y,z)
\;=\;
\sum_{\mathrm{cyc}(x,y,z)}
\bigl\langle x,\,\ell_2^{\mathrm{tr}}(y,z)\bigr\rangle
\;=\;
\sum_{\mathrm{cyc}(x,y,z)} \kappa_{\fg}\bigl(x,[y,z]\bigr).
\end{equation}
By ad-invariance of the invariant form,
$\kappa_{\fg}([a,b],c) = \kappa_{\fg}(a,[b,c])$, the three cyclic terms
are equal:
\begin{align}
\kappa_{\fg}(x,[y,z]) &= \kappa_{\fg}(x,[y,z]),
 \label{eq:cyc-term-1}\\
\kappa_{\fg}(y,[z,x]) &= \kappa_{\fg}([y,z],x) = \kappa_{\fg}(x,[y,z]),
 \label{eq:cyc-term-2}\\
\kappa_{\fg}(z,[x,y]) &= \kappa_{\fg}([z,x],y) = \kappa_{\fg}(y,[z,x])
 = \kappa_{\fg}(x,[y,z]).
 \label{eq:cyc-term-3}
\end{align}
Therefore
\begin{equation}\label{eq:affine-cubic-three-equal}
\mathfrak{C}_{\mathrm{aff}}(x,y,z)
\;=\; 3\,\kappa_{\fg}(x,[y,z]).
\end{equation}
In the complementarity potential (with the $1/3!$ normalisation),
the cubic action is
$\frac{1}{6}\,\mathfrak{C}_{\mathrm{aff}}(a,a,a)
= \frac{1}{6}\kappa_{\fg}(a,[a,a])$
\textup{(cf.\ Theorem~\textup{\ref{thm:affine-cubic-normal-form}})},
and the normalised cubic tensor is
$\mathfrak{C}_{\mathrm{aff}}(x,y,z) = \kappa_{\fg}(x,[y,z])$ when
written without the cyclic sum
\textup{(as in the theorem statement)}.
\end{proposition}

\begin{proof}
The cyclic sum formula~\eqref{eq:affine-cubic-cyclic-sum} is the
definition of the cubic shadow jet
(Definition~\ref{def:nms-shadow-jets}). Each step
in~\eqref{eq:cyc-term-2}--\eqref{eq:cyc-term-3} uses the single
identity $\kappa_{\fg}([a,b],c) = \kappa_{\fg}(a,[b,c])$, which is
ad-invariance of the normalized invariant form. The equality of
all three terms is specific to the Lie/tree phase: it reflects the
fact that $\kappa_{\fg}$ is a \emph{cyclic} pairing for the Lie bracket,
the defining condition that makes $\mathfrak{g}$ a cyclic
$L_\infty$-algebra at the transferred binary level.
\end{proof}

\subsection{Jacobi identity kills the quartic obstruction}
\label{subsec:affine-jacobi-kills-quartic}
\index{Jacobi identity!quartic obstruction vanishing}
\index{quartic obstruction!Jacobi mechanism}

\begin{proposition}[Jacobi mechanism for quartic vanishing;
\ClaimStatusProvedHere]\label{prop:affine-jacobi-quartic-vanishing}
The quartic obstruction is the Hamiltonian self-bracket of the cubic:
\begin{equation}\label{eq:quartic-obstruction-self-bracket}
\mathfrak{o}_{\mathrm{aff}}^{(4)}
\;=\;
\tfrac{1}{2}\,
\bigl\{\mathfrak{C}_{\mathrm{aff}},
\mathfrak{C}_{\mathrm{aff}}\bigr\}_{H_{\mathrm{aff}}}.
\end{equation}
Expanding on inputs $(x_1,x_2,x_3,x_4) \in \mathfrak{g}^{\otimes 4}$:
\begin{equation}\label{eq:quartic-obstruction-expanded}
\bigl\{\mathfrak{C},\mathfrak{C}\bigr\}_H
(x_1,x_2,x_3,x_4)
\;=\;
\sum_{\mathrm{channels}} P^{ij}\,
\mathfrak{C}(x_i,x_j,\cdot)\,
\mathfrak{C}(\cdot,x_k,x_\ell)
\;=\;
\frac{1}{k_{\mathrm{eff}}}
\sum_{\mathrm{channels}}
\kappa_{\fg}\bigl([x_i,x_j],\,[x_k,x_\ell]\bigr),
\end{equation}
where the sum runs over the three pairings
$(ij|k\ell) \in \{(12|34),\,(13|24),\,(14|23)\}$.
This is the \emph{Jacobiator} applied to $(x_1,x_2,x_3)$
and then paired with~$x_4$:
\begin{align}
\label{eq:jacobiator-pairing}
&\kappa_{\fg}\bigl([x_1,x_2],\,[x_3,x_4]\bigr)
+ \kappa_{\fg}\bigl([x_1,x_3],\,[x_4,x_2]\bigr)
+ \kappa_{\fg}\bigl([x_1,x_4],\,[x_2,x_3]\bigr)
\\
&\quad=\;
\kappa_{\fg}\Bigl(
 \bigl[[x_1,x_2],x_3\bigr]
 + \bigl[[x_2,x_3],x_1\bigr]
 + \bigl[[x_3,x_1],x_2\bigr],\;
 x_4
\Bigr)
\;=\; 0
\notag
\end{align}
by the Jacobi identity. Since $\kappa_{\fg}$ is nondegenerate and
$x_4$ is arbitrary, the Jacobiator vanishes identically.
Therefore $\mathfrak{o}_{\mathrm{aff}}^{(4)} = 0$.
\end{proposition}

\begin{proof}
Line~\eqref{eq:quartic-obstruction-expanded} follows from
inserting $\mathfrak{C}(x,y,z) = \kappa_{\fg}(x,[y,z])$ and contracting
two internal legs with the propagator
$P_{\mathrm{aff}} = (1/k_{\mathrm{eff}})\,\kappa_{\fg}^{-1}$, choosing a
dual basis $\{e_a\}$, $\{e^a\}$ with
$\kappa_{\fg}(e_a,e^b) = \delta_a^b$:
\[
\sum_a \kappa_{\fg}(x_i,[x_j,e_a])\cdot
\frac{1}{k_{\mathrm{eff}}}\cdot
\kappa_{\fg}(e^a,[x_k,x_\ell])
\;=\;
\frac{1}{k_{\mathrm{eff}}}\,
\kappa_{\fg}([x_i,x_j],[x_k,x_\ell]).
\]
The passage from~\eqref{eq:quartic-obstruction-expanded}
to~\eqref{eq:jacobiator-pairing} uses ad-invariance to move one
bracket factor onto the other argument, converting each channel
term $\kappa_{\fg}([x_i,x_j],[x_k,x_\ell])$ into
$\kappa_{\fg}([[x_i,x_j],x_k],x_\ell)$ (up to sign and relabelling).
The sum of three such terms is the Jacobiator paired with $x_4$,
which vanishes by the Jacobi identity.
\end{proof}

\subsection{Boundary sewing for \texorpdfstring{$\mathfrak{sl}_2$}{sl(2)}}
\label{subsec:affine-sl2-boundary-sewing}
\index{boundary shadow!sl2 explicit@$\mathfrak{sl}_2$ explicit}

\begin{proposition}[Explicit boundary quartic for
$\mathfrak{sl}_2$;
\ClaimStatusProvedHere]\label{prop:affine-sl2-boundary-quartic}
For $\mathfrak{g} = \mathfrak{sl}_2$ with standard basis
$\{e,f,h\}$, normalized invariant form satisfying
\begin{equation}\label{eq:sl2-killing-values}
\kappa_{\fg}(e,f) = 1, \quad
\kappa_{\fg}(h,h) = 2, \quad
\kappa_{\fg}(e,e) = \kappa_{\fg}(f,f) = \kappa_{\fg}(e,h) = \kappa_{\fg}(f,h) = 0,
\end{equation}
and brackets $[e,f] = h$, $[h,e] = 2e$, $[h,f] = -2f$,
the three channel contributions to the boundary quartic
$(\mathfrak{C}_{\mathrm{aff}} \star_{P_{\mathrm{aff}}}
\mathfrak{C}_{\mathrm{aff}})(x_1,x_2,x_3,x_4)$
on the input $(e,f,e,f)$ are:
\begin{align}
\textup{$s$-channel } (12|34)\colon\quad
&\frac{1}{k_{\mathrm{eff}}}\,
\kappa_{\fg}\bigl([e,f],[e,f]\bigr)
= \frac{1}{k_{\mathrm{eff}}}\,\kappa_{\fg}(h,h)
= \frac{2}{k_{\mathrm{eff}}},
\label{eq:sl2-s-channel}
\\[4pt]
\textup{$t$-channel } (13|24)\colon\quad
&\frac{1}{k_{\mathrm{eff}}}\,
\kappa_{\fg}\bigl([e,e],[f,f]\bigr)
= \frac{1}{k_{\mathrm{eff}}}\,\kappa_{\fg}(0,0)
= 0,
\label{eq:sl2-t-channel}
\\[4pt]
\textup{$u$-channel } (14|23)\colon\quad
&\frac{1}{k_{\mathrm{eff}}}\,
\kappa_{\fg}\bigl([e,f],[f,e]\bigr)
= \frac{1}{k_{\mathrm{eff}}}\,\kappa_{\fg}(h,-h)
= \frac{-2}{k_{\mathrm{eff}}},
\label{eq:sl2-u-channel}
\end{align}
where $k_{\mathrm{eff}} = k + h^\vee = k + 2$.
The sum over channels gives
\begin{equation}\label{eq:sl2-boundary-quartic-sum}
(\mathfrak{C}_{\mathrm{aff}} \star_{P_{\mathrm{aff}}}
\mathfrak{C}_{\mathrm{aff}})(e,f,e,f)
\;=\;
\frac{2}{k+2} + 0 + \frac{-2}{k+2}
\;=\; 0,
\end{equation}
confirming the Jacobi cancellation on this specific input.
More generally, on the mixed input $(e,f,h,h)$:
\begin{align}
\textup{$s$-channel}\colon\quad
&\frac{1}{k_{\mathrm{eff}}}\,
\kappa_{\fg}(h,[h,h]) = 0,
\label{eq:sl2-mixed-s}
\\[4pt]
\textup{$t$-channel}\colon\quad
&\frac{1}{k_{\mathrm{eff}}}\,
\kappa_{\fg}([e,h],[f,h])
= \frac{1}{k_{\mathrm{eff}}}\,
\kappa_{\fg}(-2e,2f)
= \frac{-4}{k_{\mathrm{eff}}},
\label{eq:sl2-mixed-t}
\\[4pt]
\textup{$u$-channel}\colon\quad
&\frac{1}{k_{\mathrm{eff}}}\,
\kappa_{\fg}([e,h],[h,f])
= \frac{1}{k_{\mathrm{eff}}}\,
\kappa_{\fg}(-2e,-2f)
= \frac{4}{k_{\mathrm{eff}}},
\label{eq:sl2-mixed-u}
\end{align}
again summing to zero as required.
\end{proposition}

\begin{proof}
Direct substitution using the $\mathfrak{sl}_2$ bracket and
the invariant-form values~\eqref{eq:sl2-killing-values}. Each channel
computes $\kappa_{\fg}([x_i,x_j],[x_k,x_\ell])/k_{\mathrm{eff}}$
as in Corollary~\ref{cor:affine-boundary-quartic}. The
cancellation in~\eqref{eq:sl2-boundary-quartic-sum}
and~\eqref{eq:sl2-mixed-s}--\eqref{eq:sl2-mixed-u} is the
Jacobi identity evaluated on specific elements.
\end{proof}

\subsection{Genus loop for \texorpdfstring{$\mathfrak{sl}_2$}{sl(2)}}
\label{subsec:affine-sl2-genus-loop}
\index{genus loop operator!sl2 explicit@$\mathfrak{sl}_2$ explicit}

\begin{proposition}[Genus loop for $\mathfrak{sl}_2$;
\ClaimStatusProvedHere]\label{prop:affine-sl2-genus-loop}
For $\mathfrak{g} = \mathfrak{sl}_2$, one has
\begin{equation}\label{eq:sl2-genus-loop-formula}
\Lambda_P(\mathfrak{C}_{\mathrm{aff}}) \;=\; 0,
\end{equation}
confirming the general vanishing of
Theorem~\textup{\ref{thm:affine-genus-loop-weyl}}. The
genus-$1$ Hessian correction for $\widehat{\mathfrak{sl}}_2$
comes entirely from the boundary quartic
$\xi^*\mathfrak{Q}_{\mathrm{aff}}
= \mathfrak{C}_{\mathrm{aff}} \star_P
\mathfrak{C}_{\mathrm{aff}}$
\textup{(}Proposition~\textup{\ref{prop:affine-sl2-boundary-quartic})}.
\end{proposition}

\begin{proof}
Specialisation of
Theorem~\ref{thm:affine-genus-loop-weyl}; alternatively,
the direct trace
$\sum_a \kappa_{\fg}(x,[e_a,e^a])
= \kappa_{\fg}(x,[e,f]+[f,e]+[h,h/2])
= \kappa_{\fg}(x,h-h+0) = 0$.\qedhere
\end{proof}

\begin{remark}[Genus-$1$ correction from boundary sewing]
\label{rem:sl2-genus-loop-sugawara}
\index{genus loop!affine vanishing}
The vanishing $\Lambda_P(\mathfrak{C}_{\mathrm{aff}}) = 0$
means that for affine Kac--Moody, the genus-$1$ correction to
the modular Hessian is generated \emph{entirely} by boundary
sewing of two cubic vertices, not by a local one-loop diagram.
Algebraically, the Lie
bracket is antisymmetric while the propagator is symmetric.
\end{remark}

\begin{corollary}[Boundary-generated quartic nonlinearity;
\ClaimStatusProvedHere]\label{cor:affine-boundary-quartic}
\index{boundary shadow!affine quartic}
Even though the local quartic vertex vanishes, the quartic
boundary shadow is nonzero:
\[
\xi^*\mathfrak{Q}_{\mathrm{aff}}
= \mathfrak{C}_{\mathrm{aff}} \star_{P_{\mathrm{aff}}}
 \mathfrak{C}_{\mathrm{aff}}.
\]
Explicitly,
\[
(\mathfrak{C}_{\mathrm{aff}} \star_{P_{\mathrm{aff}}}
\mathfrak{C}_{\mathrm{aff}})(x_1,x_2,x_3,x_4)
\;=\;
\frac{1}{k_{\mathrm{eff}}}
\sum_{\mathrm{channels}}
\kappa_{\fg}([x_i,x_j],\,[x_k,x_\ell]),
\]
where $k_{\mathrm{eff}} = k + h^\vee$ is the normalisation of
the Hessian on the strict sector, and the sum runs over the
three $s$-, $t$-, $u$-channels of the four inputs.
\end{corollary}

\begin{proof}
Apply the quartic boundary recursion
(Theorem~\ref{thm:nms-shadow-master-equations}) with
$\mathfrak{Q}_{\mathrm{aff}} = 0$. The surviving term is the
sewing of two cubic current vertices with one propagator
$P_{\mathrm{aff}} = H_{\mathrm{aff}}^{-1}
= (1/k_{\mathrm{eff}})\,\kappa_{\fg}^{-1}$. Each channel pairs one
propagator contraction between two cubic vertices, yielding the
displayed formula.
\end{proof}

\begin{conjecture}[Critical-level emergence;
\ClaimStatusConjectured]\label{conj:affine-critical-level-emergence}
\index{critical level!cubic emergence}
\index{complementarity potential!cubic leading order}
In the formal critical limit of the strict current-level cyclic
model, the quadratic Hessian vanishes:
\[
H_{\widehat{\fg}_{-h^\vee}} = 0,
\]
while the cubic shadow survives:
\[
\mathfrak{C}_{\widehat{\fg}_{-h^\vee}}(x,y,z) = \kappa_{\fg}(x,[y,z]).
\]
The complementarity potential is therefore cubic at leading order:
\[
S_{\widehat{\fg}_{-h^\vee}}(a)
= \tfrac{1}{6}\,\kappa_{\fg}(a,[a,a]) + O(a^4).
\]
\end{conjecture}

\begin{remark}[Critical scalar versus cubic tensor]
\label{rem:critical-scalar-versus-cubic-tensor}
The quadratic part of the complementarity potential is governed
by the shifted-level factor $k + h^\vee$, which vanishes at the
critical level. The simple-pole bracket remains non-degenerate,
so the first surviving homogeneous term in the Taylor expansion
is the Lie-theoretic cubic tensor $\frac{1}{6}\kappa_{\fg}(a,[a,a])$.
Thus the scalar conductor
\(\kappa(\widehat{\fg}_{-h^\vee})=0\) vanishes while the full
Lie cubic tensor \(\kappa_{\fg}(x,[y,z])\) does not.  The
conjectural step is the existence of a critical cyclic model whose
leading jet is this formal limit.
\end{remark}

\begin{conjecture}[Critical self-intersection as Maurer--Cartan;
\ClaimStatusConjectured]\label{conj:affine-critical-mc}
\index{Maurer--Cartan locus!from critical-level complementarity}
At the critical level, the derived self-intersection of the two
modular Lagrangians has leading critical equation
\[
[a,a] = 0,
\]
the Maurer--Cartan locus for the underlying Lie bracket.
This bridges complementarity geometry to critical-level
representation theory: the locus $\{a : [a,a] = 0\}$ is the
Maurer--Cartan/flat-connection locus for the underlying Lie
bracket on the formal disc. The manuscript does not identify this
locus itself with the oper space; separately, the degree-zero
completed critical bar-Ext cohomology recovers the oper algebra
$\mathfrak{z} \cong \mathrm{Fun}(\mathrm{Op}_{\fg^\vee}(D))$
by Theorem~\ref{thm:oper-bar-h0}.
\end{conjecture}

\begin{remark}[Cubic critical equation]
Differentiate the critical-level cubic action:
$dS(a) = \tfrac{1}{2}[a,a] + O(a^3)$. The cubic truncation of
$\mathrm{Crit}(S)$ is therefore $\{a : [a,a] = 0\}$, the
Maurer--Cartan locus. Any relation to critical opers is indirect
and mediated by the separate bar-cohomological theorem
Theorem~\ref{thm:oper-bar-h0}; the present conjecture does not
identify the Maurer--Cartan locus with the oper space.
\end{remark}

\begin{conjecture}[SC deformation complex at critical level;
\ClaimStatusConjectured]\label{conj:sc-deformation-critical}
\index{critical level!SC deformation complex}
\index{Swiss-cheese!critical level deformation}
At the critical level $k = -h^\vee$ for $\fg$ simple of rank~$r$
with exponents $m_1 \le \cdots \le m_r$, the SC deformation complex
governing the chiral Swiss-cheese pair
$\bigl(\ChirHoch^*(\widehat{\fg}_{-h^\vee}),\, \widehat{\fg}_{-h^\vee}\bigr)$
exhibits the following structure:
\begin{enumerate}[label=\textup{(\roman*)}]
\item \textup{(Degree-$0$ jump)}\quad
 $\ChirHoch^0(\widehat{\fg}_{-h^\vee})
 = \mathfrak{z}(\widehat{\fg}_{-h^\vee})
 \cong \mathrm{Fun}(\mathrm{Op}_{\fg^\vee}(D))
 \cong \bC[\Theta_1, \ldots, \Theta_r]$,
 jumping from $\bC$ at generic level to
 the infinite-dimensional Feigin--Frenkel center.
\item \textup{(Degree-$1$ jump)}\quad
 $\ChirHoch^1(\widehat{\fg}_{-h^\vee})$ also jumps: at generic level it equals~$\fg$
 \textup{(}Proposition~\textup{\ref{prop:chirhoch1-affine-km}}\textup{)};
 at critical level the Beilinson--Drinfeld comparison
 \textup{(}Theorem~\textup{\ref{thm:affine-periodicity-critical}}\textup{)}
 gives the full odd-degree Gelfand--Fuchs component
 $\bigoplus_{i=1}^{r} P_i \otimes \bC[\Theta_1, \ldots, \Theta_r]$,
 which is infinite-dimensional.
\item \textup{(Full ChirHoch)}\quad
 The full chiral Hochschild cohomology at critical level is
 \[
 \ChirHoch^*(\widehat{\fg}_{-h^\vee})
 \;\cong\;
 \Lambda^*(P_1, \ldots, P_r) \otimes \bC[\Theta_1, \ldots, \Theta_r],
 \]
 with $\deg P_i = 2m_i + 1$ and $\deg \Theta_i = 2(m_i + 1)$,
 by the BD comparison and the Feigin--Tsygan computation.
 This is \emph{unbounded}, with polynomial growth $O(n^{r-1})$.
\item \textup{(SC moduli singularity)}\quad
 The generic-to-critical transition in the SC moduli is a
 singular degeneration: the finite-dimensional deformation
 complex $(\bC, \fg, \bC)$ degenerates to the infinite-dimensional
 complex governed by the Feigin--Frenkel center.
 The singularity is controlled by the curved level spectral sequence
 of Corollary~\textup{\ref{cor:critical-level-spectral}}: at
 $k = -h^\vee + \epsilon$
 the $E_1$~page is $\mathrm{Fun}(\mathrm{Op}) \otimes H^*(\fg; \bC)$,
 and the deformation parameter~$\epsilon$ governs the collapse
 on the curvature-flat comparison surface that recovers the
 Koszul regime.
\end{enumerate}
\end{conjecture}

\begin{remark}[Critical SC deformation: proved inputs and obstruction]
\label{rem:sc-deformation-critical-obligation}
Parts~(i) and~(iii) are proved.  Part~(i) is the Feigin--Frenkel theorem
(Theorem~\ref{thm:critical-level-structure}),
and (iii)~is Theorem~\ref{thm:affine-periodicity-critical}
via the BD comparison.
Part~(ii) follows from~(iii) by restriction to odd
cohomological degrees.
Part~(iv) requires a separate wall-crossing theorem for the curved
level spectral sequence with deformation parameter
\(\epsilon=k+h^\vee\).  At \(\epsilon=0\) the critical differential
is square-zero; at \(\epsilon\neq0\) the perturbation differentials
and the curvature operator
\[
[\epsilon\,\mathsf C_{\fg}/(2h^\vee),-]
\]
must be shown to identify the infinite-dimensional critical
\(E_1\)-page with the finite-dimensional Koszul comparison surface
after the chosen curvature-flat or coderived totalization.  This
identification is the remaining mathematical condition; it is not a formal
consequence of Corollary~\ref{cor:critical-level-spectral} alone.

The critical-level SC moduli are compared with the moduli of
$\fg^\vee$-opers on the formal disc
$\mathrm{Op}_{\fg^\vee}(D)$:
the center jump makes the SC base space
$\ChirHoch^0 = \mathrm{Fun}(\mathrm{Op})$,
so the SC pair at critical level is a module over the oper algebra.
Whether this module structure is governed by a formal deformation
problem on $\mathrm{Op}_{\fg^\vee}(D)$ remains open.
\end{remark}

\begin{theorem}[Vanishing of the genus loop on the affine cubic;
\ClaimStatusProvedHere]\label{thm:affine-genus-loop-weyl}
\index{genus loop operator!affine cubic vanishing}
The genus loop operator applied to the cubic shadow vanishes:
\begin{equation}\label{eq:affine-genus-loop-weyl}
\Lambda_P(\mathfrak{C}_{\mathrm{aff}}) \;=\; 0.
\end{equation}
Consequently, the genus-$1$ Hessian correction for affine
Kac--Moody comes entirely from the quartic boundary shadow
$\xi^*\mathfrak{Q}_{\mathrm{aff}}
= \mathfrak{C}_{\mathrm{aff}} \star_P \mathfrak{C}_{\mathrm{aff}}$
\textup{(}Corollary~\textup{\ref{cor:affine-boundary-quartic})},
not from the local cubic vertex.
\end{theorem}

\begin{proof}
The genus loop operator
\textup{(Definition~\textup{\ref{def:nms-genus-loop-operator}})}
contracts two slots of $\mathfrak{C}_{\mathrm{aff}}^{\mathrm{sym}}$
with the propagator~$P_{\mathrm{aff}}$. The symmetrised cubic is
$\mathfrak{C}_{\mathrm{aff}}^{\mathrm{sym}}(x,y,z)
= 3\kappa_{\fg}(x,[y,z])$, symmetric in all three arguments by
ad-invariance. In coordinates:
\[
(\Lambda_P \mathfrak{C}_{\mathrm{aff}}^{\mathrm{sym}})(x)
\;=\;
\frac{3}{k_{\mathrm{eff}}}
\sum_{a,b} \kappa_{\fg}(x,[e_a,e_b])\,(\kappa_{\fg}^{-1})^{ab}.
\]
The structure constants $f^c_{ab} = [e_a,e_b]^c$ are
antisymmetric in $(a,b)$ while $(\kappa_{\fg}^{-1})^{ab}$ is
symmetric. Their contraction vanishes identically.
\end{proof}

\begin{remark}[Phase index and archetype classification]%
\label{rem:affine-phase-index}
\index{nonlinear phase index!affine}
The nonlinear phase index
(Definition~\ref{def:nms-nonlinear-phase-index}) satisfies
$\nu_{\mathrm{nl}}(\widehat{\fg}_k) = 3$ at generic level: the
cubic shadow is the first nonzero nonlinear term, placing
affine Kac--Moody in the \emph{Lie/tree phase}. At the
critical level $k = -h^\vee$, the index remains $3$ but the
Hessian degenerates; the cubic term becomes leading rather than
subleading, and one must pass to the next jet of the
complementarity potential. The nonlinear shadow calculus detects the
critical degeneration; it does not construct the
Feigin--Frenkel center.  The center is supplied by the critical
affine theorem, while the vanishing of $H_{\mathrm{aff}}$ forces the
modular geometry to compare its cubic and higher jets with the
separate formal-disc oper theorems: under the local critical bar-Ext
comparison, degree-zero completed critical bar-Ext/Hochschild
cohomology gives the Feigin--Frenkel center and the full completed
critical package gives oper differential forms
\textup{(}Theorems~\textup{\ref{thm:oper-bar-h0} and
\ref{thm:oper-bar}}\textup{)}.
\end{remark}

\subsection{Level-independence of the cubic shadow product}
\label{subsec:affine-cubic-level-independence}
\index{cubic shadow!level-independence}
\index{shadow obstruction tower!level-independence of S3 kappa@level-independence of $S_3\kappa$}

The cubic shadow coefficient $S_3$ and the modular characteristic
$\kappa$ each depend on the level~$k$. Their product does not.

\begin{proposition}[Level-independence of the cubic shadow product;
\ClaimStatusProvedHere]\label{prop:km-cubic-shadow-level-independence}
\index{cubic shadow!level-independent product}
\index{modular characteristic!cubic product}
For $\widehat{\fg}_k$ at any non-critical level $k \neq -h^\vee$,
the product
\begin{equation}\label{eq:km-cubic-level-independence}
S_3\bigl(\widehat{\fg}_k\bigr) \cdot
\kappa\bigl(\widehat{\fg}_k\bigr)
\;=\;
\frac{2h^\vee}{3}
\end{equation}
depends only on the Lie type of~$\fg$, not on the level~$k$.
Equivalently,
\begin{equation}\label{eq:km-S3-explicit}
S_3\bigl(\widehat{\fg}_k\bigr)
\;=\;
\frac{2h^\vee}{3\kappa}
\;=\;
\frac{4(h^\vee)^2}{3\,\dim(\fg)\,(k+h^\vee)}.
\end{equation}
\end{proposition}

\begin{proof}
By Theorem~\ref{thm:affine-cubic-normal-form}, the cubic shadow
is $\mathfrak{C}_{\mathrm{aff}}(x,y,z) = \kappa_{\fg}(x,[y,z])$.
The scalar projection $S_3$ onto a primary line is the cubic
coupling divided by~$\kappa$, normalised by the structure-constant
convention (the invariant-form normalization absorbs a factor
of $\dim(\fg)/(2h^\vee)$). From the defining formula
$\kappa = \dim(\fg)(k+h^\vee)/(2h^\vee)$:
\[
S_3 \cdot \kappa
\;=\;
\frac{2h^\vee}{3\kappa} \cdot \kappa
\;=\;
\frac{2h^\vee}{3},
\]
which is independent of~$k$. The explicit formula for~$S_3$
follows by dividing both sides
by~$\kappa = \dim(\fg)(k+h^\vee)/(2h^\vee)$.
\end{proof}

\begin{remark}[Shadow metric for class~$L$]%
\label{rem:affine-shadow-metric-perfect-square}
\index{shadow metric!class L perfect square}
Since $S_4 = 0$ for affine Kac--Moody
\textup{(}Corollary~\textup{\ref{cor:affine-postnikov-termination})},
the critical discriminant
$\Delta = 8\kappa S_4$
\textup{(}Definition~\textup{\ref{def:shadow-metric})} vanishes.
The shadow metric on the primary line~$L$ is therefore a
perfect square:
\begin{equation}\label{eq:km-shadow-metric-perfect-square}
Q_L(t)
\;=\;
\bigl(2\kappa + 3S_3\,t\bigr)^2
\;=\;
4\kappa^2\Bigl(1 + \frac{h^\vee\,t}{\kappa}\Bigr)^{\!2},
\end{equation}
where the second form uses~\eqref{eq:km-cubic-level-independence}.
The shadow connection $\nabla^{\mathrm{sh}} = d - Q_L'/(2Q_L)\,dt$
has a simple pole at
\begin{equation}\label{eq:km-connection-pole}
t_0
\;=\;
-\,\frac{2\kappa}{3S_3}
\;=\;
-\,\frac{\kappa^2}{h^\vee},
\end{equation}
which moves quadratically in~$\kappa$ and inversely in~$h^\vee$.
\end{remark}

The following table records the level-independent product~$S_3 \kappa$
for every simple type.

\begin{table}[ht]
\centering
\caption{The level-independent product
$S_3 \cdot \kappa = 2h^\vee/3$ for simple Lie
algebras.}\label{tab:km-S3-kappa-table}
\begin{tabular}{lrrr}
\toprule
$\fg$ & $\dim(\fg)$ & $h^\vee$ & $S_3\cdot\kappa$ \\
\midrule
$\mathfrak{sl}_2$ & 3 & 2 & $4/3$ \\
$\mathfrak{sl}_3$ & 8 & 3 & $2$ \\
$\mathfrak{sl}_4$ & 15 & 4 & $8/3$ \\
$\mathfrak{so}_5 \cong \mathfrak{sp}_4$ & 10 & 3 & $2$ \\
$\mathfrak{so}_7$ & 21 & 5 & $10/3$ \\
$\mathfrak{so}_8$ & 28 & 6 & $4$ \\
$G_2$ & 14 & 4 & $8/3$ \\
$F_4$ & 52 & 9 & $6$ \\
$E_6$ & 78 & 12 & $8$ \\
$E_7$ & 133 & 18 & $12$ \\
$E_8$ & 248 & 30 & $20$ \\
\bottomrule
\end{tabular}
\end{table}

\noindent
The table is reproduced by
\texttt{compute/lib/theorem\_class\_l\_generating\_function\_engine.py}
for all simple types at multiple levels.

\subsection{Non-simply-laced types: shadow data for
\texorpdfstring{$B_2$, $G_2$, $F_4$}{B2, G2, F4}}%
\label{subsec:nsl-shadow-data}
\index{non-simply-laced!shadow data}
\index{shadow obstruction tower!non-simply-laced types}
\index{B2@$B_2$!shadow data}
\index{G2@$G_2$!shadow data}
\index{F4@$F_4$!shadow data}

For simply-laced $\fg$ the Coxeter number $h$ equals the dual
Coxeter number $h^\vee$, and the distinction is invisible. For
non-simply-laced types it is not: $h \neq h^\vee$, and every formula
involving the level shift must use $h^\vee$. The shadow obstruction tower
data depends only on $\dim(\fg)$ and $h^\vee$, through the universal
formula $\kappa = \dim(\fg)(k+h^\vee)/(2h^\vee)$. The Coxeter
number $h$ governs the exponent spectrum and the Strange Formula
$|\rho|^2 = dh/12$, but plays no role in the Feigin--Frenkel
involution or the curvature formula.

\begin{proposition}[Non-simply-laced shadow obstruction tower;
\ClaimStatusProvedHere]%
\label{prop:nsl-shadow-tower}
\index{non-simply-laced!class L}
Let $\fg$ be a non-simply-laced simple Lie algebra with
$\dim(\fg) = d$, dual Coxeter number $h^\vee$, Coxeter number
$h \neq h^\vee$, and lacing number $r^\vee = |\alpha_{\mathrm{long}}|^2/|\alpha_{\mathrm{short}}|^2$.
Then $\widehat{\fg}_k$ at non-critical level $k \neq -h^\vee$ has:
\begin{enumerate}[label=\textup{(\roman*)}]
\item \emph{Modular characteristic:}
 $\kappa(\widehat{\fg}_k) = d(k+h^\vee)/(2h^\vee)$, identical to
 the simply-laced formula with $h^\vee$ in place of $h$.
\item \emph{Feigin--Frenkel involution:}
 $k \mapsto k' = -k-2h^\vee$.
\item \emph{Complementarity:}
 $\kappa(\widehat{\fg}_k) + \kappa(\widehat{\fg}_{k'}) = 0$ and
 $c(\widehat{\fg}_k) + c(\widehat{\fg}_{k'}) = 2d$.
\item \emph{Shadow class:}
 $\widehat{\fg}_k$ is class $L$ with $r_{\max} = 3$. The
 cubic shadow $\mathfrak{C}(x,y,z) = \kappa_{\fg}(x,[y,z])$ is
 nonzero \textup{(}from the Lie bracket\textup{)}, and the
 quartic obstruction vanishes by the Jacobi identity, exactly as
 in the simply-laced case
 \textup{(}Theorem~\textup{\ref{thm:affine-cubic-normal-form})}.
\item \emph{Critical discriminant:}
 $\Delta = 8\kappa S_4 = 0$ \textup{(}from $S_4 = 0$\textup{)},
 so the shadow metric $Q_L$ is a perfect square and the
 tower terminates.
\item \emph{$r$-matrix data:}
 the collision residue is
 $r^{\mathrm{coll}}(z)=k\Omega_{\mathrm{tr}}/z$ in trace-form
 normalization, while the KZ coefficient is
 $r^{\mathrm{KZ}}(z) =
 \Omega_{\mathrm{KZ}}/\bigl((k{+}h^\vee)\,z\bigr)$.
 At the critical level $k = -h^\vee$ the KZ denominator vanishes,
 tracking the Sugawara singularity; the trace-form collision residue
 does not have this denominator.
\end{enumerate}
\end{proposition}

\begin{proof}
All statements follow from the proofs in
Theorem~\ref{thm:affine-cubic-normal-form} and
Corollary~\ref{cor:affine-postnikov-termination}, which use only:
(a)~the OPE $J^a(z)J^b(w) \sim k\kappa_{\fg,ab}/(z-w)^2 +
f^{ab}{}_c J^c(w)/(z-w)$, (b)~ad-invariance of the invariant form,
and~(c)~the Jacobi identity. None of these depend on the root
length ratio; the structure constants $f^{ab}{}_c$ satisfy the same
Jacobi identity regardless of the lacing number. The key
observation is that $h^\vee$ (not $h$) enters the Sugawara
denominator $k + h^\vee$, the level shift $k' = -k - 2h^\vee$,
and the kappa formula $\kappa = d(k+h^\vee)/(2h^\vee)$. The
complementarity identity $\kappa + \kappa' = 0$ is
immediate: $\kappa' = d(-k - 2h^\vee + h^\vee)/(2h^\vee)
= -d(k + h^\vee)/(2h^\vee) = -\kappa$.
\end{proof}

\begin{table}[ht]
\centering
\renewcommand{\arraystretch}{1.3}
\begin{tabular}{lccccccc}
\toprule
$\fg$ & $d$ & $h$ & $h^\vee$ & $r^\vee$ &
$\kappa(\widehat{\fg}_k)$ & $k'$ & $c + c'$ \\
\midrule
$B_2 \cong C_2$ & $10$ & $4$ & $3$ & $2$ &
 $\dfrac{5(k+3)}{3}$ & $-k-6$ & $20$ \\[6pt]
$G_2$ & $14$ & $6$ & $4$ & $3$ &
 $\dfrac{7(k+4)}{4}$ & $-k-8$ & $28$ \\[6pt]
$F_4$ & $52$ & $12$ & $9$ & $2$ &
 $\dfrac{26(k+9)}{9}$ & $-k-18$ & $104$ \\
\bottomrule
\end{tabular}
\caption{Shadow data for non-simply-laced affine Kac--Moody
algebras. All are class $L$ with $r_{\max} = 3$. The
kappa formula uses $h^\vee$, the level shift uses $2h^\vee$, and
$\kappa + \kappa' = 0$ in every case. Spot checks:
$\kappa(B_2, k{=}1) = 20/3$,
$\kappa(G_2, k{=}1) = 35/4$,
$\kappa(F_4, k{=}1) = 260/9$.
The entries are reproduced by
\texttt{test\_non\_simply\_laced\_shadows.py}.}%
\label{tab:nsl-shadow-data}
\end{table}

\begin{remark}[Root-length and Serre convention]
\label{rem:nsl-root-serre-convention}
\index{Serre relations!non-simply-laced convention}
\index{root length!non-simply-laced convention}
Long roots are normalized by $(\alpha,\alpha)=2$, and the Cartan
matrix is
\[
A_{ij}=\frac{2(\alpha_i,\alpha_j)}{(\alpha_i,\alpha_i)}.
\]
The Serre relation is
\[
(\operatorname{ad}E_i)^{1-A_{ij}}(E_j)=0,
\qquad i\neq j,
\]
so it appears in bar degree $2-A_{ij}$. The bar-complex tables give
the following rank-two checks:
\[
\begin{array}{c|c|c|c}
\text{type} & \text{chosen simple roots} & \text{largest Serre relation}
& \text{bar degree}\\
\hline
B_2 & \alpha_1\ \text{long},\ \alpha_2\ \text{short}
& (\operatorname{ad}E_2)^3(E_1)=0 & 4\\
G_2 & \alpha_1\ \text{short},\ \alpha_2\ \text{long}
& (\operatorname{ad}E_1)^4(E_2)=0 & 5\\
F_4 & \text{one double edge, lacing }2
& (\operatorname{ad}E_{\mathrm{short}})^3(E_{\mathrm{long}})=0 & 4
\end{array}
\]
These degrees are presentation syzygies in the Chevalley--Serre
generators. They do not alter the shadow depth: the affine OPE still
has only the invariant-form double pole and the Lie-bracket simple
pole, the cubic tensor is \(\kappa_{\fg}(x,[y,z])\), and the quartic
shadow vanishes by Jacobi. Thus every non-critical non-simply-laced
row above remains class~\(L\) with \(r_{\max}=3\).
\end{remark}

\begin{remark}[The $h$ versus $h^\vee$ pitfall]%
\label{rem:nsl-h-vs-h-dual}
\index{dual Coxeter number!versus Coxeter number}
For non-simply-laced $\fg$, the dual Coxeter number $h^\vee$
is strictly less than the Coxeter number $h$, but there is no
universal formula relating them through the lacing number
alone. The explicit values are:
$B_n\colon h = 2n,\; h^\vee = 2n{-}1$;\quad
$C_n\colon h = 2n,\; h^\vee = n{+}1$;\quad
$G_2\colon h = 6,\; h^\vee = 4$;\quad
$F_4\colon h = 12,\; h^\vee = 9$.
The discrepancy is maximal
for $G_2$: $h/h^\vee = 3/2$, $r^\vee = 3$.
Substituting $h$ for $h^\vee$ in the kappa formula gives a
wrong answer. For $G_2$ at $k = 1$:
\[
\kappa_{\text{correct}} = \frac{7 \cdot 5}{4} = \frac{35}{4},
\qquad
\kappa_{\text{wrong}} = \frac{14 \cdot 7}{12} = \frac{49}{6};
\qquad
\text{ratio } = \frac{14}{15} \neq 1.
\]
The source of the error: $h^\vee = 1 + \sum_i a_i^\vee$ sums
the \emph{marks} (coefficients of coroots in the highest coroot),
while $h = 1 + \sum_i a_i$ sums the \emph{labels} (coefficients
of roots in the highest root). For simply-laced types, marks
equal labels; for non-simply-laced types, they differ on roots of
different length.
\end{remark}

\begin{remark}[The $B_2 \cong C_2$ isomorphism]%
\label{rem:b2-c2-isomorphism}
\index{B2@$B_2$!isomorphism with $C_2$}
The isomorphism $\mathfrak{so}_5 \cong \mathfrak{sp}_4$ implies
$B_2$ and $C_2$ have identical Lie-algebraic data: $d = 10$,
$h^\vee = 3$, $h = 4$, and therefore the same kappa formula,
central charge, and shadow obstruction tower. The lacing number is $r^\vee = 2$
in both cases. This is a rank-2 accident; for $n \geq 3$,
$B_n$ and $C_n$ are distinct ($B_n$ has $h^\vee = 2n-1$,
$C_n$ has $h^\vee = n+1$).

The bar complex computations confirming the class $L$ structure
for $B_2$ and $G_2$ are carried out in detail in
\S\ref{sec:B2-details} and \S\ref{sec:G2-details} respectively,
where the OPEs, Serre relations, and bar differentials are
computed explicitly.
\end{remark}


\subsection{Unified exceptional shadow table and the Deligne--Cvitanovi\'c series}%
\label{subsec:unified-exceptional-table}
\index{exceptional Lie algebra!unified shadow table}
\index{Deligne--Cvitanovic series@Deligne--Cvitanovi\'c series}
\index{shadow obstruction tower!exceptional types unified}

We combine the simply-laced data of
Proposition~\ref{prop:exceptional-shadow-invariants} with the
non-simply-laced data of Proposition~\ref{prop:nsl-shadow-tower}
into a single table covering all five exceptional types.

\begin{proposition}[Complete exceptional shadow data; \ClaimStatusProvedHere]%
\label{prop:complete-exceptional-shadow}
\index{exceptional Lie algebra!complete shadow data}
For every exceptional simple Lie algebra
$\fg \in \{G_2, F_4, E_6, E_7, E_8\}$ at non-critical affine
level $k \neq -h^\vee$, the shadow obstruction tower data is as follows.
\begin{center}
\renewcommand{\arraystretch}{1.4}
\begin{tabular}{lccccccc}
\toprule
$\fg$ & $\dim$ & $r$ & $h$ & $h^\vee$ & $r^\vee$
 & $\kappa(\widehat{\fg}_k)$ & Class \\
\midrule
$G_2$ & $14$ & $2$ & $6$ & $4$ & $3$
 & $7(k+4)/4$ & L \\
$F_4$ & $52$ & $4$ & $12$ & $9$ & $2$
 & $26(k+9)/9$ & L \\
$E_6$ & $78$ & $6$ & $12$ & $12$ & $1$
 & $13(k+12)/4$ & L \\
$E_7$ & $133$ & $7$ & $18$ & $18$ & $1$
 & $133(k+18)/36$ & L \\
$E_8$ & $248$ & $8$ & $30$ & $30$ & $1$
 & $62(k+30)/15$ & L \\
\bottomrule
\end{tabular}
\end{center}
Every entry has the affine Kac--Moody complementarity
$\kappa + \kappa' = 0$
\textup{(}$k' = -k - 2h^\vee$\textup{)}, $c + c' = 2\dim\fg$,
shadow depth $r_{\max} = 3$, critical discriminant
$\Delta = 8\kappa S_4 = 0$, and
$r$-matrix data
$r^{\mathrm{coll}}(z)=k\Omega_{\mathrm{tr}}/z$ and
$r^{\mathrm{KZ}}(z)=
\Omega_{\mathrm{KZ}}/\bigl((k{+}h^\vee)\,z\bigr)$
\textup{(}trace-form collision and KZ normalisations\textup{)}.
\end{proposition}

\begin{proof}
Each row is a special case of the universal affine Kac--Moody
shadow theorem (Theorem~\ref{thm:affine-cubic-normal-form}
and Corollary~\ref{cor:affine-postnikov-termination}). The
simplified $\kappa$ forms use:
$\gcd(14, 8) = 2$ for $G_2$,
$\gcd(52, 18) = 2$ for $F_4$,
$\gcd(78, 24) = 6$ for $E_6$,
$\gcd(133, 36) = 1$ for $E_7$ (already in lowest terms),
$\gcd(248, 60) = 4$ for $E_8$.
The complementarity $\kappa + \kappa' = 0$ is immediate:
$\kappa' = d(-k-2h^\vee+h^\vee)/(2h^\vee)
= -d(k+h^\vee)/(2h^\vee) = -\kappa$.
The exceptional entries are reproduced by
\texttt{test\_exceptional\_shadow\_complete.py}.
\end{proof}

\begin{remark}[Absence of cubic Casimir invariants]%
\label{rem:exceptional-no-cubic-casimir}
\index{Casimir invariant!cubic!absence for exceptionals}
\index{cubic shadow!versus cubic Casimir}
The Casimir invariant degrees (exponents $+ 1$) for each
exceptional type are:
\begin{center}
\renewcommand{\arraystretch}{1.2}
\begin{tabular}{ll}
\toprule
$\fg$ & Casimir degrees \\
\midrule
$G_2$ & $2,\, 6$ \\
$F_4$ & $2,\, 6,\, 8,\, 12$ \\
$E_6$ & $2,\, 5,\, 6,\, 8,\, 9,\, 12$ \\
$E_7$ & $2,\, 6,\, 8,\, 10,\, 12,\, 14,\, 18$ \\
$E_8$ & $2,\, 8,\, 12,\, 14,\, 18,\, 20,\, 24,\, 30$ \\
\bottomrule
\end{tabular}
\end{center}
None of the exceptional types possess a degree-$3$ Casimir
invariant (i.e.\ a symmetric cubic tensor $d_{abc}$ invariant
under the adjoint action). Among simple Lie algebras, only
$A_n$ for $n \ge 2$ (degrees include $3$) have such an invariant.
This does \emph{not} affect the cubic shadow $S_3 \neq 0$: the
cubic shadow comes from the \emph{antisymmetric} structure
constants $f^{ab}{}_c$ (the Lie bracket), not from the symmetric
$d_{abc}$. The distinction is between the Chevalley--Eilenberg
$3$-cocycle (which uses $f$ and exists for every Lie algebra) and
the third-order Casimir (which uses $d$ and exists only when
$3 \in \{\text{exponents}\} + 1$).
\end{remark}

\begin{proposition}[Anomaly ratios for exceptional principal $\mathcal{W}$-algebras;
\ClaimStatusProvedHere]%
\label{prop:exceptional-anomaly-ratios}
\index{anomaly ratio!exceptional W-algebras@exceptional $\mathcal{W}$-algebras}
\index{W-algebra@$\mathcal{W}$-algebra!exceptional types!anomaly ratio}
For the principal $\mathcal{W}$-algebra $\mathcal{W}^k(\fg)$ of
an exceptional Lie algebra $\fg$, the anomaly ratio
$\rho(\fg) = \sum_{i=1}^r 1/(m_i + 1)$
\textup{(}where $m_1, \ldots, m_r$ are the exponents\textup{)}
satisfies $\kappa(\mathcal{W}(\fg), k) = \rho(\fg) \cdot c(\mathcal{W}(\fg), k)$.
The explicit values are:
\begin{center}
\renewcommand{\arraystretch}{1.3}
\begin{tabular}{lll}
\toprule
$\fg$ & Exponents & $\rho(\fg)$ \\
\midrule
$G_2$ & $1,\, 5$ & $1/2 + 1/6 = 2/3$ \\
$F_4$ & $1,\, 5,\, 7,\, 11$ & $1/2 + 1/6 + 1/8 + 1/12 = 7/8$ \\
$E_6$ & $1,\, 4,\, 5,\, 7,\, 8,\, 11$ & $1/2 + 1/5 + 1/6 + 1/8 + 1/9 + 1/12 = 427/360$ \\
$E_7$ & $1,\, 5,\, 7,\, 9,\, 11,\, 13,\, 17$ & $\displaystyle\sum_{m} \tfrac{1}{m+1} = \tfrac{2777}{2520}$ \\[4pt]
$E_8$ & $1,\, 7,\, 11,\, 13,\, 17,\, 19,\, 23,\, 29$ & $\displaystyle\sum_m \tfrac{1}{m+1}$ \\
\bottomrule
\end{tabular}
\end{center}
\end{proposition}

\begin{proof}
Direct computation from the exponent tables. The identity
$\kappa = \rho \cdot c$ follows from the general $\mathcal{W}$-algebra
anomaly ratio formula (Theorem~D and
Chapter~\ref{chap:w-algebras}).
\end{proof}

\begin{remark}[The Deligne--Cvitanovi\'c exceptional series]%
\label{rem:deligne-cvitanovic}
\index{Deligne--Cvitanovic series@Deligne--Cvitanovi\'c series!shadow data}
The exceptional Lie algebras fit into a parametric family
along the chain $A_1 \subset G_2 \subset F_4 \subset E_6
\subset E_7 \subset E_8$, parametrized by $h^\vee$
(Deligne \cite{Deligne96}, Cvitanovi\'c \cite{Cvitanovic08}).
At level $k = 1$, the $\kappa$ values along this chain are:
\[
\kappa(\widehat{\fg}_1)
\;=\; \frac{\dim(\fg)(1 + h^\vee)}{2h^\vee}
\;=\; \frac{9}{4},\;\; \frac{35}{4},\;\; \frac{260}{9},\;\; \frac{169}{4},\;\; \frac{2527}{36},\;\; \frac{1922}{15}
\]
The progression tracks the growth of $\dim(\fg)$ with $h^\vee$;
the shadow obstruction tower structure (class~L, depth~$3$ uniformly) is
constant along the entire chain, reflecting the universal role of
the Lie bracket in determining the cubic shadow.

The Langlands self-duality of all exceptional types
($G_2^L = G_2$, $F_4^L = F_4$, $E_n^L = E_n$) implies that
the shadow obstruction tower has no Langlands asymmetry: the
level-reflected current presentation underlying the Koszul dual is
obtained by the Feigin--Frenkel involution
$k \mapsto -k - 2h^\vee$ within the same Lie type, while the dual
object itself remains the chiral CE algebra. This contrasts with the
non-exceptional pair $B_n$, $C_n$, where Langlands duality
relates algebras of different types ($B_n^L = C_n$), and the
shadow obstruction towers of $B_n$ and $C_n$ have different $h^\vee$
\textup{(}$2n-1$ versus $n+1$\textup{)} and therefore different $\kappa$
formulas.
\end{remark}

\begin{remark}[Theta function of the $E_8$ root lattice]%
\label{rem:e8-theta-function}
\index{E8 root lattice@$E_8$ root lattice!theta function}
\index{Eisenstein series!E4@$E_4$!and $E_8$ lattice}
The theta function of the $E_8$ root lattice is
\[
\Theta_{E_8}(\tau)
\;=\; \sum_{v \in \Lambda_{E_8}} q^{|v|^2/2}
\;=\; E_4(\tau)
\;=\; 1 + 240q + 2160q^2 + 6720q^3 + \cdots,
\]
the weight-$4$ Eisenstein series for $\mathrm{SL}(2,\bZ)$.
The coefficient of $q^1$ is $240 = 2 \cdot |\Delta_+(E_8)| = 2 \cdot 120$
(the number of roots of~$E_8$). This identity
$\Theta_{E_8} = E_4$ is a manifestation of the unimodularity
and evenness of $\Lambda_{E_8}$: the unique even unimodular
lattice in dimension~$8$ has its theta function equal to the
unique modular form of weight~$4$. The shadow obstruction tower of the
$E_8$ lattice VOA $V_{\Lambda_{E_8}}$ inherits this modular
property through the lattice sector decomposition
(the lattice construction chapter).
\end{remark}


%% ================================================================
\section{The affine vertex algebra as holomorphic-topological
Chern--Simons theory}
\label{sec:affine-ht-chern-simons}
\index{Chern--Simons theory!holomorphic-topological}
\index{affine vertex algebra!as gauge theory}

The gauge-theoretic interpretation used here is the boundary HT
Chern--Simons model whose classical PVA $\lambda$-bracket has the
affine form.  The comparison is made at the level of local brackets and
BRST/BV complexes; it is not an additional identification of
$A$, $\barB(A)$, $A^!$, or the derived centre.

\subsection{The affine PVA \texorpdfstring{$\lambda$}{lambda}-bracket}

\begin{definition}[Affine PVA $\lambda$-bracket]
\label{def:affine-pva-lambda-bracket}
\index{PVA lambda-bracket@PVA $\lambda$-bracket!affine}
Let $\fg$ be a simple Lie algebra with structure constants
$f^c_{ab}$ and normalized invariant form $\kappa_{ab}$. The \emph{affine PVA
$\lambda$-bracket} at level $k$ is
\begin{equation}\label{eq:affine-pva-lambda-bracket}
 \{B_a \mathbin{{}_\lambda} B_b\}
 \;=\;
 f^c_{ab}\, B_c \;+\; k\,\kappa_{ab}\,\lambda.
\end{equation}
This is the classical $\lambda$-bracket of the Poisson vertex
algebra underlying $\widehat{V}_k(\fg)$.
\end{definition}

\subsection{The 3d HT action}

\begin{proposition}[The holomorphic-topological Chern--Simons action;
\ClaimStatusProvedElsewhere]
\label{prop:affine-cs-action}
\index{Chern--Simons theory!HT action}
The 3d holomorphic-topological field theory whose boundary chiral
algebra is $\widehat{V}_k(\fg)$ has action functional
\begin{equation}\label{eq:affine-cs-action}
 S
 \;=\;
 \int \mathrm{d}z\,
 \Bigl[
 B_a\,(d_t + \bar{\partial})\,A^a
 \;+\;
 \tfrac{1}{2}\,f^c_{ab}\,B_c\,A^b A^c
 \;+\;
 \tfrac{1}{2}\,k\,\kappa_{ab}\,A^a\,\partial A^b
 \Bigr],
\end{equation}
where $(A^a, B_a)$ are the gauge and matter fields on
$\Sigma \times \mathbb{R}_t$, with $A^a$ valued in
$\Omega^{0,1}(\Sigma) \otimes \fg$ and $B_a$ valued in
$\Omega^{1,0}(\Sigma) \otimes \fg^*$.
\end{proposition}

\begin{proof}
This is the standard HT twist of 3d $\mathcal{N}=2$ Chern--Simons
theory with gauge group $G$; see \cite{Costello2013} and
\cite{CostelloGaiotto2020}. The boundary VOA is computed by
quantising the $B$-field zero-modes on $\partial(\Sigma \times
\mathbb{R}_{\ge 0})$, which recovers the affine OPE at level $k$.
\end{proof}

\begin{remark}[Gauge transformations]
\label{rem:affine-cs-gauge}
\index{gauge transformation!affine HT CS}
The infinitesimal gauge transformations generated by
$\varepsilon^a \in \Omega^0(\Sigma) \otimes \fg$ are
\begin{align*}
 \delta B_a
 &\;=\;
 -f^c_{ab}\,B_c\,\varepsilon^b
 \;+\;
 k\,\kappa_{ab}\,\partial \varepsilon^b,
 \\
 \delta A^a
 &\;=\;
 -(d_t + \bar{\partial})\varepsilon^a
 \;-\;
 f^a_{bc}\,\varepsilon^b A^c.
\end{align*}
When $k \neq 0$, the field redefinition
$A \mapsto A$, $B \mapsto B - (k/2)\partial A$ absorbs the
$B$-field kinetic term into a pure Chern--Simons form, recovering
analytically continued Chern--Simons theory $\mathrm{CS}_k(G_\mathbb{C})$
on $\Sigma \times \mathbb{R}_t$.
\end{remark}

\begin{remark}[Bar complex and BRST complex]
\label{rem:bar-as-brst-cs}
\index{bar complex!as BRST complex}
\index{BRST differential!from bar differential}
On the boundary HT Chern--Simons model, the CE/BRST differential has
the same local structure constants and level term as the affine bar
differential.  Under the boundary field dictionary, the desuspended bar
generators play the role of ghosts and the bar differential records the
BRST variations above.  This is a genus-$0$ comparison of different
models, not an equality of the bar coalgebra with the full BV complex.
Bar-cobar inversion recovers the boundary chiral algebra from its bar
coalgebra on the completed Koszul surface; it does not identify the
Koszul dual with the gauge theory itself.
\end{remark}

\begin{remark}[Ghost dictionary]
\label{rem:affine-ghost-dictionary}
\index{ghost field!affine bar complex dictionary}
\index{bar complex!ghost identification}
The identification of Remark~\ref{rem:bar-as-brst-cs} gives an
explicit dictionary between algebraic and physical generators:
\begin{center}
\renewcommand{\arraystretch}{1.3}
\begin{tabular}{lll}
\toprule
\textbf{Bar complex} & \textbf{BRST complex} & \textbf{Role} \\
\midrule
Desuspended generator $s^{-1}B_a$ & Ghost $c^a$ & Gauge parameter \\
Field $A^a$ & Gauge field $A^a$ & Connection \\
$d_{\bar{B}}(s^{-1}B_a)$ & $Q_{\mathrm{BRST}}(c^a)$ & Gauge variation \\
Bar length $n$ & Ghost number $n$ & Filtration degree \\
$d_{\bar{B}}^{\,2}=[m_0,-]$ & $Q_{\mathrm{BRST}}^2$ equals curvature/anomaly & Curved square \\
\bottomrule
\end{tabular}
\end{center}
The genus-$0$ projection of
Proposition~\ref{prop:chriss-ginzburg-structure} identifies the common
twisting morphism
$\tau = \Theta\big|_{g=0}$ of Remark~\ref{rem:five-from-one}, row~A.
The algebraic bar differential and the physical BRST differential are
two realizations of this twisting morphism in their respective
categories.
\end{remark}

\begin{remark}[Physical bulk and cochain closed sector]
\label{rem:affine-physical-bulk-cochain-sector}
\index{bulk!physical versus cochain closed sector}
\index{chiral derived center!affine boundary comparison}
The HT Chern--Simons bulk is a three-dimensional field theory with
boundary chiral algebra \(\widehat V_k(\fg)\). The cochain closed
sector attached to the boundary algebra is
\[
Z^{\mathrm{der}}_{\mathrm{ch}}(\widehat V_k(\fg))
= C^\bullet_{\mathrm{ch}}(\widehat V_k(\fg),\widehat V_k(\fg)).
\]
The boundary condition supplies a bulk-to-boundary map landing in this
Hochschild object. It does not identify the BV field complex of the
three-dimensional theory with the bar coalgebra or with the Koszul
dual. The comparison above is local and genus~\(0\): the same
structure constants and level term define the boundary BRST variation
and the affine bar differential.
\end{remark}

\subsection{Level-rank duality as a bar-cobar test}
\label{sec:level-rank-bar-cobar}
\index{level-rank duality!as bar-cobar equivalence}
\index{Koszul duality!level-rank}

Level-rank duality is an external equivalence of boundary VOAs and
conformal-block theories.  In this chapter it is used as a test for
bar-cobar inversion: an isomorphism of boundary VOAs should lift to
their bar coalgebras, while the Koszul dual still passes through the
chiral CE object.

\begin{proposition}[Level-rank duality for boundary VOAs;
\ClaimStatusProvedElsewhere]
\label{prop:level-rank-boundary-voa}
\index{level-rank duality!boundary VOA isomorphism}
\index{Costello--Dimofte--Gaiotto}
\textup{(Costello--Dimofte--Gaiotto \cite{CDG2023}, \S6.6.)}
The boundary chiral algebras of $U(N)_k$ and $U(k)_{-N}$
holomorphic-topological Chern--Simons theories, equipped with
appropriate dual boundary conditions, are isomorphic as vertex
algebras:
\begin{equation}\label{eq:level-rank-voa-iso}
 \mathrm{VOA}_{U(N)_k}
 \;\cong\;
 \mathrm{VOA}_{U(k)_{-N}}.
\end{equation}
\end{proposition}

\begin{corollary}[Bar-complex intertwining; \ClaimStatusProvedHere]
\label{cor:level-rank-bar-intertwining}
\index{bar complex!level-rank intertwining}
At the level of bar complexes, the isomorphism
\eqref{eq:level-rank-voa-iso} lifts to an equivalence of dg
coalgebras:
\begin{equation}\label{eq:level-rank-bar-equiv}
 \bar{B}^{\mathrm{ch}}\bigl(\mathrm{VOA}_{U(N)_k}\bigr)
 \;\simeq\;
 \bar{B}^{\mathrm{ch}}\bigl(\mathrm{VOA}_{U(k)_{-N}}\bigr).
\end{equation}
In particular, the bar complex of the boundary
\(\widehat{\mathfrak{gl}}_N\)-type presentation at level~$k$ is related
to the boundary \(\widehat{\mathfrak{gl}}_k\)-type presentation at
level~$-N$. This is a level-rank comparison of boundary VOAs, not an
identification of the chiral Koszul dual with an affine algebra; the
Koszul dual still passes through the corresponding chiral CE object.
\end{corollary}

\begin{remark}[Relation to Feigin--Frenkel]
\label{rem:level-rank-feigin-frenkel}
\index{Feigin--Frenkel duality!and level-rank}
The Feigin--Frenkel involution $k \mapsto -k - 2h^\vee$ is comparable
with level-rank duality, but it is not literally the
specialization $N=k$. The former is the Verdier level shift for a fixed
simple Lie algebra; the latter exchanges the boundary gauge data
$U(N) \leftrightarrow U(k)$.
\end{remark}

\begin{proposition}[Kappa anti-symmetry under Feigin--Frenkel involution]
\label{prop:kappa-anti-symmetry-ff}
\ClaimStatusProvedHere
\index{Feigin--Frenkel duality!kappa anti-symmetry}
\index{kappa@$\kappa$!Feigin--Frenkel anti-symmetry}
For any simple Lie algebra $\mathfrak{g}$ at non-critical level
$k \neq -h^\vee$:
\begin{equation}\label{eq:kappa-ff-antisymmetry}
 \kappa\bigl(\widehat{V}_k(\fg)\bigr)
 \;+\;
 \kappa\bigl(\widehat{V}_{-k-2h^\vee}(\fg)\bigr)
 \;=\; 0.
\end{equation}
The Verlinde level-rank swap
$\dim \mathcal{V}_1(\mathfrak{sl}_N, k)
= \dim \mathcal{V}_1(\mathfrak{sl}_k, N)$
is a genus-$1$ numerical symmetry for all positive integers \(N,k\).
It is not the Feigin--Frenkel involution: level-rank changes the
finite Lie algebra, while Feigin--Frenkel keeps \(\fg\) fixed and
sends \(k\) to \(-k-2h^\vee\). Full boundary VOA rank-level
statements require the usual center and \(\mathfrak u(1)\) data.
\end{proposition}

\begin{proof}
Write $d = \dim\fg$. The modular characteristic is
$\kappa(\widehat{V}_k(\fg)) = (k + h^\vee)\,d/(2h^\vee)$
\textup{(}Theorem~D\textup{)}.
Under the Feigin--Frenkel involution
$k \mapsto k' := -k - 2h^\vee$:
\[
 \kappa(\widehat{V}_{k'}(\fg))
 \;=\;
 \frac{(-k - 2h^\vee + h^\vee)\,d}{2h^\vee}
 \;=\;
 \frac{(-k - h^\vee)\,d}{2h^\vee}.
\]
Summing:
\[
 \kappa(\widehat{V}_k(\fg)) + \kappa(\widehat{V}_{k'}(\fg))
 \;=\;
 \frac{(k + h^\vee)\,d}{2h^\vee}
 \;+\;
 \frac{(-k - h^\vee)\,d}{2h^\vee}
 \;=\; 0.
\]

For the Verlinde swap: at genus~$1$,
$\dim\mathcal{V}_1(\mathfrak{sl}_N, k) = \binom{N+k-1}{N-1}$.
The identity $\binom{N+k-1}{N-1} = \binom{N+k-1}{k-1}$ is equivalent
to $\binom{n}{a} = \binom{n}{n-a}$ with $n = N+k-1$, $a = N-1$,
which gives $\dim\mathcal{V}_1(\mathfrak{sl}_N, k)
= \dim\mathcal{V}_1(\mathfrak{sl}_k, N)$. This binomial identity
changes \(N\) and \(k\), so it cannot be a specialization of
\(k\mapsto -k-2h^\vee\) for a fixed simple algebra. The extra
center/\(\mathfrak u(1)\) data belongs to the boundary VOA
rank-level theorem, not to the affine \(\kappa\)-anti-symmetry
proved above.
\end{proof}

\begin{conjecture}[Complementarity as the level-rank pairing;
\ClaimStatusConjectured]
\label{conj:level-rank-complementarity}
\index{complementarity!level-rank pairing}
\index{level-rank duality!as complementarity}
For the affine vertex algebra $\widehat{V}_k(\fg)$ at non-critical
level $k \neq -h^\vee$, complementarity (Theorem~C) specialises to
\begin{equation}\label{eq:level-rank-complementarity}
 Q_g\bigl(\widehat{V}_k(\fg)\bigr)
 \;+\;
 Q_g\bigl(\widehat{V}_{-k-2h^\vee}(\fg)\bigr)
 \;=\;
 H^*\bigl(\overline{\mathcal{M}}_g,\;
 Z(\widehat{V}_k(\fg))\bigr).
\end{equation}
When $\fg = \mathfrak{gl}_N$ and $k \in \mathbb{Z}_{>0}$, the
conjectural comparison with level-rank duality
$U(N)_k \leftrightarrow U(k)_{-N}$ requires the
$\mathfrak{u}(1)$ factor and the boundary-condition data of
Proposition~\ref{prop:level-rank-boundary-voa}.  It is not the
Feigin--Frenkel involution for a fixed simple Lie algebra, and the case
$N=k$ is not a critical-level statement.  The proposed identity is the
cross-polarization projection of the MC equation
$D\Theta_{\widehat{V}_k(\fg)} +
\tfrac{1}{2}[\Theta_{\widehat{V}_k(\fg)},
\Theta_{\widehat{V}_k(\fg)}] = 0$
(Remark~\ref{rem:five-from-one}, row~C): the $\cA$-side of the
cross-polarized graph sum operates at level~$k$, the $\cA^!$-side
at level~$-k-2h^\vee$, and the propagator $P_{\cA}$ mediates
between them.
\end{conjecture}

\begin{remark}[Test of bar-cobar inversion]
\label{rem:level-rank-thm-b-test}
\index{bar-cobar inversion!level-rank test}
Equation~\eqref{eq:level-rank-bar-equiv} is a non-trivial test of
the bar-cobar counit in Theorem~A, and of its completed Positselski
extension in Theorem~B. Under the boundary VOA isomorphism of
Proposition~\ref{prop:level-rank-boundary-voa}, the quasi-isomorphism
$\Omega(\bar{B}^{\mathrm{ch}}(A)) \xrightarrow{\;\sim\;} A$
is conjectured to intertwine the level-rank substitution
$(N,k) \leftrightarrow (k,-N)$. Concretely, the cobar differential
on $\bar{B}^{\mathrm{ch}}(\mathrm{VOA}_{U(N)_k})^\vee$ encodes
the OPE of $\mathrm{VOA}_{U(k)_{-N}}$, and bar-cobar inversion
recovers this OPE from the coalgebra structure when the boundary
duality input is supplied.
\end{remark}

\section{Synthesis}
\label{sec:km-supplement}

\begin{remark}[Kac--Moody and BKM: Hall--Drinfeld double at K3$\times E$]
\label{rem:km-K3xE}
\index{K3!Hall--Drinfeld double}
\index{Yangian!taxonomy!K3 Hall--Drinfeld double}
Beyond Kac--Moody $\widehat{\mathfrak{g}}_k$, the K3 chiral
bialgebra of Vol~III Chapter~\ref{V3-ch:k3-chiral-bialgebra-platonic}
arises from the generalized Borcherds--Kac--Moody superalgebra
\(\mathfrak{g}_{\Delta_5}\) on
\(\Lambda^{2,1}_{II} \subset \Lambda^{3,2}\). This is not a
Drinfeld Yangian. A BKM algebra has imaginary simple roots and no
finite-type Drinfeld \(J\)-presentation or RTT presentation attached
to a finite-dimensional simple \(\fg\). The quantum-group object is
the K3 chiral Hall--Drinfeld double
\[
\mathcal{D}_\hbar\!\left(
\mathcal{Y}^{\mathrm{Hall}}_\hbar(\mathrm{CoHA}_{K3 \times E})
\right),
\]
with the Schiffmann--Vasserot shuffle product and the Davison
CY-\(3\) extension. Its Cartan data form a Hall/BKM Manin-pair
geometry on the signature-\((2,1)\) lattice, not a Drinfeld Yangian
triple.
\end{remark}

\begin{remark}[Frenkel-Kac closure at K3 Fock]
\label{rem:km-frenkel-kac-K3}
At Lie/Hopf level the K3 chiral bialgebra is abelian: $24$ Miki $U_{q_1, q_2}(\widehat{\widehat{\mathfrak{gl}}}_1)$ with $\widetilde{M}_{24}$-twisted permutation. BKM non-abelianity $\mathfrak{g}_{\Delta_5}$ emerges only via Frenkel-Kac $1980$ vertex closure on $\mathcal{F}_{K3} = \mathrm{Sym}(\bigoplus_{n \ge 1} H^*(K3) \otimes t^{-n})$: $V_\alpha(z) V_\beta(w) \to V_{\alpha + \beta}$ with Borcherds $2$-cocycle gives $[e_\alpha, e_\beta] = N_{\alpha, \beta}\, e_{\alpha + \beta}$.
\end{remark}
